% 本文件由 LaTeX 排版 Agent 根据有效 translation.json 自动生成。
% author=Eusebius of Caesarea (c. 265-c. 340); work_id=2501; title=教会史

\chapter{\WorkTitle{教会史}}

% structure: p0001 source=h1 target=omitted reason=page-title-already-rendered-by-parent

\Emph{第一卷}\newline{}\Emph{第二卷}\newline{}\Emph{第三卷}\newline{}\Emph{第四卷}\newline{}\Emph{第五卷}\newline{}\Emph{第六卷}\newline{}\Emph{第七卷}\newline{}\Emph{第八卷}\newline{}\Emph{第八卷（附录）} \WorkTitle{巴勒斯坦的殉道者}\newline{}\Emph{第九卷}\newline{}\Emph{第十卷}\par

\SourceRecord{Translated by Arthur Cushman McGiffert.；Nicene and Post-Nicene Fathers, Second Series；Vol. 1.；Edited by Philip Schaff and Henry Wace.；Buffalo, NY: Christian Literature Publishing Co.,；1890.；<http://www.newadvent.org/fathers/2501.htm>.}

% structure: p0001 source=h1 target=section reason=page-title

\section{\WorkTitle{教会史（卷一）}}

% structure: p0002 source=h2 target=subsection reason=relative-heading-rank; base=h2

\subsection{第一章  写作计划}

1. 我的目的是记述圣宗徒们的传承，以及从我们救主时代至今所经历的时间；并叙述在教会历史中所发生的许多重大事件；还要提及那些在最重要的教区治理和领导教会的人，以及每一代以口头或书面宣扬天主圣言的人。\par

2. 我的目的也是要列出那些因喜爱革新而陷入最大错误的人的名字、数目和时间，他们自称为假知识（\BibleRef{弟前 6:20}）的发现者，如同凶猛的狼一般残忍地摧残基督的羊群。\par

3. 此外，我打算记述整个犹太民族因谋害我们救主而立即遭受的灾祸，记录天主圣言被外邦人攻击的方式和时机，描述那些在不同时期为圣言而流血、受酷刑的人的特征，以及在我们这个时代所作的宣认，最后还有我们的救主赐予他们所有人的仁慈而温柔的援助。既然我计划记述这一切，我将从我们的救主、主耶稣基督的救世计划的开始着手我的工作。\par

4. 但一开始，我必须为我的作品恳求知者的宽容，因为我承认我没有能力写出一部完美而完整的历史，而且由于我是第一个进入这个主题的人，我正在尝试穿越一条孤寂而无人走过的道路。我祈求天主作为我的向导，主的能力作为我的援助，因为我甚至找不到在我之前走过这条路的人留下的足迹，除了一些零碎的片段，其中一些人以一种方式，另一些人以另一种方式，为我们传下了他们所处时代的特定记述。他们像火炬一样从远处发出声音，并从某个高耸而显眼的瞭望塔呼喊，告诫我们应当如何行走，如何稳妥而安全地引导我们工作的进程。\par

5. 因此，我们从他们所提及的、分散在各处的材料中，收集了我们认为对当前工作重要的内容，并从古代作家的作品中像从草地上采摘花朵一样选取适当的段落，我们将努力将它们全部融入一部历史叙述中，如果能够保留我们救主宗徒们的传承记忆，即使不是全部，至少是那些在最为著名、至今仍受尊敬的教会中最杰出的宗徒们，我们就满足了。\par

6. 这项工作在我看来特别重要，因为据我所知，没有一位教会作家曾致力于这个主题；我希望它对于热爱历史研究的人来说将显得非常有用。\par

7. 我已经在我编纂的《年代纪要》中给出了这些内容的概要，尽管如此，我仍在本作品中致力于尽我所能写出更完整的记述。\par

8. 我的作品将如我所说，始于救主基督的救世计划——这比人类的观念更高超、更伟大——并讨论祂的天主性。\par

9. 因为既然我们甚至从基督得名，那么一位打算写作教会史的人，就必须从基督救世计划的真正起源开始，这一计划比许多人想象的更为神圣。\par

% structure: p0012 source=h2 target=subsection reason=relative-heading-rank; base=h2

\subsection{第二章  我们救主和主耶稣基督的先存及其天主性的概要}

1. 既然在基督内有两种本性，一种——就祂被视为天主而言——类似于身体的头，而另一种则可比作脚——就祂为了我们的救恩而取了与我们相同情感的人性而言——那么以下的工作只有在我们从祂全部历史中最主要、最至高的事件开始时才会完整。这样，基督宗教的古老性和天主性就会向那些认为它起源较晚且来自外邦、以为它昨天才出现的人显明。\par

2. 没有任何语言足以表达基督的起源和价值、存有和本性。因此，神圣的圣神在预言中说：\BibleQuote{谁能述说他的世代？}（\BibleRef{依 53:8}）因为除了圣子，没有人认识圣父；除了独自生了祂的圣父，也没有人能充分认识圣子。\par

3. 因为除了圣父，谁能清楚地理解那在世界之前的光、那在万世之前存在的理智而本质的智慧、那太初就与圣父同在且就是天主的活圣言、那先于一切受造之物（可见与不可见）的天主首生而独生者、那理性不朽天军的元帅、伟大谋略的使者、圣父不可言喻旨意的执行者、与圣父共同创造万物的创造者、宇宙中次于圣父的第二因、天主的真实独生子、一切受造之物的主、天主和君王、那从圣父获得了统治权和能力（连同天主性本身、权能和荣耀）的那一位？正如论及祂天主的圣经奥秘段落所说：\BibleQuote{太初有圣言，圣言与天主同在，圣言就是天主。}（\BibleRef{若 1:1}）\BibleQuote{万物是藉着他受造的；凡受造的，没有一样不是藉着他受造的。}（\BibleRef{若 1:3}）\par

4. 伟大的梅瑟也教导了这一点，作为所有先知中最古老的一位，他在神圣圣神的感动下描述了宇宙的创造与安排。他宣称世界的创造者和万物的造物主将低级受造物的制造交给了基督本身，即祂自己那清楚具有天主性且首生的圣言，并与其商议关于人的创造。\Quote{因为，}他说，\BibleQuote{天主说：让我们照我们的肖像，按我们的模样造人。}（\BibleRef{创 1:26}）\par

5. 另一位先知也证实了这事，在他所作的赞美诗中论及天主如下：\Quote{他一发言，万物造成；他一发令，万物生成。}他在此介绍\Person{圣父}和造物主为万有的统治者，以王者的点头下令；第二位是\OriginalTerm{圣言}{divine Word}，正是我们所宣告的那一位，执行\Person{圣父}的命令。\par

6. 自人类受造以来，所有被称许在正义和虔诚上卓越的人——伟大的仆人\Person{梅瑟}，以及在他之前居首位的\Person{亚巴郎}和他的子女，以及此后出现的众多义人和先知——都曾以心灵的纯洁之眼注视祂，并认出祂，向祂献上作为\OriginalTerm{圣子}{Son of God}所应得的敬拜。\par

7. 但祂绝不忽视对\Person{圣父}应有的敬畏，被委派教导他们所有人认识\Person{圣父}。例如，据说\OriginalTerm{上主天主}{Lord God}以一个普通人的形像向\Person{亚巴郎}显现，当时他正坐在\Place{玛默勒橡树}旁。\Person{亚巴郎}立刻俯伏在地，虽然肉眼看见的是一个人，却仍敬拜祂为天主，向祂献祭为主，并承认并非不知道祂的身份，当他说出：\Quote{主，全地的审判者，难道您不施行公义审判吗？}\BibleRef{创 18:25}。\par

8. 因为假若我们认为全能的天主那\OriginalTerm{无生不变的本质}{unbegotten and immutable essence}变成了人的形状，或以某一受造物的外貌迷惑了观看者的眼睛，这是不合理的；另一方面，假若我们认为圣经会虚假地捏造这样的事，也是不合理的。当那审判全地并执行审判的天主和主以人的形状被看见时，若不称之为万物的第一因，那么除了祂那唯一先在的\OriginalTerm{圣言}{Word}之外，还能称谁呢？关于祂，在\BibleBook{}{圣咏}中曾说：\Quote{祂派遣了自己的圣言，医治了他们，拯救他们脱离了灭亡。}\par

9. \Person{梅瑟}最为清楚地宣称祂是次于\Person{圣父}的\OriginalTerm{第二位主}{second Lord}，当他说：\Quote{上主从天上降下硫磺与火，降于\Place{索多玛}和\Place{戈摩辣}，是出自上主。}\BibleRef{创 19:24}神圣的圣经也称祂为天主，当祂又以人形向\Person{雅各伯}显现时，对\Person{雅各伯}说：\Quote{你的名字不再叫\Emph{雅各伯}，要叫\Emph{以色列}，因为你与天主相搏，获得了胜利。}\BibleRef{创 32:28}因此\Person{雅各伯}给那地方起名叫\Quote{天主的神视}，说：\Quote{因为我面对面见了天主，我的性命仍得保全。}\BibleRef{创 32:30}。\par

10. 也不可认为所记载的\OriginalTerm{神视显现}{theophanies}是隶属于天主的众天使和仆役的显现，因为每当这些显现于人时，圣经并不隐瞒事实，而是称呼它们的名字，不称为\Emph{天主}或\Emph{主}，而称为\Emph{天使}，这一点有无数证据可以轻易证明。\par

11. 同样，\Person{梅瑟}的继承人\Person{若苏厄}称祂为天上众天使、总领天使及超世力量的首领，并作为\Person{圣父}的代表，受托拥有第二等的统治权和统御万有的权柄，称祂为\Quote{\OriginalTerm{上主万军的统帅}{captain of the host of the Lord}}，虽然他看见祂时，仍是以人的形像和面貌出现。因为经上记载：\par

12. \Quote{当\Person{若苏厄}在\Place{耶里哥}时，他举目看见一个人，手里持着拔出的刀，站在他对面。\Person{若苏厄}走到他那里，问他说：\InnerQuote{你是帮助我们，还是帮助我们敌人？}他回答说：\InnerQuote{我如今是以上主万军的统帅而来。}\Person{若苏厄}就俯伏在地，对他说：\InnerQuote{我主，你对你的仆人有什么吩咐？}上主万军的统帅对\Person{若苏厄}说：\InnerQuote{把你脚上的鞋脱下来，因为你所站的地方是圣的。}}\par

13. 你也会从同样的话中看出，这正是那位与\Person{梅瑟}说话的天主。因为圣经用同样的话并指着同一位说：\Quote{上主看见他前来观看，上主便从荆棘中呼唤他说：\InnerQuote{梅瑟，梅瑟。}他回答：\InnerQuote{什么事？}上主说：\InnerQuote{不要走近；把你脚上的鞋脱下来，因为你所站的地方是圣地。}上主又对他说：\InnerQuote{我是你父亲的天主，\Person{亚巴郎}的天主，\Person{依撒格}的天主，\Person{雅各伯}的天主。}}\par

14. 并且，有一个在万世之前就已存在并自立的实体，它服务于万有的\Person{圣父}和天主，为了形成一切受造之物，它被称为\OriginalTerm{天主圣言}{Word of God}和\OriginalTerm{智慧}{Wisdom}。我们除了已引用的证据外，还可以从智慧本身的口中得知，她通过\Person{撒罗满}最清楚地启示了关于她自己的以下奥秘：\Quote{我，智慧，与明达和知识同住，我并呼唤了聪敏。藉着我，君王统治，王侯制定正义。藉着我，伟人被尊大，藉着我，君主治理大地。}\par

15. 她还补充说：\Quote{上主在他道路的起始便创造了我，为他的工程；在万世以前，在起初，在他造地以前，在他造深渊以前，在群山奠定之前，在一切丘陵之前，他生了我。当他预备诸天时，我与他同在；当他立定天下源泉时，我与他同在，安排一切。我是他所喜悦的；每日在他面前欢乐，时时因他完成了世界而欢乐。}\par

16. 因此，关于\OriginalTerm{圣言}{divine Word}的在先存在及其向某些人（即使不是向所有人）显现，我们已简要地说明了。\par

17. 但是，为什么\OriginalTerm{福音}{Gospel}在古代不像现在这样向所有人和所有民族宣讲，从以下考虑中可以看出。古人的生活并不容许他们接受\Person{基督}那全智全善的教导。\par

18. 因为从一开始，在最初享福的生活之后，第一个人就轻视了天主的诫命，堕入了这个有死有坏的状态，并把他原先由天主启示的享乐换成了这个受诅咒的大地。他的后代充满大地后，除了少数例外，变得更加败坏，过着一种野蛮而无法忍受的生活。\par

19. 他们既不想城市或国家，也不思艺术或科学。他们甚至连法律、正义、德行、哲学的名词也不知晓。他们如同游牧民族，在旷野中度日，像野兽一样野蛮凶猛，以过度的故意邪恶摧毁了人天生的理性，以及种植在人类灵魂中的思想和文化的种子。他们完全沉溺于各种亵渎行为，时而互相引诱，时而互相残杀，时而吃人肉，时而胆敢与神明作战，从事那些被众人传颂的巨人之战；时而计划筑垒以对抗上天，在无法无天的骄傲疯狂中准备攻击万有的天主本身。\par

20. 正因为这些事，当他们如此行事时，那无所不见的天主对他们降下了洪水与大火，如同对付蔓延全地的野生森林。祂以不断的饥荒、瘟疫、战争和天上的雷电将他们剪除，仿佛用更严厉的惩罚来遏制一些可怕而顽固的灵魂疾病。\par

21. 随后，当邪恶的泛滥几乎淹没了全人类，如同深沉的醉态，遮蔽并昏暗了人的心智时，天主首生首造的智慧，那先在的圣言自己，因祂对人超乎寻常的爱，向祂的仆人们显现，有时以天使的形象，又有时向那些蒙天主恩宠的古人显现，以祂自己作为天主的拯救大能，但非以别的方式，而只以人的形状显现，因为不可能以其他方式显现。\par

22. 并且，通过他们，虔诚的种子被播撒在许多人中间，那从希伯来人传下来的整个民族，坚持不懈地敬拜天主，祂借着先知梅瑟赐给他们——如同赐给仍被古老习俗败坏的众人——某种奥秘的安息日和割礼的预像和象征，以及其他属灵原则的初步教导，但没有赐予他们奥秘本身的完全知识。\par

23. 但当他们的律法受到尊崇，如同馨香之气弥漫在万民中时，由于他们的影响，大多数的外邦人因各地兴起的立法者和哲学家的教导而心性软化，他们野蛮粗暴的本性转变为温和，从而享受了深厚的和平、友谊与社交往来。最后，在罗马帝国开始之时，那同一位德行的导师、圣父在一切善事中的仆役、神圣而属天的天主圣言，以一个与我们的本性在实质上毫无差别的人身，再次向世上所有民族和个人显现——他们先前已得到帮助，如今已准备好接受对圣父的认识。祂做了并承受了那些已经预言的事。因为早已预言，有一位同时是人和天主者要来住在世上，行奇妙的事，并显示自己为万民教授圣父虔诚的导师。祂诞生的奇妙性质、祂的新教导、祂的奇妙作为，以及祂的死亡方式、祂从死者中的复活，以及最后祂神圣的升天，也都被预言了。\par

24. 例如，先知达尼尔在圣神感动下，看到他末世的国度，于是受默感以适合人理解的言语描述这神圣的景象：\Quote{我观看，}他说，\Quote{直到宝座设立，那位万古常存者坐上去，他的衣服洁白如雪，他的头发如纯净的羊毛；他的宝座是火焰，其轮子是燃烧的火。一条火河在他面前流出。有千千万万服侍他，有亿亿万万侍立在他面前。他设立了审判，书卷就展开了。} \BibleRef{达 7:9}\par

25. 又说：\Quote{我观看，}他说，\Quote{见有一位像人子者，乘着天上的云彩而来，他走近万古常存者，并被引到他面前；他便被赐予统治权、荣耀和国度；各民族、部落、语言都事奉他。他的统治是永远的统治，永不消逝；他的国度永不败坏。} \BibleRef{达 7:13}\par

26. 显然，这些话只能指我们的救主，那在起初就与天主同在的天主圣言，祂因最后在肉身中的显现而被称为人子。\par

27. 但既然我们已在单独的书卷中收集了有关我们救主耶稣基督的先知选段，并以更合乎逻辑的形式编排了那些关于祂所启示的事，那么以上所述对于当前已足够。\par

% structure: p0040 source=h2 target=subsection reason=relative-heading-rank; base=h2

\subsection{第三章  耶稣的名字和基督的名字从起初就已知晓，并受到受默感的先知们的尊崇。}

1. 现在是适当的地方来表明，耶稣这个名字以及基督这个名字，曾被古代蒙天主宠爱的先知们尊崇。\par

2. 梅瑟是第一个宣扬基督之名作为特别庄严荣耀之名的人。当他传授天上事物的预像、象征和奥秘图象时，根据那对他说的话：\Quote{你要留心，照你在山上所看到的样式制造所有物件。} \BibleRef{出 25:40} 他祝圣了一个人作为天主的最高司祭——就他所可能做的而言，并称他为\Emph{基督}。如此，他将最高司祭的尊位——在他眼中超越人间最尊贵的地位——附上了基督之名，以表示尊崇与荣耀。\par

3. 他深知在\Emph{基督}内具有某种神性。同一位在圣神的感动下预见\Quote{耶稣}这名，也以某种殊荣尊崇它。因为耶稣这名在梅瑟时代之前从未在人中被宣告，他首先且唯独将其应用于他所知道的一位，这位将在他死后作为预象和象征，再次接受最高权柄。\par

4. 因此，他的继承人——之前并未拥有耶稣之名，而是被称作\Person{阿乌塞斯}（这是他父母所取的另一名字）——如今他称其为耶稣，以此赐予他远超任何王冠的荣誉礼物。因为耶稣本人，\Person{纳维}的儿子，与我们的救主相似，因为唯独他，在梅瑟之后、在由他所传递的象征性敬拜完成之后，继承了真正纯洁宗教的统治权。\par

5. 这样，梅瑟便将我们救主耶稣基督的名，作为最高荣誉的标志，赐予了他时代中在德行与荣耀上超越其余百姓的两个人：即大司祭和他在统治权上的继承人。\par

6. 后来的先知们也清楚指名预言了基督，同时预言了犹太民族将对他设下的阴谋，以及通过他对外邦的召叫。例如，耶肋米亚这样说：\BibleQuote{在我们面前的神，基督上主，在他们的毁灭中被捉住了；我们曾说，在他的荫下我们要在列国中存活。}\BibleRef{哀 4:20}。达味在困惑中则说：\BibleQuote{为什么列国愤怒，万民妄想虚妄的事？地上的君王列阵，首领聚集，反对上主和他的基督。}其后，他以基督本人的口吻补充说：\BibleQuote{上主对我说：你是我的儿子，我今日生了你。你向我请求，我必将列国赐给你作产业，地极赐给你作基业。}\par

7. 在希伯来人中，不但那些受尊为大司祭、并为象征而用特制的油所傅油的人们，被冠以基督之名；而且先知们在圣神的感动下所傅油的君王，也如此被立为预象性的基督。因为他们也在自身内承载了那真而唯一的基督——即统御万有的天主圣言——的王权与主权的预象。\par

8. 我们还被告知，某些先知本人也因傅油而成为预象中的基督，因此这一切均指向真基督——受天主启示的、属天的圣言，祂是万有的唯一大司祭，一切受造物的唯一君王，圣父的唯一至高的先知之先知。\par

9. 此事的明证是：在古时象征性地受傅的众人中，无论是司祭、君王或先知，没有一位拥有像我们的救主和主耶稣——真而唯一的基督——所显示的那样伟大的启示性德能。\par

10. 至少他们中没有任何一位——无论他们在自己的民族中历世代享有多么崇高的尊荣——曾从他们自己的预象性名称\Emph{基督}那里给予其追随者\Emph{基督徒}之名。他们的臣民也未曾向他们任何人呈献神性的尊荣；在他们死后，其追随者的心态也不至于愿意为他们所尊敬的人而死。在关于那个时代的任何一个人身上，也从未在天下万民中引起如此大的骚动；因为单纯的象征在他们中间不能像我们救主所显示的真理本身那样发挥作用。\par

11. 祂虽然未从任何人领受大司祭的象征和预象，虽然并非出生于司祭的族裔，虽然未被军队护卫拥立为王，虽然不像古时的先知，虽然未在犹太人中获得任何尊荣或高位，然而圣父仍以一切——即便不是以象征，而是以真理本身——装饰了祂。\par

12. 因此，虽然祂未拥有我们所提及的那些人类似的尊荣，但祂却比他们全体更被称为\Emph{基督}。祂作为天主的真而唯一的基督，以真正威严神圣的\Emph{基督徒}之名充满了全地，不再将预象和影像交托给祂的门徒，而是交托赤裸裸的德行本身，以及在真理教义中的属天生命。\par

13. 祂不是用物质所制备的油受傅，而是如神明所宜，以圣神亲自受傅，藉分享圣父的非受生神性。这一点依撒意亚也再次教导，他仿佛以基督本人的口吻呼喊说：\BibleQuote{上主的神临于我身，因此祂给我傅了油。祂派遣我去向贫穷人传报福音，向俘虏宣告释放，向盲者宣告复明。}\par

14. 不仅依撒意亚，达味也向祂说话：\BibleQuote{天主，你的宝座永远常存；你王国的权杖是正直的权杖。你喜爱正义，憎恨邪恶；因此天主，你的天主，用喜乐之油傅你，胜过你的同伴。}此处圣经第一节称祂为天主，第二节以王室权杖尊崇祂。\par

15. 稍后，在论及神性与王权之后，圣经第三处将祂描绘为成为基督，受傅不是用物质制成的油，而是用神圣的喜乐之油。这便指明祂的特别尊荣，远超越并不同于那些古时作为预象、以更物质方式受傅的人们。\par

16. 同一位作者在别处这样论及祂：\BibleQuote{上主对我主说：你坐在我右边，等我使你的仇敌做你的脚凳。}以及\BibleQuote{从母腹中，在晨星之前，我已生了你。上主发了誓，绝不变更：你是按照默基瑟德的品位永做司祭。}\par

17. 但是，这位\Person{默基瑟德}在圣经中被介绍为至高天主的司祭，并非由任何特制的傅油所祝圣，甚至也不属于犹太人司祭职的后裔。因此，我们的救主被宣告为基督和司祭，乃是照他的品秩，而非照其他领受预象与类型之人的品秩，且是藉着誓言的呼召。\par

18. 因此，历史并未记载他曾被犹太人傅以物质之油，也不曾记载他属于司祭的谱系，而是记载他来自天主自身，在晨星之前，即在世界被造之前就存在，并且他获得了不朽不坏的司祭职，直到永远。\par

19. 然而，以下事实是关于他无形体且神圣傅油的伟大而令人信服的证据：在所有曾存在的人中，唯独他直到今日仍被全世界所有人称为基督，并以此名被宣认和见证，且受到希腊人和外邦人的纪念，甚至直到今日，他在世界各地的追随者尊他为君王，钦佩他超过先知，并荣耀他为天主真实而唯一的大司祭。除此之外，作为先存的天主圣言，他在万世之前被召叫存在，他从圣父那里接受了尊威的荣耀，并被当作天主来崇拜。\par

20. 但最奇妙的是，我们这些奉献自己给他的人，不仅以我们的声音和言语之音来尊崇他，更以灵魂的完全提升来尊崇他，以致我们宁可为向他作证，也不愿保全自己的性命。\par

21. 我不得不以这些事作为我历史的序言，以免有人因他降生成人的时期而认为我们的救主和主耶稣基督只是近日才存在。\par

% structure: p0062 source=h2 target=subsection reason=relative-heading-rank; base=h2

\subsection{第四章 他所向万国宣告的宗教，既非新，亦非怪。}

1. 但是，为免有人以为他的教义是新奇而奇怪的，仿佛是由一个近日出现、与其他人毫无分别之人所构造的，现在就让我们也简要地考虑这一点。\par

2. 众所公认，当近来我们的救主耶稣基督的出现为众人所知时，立刻出现了一个新民族；这民族诚然不小，也没有住在地的某一角落，而是万民中最众多、最虔诚的，不可毁灭、不可征服，因为它常从天主那里获得援助。这民族就这样在天主奥秘计划所指定的时间突然出现，就是被所有人以基督之名所尊荣的那个民族。\par

3. 有一位先知，当他以圣神的眼睛预见到将要发生的事时，如此惊奇，竟喊道：\Quote{谁曾听过这样的事？谁曾说过这样的话？大地岂能在一天之内生出，一个民族岂能一旦诞生？}\BibleRef{依 66:8} 同样那位先知也暗示这民族将被称为的名字，当他说：\Quote{那些事奉我的，将被称为一个新名，这名字在地上必受祝福。}\BibleRef{依 65:15}\par

4. 然而，虽然我们显然是新的，并且基督徒这个新名确实只是近来才为万国所知，但我们的生命和行为，以及我们的宗教教义，却并非我们近来所发明，而是可以说，从人的最初受造以来，就已由古时蒙天主恩宠之人的自然理解所确立。我们将以如下方式表明这一点。\par

5. 希伯来民族并不新，反而因其古老而普遍受尊敬，这是众所周知的。这民族的书籍和著作记载了古人的事迹，这些古人确实罕见且为数不多，却因虔诚、正义及一切美德而卓著。其中，有些卓越的人活在洪水之前，另一些则活在洪水之后，是\Person{诺厄}的子孙和后裔，其中包括\Person{亚巴郎}，希伯来人尊他为自己的始祖和祖先。\par

6. 如果有人断言，所有那些从\Person{亚巴郎}本人直到第一个人，曾享有正义见证的人，事实上都是基督徒，即使名义上不是，他也不会违背真理。\par

7. 因为那名所指示的——即基督徒因认识基督和基督的教导，以节制和正义、生活中的忍耐和刚毅，以及对那独一至高天主的虔诚宣认为特色——这一切，他们都不亚于我们而热切地实践了。\par

8. 他们不关心身体的割损礼，我们也不关心。他们不关心遵守安息日，我们也不。他们不回避某些种类的食物，也不在意\Person{梅瑟}最初传给后代作为预象遵守的其他分别；现今的基督徒也不做这些事。但是，他们清楚地认识了天主本身的基督；因为已经表明，他显现给\Person{亚巴郎}，他给予\Person{依撒格}启示，他与\Person{雅各伯}交谈，他与\Person{梅瑟}以及后来的先知们对话。\par

9. 因此，你会发现那些蒙天主恩宠的人以基督之名为尊荣，正如论及他们的那段经文所说：\Quote{不可触摸我的受傅者，也不可伤害我的先知。}\par

10. 因此，显然必须认为，那近期通过基督的教导而向万国宣讲的宗教，是所有宗教中最原始、最古老的，且是由那些蒙天主恩宠的人在\Person{亚巴郎}时代所发现的。\par

11. 如果有人说，\Person{亚巴郎}在很久以后才被赐予割损礼的命令，我们回答说，然而在此之前，已被宣告他因信德而获得了正义的见证；正如天主圣言所说：\Quote{亚巴郎信了天主，天主就以此算为他的正义。}\BibleRef{创 15:6}\par

12. 事实上，对于亚巴郎——他在未受割礼前已是一个成义的人——天主向他启示了自己（而这启示者就是基督自己，天主的圣言），赐给了他一个预言，关乎将来要像他一样成为义的人。这预言如下：\Quote{地上万民都要因你获得祝福。}\BibleRef{创 12:3} 又说：\Quote{他要成为一个强大而众多的人民；地上万民都要因他获得祝福。}\BibleRef{创 18:18}\par

13. 我们可以认为这话已在我们身上应验。因为他弃绝了祖辈的迷信和以往生活的错误，承认了唯一的天主，并以德行的行为，而不是以后来梅瑟所赐的法律的侍奉来敬拜他，因着信德，因着向他显现的基督——天主的圣言——而成了义。所以，对这样一个人，曾说过：地上万民和万邦都要因他获得祝福。\par

14. 然而，亚巴郎的那种敬神之道，在现今的时代，由全世界唯独的基督徒以行为——比言语更有效——重新实践了出来。\par

15. 那么，有什么能阻止我们承认：我们这些属于基督的人，与那些古时蒙天主垂爱的人，实践着同一生活方式、拥有同一宗教呢？由此可见，基督的教导所托付给我们的圆满宗教，并非新奇怪异；相反，说句实话，它是原始而真实的宗教。关于这个题目，说到这里就够了。\par

% structure: p0078 source=h2 target=subsection reason=relative-heading-rank; base=h2

\subsection{第五章 祂在人间显现的时期}

1. 现在，经过对本教会史的必要引言之后，我们可以说开始我们的旅程了，从我们的救主在肉身内显现开始。我们恳求天主，圣言的父，以及我们一直所谈论的那一位，即我们的救主和主耶稣基督自己，天上的天主圣言，作为我们的援助和同工，以叙述真理。\par

2. 是在奥古斯都在位第四十二年，埃及被征服后第二十八年，安东尼与克娄巴特拉去世（埃及的托勒密王朝随之终结）的那一年，我们的救主和主耶稣基督，按关于祂所发的预言，在犹太的白冷出生。\BibleRef{米 5:2} 祂的诞生是在第一次户口登记期间，其时基黎诺担任叙利亚总督。\par

3. 犹太历史家中最著名的弗拉维乌斯·约瑟夫也提到这次在基黎诺任内进行的户口登记。在同一记载中，他描述了当时发生的加里肋亚人起义，我们作者中的路加也在\BibleBook{宗徒}{大事录}中提到此事，说：\Quote{此后，在户口登记的时候，加里肋亚人犹达斯起来，煽动了一群百姓跟随他；他也灭亡了，所有服从他的人也都四散了。}\BibleRef{宗 5:37}\par

4. 上述的作者在他的\WorkTitle{犹太古史}第十八卷中，与这些话一致，补充了以下内容，我们照录如下：\Quote{基黎诺，一位元老院成员，曾担任过其他官职并一直升到执政官，在其他方面也极有威望，带着少数随从来到叙利亚，奉凯撒之命担任该民族的审判官并评估他们的财产。}\par

5. 稍后他写道：\Quote{但是犹达斯，一个加鲁尼人，来自一个叫加玛拉的城市，带着法利塞人撒杜克，鼓动百姓造反，两人都说纳税无异于彻底的奴役，并劝勉那民族保卫他们的自由。}\par

6. 在他的\WorkTitle{犹太战史}第二卷中，关于同一个人写道：\Quote{此时，一个名叫犹达斯的加里肋亚人，说服他的同胞造反，声称假若他们服从向罗马人纳税，并忍受除天主以外有死有坏的主人，他们就是懦夫。}这些事是约瑟夫记载的。\par

% structure: p0085 source=h2 target=subsection reason=relative-heading-rank; base=h2

\subsection{第六章 关于基督的时代，根据预言，自古以来按常规继承统治犹太民族的统治者终止了，黑落德，第一个外邦人，成了国王。}

1. 当黑落德，第一个外邦血统的统治者，成为国王时，梅瑟的预言应验了，按该预言应\Quote{犹大的权杖不离，他的脚边也不缺立法者，直到那应得者来到。}后者，他也表明，就是万民的期待。\par

2. 这个预言一直没有应验，只要他们被允许在自己的民族统治者下生活，即从梅瑟时代到奥古斯都统治时期。在后者统治下，黑落德，第一个外邦人，由罗马人赐予了犹太王国。据约瑟夫记载，他父亲是以东人，母亲是阿拉伯人。但阿非利加努斯，一位非同寻常的作家，说那些更确切了解他情况的人报告说，他是安提帕特的儿子，而安提帕特是阿斯卡隆的一个名叫黑落德的人的儿子，那人曾是阿波罗神庙的所谓仆役之一。\par

3. 这个安提帕特，童年时被以东强盗掳去，与他们生活在一起，因为他父亲穷困，无法为他付赎金。他在他们的习俗中长大，后来与犹太人的大祭司依尔卡诺结交。他的一个儿子就是那位生活在我们的救主时代的黑落德。\par

4. 当犹太王国落到这样一个人手中时，万民的期待，按预言，已在门前了。因为与他一起，从梅瑟以来按照常规统治他们的首领和总督终止了。\par

5. 在他们被掳并迁移到巴比伦之前，他们先后由撒乌耳、达味统治，在君王之前，有称为民长的领导者治理他们，这些人继梅瑟及其继承人若苏厄之后。\par

6. 他们从\Place{巴比伦}归来后，继续不间断地保持着一种贵族政体，带有寡头政治的色彩。因为司铎们一直掌管事务，直到罗马将军\Person{庞培}以武力攻取\Place{耶路撒冷}，并进入圣殿最内层的至圣所，玷污了圣地。\Person{亚里斯托布鲁斯}，他凭借古代世袭的权利，此前一直是国王兼大司祭，\Person{庞培}将他及其子女用锁链解送至\Place{罗马}；并将大司祭的职位赐给\Person{亚里斯托布鲁斯}的兄弟\Person{希尔卡诺斯}，而整个犹太民族从那时起便向罗马人纳贡。\par

7. 但\Person{希尔卡诺斯}，他是正规大司祭世系中的最后一位，不久之后被帕提亚人俘虏，而\Person{黑落德}，正如我已说过的，是第一个外邦人，由罗马元老院和\Person{奥古斯都}立为犹太民族的国王。\par

8. 在他治下，基督以肉体形态显现，万民所期盼的救恩及其蒙召也照预言应验。从那时起，犹大的首领和统治者——我指犹太民族——终结了，随之而来的是大司祭的职位，它自古以来一直以最紧密的世系代代相传，也立即陷入混乱。\par

9. 关于这些事，\Person{约瑟夫}也是一个见证，他表明当\Person{黑落德}被罗马人立为国王后，他不再从古老世系中任命大司祭，而是将这一尊荣赐予某些无名之人。\Person{黑落德}在任命司祭一事上的做法，被他的儿子\Person{阿尔赫劳}所仿效，之后罗马人接管了政权，也照此行事。\par

10. 同一位作者表明，\Person{黑落德}是第一个将大司祭的圣袍锁在自己的印鉴下，并拒绝让大司祭们为自己保管的人。他之后的\Person{阿尔赫劳}也遵循同样的做法，\Person{阿尔赫劳}之后则是罗马人。\par

11. 我们记录这些事，是为了表明另一个预言在我们救主\Person{耶稣基督}的显现中应验了。因为圣经在\WorkTitle{达尼尔书}中，\BibleRef{达 9:26}明确提及直到\Person{基督}来临的若干星期，我们在其他著作中已讨论过此事，它最清楚地预言，在这些星期完成后，犹太人中的傅油礼将完全消失。而这一点，在救主\Person{耶稣基督}诞生之时已清楚显明应验了。我们在此必须预先陈述，以证明时间的正确性。\par

% structure: p0097 source=h2 target=subsection reason=relative-heading-rank; base=h2

\subsection{第七章　关于\Person{基督}族谱的福音记载中所谓的差异}

1. \Person{玛窦}和\Person{路加}在他们的福音书中记载了\Person{基督}不同的族谱，许多人认为他们彼此矛盾。为此，每一位信徒因不明真相，都竭力想提出某种解释来调和这两段经文。请允许我们附上关于此事流传下来的记述，它出自我们刚才提到的\Person{阿弗里卡纳斯}，他在致\Person{亚里斯提德}的信函中讨论了福音族谱的和谐问题。在驳斥了他人牵强且具欺骗性的观点之后，他根据所受的传统如此陈述：\par

2. 因为，在\Place{以色列}中，世代的谱系或按血统、或按法律计算——按血统是指合法后裔的继承，按法律则是指当某人为无子而逝的兄弟立嗣时；由于复活的明确盼望尚未赐下，他们通过一种可朽的复活来象征未来的应许，以使逝者的名字得以延续——\par

3. 那么，当这份族谱表中有些人按血统继承——子承父业，而另一些人虽由某位父亲所生，却被记为另一人的名字时，就同时提及了实际上的祖先和名义上的祖先。\par

4. 因此，两部福音书都无错误，一部按血统计算，另一部按法律计算。因为从\Person{撒罗满}和从\Person{纳堂}这两个世系，由于为无嗣者立嗣和再婚而相互交织，以至同一个人有时合理地被认为属于这个世系，有时属于另一个世系；即有时属于名义上的父亲，有时属于实际上的父亲。所以，这两种记载都严格真实，并且虽然相当复杂，却仍极其准确地传至\Person{若瑟}。\par

5. 但为了阐明我所说的话，我要解释一下这些世代的交替。如果我们从\Person{达味}经由\Person{撒罗满}来推算世代，倒数第三代是\Person{玛堂}，他生了\Person{雅各伯}，即\Person{若瑟}的父亲。但如果按照\Person{路加}，我们从\Person{达味}的儿子\Person{纳堂}来推算，同样倒数第三代是\Person{默尔希}，他的儿子\Person{厄里}是\Person{若瑟}的父亲。因为\Person{若瑟}是\Person{厄里}的儿子，\Person{厄里}是\Person{默尔希}的儿子。\par

6. 因此，既然我们面对的是\Person{若瑟}，就必须说明何以两人都被记载为他的父亲——\Person{雅各伯}（出自\Person{撒罗满}世系）和\Person{厄里}（出自\Person{纳堂}世系）；首先，\Person{雅各伯}和\Person{厄里}二人如何是兄弟，其次，他们的父亲\Person{玛堂}和\Person{默尔希}虽然家族不同，却如何被宣告为\Person{若瑟}的祖父。\par

7. \Person{玛堂}和\Person{默尔希}相继娶了同一位女子，生下了同母异父的兄弟，因为法律并不禁止寡妇（无论是因休妻还是因丈夫死亡）再嫁。\par

8. 于是，通过\Person{厄斯塔}（因为据传统，这是那女子的名字），出自\Person{撒罗满}的\Person{玛堂}先生了\Person{雅各伯}。\Person{玛堂}死后，\Person{默尔希}（其世系追溯至\Person{纳堂}，同属一个支派，但家族不同）如前所述娶了她，生了儿子\Person{厄里}。\par

9. 如此我们会发现雅各伯和厄里虽属于不同的家族，却是同母兄弟。其中一位雅各伯，在其兄弟厄里无子而死后，娶了后者的妻子，并为她生了一个儿子若瑟，按自然和理性而言，这是他自己的儿子。因此经上也记载说：\BibleQuote{雅各伯生若瑟。}\BibleRef{玛 1:6} 但按照法律，他是厄里的儿子，因为雅各伯作为厄里的兄弟，为厄里立嗣。\par

10. 因此，通过他所追溯的族谱不会无效，玛窦在他的数算中这样列出：\BibleQuote{雅各伯生若瑟。}而路加则说：\BibleQuote{耶稣}按照人的看法（他也加了这一句），\BibleQuote{是若瑟的儿子，厄里的儿子，默肋希的儿子}；因为他无法更清楚地表达按法律的生嗣。而\Quote{生}这个词他在族谱表中一直省略，将血缘追溯至亚当，天主子。这个解释既不是无法证明，也不是无谓的猜测。\par

11. 因为根据肉体而言，吾主的亲属，或出于夸耀之心，或仅想陈述事实，两者都真实地传下了以下记述：有些厄东强盗袭击了巴勒斯坦的阿市刻隆城，从靠近城墙的一座阿波罗神庙中，除其他战利品外，还掳走了一个名叫安提帕特的，他是某位名为黑落德的神庙奴隶的儿子。由于司祭无法为儿子支付赎金，安提帕特便在厄东人的习俗中成长，后来受到犹太人大司祭希尔卡诺斯的庇护。\par

12. 他被希尔卡诺斯派往庞培处执行一项使命，并将被他兄弟阿黎斯托步罗侵占的王国交还给庞培，因此有幸被任命为巴勒斯坦的总督。但安提帕特被嫉妒他巨大好运的人杀害后，他的儿子黑落德接替了他，后来在安东尼和奥古斯都的统治下，通过元老院的法令被立为犹太人的王。他的儿子是黑落德和其他分封侯。这些记述也与希腊人的记载相符。\par

13. 但由于直到那时在档案中一直保存着希伯来人以及那些将血统追溯到皈依者的人的族谱，例如阿孟人阿基约尔和摩阿布人卢德，以及与以色列人混居并一同从埃及出来的人，黑落德因以色列人的血统对他毫无益处，且因意识到自己出身卑微而深感痛苦，便焚毁了所有族谱记录，以为若无人能从官方登记簿中将自己的血统追溯到族长或皈依者及与他们混居的、被称为\OriginalTerm{革奥辣}{Georae}的人，他便能显得出身高贵。\par

14. 然而，少数细心的人通过记住名字或以其他方式从登记簿中获取了私人的记录，自豪地保留着他们高贵血统的记忆。其中包括那些已经提到的、被称为\OriginalTerm{德波斯尼}{Desposyni}的人，因为他们与救主的家族有关。他们从犹大地的纳匝辣和科哈巴村庄来到世界的其他地方，尽可能忠实地从记忆和每日记录簿中绘制了上述族谱。\par

15. 不论事实是否如此，按照我个人及所有公平之人的意见，没有人能找到更清晰的解释。但愿这足以满足我们，因为我们虽无法提出任何证据支持它，但我们没有更好或更真实的说法。无论如何，\WorkTitle{福音}陈述了真理。在同一书信的末尾，他附加了这些话：\Quote{玛堂是撒罗满的后代，生了雅各伯。玛堂死后，默肋希是纳堂的后代，由同一女人生了厄里。厄里和雅各伯因此是同母异父的兄弟。厄里无子而死，雅各伯为他立嗣，生了若瑟，按自然是他自己的儿子，但按法律却是厄里的儿子。因此若瑟是两人的儿子。}\par

17. 以上是阿非利加诺斯所述。若瑟的族谱这样追溯后，玛利亚实际也显示出与他属于同支派，因为按照梅瑟的法律，不同支派之间的通婚是不允许的。法律命令人要娶同家族和同血统的人，以免产业从一支派转移到另一支派。到此为止。\par

% structure: p0114 source=h2 target=subsection reason=relative-heading-rank; base=h2

\subsection{第八章 黑落德对婴儿的残暴以及他死亡的方式}

1. 当基督按照预言在犹大地的白冷降生时，东方的贤士前来询问那位生来就是犹太人的王在哪里——他们看见了他的星，这便是他们长途跋涉的原因；因为他们热切渴望朝拜这位婴儿为天主——黑落德因此深感不安，以为自己的王国可能受到威胁；于是他询问属于犹太民族的法学士，他们期望基督诞生的地方。当他得知米该亚的预言\BibleRef{米 5:2}宣告白冷将是他的诞生地时，他通过一道谕令，下令将白冷及其四境所有两岁及以下的男婴，按照他从贤士那里准确查问的时间全部杀死，以为耶稣很可能与同龄的其他孩子遭遇同样的命运。\par

2. 但孩子提前逃脱了陷阱，被他的父母带到了埃及，他们从显现的天使那里得知了将要发生的事。这些事由\WorkTitle{福音}中的\WorkTitle{圣经}所记载。\EditorialNote{玛 2}\par

3. 此外，值得留意的是黑落德因他针对基督及其同龄婴儿的胆大罪行所得到的报应。因为立刻，毫无延迟地，神圣的惩罚在他活着时降临到他身上，并使他预先尝到了死后将要领受的滋味。\par

4. 在此不可能详述他是如何因接踵而来的家庭不幸——妻子、儿女以及其他至亲好友的被杀——而玷污了他那所谓幸福的统治。这一叙述使所有其他悲剧都黯然失色，在若瑟夫的著作中已有详细记载。\par

5. 至于他在犯罪反对我们的救主及其他婴孩之后，天主立刻降罚，催迫他走向死亡，我们最好从那位史学家的记述中得知；他在其\\WorkTitle{犹太古史}第十七卷中关于他的结局写道：\par

6. 但黑落德的病势加重，天主为惩罚他的罪行而降下灾祸。因为他体内燃起一种缓慢的火焰，触摸者不易察觉，却加剧了他内在的痛苦；他对食物有一种可怕的欲望，无法抗拒。他还患了肠道溃疡，结肠剧痛，双脚积聚了水样透明的体液。\par

7. 他的腹部也遭受了类似的痛苦。不仅如此，他的私处腐烂生蛆。他还感到呼吸极度困难，且因气味恶臭和呼吸急促而格外令人不适。\par

8. 他四肢也抽搐，产生无法控制的力气。据说，那些拥有占卜和解释这类事件智慧的人说，天主因君王的大不敬而降此惩罚。\par

9. 上述作者在所提及的著作中记述了这些事。他在其\\WorkTitle{历史}第二卷中，对同一位黑落德也有类似的描述，如下：疾病随即侵袭了他的全身，以各种折磨使之痛苦不堪。因为他有低烧，全身皮肤瘙痒难忍。他还持续患有结肠疼痛，双脚肿胀如同患水肿之人，腹部发炎，私处腐烂生蛆。此外，他只能直立呼吸，且呼吸困难，四肢抽搐，因此占卜者说他的疾病是一种惩罚。\par

10. 然而，尽管与这样的痛苦搏斗，他仍然贪恋生命，希望痊愈，并设计治疗方法。例如，他渡过约但河，在\\OriginalTerm{卡利罗}{Callirhoë}使用温泉浴，这些温泉流入\\OriginalTerm{沥青湖}{Lake Asphaltites}，但本身甘甜可饮。\par

11. 他的医生认为，他们可以通过加热的油使他的全身重新温暖。但当他们把他放进一个装满油的浴盆中时，他双眼虚弱，像死人一样上翻。但他的侍从发出喊声，他在噪音中苏醒过来；但最终，他对治愈感到绝望，下令将大约五十\\OriginalTerm{德拉克马}{drachms}分给士兵，并将巨额款项赐予他的将领和朋友。\par

12. 然后，他返回耶里哥，在那里被忧郁所袭，计划犯下一项邪恶的罪行，仿佛要向死亡本身挑战。因为，他从全犹太各城召集了所有最杰出的人，命令将他们关进所谓的\\OriginalTerm{赛马场}{hippodrome}。\par

13. 他召来了他的姊妹撒罗默和她的丈夫亚历山大，说：\\InnerQuote{我知道犹太人会为我的死而欢喜。但我可能被其他人哀悼，并举行隆重的葬礼，如果你们愿意执行我的命令。当我咽气时，立刻用士兵围住这些现在被看守的人，将他们杀死，以便全犹太和每家每户，甚至违背他们的意愿，都为我哭泣。}\par

14. 过了一会儿，若瑟夫说：\\Quote{他再次因饥饿和痉挛性咳嗽而备受折磨，以致被痛苦压倒，计划抢先命运一步。他拿起一个苹果，也要了一把刀，因为他习惯切苹果吃。然后环顾四周，看是否无人阻止，他举起右手，仿佛要刺向自己。}\par

15. 除这些事之外，同一位作者记载，他在死前又杀了自己的一个儿子，这是他下令杀的第三个儿子，随后不久他就咽了气，死时痛苦异常。\par

16. 黑落德的下场就是如此，他因屠杀白冷婴孩而遭受了应得的惩罚，那是他谋害我们救主的结果。\par

17. 此后，一位天使在埃及显现于若瑟梦中，命令他带着孩子和祂的母亲前往犹太地，向他启示那些谋害孩子性命的人已经死了。对此，福音记载道：\\BibleQuote{但因若瑟听说阿尔赫劳继承他父亲黑落德作了王，就不敢到那里去；且在梦中得了指示，便往加里肋亚境内去了。}\\BibleRef{玛 2:22}\par

% structure: p0132 source=h2 target=subsection reason=relative-heading-rank; base=h2

\subsection{第九章 比拉多时代}

1. 前面提到的史学家与福音记载一致，确认阿尔赫劳在黑落德之后继承政权。他记述了阿尔赫劳如何凭其父黑落德的遗嘱和奥古斯都凯撒的谕令获得了犹太王国，以及在统治十年后如何失去王国，而他的兄弟斐理伯与小黑落德，连同吕撒尼雅，仍治理各自的四分领地。同一位作者在\\WorkTitle{犹太古史}第十八卷中说，在提庇留（他继奥古斯都之后登基，奥古斯都统治了五十七年）统治的约第十二年，般雀比拉多受命治理犹太地，并在那里住了整整十年，直到提庇留近死之时。\par

2. 因此，那些近来流传反对我们救主的\\WorkTitle{行传}之伪造者显系说谎。因为其中所载的日期本身就暴露了其编造者的虚假。\par

3. 因为他们竟敢将关于救主受难的事放在提庇留第四次执政期间，这发生在提庇留在位的第七年；此时显然比拉多尚未统治犹太，若我们相信若瑟夫的见证，他在上述著作中清楚表明，比拉多是在提庇留在位第十二年被任命为犹太总督的。\par

% structure: p0136 source=h2 target=subsection reason=relative-heading-rank; base=h2

\subsection{\Emph{第十章 基督施教时犹太人的大司祭}}

1. 据福音记载，在提庇留在位第十五年、般雀比拉多任总督第四年，当黑落德、吕撒尼雅和斐理伯管辖犹太其余地区时，我们的救主、主、天主的基督耶稣，年约三十岁，来到若翰面前受洗，并开始宣讲福音。\par

2. \WorkTitle{圣经}又告诉我们，他整个传教时期都在大司祭亚纳斯和盖法手下度过，这表明他的全部教导时间恰好涵盖这两位司祭的任期。既然他在亚纳斯任大司祭时开始工作，并一直教导到盖法任职为止，那么整个时间尚不足四年。\par

3. 因为自那时起，法律的礼仪已被废除，按照惯例，大司祭本应世袭且终身任职，但这一制度也被废止，罗马总督任命一人又一人担任大司祭，每人任期不超过一年。\par

4. 若瑟夫记载，从亚纳斯到盖法先后有四位大司祭。他在同一部\WorkTitle{犹太古史}中写道：\Quote{瓦勒留·格拉图终止了亚纳斯的司祭职，任命法比的儿子依市玛耳为大司祭。不久后将他撤职，任命大司祭亚纳斯的儿子厄肋阿匝尔担任同一职务。一年后也撤了他，将大司祭职赐给卡米特的儿子西满。但西满任职也不超过一年，随后被称为盖法的若瑟夫继任。}由此可见，我们救主的整个传教时期尚不足整整四年，因为从亚纳斯到盖法就任，四位大司祭每人任职一年。因此，福音正确指出救主受苦时的大司祭是盖法。由此我们也看出，救主的传教时间与前述考证并无矛盾。\par

5. 我们的救主、主在传教开始后不久，召选了十二位宗徒，并唯独称他们为宗徒，以示特殊尊荣。他又另外选立了七十人，派遣他们两个两个地在他前面，往他自己将要去的各城各地。\par

% structure: p0142 source=h2 target=subsection reason=relative-heading-rank; base=h2

\subsection{\Emph{第十一章 关于若翰洗者与基督的见证}}

1. 此后不久，若翰洗者被小黑落德斩首，正如福音所记载的。若瑟夫也记载了同一事实，并指名提到黑落狄雅，说她本是黑落德兄弟的妻子，黑落德休了前妻（培特辣王阿勒达的女儿）后，娶了她，并在她丈夫还活着时，将黑落狄雅从丈夫身边夺走。\par

2. 正是为了她，黑落德杀害了若翰，并因女儿所受的羞辱而与阿勒达交战。若瑟夫记载，在这场交战中，黑落德的全军覆没，他遭受此祸正是因他对若翰所犯的罪行。\par

3. 同一位若瑟夫在此记述中承认若翰洗者是一个极其正义的人，从而与福音中关于他的记载一致。他还记载，黑落德同样因黑落狄雅失去了王位，被与她一同流放，并被判生活在高卢的维也纳。\par

4. 他在\WorkTitle{犹太古史}第十八卷中叙述了这些事，关于若翰这样写道：\Quote{有些犹太人认为，黑落德的军队被天主消灭，乃是天主极其公义地为那位被称为洗者的若翰报仇。}\par

5. \Quote{因为黑落德杀死了他，他本是一个好人，劝勉犹太人来领受圣洗，践行德行，彼此之间并对天主实行义德；因为圣洗在他眼中是可接受的，当他们用它时，不是为了赦免某些罪，而是为了洁净身体，因为灵魂已在义德中洁净了。}\par

6. \Quote{当其他人聚集在他周围（因为他们听了他的话甚感喜悦），黑落德害怕他巨大的影响力可能引起某些叛乱，因为他们似乎准备做他可能建议的任何事。因此他认为，与其在革命发生、自己陷入困境时后悔，不如趁若翰的影响尚未生出任何新事之前，抢先杀了他。由于黑落德的猜疑，若翰被捆绑送往上述的马盖鲁堡垒，并在那里被杀。}\par

7. 在叙述了这些关于若翰的事之后，他在同一部作品中提到了我们的救主，如下：\Quote{那时有一位名叫耶稣的智者，如果可以称他为人的话。因为他能行奇事，教导那些欣然接受真理的人。他吸引了许多犹太人，也吸引了许多希腊人。他就是基督。}\par

8. \Quote{当比拉多因我们首领的控告判他受十字架刑时，那些起初爱他的人并未停止爱他。因为在第三天，他又活生生地显现给他们，神圣的先知们早已预言了这些以及无数关于他的奇妙之事。此外，以他命名的基督徒族类，一直延续到今天。}\par

9. 既然一位身为希伯来人本身的历史学家，在其著作中记载了这些关于若翰洗者和我们救主的事，那么那些伪造反对他们的行为的人，还有什么理由不被证明是毫无羞耻的呢？但这里就到此为止吧。\par

% structure: p0152 source=h2 target=subsection reason=relative-heading-rank; base=h2

\subsection{\Emph{第十二章 我们救主的门徒}}

1. 我们救主宗徒的名字，众人皆知，来自福音。但七十门徒却没有名单。据说，巴尔纳伯便是其中之一，\WorkTitle{宗徒大事录}多次提及他，尤其是保禄在\WorkTitle{迦拉达书}中。\par

2. 他们说与\Person{保禄}一同致书\Place{格林多}人的\Person{索斯特乃}也是其中之一。这是\Person{克莱孟}在他的\WorkTitle{《释要》}第五卷中的记述，其中他还说\Person{刻法}是七十门徒之一，他与宗徒\Person{伯多禄}同名，而\Person{保禄}所说的那个人就是\BibleQuote{当刻法来到\Place{安提约基雅}时，我当面反对了他。} \BibleRef{迦 2:11}\par

3. 此外，\Person{玛弟亚}——他替代\Person{犹达斯}被列于宗徒之中，以及与他一同被提名的人——据说也被视为与那七十人同蒙召选。他们还说\Person{达陡}也是其中之一，关于他我将随后记述流传下来的事迹。你若查考，便会发现我们的救主有超过七十位门徒，正如\Person{保禄}所证实的，他说救主从死者中复活后，首先显现给\Person{刻法}，然后显现给那十二位，此后又同时显现给五百多位弟兄，其中有些人已经长眠；但大多数在他写信时还在世。\par

4. 以后他说他显现给\Person{雅各伯}，他是救主的所谓弟兄之一。但是，除了这些人之外，还有许多其他的人被称为宗徒，效法那十二位，就如\Person{保禄}自己，他补充说：\BibleQuote{随后显现给众宗徒。} \BibleRef{格前 15:7} 关于这些人就说到这里。至于\Person{达陡}的事迹则如下所述。\par

% structure: p0157 source=h2 target=subsection reason=relative-heading-rank; base=h2

\subsection{第十三章  关于\Place{厄德萨}王公的叙述。}

1. 我们的主和救主耶稣基督的神性因他行奇能的事迹而传扬于万民之中，他吸引了无数来自远方、远离\Place{犹太}地的人，他们怀着希望，要治愈自己的疾病和种种苦楚。\par

2. 例如\Person{阿布加尔}王，他荣耀地统治着\Place{幼发拉的河}对岸的各族，身患重疾，非人力所能医治；当他听到耶稣的名声和他的奇迹——这些奇迹为众人一致证实——便派信使送信给他，恳求他医治自己的疾病。\par

3. 但耶稣当时并未应允他的请求；然而他还是赐予他一封亲笔信，说他会派自己的一位门徒来医治他的疾病，同时应许救恩临于他和他全家。\par

4. 不久之后，他的应许便实现了。因为在他从死者中复活、升天之后，十二宗徒之一的\Person{多默}受天主感动，派遣\Person{达陡}——他也被列为基督的七十门徒之一——前往\Place{厄德萨}，作为基督教训的宣讲者和福音传道者。\par

5. 我们救主所应许的一切，都藉着他得以实现。这些事的文字证据取自\Place{厄德萨}的档案，当时该城是一座王城。因为在那里的公共记录中，记载着古代的事迹和\Person{阿布加尔}的政令，这些材料保存至今。但没有比亲自阅读这些书信更好的办法了；我们已从档案中取出这些书信，并从\OriginalTerm{叙利亚文}{Syriac}逐字翻译如下。\par

% structure: p0163 source=h3 target=subsubsection reason=relative-heading-rank; base=h2

\subsubsection{\Person{阿布加尔}王致耶稣的书信副本，由快信使者\Person{阿纳尼雅}送往\Place{耶路撒冷}。}

6. \Place{厄德萨}王\Person{阿布加尔}，致在\Place{耶路撒冷}地方显现的卓越救主耶稣，问安。我听到了关于你和你所施治愈的传闻，你无需药物或草药。据说你使瞎子看见，瘸子行走，洁净癞病人，驱逐邪灵和魔鬼，医治久病之人，并使死者复活。\par

7. 我听到这一切关于你的事，断定必是二者之一：要么你是天主，从天降下而行这些事；要么你行这些事，是天主之子。\par

8. 因此我写信给你，询问你是否愿意劳驾到我这里来，治愈我所患的疾病。因为我听说犹太人正对你不满，图谋加害于你。但我有一座很小却尊贵的城，足以容纳我们两人。\par

% structure: p0167 source=h3 target=subsubsection reason=relative-heading-rank; base=h2

\subsubsection{耶稣经由信使\Person{阿纳尼雅}给\Person{阿布加尔}王的答复。}

9. \Quote{未曾看见我就相信的人，是有福的。因为经上论及我说：看见我的人不会信我，未曾看见我的人将相信而得救。至于你写信请我到你那里去，我必须在此完成我被派遣而来的一切事，然后被接升到那派遣我者那里。但我被接升天之后，将派我的一位门徒到你那里，医治你的疾病，并赐生命给你和你的家人。}\par

% structure: p0169 source=h3 target=subsubsection reason=relative-heading-rank; base=h2

\subsubsection{后续记述}

10. 在这些信函之后，还附加了如下的叙利亚文记载。耶稣升天之后，又名\Person{多默}的\Person{犹达斯}派了\Person{达陡}——一位属于七十人的宗徒——到他那里。\Person{达陡}来到后，住在\Person{托彼特}之子\Person{托彼特}家中。他的名声传开，有人告诉\Person{阿布加尔}，耶稣的一位宗徒来了，正如耶稣曾写信给他所说的。\par

11. \Person{达陡}于是开始藉天主的权能医治各种疾病和软弱，以致众人都惊异。\Person{阿布加尔}听说他所行的伟大奇事和治愈，便开始猜想他就是耶稣曾写信所说的：\Quote{我升天之后，将派我的一位门徒到你那里，他会治愈你。}\par

12. 于是，他召来\Person{达陡}所寄居的\Person{托彼特}，说：我听说一位有权能的人来了，住在你家中。带他来见我。\Person{托彼特}来到\Person{达陡}那里，对他说：\Person{阿布加尔}王召见我，让我带你去他那里，好让你医治他。\Person{达陡}说：我去，因我就是为此被派到他那里，并带着权能。\par

13. 多俾亚于是在次日早起，带着塔德来到阿布加尔处。当他来到时，贵族们都在场，围绕阿布加尔站着。他一进来，阿布加尔立即在宗徒塔德的面容上看见了一个大异像。阿布加尔看见后，就俯伏在塔德面前，而所有在场的人都惊讶；因为他们没有看见那异像，那异像只显示给阿布加尔一人。\par

14. 于是他问塔德，他是否真是天主子耶稣的门徒；耶稣曾对他说：\Quote{我要派我的一个门徒到你那里，他将医治你，并赐给你生命。}塔德说：\Quote{因为你已强力地信了那派遣我来的，所以我被派遣到你这里。而且，你若信他，你心里的祈求将按你所信得赐给你。}\par

15. 阿布加尔对他说：\Quote{我如此信了他，以致我曾想率领军队去消灭那些钉死他的犹太人，若不是因罗马人的统治而受阻的话。}塔德说：\Quote{我们的主已完成了父的旨意，完成之后便被接升到父那里去了。}阿布加尔对他说：\Quote{我也信了他和他的父。}\par

16. 塔德对他说：\Quote{因此，我以他的名按手在你身上。}他按手之后，阿布加尔的疾病和他所受的痛苦立刻痊愈了。\par

17. 阿布加尔十分惊奇，因为他所听见关于耶稣的事，如今果然藉他的门徒塔德亲身领受了；塔德不用药物和草药就治好了他，不仅治好他，也治好了患痛风的阿布杜斯之子阿布杜斯；因为此人来到塔德跟前，俯伏在他脚前，经他按手祝福后，也得了痊愈。同一位塔德还治好了城里许多其他居民，行了奇事和异迹，并宣讲了天主的话。\par

18. 后来阿布加尔说：\Quote{塔德啊，你以天主的能力行这些事，我们很惊奇。但除此之外，我请求你告诉我关于耶稣的来临：他如何诞生；以及他的能力，是凭哪种能力行了我所听见的那些事。}\par

19. 塔德说：\Quote{现在我暂且缄默，因为我奉命公开宣讲这话。但明天请你召集所有市民，我要在他们面前宣讲，在他们当中播下天主的话，论及耶稣的来临：他如何诞生；论及他的使命：父派遣他来有何目的；论及他工作的能力，以及他在世上所宣讲的奥秘，他凭哪种能力行了这些事；论及他的新宣讲，他的贬抑和谦卑，他如何自谦，死亡并贬抑自己的神性，被钉十字架，下降阴府，打破了从永远未被打碎的栅栏，复活了死者；因为他独自下降，却与多人一同上升，这样升到他的父那里。}\par

20. 阿布加尔于是命令市民清早集合，来听塔德的宣讲。此后，他命令给塔德金银。但塔德拒绝收受，说：\Quote{我们舍弃了自己的东西，怎能拿别人的东西呢？}这些事发生在第三百四十年。\par

我已将这些事按适当顺序插在这里，从叙利亚文逐字翻译出来，并希望是出于良好用意。\par

\SourceRecord{Translated by Arthur Cushman McGiffert.；Nicene and Post-Nicene Fathers, Second Series；Vol. 1.；Edited by Philip Schaff and Henry Wace.；Buffalo, NY: Christian Literature Publishing Co.,；1890.；<http://www.newadvent.org/fathers/250101.htm>.}

% structure: p0001 source=h1 target=section reason=page-title

\section{教会史（卷二）}

% structure: p0002 source=h2 target=subsection reason=relative-heading-rank; base=h2

\subsection{引言。}

1. 我们在前一卷中讨论了教会历史中那些必须作为引言来处理的主题，并附上了简要的证明。这些主题包括救恩圣言的天主性、我们所教导的教义的古老性，以及基督徒所过的福音生活的古老性，连同与基督最近显现、祂的受难以及宗徒拣选相关的事件。\par

2. 在本卷中，让我们考察祂升天后发生的事件，其中一些由神圣圣经证实，另一些则由我们将不时引用的著作来证实。\par

% structure: p0005 source=h2 target=subsection reason=relative-heading-rank; base=h2

\subsection{第一章　基督升天后宗徒所行之道。}

1. 首先，在背叛者犹大的位置上，玛弟亚——正如已经表明的，他也是七十人之一——被选入宗徒职。另外，为了会众的服务，通过祈祷和宗徒们的覆手，选定了七位被认可的人被委任为执事，其中斯德望是之一。他在主之后，在他被按立时，首先被杀害主的人用石头砸死，仿佛他为此目的而被提升。因此，他是第一个获得冠冕的人，与他的名号相符，这冠冕属于基督的殉道者，他们堪当胜利的奖赏。\par

2. 然后，雅各伯——古人因他卓越的品德而称他为义人——据记载首先被立为耶路撒冷教会的主教。这位雅各伯被称为主的兄弟，因为他被认为是若瑟的儿子，而若瑟被当作基督的父亲，因为童贞女与他订婚后，\BibleQuote{她因圣神有孕，在他们同居之前，}\BibleRef{玛 1:18}正如神圣福音的记述所表明的。\par

3. 但克肋孟在他的\WorkTitle{沉思录}第六卷中写道：\Quote{因为据说伯多禄、雅各伯和若望在我们救主升天后，如同也被我们的主所偏爱一样，并不争夺尊荣，而是选择了义人雅各伯为耶路撒冷的主教。}\par

4. 但同一作者在同一作品的第七卷中，也记载了关于他的以下事情：\Quote{主在复活后，将知识传授给义人雅各伯、若望和伯多禄，他们又将知识传授给其余的宗徒，其余的宗徒传授给七十人，其中巴尔纳伯是之一。但有两个雅各伯：一个被称为义人，他从圣殿的檐顶被扔下，被一个漂布匠用棍棒打死；另一个被斩首。}保禄也提到了同一位义人雅各伯，他写道：\BibleQuote{我没有看见别的宗徒，只看见了主的兄弟雅各伯。}\BibleRef{迦 1:19}\par

5. 那时，我们救主对奥斯若厄尼王的应许也应验了。因为多默，在神圣的推动下，派遣了塔德乌到厄德撒，作为基督信仰的宣讲者和福音传道者，正如我们稍前根据在那里发现的文献所表明的。\par

7. 当他到达那地方时，他用基督的话治好了阿布加尔；并通过他的作为使那里所有的人都拥有正确的心态，引导他们朝拜基督的权能，使他们成为救主教义的门徒。从那时直到现在，整个厄德撒城都忠于基督之名，这为我们救主对他们的恩惠提供了非凡的证明。\par

8. 这些事摘自古代记载；但现在让我们再回到神圣圣经。当犹太人因斯德望殉道而对耶路撒冷教会发动第一次也是最大的迫害时，所有门徒（除十二位外）都分散到犹太和撒玛黎雅各地，有些，如神圣圣经所说，甚至到了腓尼基、塞浦路斯和安提约基雅，但他们还不敢将信德之言传给外邦人，因此只向犹太人宣讲。\par

9. 在这期间，保禄仍在迫害教会，进入信徒的家中，将男女拉出来投入监牢。\par

10. 斐理伯也是与斯德望一同被委任执事的人之一，他在分散的人群中，下到撒玛黎雅，充满神圣的权能，首先将圣言传给那地的居民。神圣的恩宠如此有力地与他同工，连西满术士和许多其他人都被他的话吸引。\par

11. 西满那时如此著名，并藉他的戏法对被欺骗的人施加了如此大的影响，以致他被认为是天主的大能。但此时，他对斐理伯藉神圣权能所行的奇事感到惊讶，便假装和假冒对基督的信德，甚至接受了圣洗。\par

12. 令人惊讶的是，直到今天，那些追随他最不洁异端的人也在做同样的事。因为他们仿效他们祖先的样式，像一种瘟疫和麻风病那样潜入教会，极大地折磨那些他们能够将藏在自己里面的致命而可怕的毒药注入其中的人。他们中的大多数一旦在邪恶中被抓住就被驱逐出去，如同西满本人，被伯多禄识破后，得到了应得的惩罚。\par

13. 但随着救主福音的宣讲日益进展，某种上智安排引导了一位来自埃塞俄比亚地的女王军官——因为埃塞俄比亚直到今日仍按祖制由女子统治。他因一个启示，在斐理伯那里首先接受了神圣圣言奥秘的分享，成为全世界信徒的初果，据说他回到本国后首先宣讲了宇宙天主的认识以及我们救主生命赋予的临在人间的历程；如此，通过他，那预言真正得到应验，它宣称\BibleQuote{埃塞俄比亚伸手向天主}。\par

14. 此外，保禄，那\BibleQuote{蒙选的器皿，}\BibleRef{宗 9:15}，\BibleQuote{不是从人来的，也不是借着人，乃是借着耶稣基督本身，以及使他从死者中复活的父天主的启示，}\BibleRef{迦 1:1}，被立为宗徒，因他藉着由天上启示而发出的异象和声音，堪受此召叫。\par

% structure: p0019 source=h2 target=subsection reason=relative-heading-rank; base=h2

\subsection{第二章 提庇留得知比拉多关于基督的报告后如何反应}

1. 当我们的救主奇妙的复活与升天已在四处传扬时，按照各省统治者向皇帝报告该省发生的新奇事物以免他不知道的古老惯例，本丢·比拉多将全巴勒斯坦所传扬的关于我们的救主耶稣从死者中复活的报告告知了提庇留。\par

2. 他也汇报了从他那里得知的其他奇迹，以及他死后如何从死者中复活，如今被许多人相信是一位神。据说提庇留将此案提交元老院，但元老院拒绝了，表面上是因为他们尚未先行调查（因为有一条古老法令规定，除非经元老院投票和法令，任何人不得被罗马人尊为神），但实际上是因为神圣福音的救恩教导不需要人的证实和推荐。\par

3. 虽然罗马元老院拒绝了关于我们救主的提案，但提庇留仍然保留他最初的意见，并未对基督策划任何敌对措施。\par

4. 这些事由戴尔都良记载，他是精通罗马法律的人，在其他方面也享有盛誉，是罗马城中特别杰出的人物之一。他在为基督徒辩护的著作中，用拉丁文写成，并被译成希腊文，写道如下：\par

\Quote{5. 但为了从起源说明这些法律，有一条古老的法令：任何人不得由皇帝封神，除非元老院已表示赞同。马尔谷·奥勒留在某个偶像阿尔布努斯身上就是这样做的。这对我们的教义是有利的：在你们那里，神性尊严是由人的法令赋予的。如果一位神不讨人喜欢，他就不能成为神。因此，按照这一惯例，人必须对神施恩。}\par

6. 因此，提庇留——在其治下基督之名进入了世界——当这教义从它最初开始的巴勒斯坦向他报告时，他通知了元老院，表明他赞成这教义。但元老院因尚未亲自验证此事，便拒绝了。但提庇留仍坚持自己的意见，并威胁要处死控告基督徒的人。天上的上智已明智地将此灌输给他，为的是使福音的教义在开始时不受阻碍，得以传遍世界各地。\par

% structure: p0026 source=h2 target=subsection reason=relative-heading-rank; base=h2

\subsection{第三章 基督的教义迅速传遍普世}

1. 如此，在天上力量的影响和天主的合作下，救主的教义如同太阳的光芒，迅速照亮了整个世界；并且，依照神圣的圣经，灵感充盈的福音作者和宗徒们的声音传遍了全地，他们的话语直达地极。\par

2. 在每座城和每个村庄，教会迅速建立起来，充满如同盈满的禾场般众多的人群。那些因祖先流传的错误而心智被偶像崇拜迷信的古老疾病所束缚的人，因基督通过其门徒的教导和奇妙作为所施展的权能，得以从可怕的主宰下释放，并从最残酷的枷锁中找到解脱。他们厌恶地弃绝了各种魔鬼般的多神论，承认只有一位天主，万物的创造者，并以真正虔敬的礼仪敬拜祂，这是通过我们的救主在人类中所栽植的灵感与理性的敬拜。\par

3. 如今，天主的恩宠倾注在其余的民族身上。凯撒勒雅（巴勒斯坦）的科尔乃略及其全家，藉着天主的启示和伯多禄的中介，首先接受了基督内的信德；此后，安提约基雅的其他许多希腊人也接受了信德，那些因斯德望受迫害而四散的人曾向他们宣讲了福音。当安提约基雅的教会不断增长、人数增多，且许多来自耶路撒冷的先知在场时，其中有巴尔纳伯、保禄以及许多其他弟兄，基督徒之名在那里首次出现，如同从新鲜而赋予生命的泉源涌出。\par

4. 与他们同在的一位先知阿加波预言了即将发生的饥荒，保禄和巴尔纳伯被派去援助弟兄们的急需。\par

% structure: p0031 source=h2 target=subsection reason=relative-heading-rank; base=h2

\subsection{第四章 提庇留死后，卡里古拉立阿格黎帕为犹太人的王，并将黑落德处以永久放逐}

1. 提庇留驾崩，在位约二十二年，卡里古拉继位为帝。他立即任命阿格黎帕治理犹太人，立他为斐理伯和吕撒尼雅四区的王；不久之后，他又将黑落德的四区赐予他，因黑落德（就是救主受难时的那位）和他的妻子黑落狄雅犯下众多罪行，被处以永久放逐。约瑟夫是这些事实的见证人。\par

2. 在此皇帝治下，斐洛为世所知；他不仅在我们许多人中，而且在教会之外的许多学者中都是极负盛名的人。他生为希伯来人，但在亚历山大城中，他不亚于任何身居高位的人。他如何勤勉地研读圣经和他本国学问，从他已完成的作品中人人可见。他在哲学和外国学者的自由艺术方面何等精通，无需多言，因为据记载他在研究柏拉图主义和毕达哥拉斯哲学方面超越了他所有的同代人，他特别专注于这些研究。\par

% structure: p0034 source=h2 target=subsection reason=relative-heading-rank; base=h2

\subsection{第五章 斐洛为犹太人向卡里古拉派遣使团}

1. 斐洛以五卷书记述了犹太人于加约治下所遭遇的灾祸。他同时叙述了加约的疯狂：他如何自称神，在位期间施行无数暴政；他还描述了犹太人在其治下的苦难，并报告了他本人代表亚历山大里亚同胞出使罗马的经过；当他为祖传法律向加约陈情时，所得到的只是嘲笑和讥讽，几乎丧命。\par

2. 若瑟夫在\WorkTitle{古史}第十八卷中也提及此事，原文如下：亚历山大里亚的犹太人与希腊人发生冲突，双方各选三名代表前往加约处。\par

3. 亚历山大里亚代表之一阿比翁对犹太人提出诸多指控，其中一项是说他们忽略对凯撒应有的尊崇。因为所有其他罗马臣民都为加约设立祭坛和庙宇，在各方面对待他如同对待神祇，唯有犹太人认为用雕像尊崇他、并凭他的名起誓是可耻的。\par

4. 阿比翁发表了诸多严厉指控，希望以此激怒加约——这很可能发生——此时犹太使团首领斐洛，一位各方面都著名的人，是阿拉巴耳赫的亚历山大的兄弟，且精于哲学，准备针对指控进行答辩。\par

5. 但加约阻止他，命他离开，并且极其愤怒，显然在酝酿某种严酷措施对付他们。斐洛蒙受侮辱离去，告诉与他同来的犹太人要鼓起勇气；因为加约虽然向他们发怒，实际上却已与天主争战。\par

6. 以上是若瑟夫的记载。而斐洛本人在他所著的\WorkTitle{论使团}中详细准确地描述了他当时所做的一切。但我将略去大部分，只记录那些能使读者清楚看到犹太人的灾祸是在他们冒犯基督之后不久，且因同样的缘故临到他们的事。\par

7. 他首先叙述，在提庇留执政时期，当时深受皇帝信任的色雅努斯竭力要彻底消灭犹太民族；而在犹太地区，比拉多——正是在他任内发生了对救主的罪行——曾试图在圣殿问题上做出违反犹太法律的事，当时圣殿仍在耶路撒冷，从而激起了他们极大的骚乱。\par

% structure: p0042 source=h2 target=subsection reason=relative-heading-rank; base=h2

\subsection{第六章　犹太人对基督自恃之后所遭受的灾祸}

1. 提庇留死后，加约继承帝位，除了对许多人施行无数暴政外，尤其极大地折磨了整个犹太民族。我们可以从斐洛的话中简要了解这些事，他写道：\par

2. \Quote{加约对众人的态度如此反复无常，尤其是对犹太民族。他极其憎恨他们，以至于占用了其他城市中犹太人的会堂，从亚历山大里亚开始，在其中充满自己的像和雕像（因为他允许别人竖立，实际上等于他自己竖立）。在圣城中的圣殿，此前一直未受侵扰，被视为不可侵犯的避难所，也被他改变并转变为自己的庙宇，称之为可见的朱庇特——年轻的加约之庙。}\par

3. 同一作者在另一部题为\WorkTitle{论德性}的著作中，记载了在同一皇帝统治期间亚历山大里亚犹太人遭受的无数可怕而几乎不可言喻的灾难。若瑟夫也同意这一点，指出整个民族的灾祸始于比拉多时期，以及他们对救主的胆大妄为之罪。\par

4. 请听他在\WorkTitle{犹太战记}第二卷中如何写道：\Quote{比拉多受提庇留派遣任犹太总督，夜间秘密将带有皇帝像的军旗运入耶路撒冷。次日，此事在犹太人中引起极大骚动。因为附近的人见此情景都惊慌失措，认为他们的法律被践踏了。他们不允许在城中设立任何肖像。}\par

5. 将此事与福音作者的记载相比较，你将看到不久之后，他们便因曾在同一比拉多面前喊出的那句话而受罚，那时他们高呼除了凯撒别无君王。\BibleRef{若 19:15}\par

6. 同一作者又记载此后另一场灾祸临到他们。他写道：此后，他又利用称为\OriginalTerm{科尔班}{Corban}的圣库建造一条三百斯塔狄长的水道，从而激起另一场骚动。\par

7. 民众对此极为不满，比拉多到耶路撒冷时，他们围住他的审判台大声抱怨。但比拉多预见到骚动，已在人群中散布了伪装成平民的武装士兵，命令他们不要使用刀剑，但要用棍棒殴打那些呼喊的人。他从审判台上发出预定信号，犹太人遭到棍打，许多人因受伤而死，另有许多人在逃跑中被自己的同胞践踏丧命。民众因被杀者的命运而恐惧，遂沉默不语。\par

8. 除此之外，同一作者还记载了耶路撒冷城内发生的许多其他骚乱，并指出从那时起，叛乱、战争和阴谋接连不断，在城中及整个犹太地区从未止息，直至最后韦斯巴芗的围攻使他们覆灭。如此，天主的报应临到犹太人身上，因他们胆敢对基督所犯的罪行。\par

% structure: p0051 source=h2 target=subsection reason=relative-heading-rank; base=h2

\subsection{第七章　比拉多的自杀}

值得注意的是，比拉多本人，即在我们救主时代担任总督的那位，据传在卡伊乌斯（我们现在正记述他的时代）统治下遭遇如此不幸，以至于被迫成为自己的谋杀者和刽子手；如此，正如似乎的，天主的报应并未迟迟不临到他身上。这是那些记录奥林匹亚德以及每个时期各自发生的事情的希腊历史学家所陈述的。\par

% structure: p0053 source=h2 target=subsection reason=relative-heading-rank; base=h2

\subsection{第八章　克劳狄乌斯统治时期发生的饥荒}

1. 卡伊乌斯掌握政权不到四年，继位的是皇帝克劳狄乌斯。在他的统治下，世界遭遇了饥荒，那些与我们宗教完全无关的作家也在他们的历史中记载了这一点。因此，阿加波在《宗徒大事录》中记载的预言（\BibleRef{宗 11:28}），即全世界将遭遇饥荒，得到了应验。\par

2. 路加在《宗徒大事录》中，在提到克劳狄乌斯时代的饥荒，并说明安提约基雅的弟兄们各按能力，托保罗和巴尔纳伯的手送给犹太的弟兄们后（\BibleRef{宗 11:29}），接着记述了以下内容。\par

% structure: p0056 source=h2 target=subsection reason=relative-heading-rank; base=h2

\subsection{第九章　宗徒雅各伯的殉道}

1. \Quote{\BibleRef{宗 12:1} 那时}（显然他指的是克劳狄乌斯时代）\Quote{黑落德王下手苦害教会中的一些人，用剑杀了若望的兄弟雅各伯。}\par

2. 关于这位雅各伯，克肋孟在他的《假说》第七卷中记载了一个值得提及的故事；他根据从前辈那里听来的说法讲述。他说，那个把雅各伯带到审判席的人，看见他作证，深受感动，承认自己也是一个基督徒。\par

3. 因此，他说，他们两人一同被带走；在路上，那人请求雅各伯原谅他。雅各伯稍加思索后，说：\Quote{愿你平安}，并亲吻了他。就这样，他们两人同时被斩首。\par

4. 然后，正如神圣的圣经所说（\BibleRef{宗 12:3}），雅各伯死后，黑落德见此事使犹太人喜悦，就下手拿住伯多禄，把他关进监牢，本想杀他；若非因一位天使夜间显现奇迹般地解开他的锁链，释放他继续为福音服务，他已遭杀害。这就是天主对伯多禄的眷顾。\par

% structure: p0061 source=h2 target=subsection reason=relative-heading-rank; base=h2

\subsection{第十章　阿格黎帕，又称黑落德，迫害宗徒后立即遭受天主的惩罚}

1. 国王对宗徒们的图谋不久就带来了后果；神圣正义的报复之仆在他对宗徒们施计之后立刻追上了他，正如《宗徒大事录》所记载的。因为当他前往凯撒勒雅，在一个著名的节日里，穿着华丽而威严的王服，在审判台前的高台上向民众发表演说时，全体群众欢呼喝彩，仿佛那是神的声音而非人的声音；圣经记载，主的一位天使击打了他，他被虫子咬死，断了气（\BibleRef{宗 12:23}）。\par

2. 我们必须钦佩约瑟夫关于这一奇妙事件的记载与神圣圣经的一致性；因为他明确地在他的《犹太古史》第十九卷中证实了真理，以以下话语叙述了这件奇事：\par

3. 他统治全犹太地三年后，来到凯撒勒雅，即从前称为斯特拉托塔的地方。在那里，他为凯撒举行竞技，得知这是一项为凯撒安全而举行的庆典。在这场庆典中，聚集了该省众多最高贵、最显赫的人物。\par

4. 在竞技的第二天，他于破晓时分前往剧场，身穿一件全银编织、质地奇妙的袍子。在那里，银器被初升太阳的第一缕光线反射，神奇地闪耀着光芒，如此明亮，以至于在那些注视他的人心中产生了一种恐惧和敬畏。\par

5. 立刻，那些谄媚者，有的从这边，有的从那边，提高声音，用一种对他并无益处的口吻称他为神，说：\Quote{请仁慈待我们；如果直到如今我们敬畏你如同一个人，从今以后我们承认你超越凡人的本性。}\par

6. 国王没有责备他们，也没有拒绝他们不敬的谄媚。但过了一会儿，他抬头看见一位天使坐在他的头顶上方。他立刻意识到这将是不幸的原因，正如它曾经是好运的原因，他被一种刺心的痛苦所击打。\par

7. 随即，以最剧烈的痛楚开始的痛苦攫住了他的内脏。他望着朋友们说：\Quote{我，你们的神，如今被命令离开此生；命运当场就驳斥了你们刚才关于我所说的谎言。被你们称为不死的人，如今正被带向死亡；但我们必须接受命运，正如天主所命定。我们度过的生活绝非不光彩，而是在那种被称为幸福的光辉之中。}\par

8. 他说完这话，痛苦加剧。于是他被急忙抬进王宫，消息传开，说国王必定很快死去。但民众，连同他们的妻子儿女，按祖先的习俗坐在麻布上，为国王向天主祈祷，每个地方都充满哀号和眼泪。国王躺在一间高处的房间里，看到他们伏在地上，自己也不禁流泪。\par

9. 他在肠痛中持续受苦五天后，去世了，享年五十四岁，在位第七年。他在卡伊乌斯皇帝统治下统治了四年——其中三年治理斐理伯的封地，第四年加上了黑落德的封地——在克劳狄乌斯皇帝统治下又统治了三年。\par

10. 我非常惊奇，若瑟夫在这些事以及别的事上，如此完全符合圣经。但若有人以为关于君王的名字有分歧，至少时间和事件表明所指为同一人，无论名字的改变是因抄写员的错误，还是因为他像许多人一样有两个名字。\par

% structure: p0072 source=h2 target=subsection reason=relative-heading-rank; base=h2

\subsection{第11章 伪先知特乌达斯及其追随者}

1. 路加在\WorkTitle{宗徒大事录}中记载，加玛里耳在关于宗徒的会议上说，在那时\Quote{特乌达起来，自命不凡，他被杀，所有服从他的人都四散了。} \BibleRef{宗 5:36} 因此，让我们加上若瑟夫关于此人的记载。他在上文提到的著作中记述了以下情况：\par

2. 当法杜斯任犹达省总督时，一个名叫特乌达斯的骗子说服了极多的人带着他们的财产跟随他到约但河。因为他说自己是个先知，河水会因他的命令分开，给他们一条易行的通道。\par

3. 他用这些话欺骗了许多人。但法杜斯不容他们享受愚行，派出一队骑兵袭击他们，出其不意杀了多人，活捉了许多人。他们俘获了特乌达斯本人，砍下他的头带到耶路撒冷。此外，他也提到在克劳狄统治期间发生的饥荒，如下所述：\par

% structure: p0076 source=h2 target=subsection reason=relative-heading-rank; base=h2

\subsection{第12章 奥斯若恩女皇赫楞}

1. \Quote{那时在犹达发生大饥荒，赫楞女王花费巨款从埃及购买粮食，分发给穷人。}\par

2. 你会发现这话也与\WorkTitle{宗徒大事录}相符，其中说安提约基雅的门徒\Quote{每人按自己的力量决定送救济给住在犹达的弟兄们；他们就这样做了，托巴尔纳伯和保禄的手送到长老那里去。}\par

3. 这位历史学家提到的赫楞的宏伟纪念碑至今仍显于现在称为埃利亚城的郊外。据说她是阿迪亚贝尼的女王。\par

% structure: p0080 source=h2 target=subsection reason=relative-heading-rank; base=h2

\subsection{第13章 西满术士}

1. 但我们救主和主耶稣基督的信仰既已传遍众人，那憎恨一切善事并图谋人类救恩的仇敌便设计将帝国京城据为己有。他带领上述的西满到那里，助长其欺诈伎俩，迷惑了许多罗马居民，从而将他们纳入自己的权势。\par

2. 这是我们一位杰出的作家犹斯定所陈述的，他生活于宗徒时代之后不久。我将在适当的地方谈到他。请取阅此人的著作，他在为我们的宗教写给安敦尼的第一篇\WorkTitle{护教书}中写道：\par

3. 主升天之后，魔鬼推出某些自称是神的人，你们不仅允许他们不受迫害，甚至认为他们配得尊荣。其中之一是西满，一位来自吉托村的撒玛黎雅人，在克劳狄皇帝统治期间，藉着运行在他内的魔鬼的技艺，在你们帝国京城行了某些大能的魔法，被看作神，且被你们以一座雕像尊为神，该雕像竖立在台伯河上，两桥之间，上有拉丁文铭文 \OriginalTerm{向神圣的西满神}{Simoni Deo Sancto}。\par

4. 几乎所有的撒玛黎雅人，甚至少数其他民族，都承认并崇拜他为第一个神。那时有一个名叫赫楞的女子与他同行，她原是在腓尼基提洛作妓女的；他们称她是从他发出的第一个意念。\par

5. 犹斯定叙述这些事，依勒内也在他\WorkTitle{驳斥异端}的第一部书中同意他，并记述了此人和他亵渎不洁的教导。在此引用他的叙述是多余的，因为那些想要知道各个跟随他的异端创始人的起源、生活和谬误教义，以及他们所有实践习俗的人，可以在依勒内的上述著作中找到详尽论述。\par

6. 我们得知西满是所有异端的创始人。从他的时代直到现在，跟随他异端的人假装持守基督徒克己的哲学，这哲学因其生活的纯洁而闻名于世。但他们却又拥抱他们似乎已弃绝的偶像迷信，跪拜西满本人和与他同在的上述赫楞的画像和雕像，大胆地用馨香、祭品和奠酒崇拜它们。\par

7. 至于他们所保守的更隐秘的事，据他们说，人初次听到会惊异，并用他们中间流行的一种书面说法来说，会困惑，其实是充满骇人听闻的事，癫狂和愚蠢，其性质不仅不能写成文字，就连有节制的人也无法启齿，因其过于卑鄙淫秽。\par

8. 因为凡能想象出的比最卑劣的更卑劣之事——这一切都已被这个最可憎的教派所超越，此教派由那些把真正被各种恶习压垮的可怜女子当作玩物的人组成。\par

% structure: p0089 source=h2 target=subsection reason=relative-heading-rank; base=h2

\subsection{第14章 宗徒伯多禄在罗马的宣讲}

1. 那憎恨一切善事并图谋人类救恩的邪恶势力，当时立西满为如此恶行的父亲和创始者，仿佛要使他成为我们救主伟大而蒙受灵感的宗徒的巨大对手。\par

2. 因为那与其臣仆合作的神圣而属天的恩宠，藉着他们的显现和同在，迅速熄灭了点燃的邪恶火焰，并通过他们降低了\Quote{一切高举起来反对天主之知识的傲慢事物。} \BibleRef{格后 10:5}\par

3. 因此，无论是西满的阴谋，还是当时出现的其他任何人的阴谋，都无法在那些宗徒时代成就任何事。因为一切都因真理的光辉和那神圣的圣言本身而被征服和制服，这圣言刚刚开始从天上照耀人间，并在地上兴盛，且住在宗徒们自己内。\par

4. 上述那位骗子立即被一道神圣而神奇的闪光击中了他的心灵之眼，在宗徒伯多禄于犹太揭露了他所做的恶事之后，他逃走并长途跋涉，从东方渡海到西方，以为只有这样他才能按自己的心意生活。\par

5. 他来到罗马城，由于那潜伏在那里的势力的强大合作，他的计划在短时间内非常成功，以至于那里的人通过立雕像尊崇他为神。\par

6. 但这并没有持续多久。因为立即，在克劳狄乌斯统治期间，那全善而仁慈、照管万有的天主眷顾，带领伯多禄——这位最坚强、最伟大的宗徒，因他的德行而成为所有宗徒代言人的那位——到罗马对抗这个生命的巨大败坏者。他如同一位高贵的天主的统帅，身穿神圣的盔甲，从东方将智慧的宝贵商品即光明带到住在西方的人那里，宣告光明本身和带来灵魂救恩的圣言，并宣讲天国。\par

% structure: p0096 source=h2 target=subsection reason=relative-heading-rank; base=h2

\subsection{第十五章 马尔谷福音}

1. 这样，当神圣的圣言在他们中间安家时，西满的势力立即被扑灭并与他本人一同被摧毁。虔诚的光辉如此照亮了伯多禄听众的心灵，以至于他们不满足于只听到一次，也不满足于神圣福音的无形教导，而是以各种恳求请求马尔谷（伯多禄的门徒，他的福音尚存）将他们口头领受的教导留下一份书面记录。他们坚持不懈，直到说服了此人，并因此成了那部以马尔谷命名的书面福音的缘起。\par

2. 据说伯多禄通过圣神的启示得知了所做的事后，对这些人的热诚感到满意，并且这部作品得到了他权威的认可，以便在教会中使用。克莱孟在他的\WorkTitle{学术要义}第八卷中给出了这一记载，希拉波利斯的主教帕皮亚也同意他的说法。伯多禄在他的第一封书信中提到了马尔谷，据说这封信是在罗马写成的，正如他通过比喻称该城为巴比伦所表明的那样，他在以下话中说：\Quote{在巴比伦与你们一同被选的教会问候你们，我的儿子马尔谷也问候你们。}\par

% structure: p0099 source=h2 target=subsection reason=relative-heading-rank; base=h2

\subsection{第十六章 马尔谷首先向埃及居民宣示基督教}

1. 据说这位马尔谷是第一个被派往埃及的人，他宣示了他所写的福音，并在亚历山大首先建立了教会。\par

2. 起初在那里聚集的信徒，无论男女，人数众多，他们过着最哲学和最严格的苦修生活，以至于斐罗认为值得描述他们的追求、聚会、款待和整个生活方式。\par

% structure: p0102 source=h2 target=subsection reason=relative-heading-rank; base=h2

\subsection{第十七章 斐罗论及埃及的苦修者}

1. 据说斐罗在克劳狄乌斯统治时期在罗马结识了当时在那里宣讲的伯多禄。这并非不可能，因为我们提到的那部他后来写成的著作显然包含了直到今天在我们中间遵守的那些教会规则。\par

2. 由于他尽可能准确地描述了我们苦修者的生活，因此他不仅认识，而且赞同、尊敬并赞扬了他那个时代的宗徒们，这些人似乎属于希伯来种族，因此按照犹太人的方式遵守了大部分古代习俗。\par

3. 在他题为\WorkTitle{论默观生活或论恳求者}的著作中，他首先肯定他不会在他的叙述中增添任何违背真理或自己编造的东西，然后说这些人被称为\OriginalTerm{泰拉普特}{Therapeutae}，与他们在一起的妇女被称为\OriginalTerm{泰拉普特里}{Therapeutrides}。然后他解释了这名称的理由，说它源于他们像医生一样治疗和治愈那些来到他们那里的人的灵魂，除去邪恶的情欲，或者源于他们以纯洁和真诚侍奉和崇拜神。\par

4. 无论是斐罗自己给了他们这个名称，用了一个适合他们生活方式的称号，还是他们最初就这样自称，因为基督徒这个名称尚未到处传开，我们在此无需讨论。\par

5. 然而，他首先见证了他们放弃财产。他说，当他们开始这种哲学生活方式时，他们将财物分给亲属，然后放弃所有生活的牵挂，走出城墙，住在孤寂的田野和花园中，深知与不同性格的人交往是无益和有害的。他们当时这样做，似乎是出于一种热烈而坚定的信德，效法先知们的生活方式。\par

6. 因为在被普遍承认为正典的\BibleBook{}{宗徒大事录}中记载，宗徒们的所有同伴都变卖了他们的产业和财物，按照各人的需要分给众人，以致他们中间没有一个匮乏的人。\BibleQuote{凡有田地或房屋的，都卖了，把卖得的价银拿来，放在宗徒们脚前，照每人所需要的分给每人。} \BibleRef{宗 2:45}\par

7. \Person{斐洛}为与这里所描述的极为相似的事实作证，然后补充了以下记载：在全世界都能找到这个种族。因为既已圆满地享有至善，希腊人和外邦人都理当有分。但这个种族尤其在埃及各处，在其所谓的各\OriginalTerm{诺姆}{nome}中，特别是在亚历山大附近，为数众多。\par

8. 来自各地的优秀人士，如同迁往\OriginalTerm{治疗者}{Therapeutae}父的殖民地，移居到一个非常合适的地点，它位于\Place{玛利亚湖}之上，坐落在一座低矮的山丘上，因其安全与气候温和而位置极佳。\par

9. 随后，在描述了他们的房屋类型之后，他接着谈到他们散布在各处的教堂，说：\Quote{每间房屋内都有一个称为圣所和隐修院的圣室，他们在那里独处，实践宗教生活的奥秘。他们不将任何东西带进去，既不帶饮料也不带食物，也不带任何其他满足身体需要的东西，只带法律、先知受默感的神谕、赞美诗，以及那些增进并完善他们知识与虔敬的其他事物。}\par

10. 在谈及一些其他事情之后，\Person{斐洛}说：\par

从早到晚的整个时间段，对他们而言是修行的时间。因为他们阅读圣经，并以寓意的方式解释他们祖先的哲学，将书面文字视为隐藏真理的象征，这真理以隐晦的形象传达。\par

11. 他们也有古代人的著作，那些人曾是他们宗派的创立者，并留下了许多寓意法的纪念碑。他们以这些著作为模型，并模仿他们的原则。\par

12. 这些话似乎是由一个曾听过他们解释其圣书的人所述。但极有可能，他所称的古代人的著作，乃是福音书和宗徒们的著作，或许也包括一些古代先知的注释，例如\BibleBook{希伯来}{书}和\Person{保禄}的其他许多书信中所包含的那些。\par

13. 接着，他关于他们所创作的新圣咏又写道：\Quote{所以他们不仅花时间默想，而且以各种韵律和曲调编撰赞美诗和颂词献给天主，当然，他们将其划分为超出寻常庄严的音节。}\par

14. 同一部书还记载了许多其他事物，但似乎有必要选取那些展示教会生活方式特征的事实。\par

15. 但若有人认为以上所述并非\OriginalTerm{福音政体}{Gospel polity}所特有，而可以应用于所提及者之外的人，那么让同一位作者后来的话说服他吧；如果他无偏见，将在其中找到关于此事无可争议的见证。\Person{斐洛}的话如下：\par

16. 他们将节制作为灵魂的基础奠定之后，在其上建造其他德行。他们中没有人可以在日落前进食或饮水，因为他们认为从事哲学思考是配得上光明的工作，而关注身体的需要则只适宜在黑暗中，因此将白天分配给前者，而将黑夜的一小部分分配给后者。\par

17. 但是有些人心怀极大的求知欲，三天忘记进食；有些人则如此喜悦，如此丰盛地享用智慧——它毫不吝啬地提供丰盛的教义——以致他们甚至两倍于此的时间禁食，并习惯于六天之后才勉强摄取必要的食物。我们认为\Person{斐洛}的这些陈述清楚而无可争议地指向我们共融团体中的人。\par

18. 但若在此之后，仍有人固执地否认这指向，让他摒弃不信，并由更显著的例子说服吧；这些例子只能在基督徒的\OriginalTerm{福音宗教}{evangelical religion}中找到。\par

19. 因为他们说，我们正在谈论的那些人中还有妇女，其中大多数是年长的贞女，她们保持贞洁，并非出于不得已，如同希腊人中某些女祭司那样，而是出于自己的选择，出于对智慧的热忱和渴望。她们热切地以智慧为伴侣而生活，不关注肉身的快乐，寻求的不是可朽的后裔，而是不朽的后裔——只有虔诚的灵魂才能凭自身孕育这样的后裔。\par

20. 随后不久，他更加有力地补充说：\Quote{他们以寓意的方式喻意式地解释圣经。因为这些人视整个法律如同一个活的有机体，其中话语构成身体，而隐藏于话语中的意义构成灵魂。这个隐藏的意义首先由这一派别特别研究，他们看见，正如在名字的镜子中所揭示的那样，思想的无与伦比之美。}\par

21. 何以须要再补充这些聚会、聚会中男女各自的职责，以及那些直到如今仍为我们所惯常遵守的实践，尤其是我们在\OriginalTerm{救主受难节}{feast of the Saviour's passion}惯常遵守的实践，包括斋戒、守夜和研读天主圣言。\par

22. 上述作者在他自己的著作中叙述了这些事情，指出一种直到如今仍由我们单独保存的生活方式，特别记载了与大节日相关的守夜，以及这些守夜中所做的功课，还有我们所惯常吟唱的圣咏，并描述如何：当一人按节拍规律地歌唱时，其余人静默聆听，只在圣咏结尾处加入咏唱；以及在这些日子，他们睡在地上的草垫上，用他自己的话说，\Quote{完全未尝酒和任何肉食，水是他们唯一的饮料，面包的配菜是盐和\OriginalTerm{牛膝草}{hyssop}。}\par

23. 此外，\Person{斐洛}还描述了从事教会服务者中的品级次序，提及\OriginalTerm{执事职}{diaconate}和位居所有其他品级之上的\OriginalTerm{主教}{bishop}职务。但谁若想更准确地了解这些事，可从上述历史著作中获知。\par

24. 但众人都清楚，\Person{斐洛}写这些事时，心中所指的是福音的首批传报者以及从起初由宗徒们传下来的习俗。\par

% structure: p0128 source=h2 target=subsection reason=relative-heading-rank; base=h2

\subsection{第十八章 \Person{斐洛}流传至今的著作。}

1. \Person{斐洛}语言丰富，思想广博，对圣经的观点崇高而卓越，他创作了多种多样的圣经注释。一方面，他在称为\WorkTitle{神圣律法寓意}的书中按顺序解释了\BibleBook{创}{世纪}中的事件；另一方面，他在\WorkTitle{创世纪与出世纪问答}一书中对作为探究对象的圣经章节进行了连续划分，并提出了质疑与解答。\par

2. 此外，还有他专门就某些主题撰写的论文，例如\WorkTitle{论农业}两卷，以及同数的\WorkTitle{论醉酒}；还有一些以各自内容相应的不同标题命名的作品，例如\WorkTitle{论清醒心灵所愿与所憎}、\WorkTitle{论语言混乱}、\WorkTitle{论逃亡与发现}、\WorkTitle{论为教导而集会}、\WorkTitle{论谁承受神圣事物？}、\WorkTitle{论事物的相等与不等之分}，以及进一步的\WorkTitle{论梅瑟所描述的三德与其他}。\par

3. 此外还有\WorkTitle{论名字的改变及其原因}这部作品，其中他说自己也写了\WorkTitle{论盟约}两卷。\par

4. 还有他的著作\WorkTitle{论迁移}，以及\WorkTitle{论在义德上达到完美的智者的生活}或\WorkTitle{论不成文法}；此外还有\WorkTitle{论巨人}或\WorkTitle{论天主的不变性}，以及\WorkTitle{论按梅瑟说法梦是由天主派遣}第一、二、三、四、五卷。这些便是流传至今的关于\BibleBook{创}{世纪}的著作。\par

5. 但关于\BibleBook{出}{世纪}，我们知道有\WorkTitle{问答}第一、二、三、四、五卷；还有\WorkTitle{论会幕}，以及\WorkTitle{论十诫}，以及\WorkTitle{论特别涉及十诫主要部分的法律}四卷，还有\WorkTitle{论献祭的动物}与\WorkTitle{论祭献的种类}，以及另一部\WorkTitle{论法律为善人设定的赏报与为恶人设定的惩罚和诅咒}。\par

6. 除所有这些外，还可找到某些单卷作品，例如\WorkTitle{论眷顾}，以及他撰写的\WorkTitle{论犹太人}和\WorkTitle{政治家}；还有\WorkTitle{亚历山大}，或\WorkTitle{论非理性动物拥有理性}。此外还有\WorkTitle{论每个恶人是奴隶}，其后附有\WorkTitle{论每个善人是自由的}。\par

7. 此后他撰写了\WorkTitle{论默观生活}或\WorkTitle{论恳求者}，我们从中引用了关于宗徒时代人物的生活事实；此外还有，\WorkTitle{法律与先知中希伯来名字的解释}，据说是他勤奋的成果。\par

8. 据说他在\Person{克劳狄}统治期间，曾在整个\Place{罗马元老院}面前宣读了他撰写的作品，该作品是他当\Person{加约}在位时来到\Place{罗马}时写的，关于\Person{加约}对诸神的憎恨，他讽刺性地将此书命名为\WorkTitle{论美德}。他的演讲如此受人赞赏，以致被认为值得收藏于图书馆中。\par

9. 此时，\Person{保禄}正完成他\BibleQuote{从\Place{耶路撒冷}直到\Place{依里黎果}}的行程（见：\BibleRef{罗 15:19}），\Person{克劳狄}将犹太人驱逐出\Place{罗马}；\Person{阿桂拉}和\Person{普黎史拉}与其他犹太人一同离开\Place{罗马}，来到\Place{亚细亚}，并与宗徒\Person{保禄}同住，\Person{保禄}正在坚固那个地区他新近奠定根基的教会。圣书\WorkTitle{宗徒大事录}也记载了这些事。\par

% structure: p0138 source=h2 target=subsection reason=relative-heading-rank; base=h2

\subsection{第十九章 逾越节当日在\Place{耶路撒冷}降临在犹太人身上的灾祸。}

1. \Person{克劳狄}仍为皇帝时，在\Place{耶路撒冷}的逾越节庆期中发生了如此巨大的骚乱和动荡，以至于仅被强行挤在圣殿门口的犹太人就有三万人因相互践踏而丧生。因此，这个庆节成了全体民族的哀悼之时，家家户户都有哭泣声。这些事情皆由\Person{若瑟夫}如实记述。\par

2. 但\Person{克劳狄}任命\Person{阿格黎帕}之子\Person{阿格黎帕}为犹太人的王，并派\Person{斐理斯}为整个\Place{撒玛黎雅}、\Place{加里肋亚}以及称为\Place{培勒雅}地区的总督。他在位十三年零八个月后去世，留下\Person{尼禄}继承他的帝国。\par

% structure: p0141 source=h2 target=subsection reason=relative-heading-rank; base=h2

\subsection{第二十章 \Person{尼禄}统治时期在\Place{耶路撒冷}发生的事件。}

1. \Person{若瑟夫}在他的\WorkTitle{犹太古史}第二十卷中又叙述了\Person{尼禄}统治时期，\Person{斐理斯}任\Place{犹太}总督时，司祭之间发生的争执。\par

2. 他的话如下：大司祭一方与司祭及\Place{耶路撒冷}民众首领之间发生了争执。双方各自聚集了一批最胆大妄为的人，并自任首领，每当相遇时便互相辱骂并投掷石块。无人敢于调停；这些事便恣意而行，仿佛在一座没有统治者的城中。\par

3. 大司祭的厚颜无耻和胆大妄为达到了如此地步，竟敢派仆人到打谷场强取应归司祭的什一之物；于是那些贫穷的司祭眼看着因匮乏而濒于死亡。就这样，各派的暴力凌驾于一切正义之上。\par

4. 同一位作者又叙述，大约在同一时期，\Place{耶路撒冷}出现了一种强盗，正如他所说，\Quote{他们白天在城中间杀死遇到的人。}\par

5. 尤其在节庆时，他们混入人群中，用藏在衣服里的短剑刺杀最显贵的人。当受害者倒下时，凶手自己也在那些表示愤慨的人之中。因此，由于众人对他们的信任，他们始终未被发现。\par

6. 第一个被他们杀害的是大司祭\Person{约纳堂}；此后每天有许多人被杀死，直到恐惧本身比邪恶更甚，每个人都如在战场上一样，时刻等待死亡。\par

% structure: p0148 source=h2 target=subsection reason=relative-heading-rank; base=h2

\subsection{第二十一章 埃及人——\WorkTitle{宗徒大事录}中也曾提及。}

1. 在其他事情之后，他接着写道：但犹太人遭受了比这些更严重的灾害，来自那位埃及假先知。因为那地出现了一个骗子，他使人们相信他是先知，聚集了约三万被他欺骗的人，带领他们从旷野到所谓的\Place{橄榄山}，准备用武力进入\Place{耶路撒冷}，击败罗马驻军，夺取人民的统治权，用那些与他一起攻击的人作为护卫。\par

2. 但\Person{斐理斯}预先得知他的攻击，带着罗马军团士兵出去迎战，所有民众都参与防御，因此战斗发生时，埃及人带着少数随从逃跑了，但大部分被消灭或被俘。\par

3. \Person{若瑟夫}在他的\WorkTitle{历史}第二卷中记述了这些事。但将这里关于埃及人的记载与\WorkTitle{宗徒大事录}中的记载相比较是值得的。在\Person{斐理斯}时期，当犹太民众向宗徒发动骚乱时，在\Place{耶路撒冷}的百夫长对\Person{保禄}说：\BibleQuote{你莫非就是前些日子作乱，带领四千个刺客往旷野去的那人吗？} \BibleRef{宗 21:38} 这些就是\Person{斐理斯}时期发生的事。\par

% structure: p0152 source=h2 target=subsection reason=relative-heading-rank; base=h2

\subsection{第二十二章 \Person{保禄}被捆绑从\Place{犹太}解往\Place{罗马}，辩护后获释，免受一切指控。}

1. \Person{斐斯托}被\Person{尼禄}派去接替\Person{斐理斯}。在他手下，\Person{保禄}辩护后被捆绑解往\Place{罗马}。\Person{阿黎斯塔苛}与他同去，\Person{保禄}在他某封书信中恰当地称他为自己的同囚。 \BibleRef{哥 4:10} \Person{路加}写了\WorkTitle{宗徒大事录}，他的记述到此为止，说明\Person{保禄}在\Place{罗马}作为自由的囚犯度过了整整两年，毫无阻碍地宣讲天主的圣言。\par

2. 据说他在辩护之后，又被差遣从事宣讲的职务，第二次来到同一城市时殉道。在这次囚禁中，他写了\WorkTitle{弟茂德后书}，其中提到他首次辩护和即将到来的死亡。\par

3. 请听他关于这些事的见证：\BibleQuote{在我初次辩护时，}他说，\BibleQuote{没有人支持我，众人都离弃了我；愿天主不将此罪归于他们。然而主站在我身边，坚固了我，使我藉着我能完全宣讲，使所有外邦人都能听见；我也从狮子的口中被救出。} \BibleRef{弟后 4:16}\par

4. 他在这话中清楚地表明，在先前那次，为了使他能完成宣讲，他从狮子口中被救出；这措辞很可能指\Person{尼禄}，因其残忍。他后来并未加上类似的话\Quote{他必救我脱离狮子的口}；因为他以神视预见到自己的结局不久将至。\par

5. 因此他在\BibleQuote{他从狮子口中救了我}这句话后加上：\BibleQuote{主必救我脱离各种凶恶的事，并保守我进入他的天国。} \BibleRef{弟后 4:18} 表明他即将殉道；在同一封信中他更清楚地预言此事，写道：\BibleQuote{因为我已成了奠祭，我离世的时期已经近了。}\par

6. 此外，他在\WorkTitle{弟茂德后书}中表明，他写信时\Person{路加}与他同在，但在他初次辩护时，连\Person{路加}也不在。因此很可能\Person{路加}那时写了\WorkTitle{宗徒大事录}，将他的记述延续到他与\Person{保禄}同处的时期。\par

7. 我们引证这些事，是为表明\Person{保禄}的殉道并未发生在\Person{路加}所记的那次\Place{罗马}停留期间。\par

8. 很可能，因为\Person{尼禄}起初较为温和，\Person{保禄}为自己的教义辩护较易被接受；但当他做出无法无天的胆大妄为之事后，他也使宗徒们和其他人成为他攻击的对象。\par

% structure: p0161 source=h2 target=subsection reason=relative-heading-rank; base=h2

\subsection{第二十三章 称为主的兄弟的\Person{雅各伯}殉道。}

1. 但在\Person{保禄}因向凯撒上诉而被\Person{斐斯托}解往\Place{罗马}之后，犹太人因未能用为他们所设的陷阱捉拿他而失望，便转而对付主的兄弟\Person{雅各伯}——宗徒们曾将\Place{耶路撒冷}的主教职位委托给他。他们对他采取了以下大胆的措施。\par

2. 他们把他带到人群中间，要求他在全体民众面前背弃对基督的信仰。但出乎所有人的意料，他声音清晰，以比他们预想更大的胆量，在全体群众面前宣告，我们的救主和主耶稣是天主子。他们再也无法忍受此人的见证——他因生活中所展现的卓越苦修美德和虔诚，被所有人视为最正义的人——于是便杀了他。这一暴力行为的时机是由于当时盛行的无政府状态，因为\Person{斐斯托}恰好在这时死于\Place{犹太}，该省没有了总督和首领。\par

3. \Person{雅各伯}死亡的方式已由上述\Person{克莱孟}的话指出，他记载说\Person{雅各伯}从圣殿的翼顶被扔下，并被棍棒打死。但生活于宗徒之后的\Person{赫格西普}在他的\WorkTitle{回忆录}第五卷中给出了最准确的记载。他写道：\par

4. \Person{雅各伯}，主的兄弟，与宗徒们一同接管了教会的治理。从我们救主时代直到今日，众人都称他为义者；因为有许多人名叫雅各伯。\par

5. 他从母腹中就是圣洁的；不喝酒，也不喝烈酒，不吃肉；从不剃头，不用油抹身，也不沐浴。\par

6. 唯独他被允许进入圣所；因为他穿的不是羊毛，而是细麻衣服。他常常独自进入圣殿，经常双膝跪地，为百姓求赦免，以致他的膝盖因不断屈膝敬拜天主、为百姓求赦免而变得像骆驼的膝盖一样坚硬。\par

7. 因他极其正义，被称为义者，又称\OriginalTerm{奥布利亚斯}{Oblias}，这在希腊文中意为\Quote{人民的壁垒}和\Quote{正义}，正如先知们论及他时所宣告的。\par

8. 当时在百姓中有七个派别，我在\WorkTitle{回忆录}中已提及它们。其中有些人问他：\Quote{耶稣的门是什么？}他回答说，耶稣就是救主。\par

9. 因这些话，有些人相信耶稣就是基督。但上述那些派别既不信复活，也不信那位要按各人的行为施行报应者的来临。但凡是相信的人，都是因\Person{雅各伯}而信的。\par

10. 因此，当许多统治者也都相信时，在犹太人、经师和法利塞人中就起了骚动，他们说有危险，恐怕全体百姓都要把耶稣当作基督了。于是他们一同来到\Person{雅各伯}那里，说：\Quote{我们恳求你，制止百姓；因为他们关于耶稣走迷了路，以为他是基督。我们恳求你，要说服所有来参加逾越节庆节的人关于耶稣的事；因为我们都信赖你。因为我们和全体百姓都为你作证，你是公义的，不徇情面。}\BibleRef{玛 22:16}\par

11. \Quote{因此，你要说服群众，不要被耶稣误导。因为全体百姓，包括我们所有人，都信赖你。所以你要站在圣殿的顶上，好从那高处使人清楚看见你，你的话也容易让全体百姓听见。因为所有的支派，连同外邦人，都因逾越节聚集在一起。}\par

12. 于是上述的经师和法利塞人把\Person{雅各伯}放在圣殿顶上，向他喊叫说：\Quote{你这义人，我们大家都应当信赖你，因为百姓被钉在十字架上的耶稣误导了，请告诉我们，什么是耶稣的门？}\par

13. 他大声回答说：\Quote{你们为什么问我关于耶稣、人子的事？他本人坐在天上大能者的右边，并即将乘着天上的云彩降来。}\par

14. 当许多人完全信服，因\Person{雅各伯}的见证而自豪，并说\Quote{贺三纳归于达味之子}时，这些经师和法利塞人又彼此说：\Quote{我们不该为耶稣提供这样的见证。让我们上去把他扔下去，好使他们害怕而不敢相信他。}\par

15. 他们喊叫说：\Quote{哦！哦！义人也在谬误中。}他们便应验了依撒意亚书上所写的：\BibleQuote{让我们除掉义人，因为他使我们麻烦；因此他们必吃自己行为的果实。}\par

16. 于是他们上去把义人扔下来，彼此说：\Quote{让我们砸死义人\Person{雅各伯}。}他们就开始用石头砸他，因为他未被摔死；但他转身跪下说：\Quote{我恳求你，主、天主、我们的父，宽恕他们，因为他们不知道自己在做什么。}\BibleRef{路 23:34}\par

17. 当他们这样砸他时，有一个勒加布的儿子、勒加布人的司祭（先知耶肋米亚曾提及他们）喊道：\Quote{住手！你们在做什么？义人在为你们祈祷。}\par

18. 其中有一个人，是漂布工，拿起他用来捶打衣服的棍棒，击打义人的头。他就这样殉道了。他们把他葬在圣殿旁边的原地，他的墓碑至今仍在圣殿旁。他成了向犹太人和希腊人作证的忠信证人，证明耶稣就是基督。随即，\Person{韦斯巴芗}便围困了他们。\par

19. 这些事由\Person{赫格西普}详尽叙述，他与\Person{克莱孟}看法一致。\Person{雅各伯}是如此卓越的人，因他的正义而闻名于众人，以至于连犹太人中最有见识的人都认为，这就是\Place{耶路撒冷}被围困的原因——这围困发生在他殉道之后，并非由于其他原因，而是因为他们对他所行的胆大妄为之事。\par

20. 至少\Person{约瑟夫}毫不迟疑地在著作中为此作证，他说：\Quote{这些事发生在犹太人身上，是为了报复义人\Person{雅各伯}，他是那称为基督的耶稣的兄弟。因为犹太人杀了他，尽管他是最正义的人。}\par

21. 同一位作者也在其\WorkTitle{犹太古史}第二十卷中以如下文字记载了他的死：但皇帝得知\Person{斐斯托斯}去世的消息后，便派\Person{阿尔比努斯}作犹太省总督。但年轻的\Person{亚纳奴斯}，正如我们已说过的，获得了大司祭职，他性情极其狂妄鲁莽。此外，他属于撒杜塞人的派别，正如我们已指出的，撒杜塞人在执行审判上是所有犹太人中最残忍的。\par

22. 因此，\Person{亚纳奴斯}既是这种品性，又以为因\Person{斐斯托斯}已死、\Person{阿尔比努斯}尚在途中而有了可乘之机，便召集公议会，将那个称为基督的耶稣的兄弟——名叫\Person{雅各伯}——连同一些其他人带到他们面前，指控他们违法，并判处他们用石头砸死。\par

23. 但城中那些似乎最温和、精通法律的人对此非常愤怒，并秘密派人去见王，请求他命令\Person{亚纳奴斯}停止这种行为。因为他连这第一次也没有做对。他们中有些人也去迎接从\Place{亚历山大里亚}来的\Person{阿尔比努斯}，提醒他，\Person{亚纳奴斯}在未通知他的情况下召集公议会是不合法的。\par

24. \Person{阿尔比努斯}听了他们的陈述，愤怒地写信给\Person{阿纳努斯}，威胁要惩罚他。结果，\Person{阿格黎帕}王剥夺了他担任了三个月的大司祭职，并任命了\Person{达乃乌}的儿子\Person{耶稣}。\par

25. 这些事是关于\Person{雅各伯}的记载，他被认为是第一封所谓的公函的作者。但值得注意的是，这封信的真伪有争议；至少古代的作者很少有人提到它，正如名为\Person{犹达}的书信一样，它也是七封所谓的公函之一。然而我们知道，这些书信与其他书信一起，已经在许多教会中被公开诵读。\par

% structure: p0187 source=h2 target=subsection reason=relative-heading-rank; base=h2

\subsection{二十四、\Person{安尼安}——\Person{马尔谷}之后\Place{亚历山大}教会的第一位主教}

1. 当\Person{尼禄}在位第八年时，\Person{安尼安}继圣史\Person{马尔谷}之后管理\Place{亚历山大}教区。\par

% structure: p0189 source=h2 target=subsection reason=relative-heading-rank; base=h2

\subsection{二十五、\Person{尼禄}迫害下，\Person{保禄}和\Person{伯多禄}在\Place{罗马}为信仰殉道而受荣}

1. 当\Person{尼禄}的统治已经稳固时，他开始沉溺于邪恶的追求，甚至武装自己对抗宇宙的天主的宗教。\par

2. 描述他堕落的程度不属于本著作的计划。因为确实有许多人用极其准确的叙述记录了这段历史，每个人都可以随意从中了解这人异常疯狂行为的粗鄙。在这种疯狂的影响下，他在毫无理由地毁灭了无数人之后，陷入了如此血腥的罪恶，以致于连最亲近的亲戚和最亲爱的朋友也不放过，而是用各种死法杀害了他的母亲、兄弟、妻子以及他自己家族中的许多其他人，如同对待私人和公共敌人一样。\par

3. 但尽管有所有这些罪状，他的罪行录中还缺少这一条：他是第一个表现出自己是神圣宗教之敌的皇帝。\par

4. 罗马的\Person{戴尔都良}同样为此作证。他写道：\Quote{查看你们的记录。你们会发现\Person{尼禄}是第一个迫害这一教义的人，特别是在他征服了整个东方之后，在\Place{罗马}对所有人大施暴虐。我们以有这样一个人作为我们惩罚的首领而自豪。因为凡认识他的人都明白，\Person{尼禄}所谴责的绝不是某种伟大的卓越之物。}\par

5. 这样他公开宣布自己是天主首要敌人中的第一个，于是被引导去屠杀宗徒。因此，据记载，\Person{保禄}在\Place{罗马}被斩首，\Person{伯多禄}也在\Person{尼禄}手下被钉十字架。这一关于\Person{伯多禄}和\Person{保禄}的记载，因他们的名字至今仍保存在该地的墓园中而得到证实。\par

6. 同样，教会成员\Person{卡尤斯}也证实了这一点，他是在罗马主教\Person{则斐琳}时期兴起的。他在与弗吕家异端领袖\Person{普罗克鲁斯}的公开辩论中，就上述宗徒的圣骸所安放的地点发言如下：\par

7. \Quote{但我能展示宗徒的凯旋碑。因为如果你去\Place{梵蒂冈}或\Place{奥斯蒂亚路}，你会找到那些奠定这教会基础的人的凯旋碑。}\par

8. \Place{格林多}主教\Person{狄约尼削}在致\Place{罗马}人的书信中，用以下的话说明他们两人同时殉道：\Quote{你们就这样通过这样的劝告，将\Person{伯多禄}和\Person{保禄}在\Place{罗马}和\Place{格林多}的栽种联系在一起。因为他们两人都在我们\Place{格林多}栽种并教导了我们。他们同样在\Place{意大利}一同教导，并同时殉道。}我引用这些事，是为了使历史的真实性更加确证。\par

% structure: p0198 source=h2 target=subsection reason=relative-heading-rank; base=h2

\subsection{二十六、犹太人受无数灾祸折磨，开始与罗马人的最后战争}

1. \Person{若瑟夫}在叙述了降临整个犹太民族的灾祸的许多事情之后，又记录了其他许多情况，其中提到许多最尊贵的犹太人在\Place{耶路撒冷}被\Person{弗洛鲁斯}鞭打并钉十字架。当时是\Person{尼禄}在位第十二年，战争开始点燃，他正担任\Place{犹太}行省总督。\par

2. \Person{若瑟夫}说，那时由于犹太人的叛乱，整个\Place{叙利亚}爆发了可怕的骚动，各地犹太人都被城市居民毫不留情地像敌人一样消灭，\Quote{以致人们可以看到城市里堆满了未埋葬的尸体，老人的尸骨与婴儿的尸骨散落一地，妇女甚至没有遮盖裸体的衣物，整个省份充满了难以形容的灾难，而对那些威胁之事的恐惧大于他们任何地方所受的痛苦本身。}这就是\Person{若瑟夫}的记述；当时犹太人的状况就是这样。\par

\SourceRecord{Translated by Arthur Cushman McGiffert.；Nicene and Post-Nicene Fathers, Second Series；Vol. 1.；Edited by Philip Schaff and Henry Wace.；Buffalo, NY: Christian Literature Publishing Co.,；1890.；<http://www.newadvent.org/fathers/250102.htm>.}

% structure: p0001 source=h1 target=section reason=page-title

\section{教会史（第三卷）}

% structure: p0002 source=h2 target=subsection reason=relative-heading-rank; base=h2

\subsection{宗徒在普世传扬基督的地区}

1. 犹太人的情况便是如此。与此同时，我们救主的神圣宗徒和门徒分散到了世界各地。\Place{帕提亚}分配给了\Person{多默}作为他的工场，\Place{斯基提亚}给了\Person{安德肋}，\Place{亚细亚}给了\Person{若望}，他在那里住了一段时间后，死于\Place{厄弗所}。\par

2. \Person{伯多禄}似乎曾在\Place{本都}、\Place{迦拉达}、\Place{比提尼亚}、\Place{卡帕多细亚}和\Place{亚细亚}向散居的犹太人传道。最后，他到了\Place{罗马}，被头朝下钉死在十字架上，因为他曾请求这样受苦。关于\Person{保禄}，从\Place{耶路撒冷}直到\Place{依里黎坤}传扬基督的福音，后来在\Place{罗马}在\Person{尼禄}治下殉道，我们有需要说什么吗？这些事实由\Person{奥利振}在他的\WorkTitle{创世纪注释}第三卷中叙述。\par

% structure: p0005 source=h2 target=subsection reason=relative-heading-rank; base=h2

\subsection{圣伯多禄在罗马的首位继任者}

1. \Person{保禄}和\Person{伯多禄}殉道后，\Person{里诺}首先获得了罗马教会的主教职。\Person{保禄}在从\Place{罗马}写信给\Person{弟茂德}时，在书信结尾的问候中提到过他。\par

% structure: p0007 source=h2 target=subsection reason=relative-heading-rank; base=h2

\subsection{宗徒的书信}

1. \Person{伯多禄}的一封书信，即被称为第一封的，公认是真作。古代的长老们在自己的著作中自由地使用它，视之为无可争议的作品。但我们得知，他现存第二封书信不属于正典；然而，由于它对许多人显得有益，它已与其他圣经一同被使用。\par

2. 然而，所谓的\WorkTitle{伯多禄行传}、以他命名的福音、\WorkTitle{宣讲}和\WorkTitle{默示录}，我们知道并未被普遍接受，因为没有一位教会作者，无论古今，使用过从它们引用的证据。\par

3. 但在我的历史进程中，我将小心地展示，除了正式的传承外，哪些教会作者不时使用了有争议的著作，以及他们对于正典和公认的著作，以及对于不属于这一类的著作说了什么。\par

4. 这些就是那些以\Person{伯多禄}之名流传的著作，其中我只知道一封是真的，并被古代长老所承认。\par

5. \Person{保禄}的十四封书信是众所周知且无可争议的。确实，不应忽略一个事实，即有些人拒绝了\BibleBook{}{希伯来书}，说它被罗马教会争议，理由是它不是\Person{保禄}所写。但关于这封书信，在我们之前的人所说过的话，我将在适当的地方引用。至于所谓的\WorkTitle{保禄行传}，我未发现它属于无可争议的著作。\par

6. 但正如同一位宗徒，在\BibleBook{}{罗马书}结尾的问候中，除了其他人之外，也提到了\Person{赫尔马}，那本称为\WorkTitle{牧者}的书就是归于他的，应该注意到，这本书也被一些人争议，因此不能归入公认的书目；而在另一些人看来，它被认为是相当不可或缺的，尤其是对那些需要在信仰基本原理上受教导的人。因此，我们知道，它曾在教会中公开诵读，并且我发现一些最古老的作者使用过它。\par

7. 这足以表明哪些神圣著作是无可争议的，以及哪些不是普遍公认的。\par

% structure: p0015 source=h2 target=subsection reason=relative-heading-rank; base=h2

\subsection{宗徒的首批继任者}

1. \Person{保禄}向外邦人传道，并从\Place{耶路撒冷}直到\Place{依里黎坤}四围建立了教会的基础，这从他自己的话\Quote{从耶路撒冷直到依里黎坤四处}（\BibleRef{罗 15:19}）以及\Person{路加}在\BibleBook{}{宗徒大事录}中的叙述可以清楚看出。\par

2. \Person{伯多禄}在多少省份传扬了基督，并向受割礼的人教导了新盟约的教义，这从他已提到的无可争议的书信中的话可以清楚看出，在信中他写给散居在\Place{本都}、\Place{迦拉达}、\Place{卡帕多细亚}、\Place{亚细亚}和\Place{比提尼亚}的希伯来人（\BibleRef{伯前 1:1}）。\par

3. 但那些成为宗徒真正而热忱的追随者，并被认为配得牧养他们所建立的教会的人，其数目和名字，除了\Person{保禄}著作中提到的那几位外，并不容易说出。\par

4. 因为他有数不清的同工，或如他所说的\Quote{战友}，他们中大多数人被他以不朽的纪念尊崇，因为他在自己的书信中为他们作了持久的见证。\par

5. \Person{路加}也在\BibleBook{}{宗徒大事录}中谈到他的朋友们，并提了他们的名字。\par

6. 据记载，\Person{弟茂德}是第一个获得\Place{厄弗所}教区主教职的人，\Person{弟铎}则获得了\Place{克里特}各教会的主教职。\par

7. 但\Person{路加}，出生于\Place{安提约基雅}，职业是医生，与\Person{保禄}特别亲密，并熟悉其余宗徒，他留下了两部受默感之书，作为他从他们那里学到的灵性医治技艺的证明。其中一部是\BibleBook{路加}{福音}，他证明自己写作时所依据的，是那些从起初就亲眼看见并为圣言服务的人所交付给他的，并且他声称自己从最初就准确地追溯了这一切（\BibleRef{路 1:2}）。另一部是\BibleBook{}{宗徒大事录}，他不是根据别人的叙述，而是根据自己亲眼所见写成的。\par

8. 他们说，\Person{保禄}每当好像谈论自己的某种福音，而使用\Quote{按我的福音}这些字眼时，意思是指\BibleBook{路加}{福音}。\par

9. 至于他其余的同工，\Person{保禄}见证\Person{克勒斯森}被派往\Place{高卢}；但\Person{里诺}，他在\BibleBook{弟茂德}{后书}中（\BibleRef{弟后 4:21}）提到他是他在\Place{罗马}的同伴，却是\Person{伯多禄}在那里的主教职的继任者，如上所述。\par

10. \Person{克肋孟}，被任命为罗马教会的第三位主教，也正如\Person{保禄}所见证的，是他的同工和战友。\par

11. 此外，那位名叫\Person{狄约尼削}的亚略巴古人，他在\Person{保禄}于\Place{亚略巴古}向雅典人讲道后第一个相信（如\Person{路加}在\BibleBook{}{宗徒大事录}中所记载），被另一位\Person{狄约尼削}（一位古代作家、\Place{格林多}教区的牧者）提及为\Place{雅典}教会的第一位主教。\par

12. 但与宗徒传承相关的事件，我们将在适当的时候叙述。现在，让我们继续我们历史的进程。\par

% structure: p0028 source=h2 target=subsection reason=relative-heading-rank; base=h2

\subsection{基督之后犹太人最后的围城}

1. 尼禄统治十三年，伽尔巴和奥托统治一年零六个月之后，韦斯帕芗因在对抗犹太人的战役中声名显赫，在犹太地被拥立为君主，并获得当地军队授予的皇帝称号。他随即启程前往罗马，将对抗犹太人的战事托付给他的儿子提图斯。\par

2. 因为犹太人在我救主升天之后，除了对祂所犯的罪行之外，还竭力策划尽可能多的阴谋陷害祂的宗徒。首先，斯德望被他们用石头砸死；此后，载伯德的儿子、若望的兄弟雅各伯被斩首；最后，在我救主升天之后首先获得耶路撒冷主教席位的那位雅各伯，也以已描述的方式死去。但其余的宗徒，尽管不断被密谋杀害并被逐出犹太地，却前往万邦宣讲福音，依靠基督的力量，因为祂曾对他们说：\Quote{你们要去，因我的名使万民成为门徒。}\par

3. 但耶路撒冷教会的人，在战争之前，曾通过启示（赐予那里有德之人）得到命令，要离开那城，住在培勒雅一个名叫培拉的镇上。当那些信基督的人从耶路撒冷来到那里之后，仿佛犹太人的王城和整个犹太地完全失去了圣洁的人，天主的审判终于临到那些对基督及其宗徒犯下如此暴行的人，彻底毁灭了那代不敬虔之人。\par

4. 但那时各地降在那民族身上的灾祸数量；犹太居民特别遭受的极度不幸；成千上万的男人、女人和儿童死于刀剑、饥荒和其他无数形式的死亡——所有这一切，以及针对犹太各城所进行的大规模围攻，那些逃往耶路撒冷（如同绝对安全的城市）的人所受的极大痛苦，最后整个战争的进程及其具体细节，以及先知所预言的\OriginalTerm{可憎的荒凉}{abomination of desolation}如何在长久以来备受尊崇的天主圣殿中站立（\BibleRef{达 9:27}），这座圣殿如今正等待被火彻底毁灭——凡愿意了解这些事的人，都可以在若瑟夫所著的\WorkTitle{历史}中找到精确的描述。\par

5. 但有必要指出，这位作者记载，在逾越节期间从犹太各地聚集的人数达三百万之多，被关在耶路撒冷里，\Quote{如同在监牢中，}这是他的原话。\par

6. 因为理当如此：在他们对所有人的救主和恩主、天主的基督施加苦难的那些日子里，他们被关闭\Quote{如同在监牢中，}遭受天主公义的惩罚。\par

7. 但略过他们因刀剑和其他方式所受的个别灾祸，我认为必须只叙述饥荒所带来的不幸，以便读这作品的人能多少知道，天主并未迟延为他们反对天主的基督的邪恶而施行报复。\par

% structure: p0036 source=h2 target=subsection reason=relative-heading-rank; base=h2

\subsection{第六章 压迫他们的饥荒}

1. 我们再次拿起若瑟夫的\WorkTitle{历史}第五卷，重温当时发生的悲剧事件。\par

2. \Quote{至于富有的人，}他说道，留下来同样危险。暴徒的疯狂随饥荒加剧，两种苦难日复一日地愈演愈烈。\par

3. 哪里也看不到食物；但暴徒闯入人家，彻底搜查，每当发现可吃的东西，他们就折磨主人，理由是主人曾否认有任何食物；如果什么也没发现，他们就施以酷刑，理由是主人藏得更隐蔽。\par

4. 这些可怜人有无食物，证据就在他们身上。那些仍然身体康健的，他们便假定食物充足；而那些已经消瘦的，他们则放过，因为杀害那些快要饿死的人似乎荒谬。\par

5. 的确，许多人偷偷卖掉财物，换来一量器小麦（如果属于较富有阶级）或大麦（如果较贫穷）。然后关在屋内最深处，有些因极需而生吃谷粒，另一些则根据必要和恐惧所迫而烘烤。\par

6. 哪里也未见摆桌，而是从火中夺出尚未煮熟的食物，撕成碎片。可悲的饭食，可怜的景象——强者充裕，弱者哀叹。\par

7. 的确，万恶之中饥荒最甚，它所消灭的，莫过于羞耻。因为那种在别的情况下值得尊敬的事，在饥荒中却遭藐视。于是女人从丈夫和孩子口中夺食，从父亲口中夺食，而最可悲的是，母亲从婴儿口中夺食。当他们的至亲在自己怀中日渐消瘦时，他们竟不羞于夺走支撑生命的最后几滴。\par

8. 甚至在他们如此吃的时候，也未能隐藏。暴徒到处出现，连这些食物也要抢走。因为只要他们看见一扇门关着，就认为里面的人正在吃东西。于是立刻破门而入，冲进去夺走他们正在吃的东西，几乎是从喉咙里抢出来。\par

9. 那些紧紧抓住食物的老人遭到殴打，而妇女若将食物藏在手中，她们就会为此被扯头发。无论是白发苍苍的老者还是幼小的婴儿，都没有得到怜悯；他们抱起那些紧抓着食物残渣的婴孩，把他们摔在地上。但对于那些预先察觉他们闯入、并吞下他们将要抢夺之物的，他们更加残忍，仿佛他们受了这些人的亏欠似的。\par

10. 他们还发明了最可怕的刑讯方式来找出食物，用苦菜塞住穷人的排泄道，并用尖棍刺穿他们的臀部。人们遭受了甚至听来都可怕的事情，只为了迫使他们承认藏有哪怕一个面包，或为了让他们交出一把藏起来的大麦。但施刑者自己并未挨饿。\par

11. 如果他们是迫于必需才这样做，或许还显得不那么野蛮；但他们这样做是为了发泄疯狂，并为以后多日提供食物。\par

12. 每当有人夜里偷偷出城，直到罗马人的前哨去采集野草和青草时，他们就迎上去；当那人以为已逃脱敌人时，他们抢走了他带来的东西，即使那人常常恳求他们，并呼求天主最可敬畏的名，发誓把冒生命危险得来的东西分给他们一部分，他们什么也不还给他。事实上，被抢的人若不被杀，已是幸运。\par

13. 约瑟夫在叙述其他事情后，又加上以下内容：出城的可能性断绝后，犹太人得救的一切希望都被切断了。饥荒加剧，逐户逐家地吞噬人们。房间里堆满了妇女和儿童的尸体，街上堆满了老人的尸体。\par

14. 儿童和青年，因饥荒而浮肿，像影子一样在市场游荡，在死亡之痛袭来时倒下。病人无力埋葬自己的亲人，有力气的人也因死者众多和自身命运的不确定而犹豫。许多人其实在埋葬他人时死去，许多人在死亡来临之前就已躺进自己的坟墓。\par

15. 在这些灾难中，没有哭泣，没有哀悼；饥荒窒息了天然的亲情。那些垂死的人用干涸的眼睛看着那些先他们安息的人。深沉寂静和死亡笼罩的夜色环绕着城市。\par

16. 但强盗比这些苦难更可怕；他们闯入那些已成坟墓的房屋，抢劫死者，剥掉他们身上的覆盖物，大笑着离去。他们用死尸试剑锋，有些还活着的倒在地上的人，他们刺穿以试验武器。而那些祈求他们用右手和剑杀死他们的人，他们轻蔑地丢下，让饥荒去消灭。这些人每一个都眼望圣殿而死；他们留下了叛乱者活着。\par

17. 起初，他们命令用公费埋葬死者，因为他们受不了臭味。但后来，当他们无法这样做时，就把尸体从城墙扔进壕沟。\par

18. 当提图斯巡视，看到壕沟里堆满尸体，浓烈的血水从腐烂的尸体中渗出，他大声呻吟，举起双手，呼求天主作证这不是他的作为。\par

19. 在谈及一些别的事情之后，约瑟夫继续写道：\Quote{我不能犹豫，要宣布我的感受迫使我说出的东西。我想，如果罗马人再迟些来惩罚这些罪犯，这座城市或许早已被深渊吞没，或被洪水淹没，或被像毁灭索多玛那样的雷电击中。因为它产生了一代人，比那些遭受此类惩罚的人更不虔敬。确实，正是由于他们的疯狂，整个民族才走向灭亡。}\par

20. 在第六卷中，他写道：城中死于饥荒的人数不计其数，他们所遭受的苦难难以言表。因为哪怕一点食物的影子出现在任何房屋中，就会引起争斗，最亲密的朋友彼此肉搏，互相抢夺最可悲的维生之物。\par

21. 他们甚至不相信垂死的人没有食物；强盗会搜查即将咽气的人，以免有人假装死亡而将食物藏在怀中。他们因缺乏食物而张着嘴，像疯狗一样跌跌撞撞，像醉汉一样敲打门，无力地在一小时内两次或三次冲进同一房屋。\par

22. 迫不得已，他们吃任何能找到的东西，收集并吞下连最肮脏的畜生都不该吃的东西。最后，他们甚至不放过腰带和鞋子，剥下盾牌上的皮吃掉。有的人甚至用旧干草作为食物，有的人收集秸秆，以最小的重量卖四个阿提卡德拉克马。\par

23. 但我何必提及饥荒中对无生命之物的无耻行径呢？因为我要叙述一件事，既未被希腊人也未被外邦人记载过；它叙述起来可怖，听起来难以置信。我本应乐意略过这一灾难，以免被后人视为编造离奇故事的人，如果我当时没有无数的见证者的话。此外，如果我隐瞒她所受的苦难，就是对我可怜的祖国服务不力。\par

24. 有一个名叫玛利亚的妇人，住在约旦河对岸，她的父亲是厄肋阿匝尔，来自巴特佐尔村（意即\Emph{牛膝草之家}）。她因家世和财富而显赫，与其余民众一起逃到耶路撒冷，在围城期间被困在那里。\par

25. 暴君们抢走了她从\Place{培勒雅}带入城中的其余财产。她的残余物品以及任何可见的食物，卫兵们每日冲进来夺走。这使那妇人极其愤怒，她频频责骂咒诅，激起了那些贪婪恶棍对她的怒气。\par

26. 但没有人因愤怒或怜悯而杀死她；她已厌倦为别人寻找食物。搜寻也到处变得困难，饥荒刺痛她的五脏六腑，而怨恨比饥荒更猛烈地燃烧。于是，她以愤怒和必需为参谋，做出了一件最违背天性的事。\par

27. 她抓住正在吃奶的男孩，说道：\Quote{噢，不幸的孩子，在战争、饥荒、叛乱中，我为何要保全你？即使罗马人容许我们活下去，我们也将成为他们的奴隶。但就连奴隶制也被饥荒抢先，而暴徒比两者都残忍。来，成为我的食物，成为这些暴徒的狂怒，成为世人的笑柄，因为只缺这一件就成全了犹太人的灾祸。}\par

28. 说完这话，她杀了儿子；烤熟后，自己吃了一半，将剩余的部分覆盖保存。不久暴徒们来到现场，闻到那罪恶的气味，威胁要立刻杀死她，除非她交出所预备的东西。她回答说为他们留了一份上好的肉，并揭开了孩子的残骸。\par

29. 他们立即惊恐万状，呆立当场。但她说：\Quote{这是我自己的儿子，这行为是我的。吃吧，因为我也吃了。不要比女人更慈悲，也不要比母亲更怜恤。但如果你们过于虔诚，回避我的祭品，我已经吃了；剩下的也留给我吧。}\par

30. 听了这话，这些人战栗地离去，这一次他们感到恐惧；然而他们还是勉强把那食物留给了母亲。顷刻间，全城充满了这可怕的罪行，每个人都仿佛亲眼目睹那可怕的行为，战栗如亲手所为。\par

31. 那些忍受饥荒的人如今渴望死亡；那些在听闻和目睹如此悲惨之前就已死去的人是有福的。\par

32. 这就是犹太人因对天主的基督所行的邪恶和不虔敬而得的报应。\par

% structure: p0069 source=h2 target=subsection reason=relative-heading-rank; base=h2

\subsection{第七章 基督的预言。}

1. 在这些叙述之后，宜补充我们救主的真实预言，祂曾预言了这些事件。\par

2. 祂的话如下：\BibleQuote{那些日子怀孕的和哺乳的有祸了！你们要祈求逃走不在冬天，也不在安息日。因为那时必有大灾难，从世界的开始直到如今，未曾有过，将来也绝不会再有。}\par

3. 史学家计算被杀的总数，说有一百一十万人死于饥荒和刀剑，其余的暴徒和强盗在城被攻取后互相出卖而被杀。最英俊的青年被保留用于凯旋式。其余群众中，十七岁以上的被送到埃及的工场囚禁劳动，还有更多的人被分散到各省，在剧场中死于刀剑和野兽。十七岁以下的被掳去卖为奴隶，仅这些人的数目就达到九万。\par

4. 这些事发生在\Person{韦斯巴芗}统治的第二年，符合我们的主和救主耶稣基督的预言，祂以天主的能力预见了它们，仿佛已然发生，并为此哭泣哀恸，正如圣史所记，他们记载了祂所说的话，祂仿佛对耶路撒冷说话时说：\par

5. \BibleQuote{巴不得你在这日子知道关系你平安的事！无奈这事现在是隐藏的，叫你的眼看不出来。因为日子将到，你的仇敌必筑起土垒，围绕你，四面困住你，并要扫灭你和你里的儿女。}\par

6. 然后，祂仿佛论及百姓说：\BibleQuote{那时必有大灾难在这地方，也有震怒临到这百姓。他们要倒在刀下，又被掳到各国去。耶路撒冷要被外邦人践踏，直到外邦人的日期满了。} 又说：\BibleQuote{你们看见耶路撒冷被军队围困，就知道它荒凉的日子近了。}\par

7. 若有人将我们救主的话与史学家关于整个战争的记载相比较，怎能不惊叹，并承认我们救主的预知和预言真是神圣而奇妙非凡。\par

8. 因此，关于救主受难后，以及犹太群众要求释放强盗和杀人犯、却请求将生命之主从他们中间夺走之后，整个犹太民族所遭遇的灾祸，无需在史学家的记载之外添加任何内容。\par

9. 但或许也宜提及那些事件，它们彰显了那至善上智的恩慈——在他们对基督犯罪之后，仍将毁灭延迟了整整四十年；在此期间，许多宗徒和门徒，以及\Person{雅各伯}本人——当地的第一位主教，被称为主的兄弟——仍活着，并住在耶路撒冷，成为那里最坚固的堡垒。上智眷顾仍长久容忍他们，要看他们是否因所行而悔改，以获得宽恕与救恩；除了如此长久的容忍，上智眷顾还提供了奇妙的征兆，显示若他们不悔改即将发生的事。\par

10. 既然这些事被前述史家认为值得提及，我们也最好把它们叙述出来，以裨益本书的读者。\par

% structure: p0080 source=h2 target=subsection reason=relative-heading-rank; base=h2

\subsection{第八章  战争前夕的预兆}

1. 那么，且取此作者的作品，阅读他在\WorkTitle{历史}第六卷中所记载的。他的话如下：\Quote{就这样，当时那些可怜的人被骗子与假先知所蒙骗；但他们并未留意也不相信那预告即将来临的荒凉的异象和征兆。相反，他们如同被闪电击中，如同既无眼睛又无理智，轻看了天主的宣告。}\par

2. 有一次，一颗形状如剑的星悬在城上，还有一颗彗星持续了整整一年；又在叛乱之前，在导致战争的那些骚动之前，当民众聚集庆祝无酵节时，在桑西克斯月第八日的夜间第九时辰，祭台和圣殿周围照耀着如此大的光，以至如同白昼；这光持续了半小时。这在无知者看来是个好兆头，但被圣书士解释为预示那些很快发生的事件。\par

3. 在同一节期，一头由大司祭牵去献祭的母牛，在圣殿中间生了一只羔羊。\par

4. 内殿的东门是铜制的，极其沉重，傍晚需二十个人费力关闭，它靠在铁包的木梁上，门闩深深插在地里，但在夜间第六时辰，它自动打开了。\par

5. 节期过后不久，在阿特米西翁月第二十一日，看见了一个奇妙的异象，令人难以置信。若非由目睹者所述，且随后降临的灾祸与这样的征兆相称，这奇事或许会被视为虚构。因为在日落之前，整个区域半空中可见战车和武装部队在云中盘旋，环绕着城市。\par

6. 在称为五旬节的节期，当司祭们在夜间照例进入圣殿举行礼仪时，他们说起初感到动静和响声，随后有如一大群人的声音说：\Quote{我们离开这里吧。}\par

7. 但接下来的事更加可怕：在战争前四年，当城市特别繁荣和平时，有一个名叫耶稣、阿纳尼雅的儿子、平凡的乡下人，来到节期（其时所有人都照例在圣殿搭棚以光荣天主），突然喊道：\Quote{有声音来自东方，有声音来自西方，有声音来自四方，有声音反对耶路撒冷和圣殿，有声音反对新郎和新娘，有声音反对全体百姓。}他日夜穿过所有街巷这样喊着。\par

8. 但一些显要的市民，因这不祥的喊声而恼怒，抓住那人，用许多鞭子打他。但他没有为自己说一句话，也没有对在场的人说什么特别的话，继续喊着同样的话。\par

9. 官长们认为（确实如此）那人被一种更高的力量所驱动，便把他带到罗马总督面前。然后，尽管他被鞭打到骨头，但他既未求情也未流泪，反而以最悲惨的声调，每一鞭都回答说：\Quote{祸哉，祸哉，耶路撒冷。}\par

10. 同一史家还记载了另一件比这更奇妙的事。他说在他们的圣书中发现了一个神谕，宣告那时将有一个人从他们的国家出去统治世界。他本人认为这应验在韦斯巴芗身上。\par

11. 但韦斯巴芗并未统治全世界，而只统治了罗马人统治的那部分。更合适的是把它应用于基督；圣父曾对他说：\BibleQuote{你向我请求，我就将外邦人赐你作产业，将地极赐你作田产。}实际上，就在那时，他圣宗徒的声音\BibleQuote{传遍天下，他们的话语传到地极。}\par

% structure: p0092 source=h2 target=subsection reason=relative-heading-rank; base=h2

\subsection{第九章  若瑟夫及其遗留的著作}

1. 在此之后，我们应当知道一些关于若瑟夫的出身和家族的事，他对本历史贡献甚多。他本人这样告诉我们：\Quote{若瑟夫，玛塔提雅的儿子，是耶路撒冷的一位司祭，他起初曾与罗马人作战，后来被迫目睹了随后发生的事。}\par

2. 他是当时所有犹太人中最著名的，不仅在本族中，也在罗马人中，以致在罗马为他立了雕像，他的作品被认为值得收藏在图书馆中。\par

3. 他写了整部\WorkTitle{犹太古史}共二十卷，以及他当时发生的与罗马人战争的\WorkTitle{历史}共七卷。他本人证实后者不仅用希腊文写成，而且由他本人译成本族语言。他在此事上值得信任，因为他在其他事上是诚实的。\par

4. 他的另外两本书也存世，值得一读。它们论述犹太人的古老，其中他答复了文法家阿皮翁，此人当时写了反对犹太人的论文，也答复了其他试图诋毁犹太人祖传制度的人。\par

5. 在第一本书中，他给出了所谓的旧约正典书卷的数目。他显然根据古传，指出希伯来人中哪些书卷是毫无争议地接受的。他的话如下。\par

% structure: p0098 source=h2 target=subsection reason=relative-heading-rank; base=h2

\subsection{第十章  若瑟夫提及神圣书卷的方式}

\Quote{1. 因此，我们并没有许多互相矛盾冲突的书；我们只有二十二卷，它们包含了全部时间的记录，并公正地被视为神圣的。}\par

\Quote{2. 其中五卷是梅瑟的，包含法律和关于人类起源的传统，并把历史延续到他自己的死亡。这一时期约有三千年。}\par

3. 从梅瑟逝世到继薛西斯为波斯王的阿塔薛西斯逝世，继梅瑟之后的先知们将他们所处时代的历史写成了十三卷书。另外四卷书包含对天主的颂歌，以及规范人类生活的诫命。\par

4. 自阿塔薛西斯时代至今日，所有事件均有记载，但这些记载不值得像先前那些记载一样值得信任，因为在此期间没有准确的先知继承。\par

5. 我们对自己著作的珍视程度，从我们对待它们的方式可以清楚看出。因为尽管如此漫长的岁月已经过去，却无人敢在其中增添或删减；相反，所有犹太人从一出生就根深蒂固地视它们为天主的教训，并坚守它们，如有必要，甘愿为之舍命。\par

我以为在此处引述这位史学家的这些评论是有益的。\par

6. 同一位作者还写了一部颇有价值的作品\WorkTitle{论理性的至上}，有些人称之为\WorkTitle{玛加伯书}，因为它记载了那些为真道奋勇作战的希伯来人的斗争，正如称为\WorkTitle{玛加伯书}的书卷中所记述的。\par

7. 在\WorkTitle{犹太古史}第二十卷的末尾，约瑟夫斯本人表示，他曾计划写一部四卷本的著作，论及天主及其存有，按照犹太人的传统见解，以及论及法律，为何允许一些事而禁止另一些事。同一位作者还在自己的作品中提到了他自己所写的其他书籍。\par

8. 除了这些，还应引述他在\WorkTitle{犹太古史}结尾所写的话，以证实我们从他的叙述中得出的见证。在那里，他攻击提庇里亚的犹斯托，此人像他一样试图撰写一部当代史，理由是他没有如实记述。在对他提出诸多指控之后，他接着说了这样一些话：\par

9. 我确实不像你那样害怕我的著作，相反，当事件几乎还在人们眼前时，我就把我的书呈献给皇帝们本人。因为我知道我在记述中保持了真相，所以没有失望地得到了他们的认证。\par

10. 我还把我的历史书呈献给许多其他人，其中有些人曾参与过这场战争，例如阿格里帕王和他的几位亲属。\par

11. 因为提托皇帝非常希望人们对事件的了解只通过我的历史书传达，所以他亲手签署了这些书，并命令出版。阿格里帕王写了六十二封信，证明我的记述的真实性。约瑟夫斯附上了其中两封。但这些关于他的事就足够了。现在让我们继续我们的历史。\par

% structure: p0111 source=h2 target=subsection reason=relative-heading-rank; base=h2

\subsection{十一、西默盎在雅各伯之后治理耶路撒冷教会。}

1. 据说在雅各伯殉道以及随即发生的耶路撒冷被攻陷之后，当时仍然存活的宗徒和主的门徒们从各处聚集，连同那些按肉身与主有亲属关系的人（因为他们中的大多数也还活着），商议谁配继承雅各伯。\par

2. 他们全体一致同意，认定克罗帕的儿子西默盎（福音中也曾提及他）配得上那教区的主教职位。据说他是救主的堂兄弟。因为赫格西普斯记载说，克罗帕是若瑟的兄弟。\par

% structure: p0114 source=h2 target=subsection reason=relative-heading-rank; base=h2

\subsection{十二、苇斯巴芗下令搜寻达味的后裔。}

他还记述说，苇斯巴芗在攻陷耶路撒冷之后，下令搜寻所有属于达味家族的人，以免任何王族血统留在犹太人中间；因此，一场极其可怕的迫害再次降临在犹太人身上。\par

% structure: p0116 source=h2 target=subsection reason=relative-heading-rank; base=h2

\subsection{十三、阿纳克肋托，罗马第二位主教。}

苇斯巴芗在位十年后，其子提托继位。在他执政的第二年，担任罗马教会主教十二年的理诺将职位交给了阿纳克肋托。但提托在位两年零两个月后，由其弟多米提安继位。\par

% structure: p0118 source=h2 target=subsection reason=relative-heading-rank; base=h2

\subsection{十四、阿彼略，亚历山大里亚第二位主教。}

多米提安在位第四年，亚历山大里亚教区第一位主教阿尼雅诺在任职二十二年后去世，由第二位主教阿彼略继任。\par

% structure: p0120 source=h2 target=subsection reason=relative-heading-rank; base=h2

\subsection{十五、克肋孟，罗马第三位主教。}

同一位皇帝在位第十二年，阿纳克肋托在担任罗马教会主教十二年后，由克肋孟继任。宗徒在致斐理伯人书中告诉我们，这位克肋孟是他的同事。他的话如下：\Quote{与克肋孟以及其余与我一同工作的同伴，他们的名字已写在生命册上。}\par

% structure: p0122 source=h2 target=subsection reason=relative-heading-rank; base=h2

\subsection{十六、克肋孟书信。}

有一封克肋孟的书信流传至今，被公认为真作，篇幅颇长，价值非凡。他是以罗马教会的名义写给格林多教会的，当时格林多教会中发生了纷争。我们知道，这封书信无论在古代还是在我们时代，都在许多教会中公开使用过。至于当时格林多教会确实发生了纷争，赫格西普斯是一位可靠的见证人。\par

% structure: p0124 source=h2 target=subsection reason=relative-heading-rank; base=h2

\subsection{十七、多米提安下的迫害。}

多米提安对许多人残暴至极，在罗马不义地处死了不少出身高贵、声名显赫的人，又无缘无故地流放了许多其他杰出人士并没收其财产，最终在仇恨和敌视天主方面成了尼禄的继承人。事实上，他是第二个对我们发动迫害的人，尽管他的父亲苇斯巴芗没有采取任何对我们不利的行动。\par

% structure: p0126 source=h2 target=subsection reason=relative-heading-rank; base=h2

\subsection{十八、宗徒若望与默示录。}

1. 据说在这次迫害中，仍然活着的宗徒兼福音作者若望，因他为天主圣言作证，被判处居住在帕特摩岛上。\par

2. \Person{依勒内}在其著作《\WorkTitle{驳斥异端}》第五卷中，讨论《所谓\WorkTitle{若望默示录}》所载假基督之名的数目时，论及他说：\par

3. \Quote{倘若必须在这时代公开宣告他的名字，自当由那看见启示者宣布。因为这启示并非很久以前才被看见，而是在几乎我们这一代、\Person{多米提安}统治末期。}\par

4. 我们的信德道理在那时确实如此兴盛，以致那些远离我们宗教的作家也毫不迟疑地在他们的史书中提及当时发生的迫害与殉道。\par

5. 他们确实准确指出了时间。因为他们记载，\Person{多米提安}在位第十五年，\Person{弗拉维·克肋孟}之姊妹的女儿\Person{弗拉维雅·多米提拉}，与许多别的因对基督作证而被放逐至\Place{蓬提亚岛}的人一起被流放。\par

% structure: p0132 source=h2 target=subsection reason=relative-heading-rank; base=h2

\subsection{第19章　\Person{多米提安}命令杀害\Person{达味}的后裔}

当这同一个\Person{多米提安}命令杀害\Person{达味}的后裔时，古代传统说，有些异端者控告\Person{犹达}的后裔——据说\Person{犹达}是按肉身而论之救主的兄弟——说他们是\Person{达味}的后裔，并与基督本人有亲缘。\Person{赫格西普}记载了这些事实，措辞如下。\par

% structure: p0134 source=h2 target=subsection reason=relative-heading-rank; base=h2

\subsection{第20章　我们救主的亲属}

1. 主的家族中，尚存活着\Person{犹达}的孙辈。据说\Person{犹达}是主按肉身的兄弟。\par

2. 有人报告说他们属于\Person{达味}家族，于是被一位\OriginalTerm{招集官}{Evocatus}带到\Person{多米提安}皇帝面前。因为\Person{多米提安}惧怕基督的来临，如同\Person{黑落德}也曾惧怕。他便问他们是否为\Person{达味}后裔，他们承认了。然后他问他们有多少财产或拥有多少银钱。两人回答说，他们仅有九千\OriginalTerm{银币}{denarii}，每人各得一半。\par

4. 这财产并非银子，而是一块仅三十九英亩的土地，他们从中纳税，并靠自己的劳动维持生计。\par

5. 然后他们伸出手，展示他们身体的粗硬和因持续劳动而产生的茧子，作为自己劳作的证据。\par

6. 当被问及基督和他的王国——它是什么样的，在哪里，何时显现——他们回答说，那不是现世的或地上的王国，而是天上的、属天使的王国，将在世界末日，当他带着光荣来临审判生者死者，并按每个人的行为给予各人时显现。\par

7. \Person{多米提安}听了这话，并未判决他们，反而轻视他们为不足挂齿，释放了他们，并颁布一道谕令，停止了教会的迫害。\par

8. 但他们被释放后，因作证者并因是主的亲属，管理着各教会。和平建立后，他们一直活到\Person{图拉真}时代。这些事由\Person{赫格西普}记述。\par

9. \Person{戴尔都良}也曾以下列词语提及\Person{多米提安}：\Quote{\Person{多米提安}也有份于\Person{尼禄}的残忍，曾一次试图效仿其所为。但我想他还有几分理智，便很快停止，甚至召回了那些他曾放逐的人。}\par

10. 但\Person{多米提安}执政十五年后，\Person{内尔瓦}继位帝国，根据记载当时历史的作者，罗马元老院投票决定撤销\Person{多米提安}的荣誉，并让那些被不义放逐的人返回家园，恢复他们的财产。\par

11. 就在这时，宗徒\Person{若望}从他在岛上的放逐中归来，并根据古代基督徒传统，定居在\Place{厄弗所}。\par

% structure: p0145 source=h2 target=subsection reason=relative-heading-rank; base=h2

\subsection{第21章　\Person{克尔东}成为\Place{亚历山大}教会第三任统治者}

1. \Person{内尔瓦}执政稍过一年后，由\Person{图拉真}继位。在他统治的第一年，管理了\Place{亚历山大}教会十三年的\Person{阿彼略}，由\Person{克尔东}继任。\par

2. 他是继第一任\Person{安尼安}之后，主持该教会的第三人。那时\Person{克肋孟}仍治理罗马教会，他也是继\Person{保禄}和\Person{伯多禄}之后，在那里担任主教职务的第三人。\par

3. \Person{利诺}是第一任，其后是\Person{阿内克雷图}。\par

% structure: p0149 source=h2 target=subsection reason=relative-heading-rank; base=h2

\subsection{第22章　\Person{依纳爵}，\Place{安提约基亚}第二任主教}

当时\Person{依纳爵}被誉为\Place{安提约基亚}第二任主教，\Person{厄沃狄}是第一任。\Person{西默盎}那时也是\Place{耶路撒冷}教会第二任统治者，我们救主的兄弟是第一任。\par

% structure: p0151 source=h2 target=subsection reason=relative-heading-rank; base=h2

\subsection{第23章　关于宗徒\Person{若望}的记述}

1. 那时，那为\Person{耶稣}所爱的宗徒兼福音作者\Person{若望}，仍生活在\Place{亚细亚}，治理该地区的诸教会，他是在\Person{多米提安}死后从放逐之岛返回的。\par

2. 他那时仍活着，可由两位见证人的见证证明。他们应是可信的，因他们持守了教会的正统信仰；\Person{依勒内}和\Place{亚历山大}的\Person{克肋孟}正是如此。\par

3. 前者在其著作《\WorkTitle{驳斥异端}》第二卷中写道：\Quote{所有在\Place{亚细亚}与主之门徒\Person{若望}交往的长老都证明，\Person{若望}把这些传授给了他们。因为他在他们中间一直逗留到\Person{图拉真}时代。}\par

4. 在同一著作的第三卷中，他以如下词句证实同一件事：\Quote{\Place{厄弗所}教会也是由\Person{保禄}建立、\Person{若望}在那里逗留到\Person{图拉真}时代的，它是宗徒传统的忠实见证者。}\par

5. \Person{克肋孟}在其题为《\WorkTitle{什么富人可得救？}》的书中也指出了时间，并附上了一段叙述，对那些喜爱聆听优美而有益事物的人极具吸引力。请阅读以下记载：\par

6. 请听一个故事，这并非普通传说，而是关于宗徒若望的叙述，世代相传，珍藏于记忆中。当暴君死后，他从帕特摩斯岛返回厄弗所，应邻邦外邦人之请前往，在有些地方任命主教，在另一些地方整顿整个教会，在其他地方则拣选圣神所指明的人服务于圣职。\par

7. 他来到一座不远城市（有人给出了其名称），在安慰了弟兄们的其他事务之后，最后转向被任命的主教，看见一位体格强健、容貌俊美、性情热烈的青年，便说：\Quote{在教会面前，并以基督为见证，我郑重地将这人托付给你。}主教接受了委托并承诺一切后，他再次以同样的见证重申同一嘱托，然后启程往厄弗所去了。\par

8. 但那位长老将托付给他的青年领回家中，养育、看顾、爱护，最终为他施了圣洗。此后，他放松了更严格的看管和警惕，以为给他印上主的印记便是给了他完美的保护。\par

9. 然而，一些同龄的游手好闲、放荡成性、惯于行恶的青年，在他过早摆脱管束时腐蚀了他。起初，他们用奢华的宴饮引诱他；然后，当他们夜间出去抢劫时，便带他同行；最后，他们要求他参与更大的罪行。\par

10. 他逐渐习惯了这些行径，因他性格刚烈，偏离正路，像一匹倔强有力的马咬紧嚼环，更加猛烈地冲向深渊。\par

11. 最终，因对天主的救恩绝望，他不再考虑小事，犯下大罪后，既然彻底失落，便预期与他人同遭厄运。于是，他聚集他们，组成一伙强盗，成为大胆的匪首，众人中最凶暴、最残忍、最冷酷的一个。\par

12. 时光流逝，因某种需要，他们派人去请若望。他安排好了前来的其他事务后，说：\Quote{主教啊，请将我和基督托付给你的交托归还我们，你所主持的教会作证。}\par

13. 主教起初惊慌，以为在钱财上被诬告，因他未曾收受，既不能相信这指责，又不能不信若望。但当若望说：\Quote{我要那青年和弟兄的灵魂。}老人深深叹息，同时泪流满面，说：\Quote{他死了。}若望问：\Quote{怎么死的？什么死？}老人说：\Quote{他向天主死了，}他说，\Quote{因为他变得邪恶放荡，最终成了强盗。如今，他不去教会，反而与一群同类盘踞山间。}\par

14. 宗徒撕裂自己的衣服，捶头痛哭，说：\Quote{我为弟兄的灵魂留下了一个多好的守卫啊！但给我备一匹马，让人引路。}他照旧骑马离开教会，来到那地方，被强盗的哨兵俘虏。\par

15. 他却不逃跑，也不恳求，只喊道：\Quote{我正是为此而来，带我去见你们的头目。}\par

16. 那时，后者正全副武装地等待着。但当他认出若望走近时，羞愧地转身要逃。\par

17. 但若望忘记了自己的年岁，全力追赶他，喊道：\Quote{我儿，为什么逃避你年迈无武装的父亲？可怜我吧，我儿，不要怕；你还有生命的希望。我愿为你在基督面前交账。若有必要，我情愿替你死，如同主为我们舍命一样。我愿为你舍弃我的生命。站住，相信吧；基督派遣了我。}\par

18. 他听见后，先是停步低头；然后扔掉武器，颤抖着痛哭。当老人走近时，他拥抱了他，尽己所能以哀痛告解，用泪水为自己再次施行圣洗，只隐藏了右手。\par

19. 但若望向他保证并起誓他将在救主那里获得宽恕，恳求他，跪下，亲吻他那仿佛已被悔改净化了的右手，领他回到教会。若望为他恳切祈祷，与他一同持守斋戒，用各种言辞驯服他的心思，据说直到将他恢复在教会中才离开，提供了真实悔改的伟大榜样和重生的明证，一座可见复活的凯旋碑。\par

% structure: p0171 source=h2 target=subsection reason=relative-heading-rank; base=h2

\subsection{第二十四章　福音书的次序}

1. 我插入这段克莱孟的引文，是为历史和读者的益处。现在让我们指出这位宗徒无可争议的著作。\par

2. 首先，他的福音为普世万民教会所知，必须承认为真作。古人将其置于其他三部福音之后，列于第四，这是有充分理由的，可由以下方式显明。\par

3. 那些伟大而真正属神的人，即基督的宗徒，他们在生活中被净化，以灵魂的每一种美德为装饰，但在言辞上未经雕琢。他们固然信赖救主所赐予的神圣而奇能的力量，却不知道，也不试图用修饰且讲究的语言宣讲老师的教义，而只运用与他们同工的圣神的证明，以及藉他们显明的基督的奇能，向普世宣扬天国的知识，极少关注书面作品的撰写。\par

4. 他们之所以如此行，是因为他们在职务上得到了超乎凡人的帮助。例如，\Person{保禄}，他在表达的力度和思想的丰富上超越众人，却只写下了最简短的书信，尽管他有无数的奥秘之事要传达，因为他曾达到第三层天的景象，被带到天主的乐园，并蒙恩听见不可言传的话。\par

5. 我们救主的其余门徒——十二宗徒、七十门徒以及无数其他人——并非对这些事无知。然而，在主的所有门徒中，只有\Person{玛窦}和\Person{若望}给我们留下了书面的记述，而据传统所言，他们只是在迫不得已的压力下才动笔的。\par

6. 因为\Person{玛窦}起初向希伯来人宣讲，当他即将前往其他民族时，便用本乡话写下了他的福音，以此补偿那些他不得不离开的人因失去他的面见所受的损失。\par

7. 当\Person{马尔谷}和\Person{路加}已经出版他们的福音之后，据说\Person{若望}——他一直以来都口传福音——最终出于以下原因动笔写作。上述三福音既已传到所有人手中，也传到了他手中，据说他接受了它们并为它们的真实性作证；但其中缺少了基督在传教初期所行的事迹。\par

8. 这确实是事实。因为显然三位福音作者只记录了救主在洗者若翰被囚禁后一年内所行的事迹，并且在他们的记述开头指明了这一点。\par

9. 因为\Person{玛窦}在四十天禁食和随后的试探之后，这样指出他作品的时间顺序：\Quote{他听说若翰被交付后，就从犹太退避到加里肋亚。}\BibleRef{玛 4:12}\par

10. \Person{马尔谷}同样说：\Quote{若翰被交付后，耶稣来到加里肋亚。}\BibleRef{谷 1:14} 而\Person{路加}在开始记述耶稣事迹前，同样指明时间，说他提到\Person{黑落德} \Quote{又在这一切恶事上加了囚禁若翰于监中。}\BibleRef{路 3:20}\par

11. 因此，据说宗徒\Person{若望}为此缘故受托，在他的福音中记述了先前福音作者所省略的那一时期，以及救主在那段时期所行的事迹；也就是说，在洗者若翰被囚禁之前所行的事。他们说他用以下话语指明这一点：\Quote{这是耶稣所行的第一个神迹}；此外，当他在耶稣事迹中间提及洗者若翰时，说若翰仍在\Place{厄农}靠近\Place{撒林}的地方施洗，\BibleRef{若 3:23} 他清楚地说：\Quote{因为若翰那时还没有被投进监狱。}\par

12. 因此，\Person{若望}在他的福音中记载了基督在洗者若翰被囚禁之前所行的事迹，而其他三位福音作者提到的是此后发生的事件。\par

13. 明白了这一点，就不会再认为诸福音互有矛盾，因为根据\Person{若望}的福音包含基督最初的事迹，而其他福音记述他后来的生活。至于我们救主按血统的族谱，\Person{若望}很自然地省略了，因为\Person{玛窦}和\Person{路加}已经给出，而他以关于他神性的教义开始，这似乎是由圣神为他保留的，作为更卓越的一位。\par

14. 以上我们关于\Person{若望}福音所说的话或许足够了。关于\Person{马尔谷}福音的成书原因，我们已经陈述过。\par

15. 至于\Person{路加}，他在福音开头自己陈述了促使他写作的理由。他说，既然许多其他人更冒失地着手编写那些他已知之甚详的事件的叙述，他感到有必要使我们摆脱他们不确定的意见，因此在他的福音中确切地传达了那些他已得知全部真相的事件，这得益于他与\Person{保禄}的亲密关系和停留，以及他对其他宗徒的认识。\par

16. 以上是关于这些事我们自己的叙述。在更合适的地方，我们将尝试引用古人的话，表明其他人对此所说的。\par

17. 至于\Person{若望}的著作，不仅他的福音，而且他的第一封书信，无论在现今还是古代，都被不加争议地接受。但另外两封书信是有争议的。\par

18. 关于\OriginalTerm{默示录}{Apocalypse}，大多数人的意见仍有分歧。但在适当的时候，这个问题也将根据古人的见证来定断。\par

% structure: p0190 source=h2 target=subsection reason=relative-heading-rank; base=h2

\subsection{第二十五章 被接受的和不被接受的圣经}

1. 既然我们在讨论这个主题，宜总结一下已经提到的新约著作。首先必须列为\WorkTitle{四福音}；接着是\WorkTitle{宗徒大事录}。\par

2. 此后应算上\WorkTitle{保禄书信}；其次，现有的\WorkTitle{若望前书}，同样\WorkTitle{伯多禄书信}，必须持守。之后，如果确实合适，应放置\WorkTitle{若望默示录}，关于它我们将在适当时候给出不同意见。这些属于被接受的著作。\par

3. 在有争议的著作中——尽管仍为许多人所承认——现存的有所谓\WorkTitle{雅各伯书}和\WorkTitle{犹达书}，还有\WorkTitle{伯多禄后书}，以及那些称为\par

\SourceRecord{Translated by Arthur Cushman McGiffert.；Nicene and Post-Nicene Fathers, Second Series；Vol. 1.；Edited by Philip Schaff and Henry Wace.；Buffalo, NY: Christian Literature Publishing Co.,；1890.；<http://www.newadvent.org/fathers/250103.htm>.}

% structure: p0001 source=h1 target=section reason=page-title

\section{\WorkTitle{教会史（第四卷）}}

% structure: p0002 source=h2 target=subsection reason=relative-heading-rank; base=h2

\subsection{第一章　图拉真在位期间罗马与亚历山大诸主教}

1. 约在图拉真在位第十二年，上述亚历山大教会的主教去世，由从宗徒算起第四任的普黎慕斯被选任此职。\par

2. 其时，自伯多禄与保禄算起第五任的亚历山大，在罗马接受主教之职，此前厄瓦雷斯图斯已任职八年。\par

% structure: p0005 source=h2 target=subsection reason=relative-heading-rank; base=h2

\subsection{第二章　图拉真在位期间犹太人的灾难}

1. 我们救主的教导和教会日益兴旺，日日进步；但犹太人的灾难却加剧，他们遭受持续不断的祸患。在图拉真在位第十八年，犹太人再次发生骚乱，其中大批人丧生。\par

2. 因为在亚历山大和埃及其余地区，以及在基勒乃，仿佛被某种可怕而好斗的精神煽动，他们向同住的希腊居民发起叛乱。叛乱规模大增，次年，当路普斯担任全埃及总督时，它发展成一场规模不小的战争。\par

3. 在首次攻击中，他们碰巧战胜了希腊人；希腊人逃往亚历山大，将城内的犹太人监禁并杀害。但基勒乃的犹太人虽失其援助，仍继续劫掠埃及土地，蹂躏其地区，由路夸斯领导。皇帝派遣马尔基乌斯·图尔博率步兵、海军及骑兵前去对抗他们。\par

4. 他长时间与他们作战，多次交战，击杀成千上万的犹太人，不仅包括基勒乃的，也包括居住在埃及并前来支援其王路夸斯的人。\par

5. 但皇帝担心美索不达米亚的犹太人也攻击该地居民，便命路基乌斯·昆图斯清剿该省的犹太人。他进军击杀该地大量居民；因之成功，被皇帝任命为犹太省总督。这些事件也以同样的措辞记载于撰写那些时期历史的希腊史家笔下。\par

% structure: p0011 source=h2 target=subsection reason=relative-heading-rank; base=h2

\subsection{第三章　哈德良在位期间为信仰辩护的护教者}

1. 图拉真统治十九年半后，埃利乌斯·哈德良继位为帝。夸德拉图斯向他呈上一篇为我们的宗教辩护的演说，因为某些恶人曾企图扰乱基督徒。该作品至今仍为许多弟兄所持有，也同样为我们所有，并清楚显明此人的才智及其宗徒般的正统信仰。\par

2. 他本人以下列话语揭示其生活的早期年代：\Quote{但我们救主的作为始终显现，因为它们是真实的：——那些得医治的，那些从死里复活的，他们不仅在被医治和被复活时被看见，而且始终显现；不仅救主在世时如此，在他死后，他们也活了相当长的时间，以致其中有些人甚至活到我们这个时代。}夸德拉图斯便是如此。\par

3. 阿里斯提德斯也是一位热诚于我们宗教的信徒，他如同夸德拉图斯一样，留下一篇向哈德良进言的护教文。他的作品也直到今日为许多人保存。\par

% structure: p0015 source=h2 target=subsection reason=relative-heading-rank; base=h2

\subsection{第四章　同一位皇帝统治下罗马与亚历山大诸主教}

同一位皇帝在位第三年，罗马主教亚历山大去世，任职十年。其继任者为西斯图斯。大约同时，亚历山大主教普黎慕斯在其主教任第十二年去世，由犹斯图斯继任。\par

% structure: p0017 source=h2 target=subsection reason=relative-heading-rank; base=h2

\subsection{第五章　从我们救主时代至所述时期耶路撒冷诸主教}

1. 耶路撒冷诸主教的年代顺序我从未在任何文献中找到记载；因为传统称他们皆享年不久。\par

2. 但我从文献中得知此事：直到哈德良治下犹太人的围城战，那里先后有十五位主教，据说他们全是希伯来人后裔，并纯正地接受了基督的知识，因而受到能判断此事之人的认可，被认为堪当主教之职。因为他们的整个教会那时由信奉基督的希伯来人组成，他们从宗徒时代延续到此时发生的围城战；在那次围城中，犹太人再次背叛罗马人，在激战之后被征服。\par

3. 但由于受割礼的主教此时终止，故宜在此列出他们从起初开始的名单。第一位是雅各伯，即所谓主的兄弟；第二位是西默盎；第三位是犹斯图斯；第四位是匝凯；第五位是多俾亚；第六位是本雅明；第七位是若望；第八位是玛提亚；第九位是斐理伯；第十位是塞内卡；第十一位是犹斯图斯；第十二位是肋未；第十三位是厄弗赖；第十四位是若瑟；最后第十五位是犹达。\par

4. 这些是在宗徒时代与所述时期之间生活在耶路撒冷的主教，全都属于受割礼者。\par

5. 哈德良在位第十二年，西斯图斯完成其主教任期第十年后，由特勒斯福鲁斯继任，他是自宗徒算起第七位。与此同时，一年零数月后，第六任的欧默乃斯接替亚历山太教会之领导，其前任任职十一年。\par

% structure: p0023 source=h2 target=subsection reason=relative-heading-rank; base=h2

\subsection{第六章　哈德良治下犹太人最后一次围城}

1. 由于此时犹太人的叛乱更为严重，犹太总督鲁福斯在皇帝派来援军后，以他们的疯狂为借口，毫不留情地攻击他们，不分男女幼童，任意杀戮，并按战争法使该国完全臣服。\par

2. 当时犹太人的首领是一个名叫巴尔科赫巴（意为\Quote{星}）的人，他具有强盗和凶手的品性，然而凭借自己的名字，向犹太人夸口，仿佛他们是奴隶一般，声称自己拥有奇异的权能；他假装自己是从天上降临到他们中间的一颗星，要在他们的苦难中带给他们光明。\par

3. 战争在哈德良（Adrian）统治第十八年最为激烈，发生在比塔拉（Bithara）城，那是一座非常坚固的堡垒，离耶路撒冷不远。围城持续了很长时间，叛军因饥渴而陷入绝境，叛乱的煽动者也受到了应有的惩罚，此后整个民族便因哈德良的法令和命令被禁止再上耶路撒冷地区。因为皇帝下令，他们甚至不得从远处眺望他们祖先的土地。这是佩拉人阿里斯托的记述。\par

4. 这样，当城市因犹太民族被清空、其古代居民完全毁灭之后，它被另一个民族殖民，随后兴起的罗马城市改变了名称，为纪念埃利乌斯·哈德良（Ælius Adrian）皇帝而被称为\OriginalTerm{埃利亚}{Ælia}。由于那里的教会如今由外邦人组成，第一位在割礼的主教之后治理该教会的是马尔库斯（Marcus）。\par

% structure: p0028 source=h2 target=subsection reason=relative-heading-rank; base=h2

\subsection{第七章 当时成为所谓的假知识领袖的人物}

1. 当时，全世界的教会像最明亮的星辰般闪耀，对我们救主和主耶稣基督的信德在全人类中兴旺，那仇恨一切美善、总是敌对真理、最猛烈反对人类得救的恶魔，使出所有手段攻击教会。起初，他借助外在的迫害武装自己攻击教会。\par

2. 但现在，由于无法使用这些手段，他谋划各种计策，采用其他方法来与教会争斗，利用卑鄙诡诈的人作为毁灭灵魂的工具和破坏的仆役。受他煽动，骗子们假借我们宗教之名，将所能赢得的信徒带入毁灭的深渊，同时，通过他们所行的事，使那些不认识信德的人远离通向救恩之道的道路。\par

3. 因此，从那我们已提及的梅南德（Menander）——西蒙（Simon）的继承人——出来了一种蛇般的势力，双舌双头，产生了两种不同异端的领袖：安提约基雅人撒都尼诺（Saturninus）和亚历山大里亚人巴西利德（Basilides）。前者在叙利亚设立了不敬神的异端学派，后者在亚历山大里亚。\par

4. 依勒内（Irenæus）指出，撒都尼诺的虚假教导在大多数方面与梅南德一致，但巴西利德以不可言传的奥秘为借口，编造了怪诞的寓言，将其不敬异端的虚构推到了极端。\par

5. 但当时教会中有许多成员为真理而战，以非凡的口才捍卫宗徒的和教会的教义，也有一些通过著作为后人提供了防范上述异端的自卫手段。\par

6. 在这些著作中，我们有一篇由阿格利帕·卡斯特（Agrippa Castor）——当时最著名的作家之一——所写的对巴西利德的极其有力的驳斥，揭示了此人可怕的欺骗。\par

7. 他在揭露巴西利德的奥秘时说，巴西利德写了二十四本关于福音的书，为自己编造了名叫巴尔卡巴斯（Barcabbas）和巴尔科夫（Barcoph）的先知和其他不存在的人物，并给他们起了野蛮的名字，以眩惑那些惊叹此类事物的人；他还教导说，吃祭偶像的肉和在迫害时期不加谨慎地背弃信德都是无关紧要的事；并像毕达哥拉斯（Pythagoras）那样，命令其追随者保持五年的沉默。\par

8. 上述作者记录了有关巴西利德的其他类似事情，并有能地揭露了他异端的错误。\par

9. 依勒内还写道，卡尔波克拉特（Carpocrates）是这些人的同时代人，他是另一种异端——称为\OriginalTerm{诺斯替派}{Gnostics}的异端——之父，他们不再像西蒙那样秘密传授魔法技艺，而是公开进行。因为他们夸口——好像是什么了不起的事——他们精心准备的情药，以及某些派送梦境并提供保护的恶魔，以及其他类似的能力；据此，他们教导说，凡希望完全进入他们奥秘——或不如说是可憎之事——的人，必须实行一切最恶劣的邪恶，理由是，除非通过无耻行为履行他们对所谓\Quote{宇宙权能}的义务，否则无法逃脱它们。\par

10. 于是，那恶毒的恶魔利用这些仆役，一方面奴役那些被他可怜地引入歧途的人，使之走向毁灭；另一方面，他向不信的外邦人提供了大量诽谤天主圣言的机会，因为这些人的声誉给整个基督徒群体带来了耻辱。\par

11. 因此，在那时代的不信者中，流传着关于我们的可耻而最荒谬的猜疑，认为我们与母亲和姐妹行非法之事，并享受不敬神的筵席。\par

12. 然而，他并未在这些诡计中长期得逞，因为真理确立了自己，并随着时间的推移大放光芒。\par

13. 敌人阴谋因教会的势力而被驳斥，很快便消失了。一个又一个异端接连兴起，而先前的异端总是消逝，如今时而这样，时而那样，以各种方式和形态消失在五花八门的思想中。但大公且唯一真实教会的荣光，始终如一，其伟大和力量不断增长，并将其虔敬、纯朴、自由，以及其神感生活与哲学之谦逊和纯洁，反映给希腊人和外邦人每一民族。\par

14. 同时，针对整个教会的诽谤控告也消失了，唯我们的教导留存，它已胜过一切，并被公认为在尊严、节制以及神圣与哲学教义方面优于一切。因此，现在无人敢于对我们的信仰加以卑劣的诽谤，或像古代敌人曾喜欢散布的谗言。\par

15. 然而，在那时，真理再次召唤许多战士为其辩护，反对不敬神的异端，不仅以口头，也以书面论据驳斥他们。\par

% structure: p0044 source=h2 target=subsection reason=relative-heading-rank; base=h2

\subsection{第八章 教会作家}

1. 在这些战士中，\Person{赫格西普斯}是著名人物。我们已经多次引用他的话，根据他的记述，叙述了宗徒时代发生的事件。\par

2. 他以最简洁的风格撰写了五卷书，记录了宗徒教义的真传承，并在他关于最初设立偶像者的话中表明了他活跃的时期，他写道：\Quote{他们为他们建造了空墓和庙宇，就像至今仍做的那样。其中也有\Person{安提诺乌斯}，他是\Person{哈德良}皇帝的奴隶，为纪念他，也举行了安提诺乌斯运动会，这运动会是在我们的时代设立的。因为他（指\Person{哈德良}）还建立了一座以\Person{安提诺乌斯}命名的城市，并任命了先知。}\par

3. 同时，真正的哲学之真诚爱好者\Person{犹斯定}仍在继续研究希腊文学。他在致\Person{安东尼}的\WorkTitle{护教书}中表明了这个时代，他写道：\Quote{我们认为在这里提及\Person{安提诺乌斯}也并非不当，他生活在我们时代，所有人都因恐惧被迫将他敬奉为神，尽管他们知道他是谁，从何而来。}\par

4. 同一作者在谈到那时发生的犹太战争时补充道：\Quote{因为在最近的犹太战争中，犹太叛军的领袖\Person{巴尔科赫巴}命令，除非基督徒否认并亵渎\Person{耶稣基督}，否则唯独他们应遭受可怕的刑罚。}\par

5. 在同一作品中，他表明他从希腊哲学皈依基督教并非没有理由，而是经过深思熟虑的结果。他的话如下：\Quote{因为我自己，当我陶醉于\Person{柏拉图}的学说，听到基督徒被诽谤，看到他们既不怕死亡，也不怕通常视为可怕的其他事物时，我断定他们不可能生活在邪恶和享乐中。因为哪个贪图享乐或放纵的人，哪个认为以人肉为乐是好事的人，会欢迎死亡以致被剥夺享乐，而不愿努力继续他目前的生活，并躲避统治者的注意，反而献上自己被处死呢？}\par

6. 此外，同一作者还叙述说，\Person{哈德良}收到一位最杰出的总督\Person{塞伦尼乌斯·格拉尼亚努斯}为基督徒辩护的信函，信中陈述仅仅为了满足民众的呼喊，在没有正式控告和审判的情况下处死基督徒是不公正的；于是\Person{哈德良}给亚细亚代行执政官\Person{米努奇乌斯·丰达努斯}发了一道复文，命令他未经起诉和确凿指控不得定任何人的罪。\par

7. 并且他抄录了这封信，保留了原文的拉丁文，并在前面加了以下的话：\Quote{虽然从最伟大、最杰出的皇帝、你的父亲\Person{哈德良}的信中，我们有充分理由要求你按照我们的愿望进行审判，但我们提出这个请求并非因为是\Person{哈德良}命令的，而是因为我们知道我们所求的是公正的。我们附上了\Person{哈德良}书信的副本，好让你知道我们在这件事上也说的是真话。这是副本。}\par

8. 在这些话之后，上述作者给出了拉丁文的复文，我们尽可能准确地翻译成希腊文。内容如下：\par

% structure: p0053 source=h2 target=subsection reason=relative-heading-rank; base=h2

\subsection{第九章 \Person{哈德良}的诏书，规定我们不应不经审判而受罚}

1. 致\Person{米努奇乌斯·丰达努斯}。我收到了最杰出的人\Person{塞伦尼乌斯·格拉尼亚努斯}（你已接替他）写给我的信。我认为此事不经调查就放过似乎不妥，以免这些人受到骚扰，并为告密者提供作恶的机会。\par

2. 因此，如果该省的居民能够明确地支持这项针对基督徒的请愿，以便在法庭上作出答辩，就让他们只走这条路，但不要诉诸人们的请愿和呼喊。因为，如果有人想提出控告，由你来调查此事要合适得多。\par

3. 因此，如果有人控告他们，并表明他们做了任何违法的事，你就根据罪行的轻重来判决。但是，凭\Person{赫拉克勒斯}发誓！如果有人纯粹出于诽谤而提出控告，你要判定他的罪行，并确保你施加惩罚。\par

此类乃\Person{哈德良}复文的内容。\par

% structure: p0058 source=h2 target=subsection reason=relative-heading-rank; base=h2

\subsection{第十章 \Person{安东尼努斯}在位期间罗马与亚历山大的主教}

\Person{哈德良}在位二十一年后去世，由\Person{安东尼努斯}（号称\Quote{虔敬者}）继承罗马帝国的统治。在他统治的第一年，主教\Person{德肋斯福禄}在任第十一年去世，\Person{希吉诺}成为罗马主教。\Person{依勒内}记载，\Person{德肋斯福禄}的死亡因殉道而光荣，并在同一语境中说明，在上述罗马主教\Person{希吉诺}时期，异端创始人\Person{瓦伦提诺}和\Person{马尔基昂}谬误的始作俑者\Person{克尔东}都在罗马知名。他写道：\par

% structure: p0060 source=h2 target=subsection reason=relative-heading-rank; base=h2

\subsection{第十一章 该时代的异端魁首}

1. \Quote{因为瓦伦提诺在希吉诺治下来到罗马，在比约时兴旺，并留到亚尼塞托。马西昂的前驱塞尔东也在希吉诺，第九位主教，的时代进入教会，做了宣认，并继续这样，时而秘密教导，时而又宣认，时而因腐败教义受谴责并从弟兄们的集会中退出。}\par

2. 这些话见于\WorkTitle{驳异端}第三卷。又在第一卷中他论及塞尔东说：\Quote{有一个叫塞尔东的人，他从西满的门徒那里接受了体系，在希吉诺，按宗徒继承的第九位主教，治下来到罗马，教导说法律和先知所宣告的天主不是我们的主耶稣基督的父。因为前者是已知的，后者则是未知的；前者是公义的，后者是良善的。本都的马西昂继承了塞尔东，发展了他的教义，口出无耻的亵渎之言。}\par

3. 同一位依勒内以极大的活力揭露了瓦伦提诺关于物质的错谬的莫测深渊，并揭示了他的邪恶，像藏匿在窝中的蛇一样秘密隐藏。\par

4. 除了这些人，他还说当时还有另一个名叫马尔库斯的人，精通魔法。他也以如下话语描述了他们的亵圣入教仪式和可憎秘仪：\par

5. \Quote{因为他们中有些人预备婚床，用某些表达形式向被领入教者行神秘礼仪，并说这是由他们举行的属灵婚姻，效仿天上的婚姻。另一些人则领他们到水边，当给他们施洗时，念出以下话语：归于宇宙的未知之父之名，归于真理，万有之母，归于那降临在耶稣身上的那位。另一些人则念希伯来名字，为了更好地迷惑被领入教者。}\par

6. 但希吉诺在主教任第四年结束时去世，比约继任管理罗马教会。在亚历山大，欧默内担任牧者十三年后，马尔库斯被任命为牧者。马尔库斯任职十年后去世，塞拉迪翁继任管理亚历山大教会。\par

7. 在罗马，比约在主教任第十五年去世，亚尼塞托在那里担任基督徒的领袖。赫格西普记载他本人当时在罗马，并一直留到艾琉德若任教宗期间。\par

8. 但犹斯定在那些日子特别突出。他以哲学家的姿态宣讲天主圣言，并以著作捍卫信德。他也写了一部\WorkTitle{驳马西昂}的作品，其中指出他写作时马西昂仍在世。他如下说：\par

9. 他如下说：\Quote{有一个本都的马西昂，他甚至至今仍教导他的跟随者认为有另一位大于造物主的天主。他借着魔鬼的帮助，说服了各族许多人出言亵渎，否认这宇宙的创造者是基督的父，而承认另有一位大于他的创造者。所有跟随他们的人，如我们所说，被称为基督徒，正如哲学的名号也给予哲学家一样，虽然他们可能没有共同的教义。}\par

10. 他补充说：\Quote{我们也写了一部反对所有存在过的异端的作品，如果你愿意读，我们会给你。}\par

11. 但同一位犹斯定最成功地与希腊人争辩，并向被称为比约的安敦尼皇帝和罗马元老院发表了包含为我们信仰辩护的演说。因为他住在罗马。他在其\WorkTitle{护教书}中以如下话语表明他是谁，来自何处。\par

% structure: p0072 source=h2 target=subsection reason=relative-heading-rank; base=h2

\subsection{第十二章 犹斯定向安敦尼呈递的\WorkTitle{护教书}}

1. \Quote{致皇帝提图斯·埃利乌斯·哈德良·安敦尼·比约·凯撒·奥古斯都，和他的儿子、哲学家维里西穆斯，以及哲学家路济乌斯，凯撒的亲生子、比约的养子、爱好学问者，并致神圣元老院和全体罗马人民，我，犹斯定，普里斯库斯之子、巴基乌斯之孙，来自叙利亚巴勒斯坦的弗拉维亚·尼阿波利斯，代表那些被不义地仇恨和迫害的各民族之人呈上这篇护辞和请愿，我本人也是其中之一。}同样这位皇帝，从亚细亚的其他弟兄那里得知该省居民给他们施加的各种伤害，认为适合向亚细亚公会发出以下谕令。\par

% structure: p0074 source=h2 target=subsection reason=relative-heading-rank; base=h2

\subsection{第十三章 安敦尼关于我们教义致亚细亚公会的信函}

皇帝凯撒·马尔库斯·奥勒留·安敦尼·奥古斯都，亚美尼库斯，大祭司长，第十五次任保民官，第三次任执政官，致亚细亚公会，敬启。\par

2. 我知道众神也会留心使这样的人不能逃脱察觉。因为他们宁可惩罚那些不肯敬拜他们的人，而非你们。\par

3. 但你们使他们陷入混乱，当你们指控他们无神时，反而坚定了他们持守的意见。对他们而言，被控告时，为自己的天主而死，比活着更可取。因此，当他们舍弃生命而不服从你们的命令时，他们反而取得了胜利。\par

4. 关于已经发生和仍在发生的地震，劝诫你们是合适的，你们每当地震发生时便丧失勇气，然而却惯于将自己的行为与他们相比。\par

5. 他们确实对天主更加自信，而你们却始终在明显的无知中忽视其他神明和不朽者的崇拜，并且压迫和迫害——甚至处死——那些崇拜他的基督徒。\par

6. 但关于这些人，许多行省的总督也写信给我们至神圣的父亲，他回信说，他们不应骚扰这些人，除非他们似乎企图做出影响罗马政府的事。也有许多人给我寄来关于这些人的信函，但我回复他们时，与我父亲所行的一样。\par

7. 但若有人仍执意控告此类人等，则被告虽显系其中一员，亦应免于罪责，而原告则应受罚。此令发布于\Place{厄弗所}亚细亚全体会议。\par

8. 关于这些事，\Person{撒尔德}教会之主教\Person{梅利托}，当时著名之人，可为见证，由其致\Person{维鲁斯}皇帝的《护教书》中为吾道辩护之语可知也。\par

% structure: p0083 source=h2 target=subsection reason=relative-heading-rank; base=h2

\subsection{第十四章　宗徒之友波利卡普的相关事迹}

1. 此时，\Person{阿尼塞托}为\Place{罗马}教会之首，\Person{依勒内}记述，尚在人世的\Person{波利卡普}曾至\Place{罗马}，与\Person{阿尼塞托}就逾越节庆日问题举行会谈。\par

2. 同一作者另有一则关于\Person{波利卡普}的记述，我觉有必要将其添入已述之事。此记述取自\Person{依勒内}之著作\WorkTitle{驳斥异端}第三卷，内容如下：\par

3. 波利卡普不仅受教于宗徒，且与许多曾见过基督者相熟，更在亚细亚被宗徒们立为\Place{士每拿}教会之主教。\par

4. 我们亦在年少时见过他；因他长寿，至耄耋之年，以光荣至极的殉道者之死辞世，终身教授从宗徒所学、教会所传、且惟独真实之事。\par

5. 所有亚细亚教会皆为此作证，直至今日继承\Person{波利卡普}者也如此，他是比\Person{瓦伦提诺}、\Person{马西昂}及诸异端者远为可信且确凿的真理见证人。他亦在\Person{阿尼塞托}之时在\Place{罗马}，使多人转离上述异端者而归向天主的教会，宣告他从宗徒所领受的，乃是教会所传的这唯一真理体系。\par

6. 且有人从他亲耳听闻，主的门徒\Person{若望}去\Place{厄弗所}沐浴，见\Person{塞林托}在内，便未洗澡跑出浴室，喊道：\Quote{我们快逃，免得浴室也塌了，因为真理之敌\Person{塞林托}在里面。}\par

7. 波利卡普本人，当\Person{马西昂}有一次遇见他，说：\InnerQuote{你认得我们吗？}他回答：\InnerQuote{我认得撒殚的长子。}宗徒及其门徒如此谨慎，甚至不与任何歪曲真理者交谈；正如\Person{保禄}也说：\BibleQuote{异端者，第一次第二次警戒后，就可弃绝他。知道这样的人已经背道，犯罪，自定己罪。} \BibleRef{铎 3:10}\par

8. 此外还有一封极具力量的\Person{波利卡普}的\WorkTitle{致斐理伯人书}，凡愿者及关心自身救恩者，皆可从中认识其信德品行与真理宣讲。\Person{依勒内}的记述如此。\par

9. 波利卡普在其上述致\Place{斐理伯人}的书信中（至今尚存），引用了一些取自\BibleBook{伯多禄}{前书}的见证。\par

10. 当号称\Quote{比约}的\Person{安东尼}完成其统治第二十二年后，其子\Person{马尔库斯·奥勒留·维鲁斯}（亦号称\Person{安东尼}）与其弟\Person{路济乌斯}一同继位。\par

% structure: p0094 source=h2 target=subsection reason=relative-heading-rank; base=h2

\subsection{第十五章　维鲁斯治下，波利卡普与诸人在士每拿殉道}

1. 此时，最大规模的迫害扰动着亚细亚，\Person{波利卡普}以殉道终结其生命。然我以为至关重要者，即将其死亡（至今仍有文字记述）载入此历史。\par

2. 有一封以他本人所主持的教会之名致\Place{本都}各堂区的信，中述其遭遇之事，文字如下：\par

3. \Quote{住在\Place{非罗美利}的天主教会，并致各地圣而公教会的所有堂区：愿仁慈、平安与爱德，从天主圣父多多地赐予你们。弟兄们，我们写信给你们，叙述那些殉道者及有福的\Person{波利卡普}所遭遇之事；他终止了迫害，仿佛以他的殉道将其封印。}\par

4. 在此等话语之后，在叙述\Person{波利卡普}之前，他们记录其余殉道者所遇之事，描述他们在苦痛中表现出的极大坚毅。他们称，旁观者见他们被鞭笞至最深处之血管脉络，以致身体内部隐藏的脏腑及四肢暴露无遗；继而置于贝壳与某些尖尖的铁叉上，遭受各种刑罚与折磨，最后被抛给野兽为食，无不惊骇。\par

5. 他们记录，最高贵的\Person{日耳曼尼库斯}尤其出众，因天主的恩宠克服了本性所植的肉身死亡恐惧。当总督想劝说他，以他年轻为由敦促他，恳求他（因他分外年轻壮硕）怜惜自己时，他毫不迟疑，反而急切地引诱野兽扑向自己，几乎逼迫并激怒它，为能更早脱离他们不义而不法的生命。\par

6. 他光荣死后，全体群众惊叹于天主所爱的殉道者之勇毅及全体基督徒种族之刚毅，突然开始呼喊：\Quote{除掉无神论者；搜寻波利卡普。}\par

7. 因呼喊声引起了极大的骚动，有一夫黎基雅人，名\Person{昆图斯}，刚从\Place{夫黎基雅}来，见到野兽与加倍的酷刑，被怯懦击中，放弃了获得救恩。\par

8. 但上述书信表明，他过于仓促、未经审慎，与其他人一同冲到审判台前，但被捕后却向众人提供了明证：如此之人，不宜轻率鲁莽地置身危险。他们之事，遂如此告终。\par

9. 然而，最可敬的\Person{波里加布}最初听到这些事时，仍然保持平静，持守安详而坚定不移的心，决意留在城内。但因朋友们恳求并劝他秘密退避，他被说服，便离开到离城不远的一个农场，与少数同伴住在那里，日夜不停地以祈祷与主角力，恳求并祈求普世各教会得享平安。这是他素常的习惯。\par

10. 在他被捕前三天，当他正在祈祷时，夜间的异象中看见他头下的枕头突然被火抓住并烧尽；醒来后他立即向在场的人解释这异象，几乎预言了将要发生的事，并向那些与他同在的人清楚宣告，他必须为基督的缘故死于火中。\par

11. 于是，当搜捕他的人加紧追查时，据说他因弟兄们的关怀与爱德再次被迫前往另一农场。不久，追捕者来到那里，抓住了两个仆人，并拷打其中一个，以从他口中得知\Person{波里加布}的藏身之处。\par

12. 他们在傍晚时分来到，发现他正躺在楼上的一间房里，他本可以转到另一间屋子去，但他不肯，却说：\Quote{愿天主的旨意成就。}\par

13. 按记载所述，当他得知他们来到时，便下楼以极其愉快而温和的面容对他们说话，以致那些原本不认识他的人，见他年事已高，举止庄重坚定，都以为看见了奇迹，并惊叹竟要如此费力来捉拿这样一个人。\par

14. 但他没有迟疑，立即吩咐为他们摆上桌子，然后邀请他们享用丰盛的餐食，并请求给他们一小时的时间让他安静祈祷。他们应允后，他便站起来祈祷，充满主的恩宠，以致那些在场听见他祈祷的人都感到惊讶，多人此刻懊悔，如此可敬而虔诚的老人竟要被处死。\par

15. 此外，关于他的记述还包含以下内容：当他终于结束祈祷，记念了一切曾与他接触的人，无论大小、有名无名，以及普世的天主教会后，离去的时刻到了，他们便把他驮在驴上，带往城中，那天正是大安息日。警察队长\Person{黑落德}和他的父亲\Person{尼格特}遇见他，便把他扶上他们的马车，坐在他旁边，试图说服他，说：\Quote{有什么害处呢？说\InnerQuote{主凯撒}并献祭以保全性命。}起初他没有回答；但他们坚持时，他说：\Quote{我不会做你们所建议的事。}\par

16. 他们见无法说服他，便说出可怕的话，并粗暴地将他推下，以致他从马车下来时擦伤了胫骨。但他没有回头，继续迅速而敏捷地前行，仿佛什么也没有发生，被带到了竞技场。\par

17. 但竞技场里一片嘈杂，以致很少人听到从天上传来的声音，这声音在\Person{波里加布}进入时临到他：\Quote{波里加布，你要刚强，要作大丈夫。}没有人看见说话者，但我们许多人都听见了那声音。\par

18. 当他被带上前时，众人听说\Person{波里加布}已被捕，引起一阵大骚动。最后，他走上前来，总督问他是否就是\Person{波里加布}。他承认了，总督便设法劝他背弃信仰，说：\Quote{要顾及你的年纪}以及其他类似的话，这是他们惯常说的：\Quote{指着凯撒的护守神起誓；悔改，并说\InnerQuote{除掉无神派}。}\par

19. 但\Person{波里加布}以庄严的目光环视聚集在竞技场中的全体人群，向他们挥手，叹息着，举目望天，说：\Quote{除掉无神派。}\par

20. 长官逼迫他说：\Quote{起誓吧，我就释放你；咒骂基督。}\Person{波里加布}说：\Quote{八十六年来我一直事奉他，他从未亏负过我；我怎能亵渎那拯救我的君王呢？}\par

21. 当他再次坚持说：\Quote{指着凯撒的护守神起誓}时，\Person{波里加布}回答说：\Quote{你若妄自以为我会照你所说的指着凯撒的护守神起誓，假装不知道我是谁，那就请听清楚：我是基督徒。但若你愿意学习基督信仰的道理，就定个日子来听吧。}\par

22. 总督说：\Quote{你去说服民众。}但\Person{波里加布}说：\Quote{至于你，我认为你堪受解释；因为我们受教，对由天主设立的君长和当权者，要予以应得的尊敬，只要不伤害我们；但对于这些人，我认为他们不配听我的答辩。}\par

23. 总督说：\Quote{我有野兽，你若不肯悔改，我就把你扔给它们。}但他说：\Quote{叫它们来吧；因为从较好转向较差，是我们不能作的改变。但从邪恶转向正义，却是高尚的事。}\par

24. 他又对他说：\Quote{你若藐视野兽，我就用火烧死你，除非你悔改。}但\Person{波里加布}说：\Quote{你所威胁的火只烧一时，不久便熄灭；因为你不知道那为不虔敬者所存留的将来审判和永恒惩罚的火。你为何迟延呢？任意而行吧。}\par

25. 说了这些话以及其他言语之后，他充满勇气和喜乐，面容洋溢着恩宠，以致他不仅没有被那些对他所说的话吓倒或惊慌，反而总督感到惊奇，并派传令官在竞技场中央三次宣告：\Quote{波里加布已承认自己是基督徒。}\par

26. 当传令官宣告此事后，住在\Place{斯米纳}的全群外邦人和犹太人都以无法抑制的愤怒大声喊叫：\Quote{这是亚细亚的教师，基督徒的父，我们神祇的推翻者，他教导许多人不可献祭也不可敬拜。}\par

27. 他们说了这话，就喊叫着请求亚细亚行省总督\Person{腓力}放出一头狮子扑向\Person{波利卡普}。但他说不可以这样做，因为他已经结束了竞技。于是他们一致认为应该高声呼喊，将\Person{波利卡普}活活烧死。\par

28. 因为必须应验那曾向他启示的关于他枕头的异象，当时他正在祈祷，看见枕头燃烧，就转过身来预言性地对同他在一起的信友说：\Quote{我必须被活活烧死。}\par

29. 这些事进行得极为迅速——比说话还快——群众立刻从工场和浴场收集木料和柴捆，犹太人尤其热心于这项工作，正如他们惯常所做的那样。\par

30. 当柴堆备好时，他脱去所有外衣，解下腰带，也试图脱下鞋子，虽然他以前从未这样做过，因为每位信徒总是努力首先接触到他的皮肤；他在白发未生之前就已因其德行备受尊敬。\par

31. 随即他们把为柴堆准备的材料围放在他身边；当他们也准备将他钉在木桩上时，他说：\Quote{就让我这样吧；因为那位赐我力量忍受火焰的，也会赐我力量在火中安然不动，无需你们用钉子固定。}于是他们没有钉他，而是绑了他。\par

32. 他双手被绑在背后，如同一只从大羊群中取出的高贵公羊，作为可悦纳的全燔祭献给全能的天主，说道：\par

33. \Quote{父啊，你所爱的、蒙祝福的圣子耶稣基督的父，藉着他我们得以认识你，你是天使、万有、一切受造物以及活在你面前的整个义族的天主，我颂扬你，因你使我配得上这一天和这个时辰，使我能在殉道者的行列中、在基督的杯中得到一份，以获致灵魂和身体的永生复活，并因圣神的不朽。}\par

34. \Quote{愿我今天在你面前被收纳于他们之中，作为丰盛而可悦纳的祭献，正如你，信实而真实的天主，所预先预备和启示，并已成就的。}\par

35. \Quote{因此，我也为一切赞美你；我颂扬你，我光荣你，藉着永恒的至高司祭，耶稣基督，你所爱的圣子；藉着他，同他，并在圣神内，愿光荣归于你，从现在直到永远，阿们。}\par

36. 当他献上阿们并结束祈祷后，点火者点燃了火；当大火燃起时，我们这些被赐予看见的人，看到了一个奇事，并被保存下来以便将所发生的事传给其他人。\par

37. 因为火焰呈现出拱顶的形状，如同风鼓满的船帆，在殉道者的身体周围形成了一道墙；他的身体在中间，不像肉在燃烧，倒像金银在熔炉中被精炼。因为我们闻到了一种芬芳的香气，如同乳香或其他珍贵香料的烟雾。\par

38. 最后，那些无法无天的人见身体不能被火焚毁，就命令一名刽子手上前用剑刺他。\par

39. 当他这样做时，流出了大量的血，以致熄灭了火；全体群众都感到惊奇，不信者与蒙选者之间竟有如此大的差别，而此人正是其中之一，是我们时代最奇妙的教师，\OriginalTerm{宗徒的和先知的}{apostolic and prophetic}，他曾是\Place{士每拿}天主教会的主教。因为从他口中说出的每一句话都成就了，并将成就。\par

40. 但那位嫉妒而满怀恶意的\OriginalTerm{邪恶者}{Evil One}，义族的仇敌，见他殉道的伟大，以及他从起初开始的无可指摘的生活，并见他戴上了不朽的冠冕，赢得了无可争议的奖赏，就设法不让我们取走他的身体，尽管许多人渴望这样做，并与他的圣体共融。\par

41. 于是，某些人暗中向\Person{尼塞特斯}——\Person{希律}的父亲、\Person{阿尔克}的兄弟——建议，要他请求长官不要交出他的身体，\Quote{免得，}他们说，\Quote{他们抛弃被钉者，开始敬拜这个人。}他们这样说，是出于犹太人的唆使和推动；那些犹太人也监视着我们，当我们正要从火中取走他的身体时，他们不知道我们绝不会离弃为拯救那得救的整个世界而受苦的基督，也不会敬拜任何其他人。\par

42. 因为我们敬拜那为天主圣子者，但对于殉道者，我们作为主的门徒和效法者，以他们应得的爱来爱他们，因他们对自身君王和师傅的无比挚爱。愿我们也能与他们一同有分，并作同门徒。\par

43. 于是百夫长见犹太人所表现的争执，就将他放在中间烧了，这是他们的惯例。之后我们收拾他的骨骸，比宝石更珍贵，比金子更可宝爱，安放在一个合适的地方。\par

44. 在那里，主将允许我们尽所能地聚集，在喜乐和欢欣中庆祝他殉道的诞辰，以纪念那些已经作战的人，并训练和预备那些将来也要这样做的人。\par

45. 这就是发生在有福的\Person{波利卡普}身上的事，他与来自\Place{非拉铁非}的十一位弟兄在\Place{士每拿}殉道。只有他一人比其他人更被众人纪念，以致连外教人在各处也谈论他。\par

46. 可敬而宗徒的\Person{波利卡普}被认为配得上这样的结局，正如\Place{士每拿}教会的弟兄们在我们已提及的书信中所记载的。在同一卷关于他的书中，还附录了其他殉道事迹，这些事迹发生在同一座城市\Place{士每拿}，大约与\Person{波利卡普}殉道同时期。其中\Person{梅特罗多鲁斯}，似乎是\OriginalTerm{马西昂派}{Marcionitic sect}的皈依者，也死于火刑。\par

47. 当时有一位著名的殉道者，名叫\Person{比奥尼乌斯}。凡愿了解他多次的宣认、他在民众和统治者面前发言的胆量、他为信仰所作的辩护、他有教益的讲话，以及他向在迫害中屈服于诱惑者的问候、他在监狱中前来探望的弟兄们所说的鼓励之言，此外他所受的酷刑、钉刑、在火刑柱上的坚定，以及经历所有这些非凡考验后的死亡——我们请他们参阅我们收录在\WorkTitle{古代殉道录}中的那封书信，其中对他有十分详尽的记述。\par

48. 此外，还有其他殉道者的记录留存，他们在亚细亚的\Place{培尔加摩}城殉道——即\Person{卡尔普斯}、\Person{帕皮鲁斯}和一位名叫\Person{阿加托尼刻}的妇女，她在许多光辉的见证之后，光荣地结束了生命。\par

% structure: p0143 source=h2 target=subsection reason=relative-heading-rank; base=h2

\subsection{第十六章　哲学家\Person{犹斯定}在\Place{罗马}宣讲基督的圣言并殉道}

1. 大约此时，上文提及的\Person{犹斯定}在向已提名的统治者递交了第二篇为我们的教义辩护的著作之后，因\Person{克雷森斯}——一位名字上与犬儒学派相似、效法其生活方式和习惯的哲学家——设计陷害，获得了神圣的殉道冠冕。\Person{犹斯定}多次在公开辩论中驳斥他，最终以殉道赢得了胜利的奖赏，为他所宣讲的真理而死。\par

2. 他本人，一位在真理方面最为博学的人，在他已提及的《护教书》中明确预言了此事将如何发生在他身上，尽管那时尚未发生。\par

3. 他的话如下：\Quote{因此，我也预料会被我提到的人之一，或者被\Person{克雷森斯}——那个不配哲学之名、虚荣自负的人——所陷害和锁上足枷。因为此人公开指控他根本不知情的事，为要取悦和讨好群众，宣称基督徒是无神和不敬虔的。}\par

4. \Quote{这样做，他大错特错。因为如果他未曾读过基督的教导就攻击我们，那么他完全堕落，比那些常常提防讨论和妄证自己不懂之事的人坏得多。而如果他读过却不明白其中的伟大，或者明白却为了不被人怀疑是信徒而做这些事，那么他更为卑鄙，完全堕落，被世俗的掌声和非理性的恐惧所奴役。}\par

5. \Quote{因为我愿你们知道，当我向他提出这类问题并征询他的意见时，我得知并证明他确实一无所知。为表明我说的是真话，如果这些辩论尚未报告给你们，我愿在你们面前重新讨论这些问题。这确实是配得上皇帝的行为。}\par

6. \Quote{但如果我的提问和他的回答已为你们所知，那么你们显然明白他对我们的事一无所知；或者如果他知情，却因听众而不敢发言，则他表明自己，正如我已说过的，不是哲学家，而是一个虚荣自负的人，甚至不尊重\Person{苏格拉底}那句最可称赞的话。}以上是\Person{犹斯定}的话。\par

7. 关于他正如所预言的，因\Person{克雷森斯}的阴谋而遇害，此事由\Person{塔提安}陈述。他早年讲授希腊科学，获得不小声誉，并留下了许多著作作为纪念。他在其著作\WorkTitle{反对希腊人}中记载了此事，写道：\Quote{那位最可敬佩的\Person{犹斯定}真实地宣称上述之人如同强盗。}\par

8. 然后，在谈论哲学家之后，他接着写道：\Quote{\Person{克雷森斯}确实在那大城中筑巢，在反自然的欲望上超过众人，完全贪爱钱财。}\par

9. \Quote{而他，教导人应当藐视死亡，自己却对死亡如此恐惧，以至于试图将死亡当作大恶加于\Person{犹斯定}，因为后者在宣讲真理时证明了哲学家们是贪食者和骗子。这就是\Person{犹斯定}殉道的原因。}\par

% structure: p0153 source=h2 target=subsection reason=relative-heading-rank; base=h2

\subsection{第十七章　\Person{犹斯定}在其著作中提及的殉道者}

1. 同一位\Person{犹斯定}在其冲突之前，在第一篇《护教书》中提到了在他之前殉道的其他人，并非常恰当地记录了以下事件。\par

2. 他如此写道：\Quote{有一位妇女与一个放荡的丈夫同住；她本人原先也如此。但当她得知基督的教导后，变得节制，并努力劝她的丈夫同样节制，反复讲述教导，并宣告那些不按正当理性节制生活的人将受永火的惩罚。}\par

3. \Quote{但他仍沉溺于同样的放纵，其行为使妻子疏远。因为她最终认为，与一个违反自然和正当之理、千方百计寻求享乐的男人作为妻子同居是错误的，想要与他离婚。}\par

4. \Quote{当她的朋友们劝她仍与他同住，认为她丈夫或许有时会改过自新时，她强逼自己留了下来。}\par

5. \Quote{但当她的丈夫去了亚历山大里亚，并据说行为更糟时——为了不因继续婚姻、同席同床而参与他的不法和不敬——她给了我们所谓的休书，离开了他。}\par

6. \Quote{但她那个高贵而杰出的丈夫——本应高兴，如他所应当的，因为她放弃了从前与仆人和雇工们轻率犯下的行为（那时她沉溺于醉酒和各种恶习），并希望他也同样放弃——却因她违背他的意愿离开而控告她，宣称她是基督徒。}\par

7. 她向皇帝您恳求，允许她先处理自己的事务，待事务处理完毕后，再针对指控进行辩护。而您允准了此事。\par

8. 然而，那个曾是她丈夫的人，因无法再控告她，便将攻击转向一位名叫\Person{托勒密}的基督徒教义导师，此人曾受到\Person{乌尔比基乌斯}的惩罚。他对其采取如下方式：\par

9. 他劝说一位与他交好的百夫长，将\Person{托勒密}投入监狱，并带他前来，只问这一个问题：你是否是基督徒？当\Person{托勒密}——他热爱真理，并非诡诈虚伪之人——承认自己是基督徒后，百夫长就将他捆绑，在监狱中折磨了很长时间。\par

10. 最后，当此人被带到\Person{乌尔比基乌斯}面前时，也只被问了同样的问题：你是否是基督徒？他再次意识到自己因基督的教导所蒙受的恩惠，便承认自己领受了神圣德性的教育。\par

11. 因为凡否认自己是基督徒的，或因鄙视基督教而否认，或自觉不配、与基督教无份而不敢承认；但真正的基督徒绝非如此。\par

12. 当乌尔比基乌斯下令将他带去受刑时，一个同样身为基督徒的\Person{卢基乌斯}，见审判如此不公，便对\Person{乌尔比基乌斯}说：\Quote{你为何惩罚这人？他并非奸夫、淫乱者、凶手、小偷、强盗，也未犯任何罪行，只是承认自己基督徒的名号。乌尔比基乌斯啊，你的审判不符皇帝皮乌斯、哲人恺撒之子或神圣元老院的风范。}\par

13. 乌尔比基乌斯没有其他回答，便对\Person{卢基乌斯}说：\Quote{我看你也是其中之一。}卢基乌斯说：\Quote{当然。}他再次下令将他带去受刑。但卢基乌斯表示感恩，因为他说，他已从如此邪恶的统治者手中得解脱，正前往良善的父与君王——天主那里。接着又有第三人挺身而出，被判受刑。\par

14. 对此，\Person{尤斯定}恰如其分且一贯地补充了我们上文所引用的话语，说道：\Quote{因此，我也预料将受到我所提名的那些人中某位的暗算。}等等。\par

% structure: p0168 source=h2 target=subsection reason=relative-heading-rank; base=h2

\subsection{第十八章　\Person{尤斯定}流传至今的著作}

1. 这位作家留下了大量心智受过神圣事物训练和操练的纪念碑，充满了各种有益的内容。我们将把勤勉的学生引向它们，并在叙述过程中记下我们所知的那些作品。\par

2. 他有一篇护教言论，呈递给绰号\Quote{大善}的\Person{安敦尼}及其儿子们，以及罗马元老院。另一部作品包含他的第二篇\WorkTitle{护教书}，是为了我们的信仰而呈递给上述皇帝的那位继任者，他与那位皇帝同名，即\Person{安敦尼·维鲁斯}，我们正在记述他的时代。\par

3. 还有另一部针对希腊人的著作，他在其中详细论述了我们与希腊哲学家之间的大部分争议问题，并讨论了魔鬼的本性。我无需在此添加这些内容。\par

4. 他的另一部反对希腊人的著作也流传至今，他命名为\WorkTitle{驳斥}。此外还有另一部\WorkTitle{论天主的至高主权}，他不仅从我们的圣经，也从希腊人的著作中加以论证。\par

5. 还有一部题为\WorkTitle{琴师}的著作，以及另一篇\WorkTitle{论灵魂}的论辩，他在其中提出了关于所讨论问题的各种疑问，并列出了希腊哲学家的意见，承诺在另一部著作中加以驳斥并提出自己的观点。\par

6. 他还写了一部反对犹太人的对话录，是在\Place{厄弗所}城与\Person{特黎丰}（当时希伯来人中一位非常杰出的人物）进行的。在其中他展示了天主的恩宠如何催促他信奉信仰的教义，他以前如何热切地追求哲学研究，以及他如何炽烈地寻求真理。\par

7. 他在同一部著作中记载犹太人曾密谋反对基督的教导，并对\Person{特黎丰}说了同样的话：\Quote{你们不但没有悔改自己所行的恶事，反而在当时挑选了一些人选，从耶路撒冷派他们走遍各地，宣告基督徒的无神异端已经出现，并控告他们所有那些对我们无知的人所指责的事，以致你们不仅成了自身不义的原因，也成了所有其他人的不义的原因。}\par

8. 他还写道，甚至直到他的时代，先知之恩仍在教会中闪耀。他提到\Person{若望}的\BibleBook{若望}{默示录}，明确说那是宗徒的作品。他还引用某些先知预言，指责\Person{特黎丰}，因为犹太人将那些预言从圣经中删除了。他的许多其他著作至今仍存于许多弟兄手中。\par

9. 此人的言论甚至被古代作者认为非常值得研究，以至于\Person{依勒内}引用他的话：例如，在其著作\WorkTitle{驳斥异端}的第四卷中写道：\Quote{尤斯定在其反对\Person{马尔基昂}的著作中说得很好：即使主亲自宣讲另一位神而非造物主，他也不会相信。}在同一部著作的第五卷中又说：\Quote{尤斯定说得很好：在主来临之前，撒殚从不敢亵渎天主，因为他尚未知道自己的刑罚。}\par

10. 我认为有必要说这些，以激励勤勉的学生仔细研读他的著作。关于他就说这些。\par

% structure: p0179 source=h2 target=subsection reason=relative-heading-rank; base=h2

\subsection{第十九章　\Person{维鲁斯}统治时期罗马教会和亚历山大里亚教会的统治者}

1. 在上述统治的第八年，\Person{索特尔}继\Person{阿尼塞特}之后成为罗马教会的主教，\Person{阿尼塞特}在职共十一年。但当\Person{克拉迪翁}治理亚历山大里亚教会十四年后，\Person{阿格里皮努斯}继任。\par

% structure: p0181 source=h2 target=subsection reason=relative-heading-rank; base=h2

\subsection{第二十章　安提约基雅教会的统治者}

1. 那时，在\Place{安提约基雅}教会中，\Person{德奥斐禄}作为自宗徒以来的第六位主教而闻名。因为继\Person{希洛}之后的\Person{科尔乃略}是第四位，之后\Person{厄洛}按顺序是第五位，他们担任了主教职务。\par

% structure: p0183 source=h2 target=subsection reason=relative-heading-rank; base=h2

\subsection{第二十一章 当时兴盛的教会作家。}

1. 那时，在教会中活跃的有\Person{赫格西卜}（我们从先前已知），以及\Place{科林托}的主教\Person{狄约尼削}，另一位\Place{克里特}的\Person{皮尼图斯}主教，此外还有\Person{斐理伯}、\Person{阿波利纳里}、\Person{梅利托}、\Person{穆萨奴}、\Person{莫德斯图}，最后是\Person{依勒内}。他们以笔传下了宗徒传统所接受的纯正而正统的信德。\par

% structure: p0185 source=h2 target=subsection reason=relative-heading-rank; base=h2

\subsection{第二十二章 赫格西卜及其所提及的事件。}

1. \Person{赫格西卜}在流传于我们的五卷《回忆录》中，留下了他自己见解的最完整的记录。在其中，他声明在前往\Place{罗马}的旅途中，他会晤了众多主教，并从他们所有人那里领受了相同的道理。在提到\Person{克莱孟}致\Place{科林托}人书的一些评论后，听听他说什么是适宜的。\par

2. 他的话如下：\Quote{科林托的教会一直持守真信德，直到\Person{普里慕斯}在\Place{科林托}担任主教。我在前往罗马的路上与他们交谈，并和科林托人一起住了许多天，在此期间，我们在真道中彼此获得慰藉。}\par

\Quote{3. 当我来到\Place{罗马}时，我留在那里直到\Person{阿尼塞特}，其执事是\Person{厄勒乌特鲁}。\Person{阿尼塞特}由\Person{索特尔}继任，\Person{索特尔}又由\Person{厄勒乌特鲁}继任。在每一次继承中，并在每一座城市里，\InnerQuote{由法律、先知和主所宣讲的}被持守。}\par

4. 同一位作者也以下列话语描述了他时代兴起的异端之开端：\Quote{在义人雅各伯殉道之后——主也因同样的缘由殉道——主的叔父克罗帕之子西默盎被任命为下一任主教。所有人都推举他作为第二位主教，因为他是主的表亲。}\par

\Quote{因此，他们称教会为童贞女，因为她尚未被虚妄的言谈所败坏。}\par

\Quote{5. 但\Person{特布提}因未被选为主教，便开始败坏她。他也出自民众中的七个派别，如\Person{西蒙}（从中产生了\OriginalTerm{西蒙派}{Simonians}）、\Person{克利奥伯}（从中产生了\OriginalTerm{克利奥伯派}{Cleobians}）、\Person{多西特}（从中产生了\OriginalTerm{多西特派}{Dositheans}）、\Person{戈尔特}（从中产生了\OriginalTerm{戈尔特派}{Goratheni}）、\Person{玛斯博特}（从中产生了\OriginalTerm{玛斯博特派}{Masbothaeans}）。从他们产生了\OriginalTerm{梅南德派}{Menandrianists}、\OriginalTerm{马尔基翁派}{Marcionists}、\OriginalTerm{卡尔波克拉特派}{Carpocratians}、\OriginalTerm{瓦伦廷派}{Valentinians}、\OriginalTerm{巴西里德派}{Basilidians}和\OriginalTerm{萨图尼努斯派}{Saturnilians}。每个人都私下并分别引入了他自己的独特意见。从他们产生了假基督、假先知、假宗徒，他们以敌对天主及其基督的败坏道理分裂了教会的合一。}\par

6. 同一位作者也以下列话语记录了犹太人中兴起的古代异端：\Quote{此外，在以色列子民中，割礼派内存在各种意见。以下这些是与犹大支派及基督对立的：\OriginalTerm{厄色尼派}{Essenes}、\OriginalTerm{加里肋亚派}{Galileans}、\OriginalTerm{每日洗者派}{Hemerobaptists}、\OriginalTerm{玛斯波泰派}{Masbothaeans}、\OriginalTerm{撒玛黎雅派}{Samaritans}、\OriginalTerm{撒杜塞人}{Sadducees}、\OriginalTerm{法利塞人}{Pharisees}。}\par

7. 他还写了其他许多事，我们已部分提及，按适当顺序引入叙述。他从\WorkTitle{依据希伯来人语的叙利亚福音}中引用了希伯来文的若干段落，显示他是希伯来人中的皈依者，并提及取自犹太人不成立传统的其他事。\par

8. 不仅是他，还有\Person{依勒内}以及所有古代的作者，都称\WorkTitle{撒罗满箴言}为\Emph{全德智慧}。在谈到被称为\OriginalTerm{伪经}{Apocrypha}的书卷时，他记录其中一些是在他时代由某些异端者编纂的。但让我们现在转向另一位。\par

% structure: p0195 source=h2 target=subsection reason=relative-heading-rank; base=h2

\subsection{第二十三章 科林托的主教狄约尼削及其所写的书信。}

1. 首先我们必须谈谈\Person{狄约尼削}，他被任命为\Place{科林托}教会的主教，并慷慨地将他受感动的劳作不仅分享给自己的人，也分享给外地的人，并通过他写给各教会的公函为所有人提供了极大的服务。\par

2. 其中包括致\Place{拉刻代蒙}人的一封，包含正统信德的教导以及对和平与合一的劝勉；还有致\Place{雅典}人的一封，激励他们信奉信德并过福音所规定的生活，他责备他们几乎轻视这些，仿佛自从统治官\Person{普布利乌斯}在当时的迫害中殉道以来，他们几乎背弃了信德。\par

3. 他也提及\Person{夸德拉图}，声明他在\Person{普布利乌斯}殉道后被任命为他们的主教，并见证通过他的热心，他们重新聚集，信德得以复兴。此外，他记载\Place{阿勒约帕哥}的\Person{狄约尼削}——按照\BibleBook{宗徒}{大事录}的陈述，由宗徒\Person{保禄}皈依信德——首先担任了\Place{雅典}教会的主教职务。\par

4. 还有另一封他致\Place{尼科默迪亚}人的书信留存，其中他攻击\Person{马尔基翁}的异端，并坚定持守真理的准则。\par

5. 他写信给在\Place{戈尔提纳}的教会，连同\Place{克里特}的其他堂区，称赞他们的主教\Person{斐理伯}，因为在他领导下的教会以诸多坚毅行为著称，并警告他们要提防异端者的偏差。\par

6. 他在致\Place{阿马斯特里斯}的教会（连同\Place{本都}的教会）的信中，提到\Person{巴库利德}和\Person{厄尔皮斯图斯}曾敦促他写信，并添加了对神圣圣经段落的解释，并点名提到他们的主教\Person{帕尔玛}。他就婚姻和贞洁给予了他们许多建议，并命令他们接纳任何跌倒后（无论是失足还是异端）回头的人。\par

7. 在这些信中还包括一封致\Place{诺索斯}人的信，其中他劝勉该堂区的主教\Person{皮尼图斯}，不要在贞洁方面给弟兄们施加沉重和强制性的负担，而要顾及众人的软弱。\par

8. \Person{皮尼图斯}回复这封信，钦佩并称赞\Person{狄约尼削}，但反过来劝勉他再写信时，适时分给更坚实的食粮，用更进深的教导牧养他手下的人，使他们不持续被这些乳般的道理喂养，在不自觉中因着适于儿童的训练而老去。在这封信中也揭示了\Person{皮尼图斯}信德的正统、他对属下的关爱、他的学识以及对神圣事物的理解，如在一幅最完美的肖像中。\par

9. 此外还存在另一封由\Person{狄约尼削}写给罗马人、并致当时的主教\Person{索泰尔}的信。我们最好附上这封信中的一些段落，他在其中称赞罗马人的做法——这做法一直保持到我们本时代的迫害。他的言辞如下：\par

10. \Quote{因为从起初，你们就习惯于以各种方式善待所有弟兄，并向各城中的许多教会送去捐助。你们罗马人藉着从起初就送去的礼物，如此救济贫乏者的需要，并为在矿中的弟兄作准备，你们保持了罗马人的祖传习俗，而你们可敬的主教\Person{索泰尔}不仅维持了它，而且予以增添，向圣者们提供丰富的供应，并以祝福的话语鼓励来自外地的弟兄，如同慈父对待儿女一般。}\par

11. 在同一封信中，他也提到\Person{克肋孟}写给格林多人的信，表明从起初就有在教会中宣读它的习惯。他的言辞如下：\Quote{今天我们度过了主的神圣之日，在这一天我们读了你们的信。从其中，无论何时我们读它，我们总能汲取劝告，就如同从先前那封经由\Person{克肋孟}写给我们的信一样。}\par

12. 同一位作者也如下谈论他自己的书信，声称它们曾被篡改：\Quote{弟兄们要我写信，我就写了。而这些书信，魔鬼的宗徒们用莠子填满了它们，删掉了一些东西，添加了另一些。为他们保留着祸哉。因此，如果有人也试图篡改主的著作，就不足为奇了，因为他们甚至对不那么重要的著作也图谋不轨。}\par

此外，还有另一封\Person{狄约尼削}写给最忠信姊妹\Person{克鲁索福拉}的信。他在信里写了恰当的话，也分给她合适的属灵食粮。关于\Person{狄约尼削}，就谈到这里。\par

% structure: p0209 source=h2 target=subsection reason=relative-heading-rank; base=h2

\subsection{第二十四章 \Place{安提约基雅}主教\Person{德奥斐禄}}

1. 关于我们曾提及的\Place{安提约基雅}教会主教\Person{德奥斐禄}，现存有致\Person{奥托利库}的三部基础著作；另有题为\WorkTitle{驳\Person{赫尔谟根}异端}的另一著作，其中他引用了\BibleBook{若望}{默示录}的见证；最后还有一些其他的教理书籍。\par

2. 由于异端者那时不亚于其他时候，像莠子一样毁坏宗徒教导的纯净庄稼，各地教会的牧者急忙约束他们，如同从基督的羊圈赶走野兽，时而通过劝诫和训勉弟兄们，时而通过口头讨论和驳斥更公开地与他们争辩，时而又通过书面著作以最精确的论证纠正他们的意见。\par

3. 而且\Person{德奥斐禄}也与其他人士一起与他们争辩，这从他写的一篇反对\Person{玛尔基昂}的不凡价值的讲演中可以显明。这部著作，连同我们所说过的其他著作，也保存至今。\par

\Person{马克西米努}，宗徒后的第七位，继他任\Place{安提约基雅}教会的主教。\par

% structure: p0214 source=h2 target=subsection reason=relative-heading-rank; base=h2

\subsection{第二十五章 \Person{斐理伯}与\Person{莫德斯图}}

\Person{斐理伯}，正如我们从\Person{狄约尼削}的话语中所知，是\Place{戈尔廷}教区的主教，也写了一部极其精详的反对\Person{玛尔基昂}的著作，\Person{依勒内}和\Person{莫德斯图}也同样写了。后者比其余人更清楚地向众人揭露了那人的错误。还有许多人，他们的著作至今仍被许多弟兄们呈现出来。\par

% structure: p0216 source=h2 target=subsection reason=relative-heading-rank; base=h2

\subsection{第二十六章 \Person{默里托}与他所记载的情况}

1. 在那些日子里，\Place{撒尔德}教区的主教\Person{默里托}和\Place{希拉波利斯}的主教\Person{阿波利纳里}也享有盛名。他们各自代表信德向当时在位的那位罗马皇帝呈递了护教书。\par

2. 以下这些作者的著作已为我们所知。\Person{默里托}的著作：\WorkTitle{论逾越节}两卷，\WorkTitle{论生活行为与先知}一卷，\WorkTitle{论教会}讲演一篇，\WorkTitle{论主日}一卷，此外\WorkTitle{论人的信德}一卷，\WorkTitle{论他的创造}一卷，\WorkTitle{论信德的服从}另一卷，\WorkTitle{论感官}一卷；此外有\WorkTitle{论灵魂与身体}的著作，\WorkTitle{论圣洗}之作，\WorkTitle{论真理}之作，\WorkTitle{论基督的受造与受生}之作；他的讲演\WorkTitle{论预言}和\WorkTitle{论款待}；此外还有\WorkTitle{钥匙}以及\WorkTitle{论魔鬼与若望默示录}诸书，\WorkTitle{论天主的形体性}之作，最后是致\Person{安敦尼}的书。\par

3. 在\WorkTitle{论逾越节}诸书中，他指明了写作的时间，以这些话开始：\Quote{当\Person{塞维利乌斯·保禄}担任\Place{亚细亚}总督时，在\Person{萨加里斯}殉道的时候，\Place{劳狄刻雅}发生了一场关于逾越节的大争论，该节日按规则落在那些日子；这些书就是那时写的。}\par

4. \Person{亚历山大里的克肋孟}在他自己的\WorkTitle{论逾越节}讲演中引用了这部著作，他说，他是因\Person{默里托}的著作而写的。\par

5. 但他在致皇帝的书信中记载了在他治下发生在吾人身上的事：\Quote{因为，前所未有地，虔敬的族类如今正遭受迫害，在亚细亚被新法令驱赶。无耻的告密者和觊觎他人财产之人，借法令之机，公然日夜进行掠夺，剥夺无罪之人的财物。} 稍后他又说：\Quote{若这些事是奉你的命令而行，那也好。因为公正的统治者绝不会采取不义的措施；我们确实欣然接受如此死亡的光荣。}\par

6. 但我们只向你提出这一请求：愿你亲自首先审查这些纷争的肇始者，并公正判断他们究竟该死该罚，还是该得安全与安宁。但反之，若此计谋和这新法令（甚至对蛮族敌人都不可执行）并非出自你，我们更恳求你，不要让我们暴露于民众如此无法无天的掠夺之下。\par

7. 他又补充道：\Quote{因为我们的哲学从前在蛮族中昌盛；但它在你的先祖\Person{奥古斯都}的伟大统治期间，在您治下的各民族中兴起，对你的帝国而言，它尤其成为吉兆的祝福。因为从那时起，罗马人的力量在伟大和辉煌中增长。你作为所渴望的继承者继承了这一力量，并且你将与你的儿子一同如此继续，只要你守护那与帝国一同成长、与奥古斯都一同出现的哲学；这一哲学你的先祖也与其他宗教一同尊崇。}\par

8. 最令人信服的证据表明，我们的教义为一个顺利开始的帝国之益处而昌盛，即：自奥古斯都统治以来，没有发生任何灾祸，反而万事辉煌荣耀，正如所有人的祈祷所愿。\par

9. 只有\Person{尼禄}和\Person{多米提安}，受某些诽谤者怂恿，曾想诋毁我们的教义；由此，这谎言得以流传，这是由于一种不合理的习俗，即针对基督徒提起诽谤性控告。\par

10. 但你们虔敬的先父纠正了他们的无知，多次以书面斥责那些胆敢对他们采取新措施的人。其中，你的祖父\Person{哈德良}似乎曾写信给许多人，也给亚细亚行省总督\Person{丰达努斯}。而你的父亲，当时你也与他一同统治，曾写信给各城，禁止他们对我们采取任何新措施；其中包括拉里萨人、得撒洛尼人、雅典人以及所有希腊人。\par

11. 至于你——既然你对基督徒的意见与他们相同，甚至更为仁慈和富哲学智慧——我们更加确信你会应允我们所求的一切。这些话见于上述著作。\par

12. 但在同一作者所著的\WorkTitle{《摘录》}中，他在引言开头给出了旧约公认书卷的目录，有必要在此引用。他写道：\par

13. \Person{梅利托}致其兄弟\Person{敖乃息摩}，安好：由于你既时常因热爱圣言，表达希望从法律和先知中摘录关于救主以及我们全部信仰的经文，又希望获得古代书卷数量与顺序的准确陈述，我努力成就此事，深知你对信仰的热忱、你对认识圣言的渴望，也知道你在渴慕天主中将这一切置于一切之上，奋力争取永恒的救恩。\par

14. 因此，当我前往东方，来到这些事被宣讲和成就的地方时，我准确得知了旧约的书卷，并将其名录如下寄给你。其名称如下：梅瑟五书：创世纪、出谷纪、户籍纪、肋未纪、申命纪；\BibleBook{若苏厄}{书}、\BibleBook{民长}{纪}、\BibleBook{卢德}{传}；列王纪四卷；编年记两卷；\BibleBook{圣咏}{集}、\BibleBook{箴言}{篇}，又有\BibleBook{智慧}{篇}、\BibleBook{训道}{篇}、\BibleBook{}{雅歌}、\BibleBook{约伯}{传}；先知书：\BibleBook{依撒意亚}{}、\BibleBook{耶肋米亚}{}；\BibleBook{}{十二先知书}一卷；\BibleBook{达尼尔}{}、\BibleBook{厄则克耳}{}、\BibleBook{厄斯德拉}{}。我从其中做了摘录，分为六卷。以上为\Person{梅利托}的话。\par

% structure: p0231 source=h2 target=subsection reason=relative-heading-rank; base=h2

\subsection{第二十七章  希拉波利教会主教阿波利纳里}

许多\Person{阿波利纳里}的著作被多人保存，以下为我们所获：致上述皇帝的演说，五卷\WorkTitle{《驳希腊人》}，\WorkTitle{《论真理》}第一、二卷，以及其后所著反对弗里吉亚人异端（此异端不久后以其新说出现，当时尚处于萌芽状态，因为\Person{孟他努}与其伪女先知正奠定其错误的基础）的著作。\par

% structure: p0233 source=h2 target=subsection reason=relative-heading-rank; base=h2

\subsection{第二十八章  穆萨努斯及其著作}

至于我们先前提到的\Person{穆萨努斯}，存有一篇极为优美的论述，是他写给某些转投所谓恩克拉提派异端的弟兄的。此异端最近兴起，引入了陌生而有害的错误。据说\Person{塔提安}是这一错误教义的首倡者。\par

% structure: p0235 source=h2 target=subsection reason=relative-heading-rank; base=h2

\subsection{第二十九章  塔提安的异端}

1. 他就是我们刚才就那位杰出人物\Person{犹斯定}所引用的话中提到的那一位，我们曾说他是一位殉道者的门徒。\Person{依勒内}在其著作\WorkTitle{《驳斥异端》}第一卷中证实了这一点，他在那里关于他和他的异端写道：\par

2. 那些被称为恩克拉提派、源自\Person{撒图尔尼努斯}和\Person{马尔基昂}的人，宣讲独身生活，废弃天主原本的安排，并暗中指责那位造男造女以繁衍人类的天主。他们还引入戒食被称为有生命之物，从而对创造万有的天主忘恩负义。他们也否认第一个人的救恩。\par

3. 但这只是他们最近才发现的，某个叫\Person{塔提安}的人首先引入了这种亵渎。他曾是\Person{犹斯定}的听众，与\Person{犹斯定}在一起时并未表达过这样的意见，但在后者殉道后，他离开了教会，自高自大，以为自己是教师，并自认优越于他人，便建立了自己独特的教义类型，捏造了一些无形的\OriginalTerm{永世}{æons}，如同\Person{瓦伦提努斯}的追随者一样；同时，如同\Person{马尔基翁}和\Person{萨图尔尼努斯}，他宣称婚姻是败坏和奸淫。然而，关于亚当是否得救的论点，则是他自己杜撰的。\Person{依勒内}当时这样写道。\par

4. 但稍后，一个名叫\Person{塞维鲁}的人给上述异端注入了新力量，从而使得那些源于此异端的人，根据他被称为\OriginalTerm{塞维鲁派}{Severians}。\par

5. 他们确实使用《法律》和《先知书》以及《福音书》，但以他们自己的方式解释圣经中的言辞。他们还辱骂宗徒\Person{保罗}，拒绝他的书信，甚至不接受《\WorkTitle{宗徒大事录}》。\par

6. 但他们的最初创立者\Person{塔提安}，不知如何，将福音书编排并汇集成一册，称之为\OriginalTerm{四史合编}{Diatessaron}，这册书至今仍存于某些人手中。但有人说，他竟敢改写宗徒的某些言语，以改良其风格。\par

7. 他留下了大量著作。其中许多人最常使用的是他那著名的《\WorkTitle{致希腊人书}》，这似乎也是他所有作品中最好、最有用的。在这部书中，他论述了最古老的时代，并指出\Person{梅瑟}和希伯来先知比希腊所有著名人物都要早。关于这些人就说这么多。\par

% structure: p0243 source=h2 target=subsection reason=relative-heading-rank; base=h2

\subsection{第三十章 叙利亚人\Person{巴尔德撒尼斯}及其现存著作。}

1. 在同一朝代，当异端在两河流域泛滥时，有一个名叫\Person{巴尔德撒尼斯}的人，他是一位非常能干的人，精通叙利亚语，擅长辩论。他撰写了对话录，反对\Person{马尔基翁}的追随者以及其他一些持有各种意见的人，并用他自己的语言将这些对话录连同许多其他著作写了下来。他的门徒众多（因为他是一位强大的信仰捍卫者），将这些作品从叙利亚语翻译成了希腊语。\par

2. 其中还有他那篇非常有能力的对话录《\WorkTitle{论命运}》，致\Person{安托尼努斯}，以及据说他因当时发生的迫害而撰写的其他著作。\par

3. 他起初确实是\Person{瓦伦提努斯}的追随者，但后来，他摒弃了\Person{瓦伦提努斯}的教导，驳斥了其大部分虚构之说，自以为已转向了更正确的意见。然而，他并未完全洗掉旧异端的污秽。\par

大约在这时，罗马教会的主教\Person{索特尔}也去世了。\par

\SourceRecord{Translated by Arthur Cushman McGiffert.；Nicene and Post-Nicene Fathers, Second Series；Vol. 1.；Edited by Philip Schaff and Henry Wace.；Buffalo, NY: Christian Literature Publishing Co.,；1890.；<http://www.newadvent.org/fathers/250104.htm>.}

% structure: p0001 source=h1 target=section reason=page-title

\section{《教会史（第五卷）》}

% structure: p0002 source=h2 target=subsection reason=relative-heading-rank; base=h2

\subsection{引言。}

1. 罗马教会的主教索特尔，在任职八年后逝世，由宗徒以来的第十二位主教埃莱乌泰里乌斯继任。在皇帝安敦尼·维鲁斯在位的第十七年，由于城市中民众的骚动，对我们的子民的迫害在某些地区更猛烈地重新燃起；仅从一个民族的人数来看，就有数以万计的人在世界各地殉道。这一记载为后世而写，确实值得永远铭记。\par

2. 我们编纂的殉道者列传集提供了关于该主题最可靠信息的完整记述，它既是一部具有教育意义的历史叙事。我将在此重复该记述中为当前目的所必需的部分。\par

3. 其他历史学家记录战争的胜利、从敌人那里获得的战利品、将领的谋略和士兵的勇敢，这些沾满了鲜血和无数杀戮，为的是子女、国家和其它财产。\par

4. 但我们对天主治理的叙述，将以不可磨灭的文字记录为灵魂的和平而进行的最和平的战争，并将讲述那些为真理而非祖国、为虔诚而非最亲爱的朋友而勇敢行事的人。它将把宗教运动员的纪律和备受考验的坚韧、从魔鬼那里获得的战利品、对无形敌人的胜利以及加在他们所有人头上的冠冕，永远铭记。\par

% structure: p0007 source=h2 target=subsection reason=relative-heading-rank; base=h2

\subsection{第一章 在维鲁斯统治下高卢为宗教而战的人数及其冲突的性质。}

1. 为她们准备竞技场的国家是高卢，其中里昂和维也纳是主要和最著名的城市。罗讷河流经这两座城市，以宽阔的水流贯穿整个地区。\par

2. 该国最著名的教会向亚细亚和弗里吉亚的教会发送了关于这些见证者的记述，以下述方式叙述了她们中间所发生的事。\par

我将引用她们的原话。\par

3. \Quote{在高卢的维也纳和里昂的基督的仆人，致亚细亚和弗里吉亚各地拥有同一信仰和救赎希望的弟兄们：愿天主圣父和我们的主耶稣基督赐予平安、恩宠和荣耀。}\par

4. 在叙述了其他一些事情之后，她们这样开始她们的记述：此地的巨大磨难、外邦人对圣人的愤怒以及蒙福的见证者的苦难，我们不能精确地叙述，甚至根本无法记录。\par

5. 因为那仇敌竭尽全力攻击我们，让我们预先尝到它未来无节制的活动。它想方设法训练和操练它的仆人来反对天主的仆人，不仅将我们赶出房屋、浴场和市场，而且禁止我们在任何地方出现。\par

6. 但天主的恩宠带领了对抗它的战斗，拯救了弱者，将他们立为坚固的柱子，能够以忍耐承受那恶者的一切愤怒。他们与它交战，忍受各种羞辱和伤害；视自己的巨大苦难为小事，他们奔向基督，真正彰显出\BibleQuote{现今的苦楚，不值得我们后来所将要显示给我们的光荣}\BibleRef{罗 8:18}。\par

7. 首先，他们勇敢地忍受了群众加给他们的伤害：喧嚷、殴打、拖曳、抢劫、石击、监禁，以及暴怒的群众乐于对敌人和对手施加的一切。\par

8. 然后，他们被千夫长和城里的官长带到广场，在全众面前受审，承认之后，被关押直到总督到来。\par

9. 后来，当他们被带到总督面前，他以最残酷的方式对待我们时，一位弟兄维提乌斯·埃帕加图斯，一个充满对天主和邻人之爱的人，出面干预。他的生活如此一贯，以致虽然年轻，却已获得了与年长的匝加利亚同等的声誉：因为他\BibleQuote{遵行主的一切诫命和典章，无可指摘}\BibleRef{路 1:6}，并为了邻人而不懈地行一切善工，对天主热忱，在心神上火热。这样的品性，使他不能忍受对我不合理的判决，而是充满义愤，请求允许他为他的弟兄们作证，证明我们中间没有任何不敬或亵渎的事。\par

10. 但审判席周围的人向他呼喊，因为他是一个显要人物；总督拒绝了他正当的请求，只是问他是否也是基督徒。他大声承认后，自己也被列入了见证者的行列，被称为\OriginalTerm{基督徒的辩护者}{Advocate of the Christians}，但他自己里面有那辩护者，即圣神，比匝加利亚更丰富。他以此表明了他爱的满全，甚至乐于为保卫弟兄们而舍弃生命。因为他过去和现在都是基督的真门徒，\BibleQuote{羔羊无论到哪里去，他都跟随}\BibleRef{默 14:4}。\par

11. 于是其余的人被分开了，初期的见证者们显然已经预备好，并极其热切地完成了他们的宣认。但有些人显得没有预备好和未受训练，仍然软弱，无法承受如此大的冲突。大约有十个人成了流产者，给我们带来了无法估量的巨大忧伤和悲痛，并削弱了其他尚未被逮捕者的热忱，他们虽然忍受各种苦难，却始终与见证者同在，没有离弃他们。\par

12. 那时我们所有人都非常惧怕，因为对他们的宣认不确定；不是因为我们害怕将要忍受的苦难，而是因为我们着眼结局，害怕他们中有些人会跌倒。\par

13. 但那些堪当的人一天天被逮捕，凑满了他们的数目，以致所有热忱的人，尤其是那些使我们事务得以建立的人，都从两座教会中聚集在一起。\par

14. 我们的一些外邦奴仆也被抓了，因为总督下令公开审讯我们所有人。他们被撒殚诱惑，害怕自己遭受他们所见圣人忍受的酷刑，又受到士兵的煽动，便诬告我们举行提厄斯忒斯的宴会和俄狄浦斯的交合，以及那些不仅我们不可言说或思考，甚至无法相信曾有人做过的事。\par

15. 这些指控被报告后，全体民众如同野兽般向我们发怒，以致即使有人先前因友情而温和，现在也极其愤怒，向我们咬牙切齿。我们主所说的话应验了：\BibleQuote{时候将到，凡杀你们的，就以为是侍奉天主。}\BibleRef{若 16:2}\par

16. 最后，圣洁的证人忍受了难以言喻的痛苦，撒殚竭力想让他们也说出一些诽谤之言。\par

17. 但全体民众、总督和士兵的愤怒都极其猛烈地转向\Person{Sanctus}，从\Place{Vienne}来的执事，\Person{Maturus}，一位新归信者，却是一位高贵的战士，以及\Person{Attalus}，来自\Place{Pergamos}本地人，他在那里一直是柱石和根基，还有\Person{Blandina}，通过她，基督显明了那些在人看来卑微、渺小、可鄙的事物，在天主眼中却因向他所显的爱和能力、不凭外貌夸耀，乃是极大的荣耀。\BibleRef{格前 1:27}\par

18. 当我们所有人都战栗时，她在地上的主母（她自己也是证人之一）担心她因身体软弱而无法大胆宣认，但\Person{Blandina}被充满了如此大的能力，以致她从早晨到傍晚轮流折磨她的那些人中解脱出来，超越了他们，他们承认自己被征服了，对她无能为力。他们对她的耐力感到惊讶，因为她的全身被撕裂、破碎；他们作证说，其中任何一种酷刑都足以致命，更不用说如此众多而巨大的痛苦了。\par

19. 但这蒙福的妇人，如同高贵的运动员，在宣认中更新了力量；她的安慰、消遣和从痛苦中解脱，在于呼喊说：\Quote{我是基督徒，我们并没有做什么卑劣的事。}\par

20. \Person{Sanctus}也奇妙地、超乎常人地忍受了他所受的一切凌辱。当恶人希望借着持续和严酷的折磨从他口中逼出一些不该说的话时，他以如此坚定对抗他们，以致他甚至不肯说出自己的名字、所属的民族或城市，或他是为奴还是自由，而是用拉丁语回答他们所有的问题：\Quote{我是基督徒。}他以此代替名字、城市、种族及其他一切来宣认，人们从他口中听不到别的话。\par

21. 因此，总督和折磨他的人产生了极大的愿望要征服他；但既然对他无计可施，他们最后将烧红的铜板贴在他身体最柔嫩的部分。\par

22. 这些部位确实被烧焦了，但他依然坚定不移，毫不屈服，坚定于他的宣认，并由从基督怀中涌出的天上生命之泉所更新和刚强。\par

23. 他的身体是他受苦的见证，完全是一处创伤和伤痕，扭曲变形，完全不像人形。基督在他内受苦，彰显了他的荣耀，从他仇敌手中解救了他，使他成为他人的榜样，表明哪里有天父的爱，就没有什么可畏惧的；哪里有基督的荣耀，就没有什么痛苦。\par

24. 因为当恶人几天后第二次折磨他时，以为他的身体已经肿胀发炎到连手碰触都不能忍受的程度，如果他们再次使用同样的刑具，就会制伏他，或者至少因他受苦而死，使其他人害怕。结果不仅没有发生，反而出乎所有人预料，他的身体在后续的折磨中起来直立，恢复了原来的形状和肢体的功能，以致因基督的恩宠，这第二次的痛苦对他而言，不是折磨，而是医治。\par

25. 但魔鬼以为它已经吞噬了\Person{Biblias}（她曾是那些否认基督者之一），想要藉着她说出亵渎的话来加重对她的定罪，便再次将她带到刑讯前，强迫她——已虚弱无力——报告关于我们的不敬之事。\par

26. 但她却在痛苦中恢复了意识，仿佛从沉睡中醒来，并被当前的痛苦提醒地狱中的永罚，她反驳了那些亵渎者。她说：\Quote{那些连无理性动物的血都认为不可尝的人，怎能吃孩子呢？}从此她宣认自己是基督徒，并被列入证人的行列。\par

27. 但由于基督藉着蒙福者的忍耐使暴虐的酷刑无效，魔鬼便发明了其他手段——囚禁在监狱中最黑暗、最污秽的地方，将脚拉开到木枷的第五个孔，以及其他暴怒、充满魔鬼的仆役们惯常对囚犯施加的凌辱。许多人在监狱中被闷死，他们被主选为这种死亡方式，为要在他们身上彰显他的荣耀。\par

28. 因为有些人虽然被如此残酷地折磨，似乎即使最细心的照料也难以存活，但他们无人照料，留在监狱中，却被主坚固，身体和灵魂都得到力量；他们鼓励和劝勉其余的人。但那些年轻、最近被捕、身体尚未适应酷刑的人，无法忍受监禁的严酷，便死在狱中。\par

29. 真福\Person{波蒂努斯}受委托管理\Place{里昂}教区，被拖到审判席。他已逾九十高龄，身体极其虚弱，因体力衰弱几乎无法呼吸；但因对殉道的热切渴望，他在属神的热忱中得着力量。虽然他的身体因年老和疾病而衰败，他的生命得以保全，为使\Person{基督}在其中得胜。\par

30. 当兵士将他带到法庭，同时有民政官和一大群人随同，他们以各种方式向他喊叫，仿佛他就是\Person{基督}自己，他作了崇高的见证。\par

31. 总督问他基督徒的天主是谁，他回答说：\Quote{你若有资格，就会知道。}然后他被粗暴地拖走，遭受各种打击。近旁的人用手脚打他，不顾他的年纪；远处的人则抓住什么就向他投掷什么；他们都认为，若省略了任何可能的凌辱，自己就犯了极大的邪恶和不敬。因为他们以为这样就能为自己的神明复仇。他几乎无法呼吸，被投入监狱，两天后去世。\par

32. 那时，发生了一件天主伟大的上智安排，\Person{耶稣}的怜悯显现得超乎度量，其方式在弟兄中罕见，但并未超越\Person{基督}的能力。\par

33. 因为在初次被捕时背弃信仰的人，与其他一同被囚禁，并忍受可怕的痛苦，以致他们的否认连眼前也毫无益处。但那些承认自己身份的人，作为基督徒被囚禁，没有别的指控加于他们。然而，前者后来被当作杀人犯和污秽之人对待，受到的惩罚是其他人的两倍。\par

34. 因为殉道的喜乐、对预许的望德、对\Person{基督}的爱德以及圣父的圣神支持着后者；但前者因良心极度痛苦，以致他们被带出时，仅凭面容就极易与其余的人区分。\par

35. 因为前者怀着喜乐出去，光荣与恩宠交织在他们的面容上，以致他们的锁链也像美丽的装饰，如同新娘佩戴着各色金穗；他们散发着\Person{基督}的馨香，以致有人以为他们被涂了世俗的香膏。但后者垂头丧气、谦卑沮丧，充满各种耻辱，并被外教人斥为卑贱懦弱，背负杀人犯的指控，失去了那唯一可敬、光荣与赋予生命的圣名。其余的人看到这些，就得了力量，在被捕时毫不犹豫地承认，毫不理会魔鬼的劝诱。\par

36. 在说了某些其他话之后，他们继续说道：\par

此后，最终，他们的殉道被分为各种形式。因为他们编织了一顶各色各样的花朵制成的冠冕，将它献给圣父。因此，高贵的运动员在经历了多重争战并光荣得胜后，理当领受那伟大而不朽的冠冕。\par

37. 因此，\Person{马图鲁斯}、\Person{桑克图斯}、\Person{布兰迪娜}和\Person{阿塔鲁斯}被带到竞技场，暴露给野兽，为外邦公众提供一场残酷的表演，一个斗兽的日子特地因我们的人而设立。\par

38. \Person{马图鲁斯}和\Person{桑克图斯}两人在竞技场中再次经历各种酷刑，仿佛他们先前没有受过任何苦，或者更确切地说，仿佛已在多次竞赛中战胜了对手，如今正为冠冕本身而奋斗。他们再次忍受了惯常的夹道鞭打和野兽的猛烈攻击，以及愤怒的人群所呼求或渴望的一切，最后是铁椅，他们的身体在其中被烤炙，以烟雾折磨他们。\par

39. 迫害者并不以此罢休，反而更加疯狂地对付他们，决心要挫败他们的忍耐。但即便如此，他们从\Person{桑克图斯}口中没有听到任何话，除了他从起初所宣认的信仰。\par

40. 这些人，在生命因着伟大的斗争而延续了很长一段时间之后，终于被献祭，在那一天成了世界的一场景象，代替了通常的各类战斗。\par

41. 但\Person{布兰迪娜}被悬挂在木桩上，暴露给将要攻击她的野兽吞食。因为她看上去仿佛挂在十字架上，又因她恳切的祈祷，她激发了战斗者们的极大热忱。因为他们注视着她战斗，用肉眼在他们姐妹的身上看见了为他们被钉在十字架上的那一位，为要说服那些信从他的人：凡为\Person{基督}的光荣而受苦的人，永远与永生天主共融。\par

42. 因当时没有一只野兽触碰她，她就被从木桩上取下，再次投入监狱。她这样被保留下来，为另一场竞赛，以便在更多冲突中得胜后，使那弯曲之蛇的惩罚不可逆转；并且，虽然她渺小、软弱、受人轻视，却身穿\Person{基督}这位大能得胜的运动员，可能激发弟兄们的热忱，并在多次战胜仇敌后，借着自己的斗争领受那不朽的冠冕。\par

43. 但\Person{阿塔鲁斯}被群众高声呼喊，因为他是一位显赫人物。他因着良心的无愧和在基督徒操练中的真诚实践，欣然投入竞赛，并且他一直是我们中间真理的见证者。\par

44. 他被带着绕竞技场一周，前面举着一块牌子，上面用罗马文字写着\Quote{此乃阿塔鲁斯，基督徒}，民众对他充满愤怒。但当总督得知他是罗马人时，便下令将他与其他在狱中的人一同带回去，关于这些人他已写信给\Person{凯撒}，正在等待回音。\par

45. 但中间的时间对他们并非虚度或徒然；因为藉着他们的忍耐，基督无量的慈悲彰显了出来。因为藉着他们的继续存活，死者得以复生，见证人向那些未能见证的人施予了恩惠。那童贞母亲在领受她曾作为死者所生的那些活人时，充满了极大的喜乐。\par

46. 因为藉着他们的影响，许多曾背弃的人得以复原，重获新生，重新燃起生命之火，并学会了宣认。当他们活过来并坚强起来后，他们走到审判台前，再次接受总督的讯问；天主不愿罪人死亡\BibleRef{则 33:11}，却仁慈地召叫他们悔改，以仁慈对待他们。\par

47. 因为凯撒命令将他们处死，但任何背弃者均可获释。因此，在当地的公共节庆开始时，来自各民族的人群聚集，总督将那些蒙福者带到审判台前，将他们作为展示和景观给众人看。因此他又审问他们，将那些看来拥有罗马公民身份的人斩首，其余的人则投入野兽。\par

48. 而基督在那些从前背弃他的人身上大受光荣，因为出乎外邦人的意料，他们宣认了。因为他们被单独审问，似乎将要获释；但宣认后，他们被列入见证者的行列。但有些人仍留在外面，他们从未有过一丝信德的痕迹，也未理解婚宴礼服\BibleRef{玛 22:11}，更未领悟对天主的敬畏；反而作为丧亡之子，以他们的背教亵渎了这道。\par

49. 但其余的人都加入了教会。在他们受审时，有一位亚历山大，本为弗里吉亚人，职业是医生，曾在高卢居住多年，因他对天主的爱和言论的大胆而为众人所熟知（因为他并非没有宗徒恩宠的一份），站在审判台前，用动作鼓励他们宣认，在旁观者看来，他仿佛在经历产痛。\par

50. 但众人因那些从前背弃的人现在宣认而大怒，向亚历山大呼喊，仿佛他是这一切的原因。于是总督召见他，问他是什么人。当他回答说自己是一名基督徒时，总督非常愤怒，判他投入野兽。翌日，他与阿塔鲁斯一同进入斗兽场。因为为了取悦民众，总督下令将阿塔鲁斯再次投入野兽。\par

51. 他们在圆形剧场中受尽为此目的而设计的一切刑具的折磨，历经极大的搏斗后，终于被祭杀。亚历山大既没有呻吟也没有抱怨任何事，而是在心中与天主交谈。\par

52. 但当阿塔鲁斯被放置在铁椅上，他燃烧的身体发出的烟雾升起时，他用罗马语对民众说：\Quote{看！你们所做的这事是吞噬人；但我们不吞噬人，也不做任何其他邪恶的事。}当被问及天主叫什么名字时，他回答说：\Quote{天主不像人那样有名字。}\par

53. 在所有这些之后，在竞赛的最后一天，布兰迪娜再次被带进来，同来的还有蓬提库斯，一个约十五岁的男孩。他们每天都被带来观看其他人的受苦，并被逼迫向偶像起誓。但因他们坚定不移且鄙视偶像，民众变得狂暴，以致他们对男孩的年龄毫无怜悯，对女子的性别也不尊重。\par

54. 因此，他们让这两个人遭受所有可怕的痛苦，使他们历经全部折磨，反复逼迫他们起誓，但未能成功；因为蓬提库斯受到他姐姐的鼓励，连外邦人都看出她在坚定和加强他，他高贵地忍受了每一种折磨后，交出了灵魂。\par

55. 但蒙福的布兰迪娜最后，如同一位高贵的母亲，鼓励了她的儿女并使他们胜利地先行到达君王面前，自己承受了他们所有的搏斗，随后赶上他们，为她的离去而喜乐欢欣，仿佛被召赴婚宴，而不是被投入野兽。\par

56. 在鞭打、野兽、烤椅之后，她最终被裹在网中，扔到一头公牛面前。她被那动物抛来抛去，但因她的望德和对所交托之事的坚握以及与基督的共融，她对发生在自己身上的事毫无感觉，她也被祭杀了。连外邦人也承认，他们中间从未有一个女子忍受过如此众多而可怕的折磨。\par

57. 但即使如此，他们对圣者的疯狂和残忍仍未满足。因为被那野兽\OriginalTerm{野兽}{Bestia}煽动，野蛮的部落并不轻易平息，他们的暴力又在尸体上找到了另一个特有的机会。\par

58. 因为他们缺乏理性的男子气概，被征服的事实并未使他们羞耻，反而更加点燃了他们的怒火，如同野兽一般，激起总督和民众的仇恨，不公正地对待我们；好使圣经得以应验：\Quote{那不义的，叫他仍旧不义；那正义的，叫他仍旧正义。}\par

59. 因为他们将那些在狱中窒息而死的人扔给狗，日夜严密看守，免得有人被我们埋葬。他们将野兽和火留下的残骸，残缺不全、烧焦的，暴露在外，并将其他人的头与身体放在一起，同样由士兵看守多日，不许埋葬。\par

60. 有些人向他们发怒，咬牙切齿，想要对他们施行更严厉的报复；另一些人则嘲笑讥讽他们，夸耀自己的偶像，将基督徒所受的惩罚归咎于他们。甚至那些较为理智、似乎有些同情的人，也常常责问他们，说：\Quote{他们的天主在哪里？他们所选择的超过生命的宗教，对他们有什么益处呢？}\par

61. 他们对我们如此行径各异；但我们深陷悲痛，因为无法埋葬尸体。黑夜未能助我们成事，金钱无法说服，恳求也未能打动怜悯之心；他们用尽各种方式看守，仿佛阻止埋葬对他们有莫大的好处。\par

此外，他们在其他话之后又说：\par

62. 殉道者的尸体就这样以各种方式被展示并暴露了六天，之后被恶人焚烧成灰，扫入\Place{罗讷河}，使地上不留丝毫痕迹。\par

63. 他们这样做，仿佛能战胜\Person{天主}，阻止他们重生；\Quote{他们说：\InnerQuote{这样，他们便毫无复活的盼望，而正是因着这盼望，他们向我们传入这外来的新宗教，轻视可怕之事，甚至乐意赴死。现在让我们看看他们是否真能复活，他们的\Person{天主}是否能帮助他们，并把他们从我们手中救出。}}\par

% structure: p0074 source=h2 target=subsection reason=relative-heading-rank; base=h2

\subsection{第二章 蒙天主钟爱的殉道者仁慈地服事在迫害中跌倒的人}

1. 在上述皇帝统治下，基督的教会遭遇了这些事；由此我们可以合理推测其他省份发生的情况。适宜从同一封信中再增选几段，其中记录了这些见证人的温良与慈悲，措辞如下：\par

2. 他们如此热切地效法基督——\BibleQuote{他虽具有天主的形体，却没有以自己与天主同等为应得之物}\BibleRef{斐 2:6}——尽管他们已获得如此荣耀，并曾不止一次，而是多次作证——从野兽口中被带回监狱，满身烧伤、疤痕与伤口——却仍不自称为见证人，也不许我们以此称呼他们。若我们中有人在书信或谈话中称他们为见证人，他们便严厉责备他。\par

3. 因为他们欣然将见证人的称号归于基督——\BibleQuote{忠信而真实的见证}\BibleRef{默 3:14}，\BibleQuote{死者中的首生者}\BibleRef{默 1:5}，以及天主生命的元首；他们提醒我们那些已离世的见证人，说：\Quote{那些已在宣认中被基督认为配得接去，并以离世印证其见证的人，他们已是见证人；但我们只是卑微谦逊的宣认者。}他们含泪恳求弟兄们，应献上热切的祈祷，使他们得以达至圆满。\par

4. 他们以实际行动彰显了\OriginalTerm{见证}{testimony}的力量，对众弟兄显出极大的胆量，并以忍耐、无畏和勇气彰显其高贵，却拒绝接受见证人的称号，以免与弟兄们分离；他们满怀着对\Person{天主}的敬畏。\par

5. 稍后他们又说：\Quote{他们在全能的手下自谦自卑，如今却大受高举。他们为众人辩护，却不控告任何人；他们赦免众人，却不捆绑任何人；他们为那些虐待他们的人祈祷，正如完美的见证\Person{斯德望}说：\BibleQuote{主，不要把这项罪债归于他们。}\BibleRef{宗 7:60} 但如果他为用石头砸他的人祈祷，更何况为弟兄们呢！}\par

6. 之后在提到其他事后，他们又说：\par

\Quote{因着他们真诚的爱，他们与那野兽的最大争战，就是使他窒息，将他以为已吞下的人活着吐出来。因为他们不以跌倒者为耻，反而以自己所丰裕的帮助他们的需要，怀有母亲的慈心，在\Person{圣父}面前为他们流了许多眼泪。}\par

\Quote{他们求得了生命，\Person{天主}就赐予他们，他们便与邻人共享。战胜一切后，他们归往\Person{天主}那里。他们素来喜爱和平，并将和平托付给我们，便平安地归往\Person{天主}，没有给他们的母亲留下忧伤，也没有给弟兄们带来分裂或争执，而是喜乐、平安、和睦与爱。}\par

8. 这些有福之人对跌倒的弟兄所怀的深情记录，值得添上，因为此后有些人以不人道、不仁慈的态度，毫不留情地对待基督的肢体。\par

% structure: p0084 source=h2 target=subsection reason=relative-heading-rank; base=h2

\subsection{第三章 显现给见证人阿塔鲁的梦中异象}

1. 上述见证人的同一封信中，还包含另一件值得纪念的事。任何人不会反对我们将此事告知读者。\par

2. 信中说：有一位叫\Person{阿尔基比亚德}的，是他们中的一员，过着极为刻苦的生活，除了面包和水，什么都不吃。当他想在狱中继续这种生活方式时，在斗兽场第一次战斗后，异象启示给\Person{阿塔鲁}，说\Person{阿尔基比亚德}这样拒绝\Person{天主}所造的万物，并给他人设置绊脚石，并不妥当。\par

3. 于是\Person{阿尔基比亚德}听从了，毫无节制地享用万物，感谢\Person{天主}。因为他们并未被剥夺\Person{天主}的恩宠，反而有\Person{圣神}作他们的顾问。这些事就到此为止。\par

4. 此时，在\Place{夫里基亚}的\Person{蒙达诺}、\Person{阿尔基比亚德}和\Person{德奥多托}的追随者，开始广泛传播他们关于预言的臆测——因为藉着\Person{天主}的恩赐，在各教会中仍在发生的许多其他奇迹，使他们的预言轻易被许多人接受——当为此发生争论时，\Place{高卢}的弟兄们就此事发表了他们审慎且最正统的判断，并公布了若干在他们中间被处死的见证人所写的书信。他们尚在狱中时，就将这些书信寄给了\Place{亚细亚}和\Place{夫里基亚}各地的弟兄们，也寄给了当时\Place{罗马}的主教\Person{厄勒乌特鲁}，为教会的和平进行协商。\par

% structure: p0089 source=h2 target=subsection reason=relative-heading-rank; base=h2

\subsection{第四章 见证人在书信中赞扬依肋内}

1. 同一批见证人也向上述\Place{罗马}主教推荐了当时已是\Place{里昂}教会司铎的\Person{依肋内}，在信中说了许多赞美他的话，正如以下摘录所示：\par

\Quote{我们祈求，厄莱乌泰鲁斯神父，愿你在一切事上常常因天主而喜乐。我们已委托我们的弟兄与同伴依勒内将这封信带给你，我们请求你重视他，因为他热切于基督的盟约。因为如果我们认为职务能使人成义，我们应将他列在最先，作为教会的长老，这正是他的身份。}\par

我们何必抄写那封已提及的信函中所列的证人名单，其中有的被斩首，有的被抛给野兽，有的在狱中安眠，或列举当时仍活着的宣认者的人数？因为凡愿意的人，只要查阅这封信本身，便可轻易找到完整的记载，而这封信，如我所说，记录在我们的\WorkTitle{殉道者集}中。这就是在安敦尼治下所发生的事。\par

% structure: p0093 source=h2 target=subsection reason=relative-heading-rank; base=h2

\subsection{第五章　天主从天降雨为马尔谷·奥勒略·凯撒，回应我们百姓的祈祷。}

据说，安敦尼的兄弟马尔谷·奥勒略·凯撒，将要与日耳曼人和萨尔马提亚人作战时，因他的军队遭受干渴而极其苦恼。但那所谓的默利泰内军团的士兵，藉着那从那时至今赐予力量的信德，当他们列阵于敌人面前时，跪在地上，正如我们在祈祷中的习惯，并向天主发出恳求。\par

这对敌人而言确实是一个奇异的景象，但据说紧接着发生了一件更奇异的事。闪电驱散敌人使其溃败灭亡，而一场甘霖滋润了那支呼求天主的军队，他们全体都曾因干渴而濒临死亡。\par

这个故事由非基督徒作家讲述，他们乐于记述所提到的时代；它也被我们自己的百姓记录下来。那些与信仰无分的历史学家提到这奇事，却不承认这是对我们祈祷的回应。而我们自己的百姓，作为真理之友，以淳朴无华的方式叙述了此事。\par

在这些人中有阿波利纳里，他说从那时起，那一支因其祈祷而发生奇事的军团，从皇帝那里得到了一个与事件相称的称号，按罗马人的语言称为雷霆军团。\par

戴尔都良是这些事的可靠见证人。在他向罗马元老院呈递的\WorkTitle{为信仰辩护}一书中——我们已提及这部作品——他以更大更确凿的证据证实了这一历史。\par

他写道，最明智的皇帝马尔谷的书信仍然存在，他在信中作证说，他在日耳曼的军队曾因干渴而濒临死亡，却因基督徒的祈祷而得救。他还说，这位皇帝曾以死亡威胁那些控告我们的人。\par

他接着写道：\par

\Quote{那些不虔敬、不义、残忍的人，对我们单独使用的，是怎样的法律？这些法律，韦斯帕芗虽征服了犹太人，却未予理会；图拉真部分废除，禁止搜寻基督徒；哈德良，虽对一切事务都好奇，也未予批准；那位被称为庇护的，也未予以认可。} 但任凭各人随己意看待这些事；我们必须继续讲述下文。\par

波提努斯与其他在高卢的殉道者一同去世，享年九十岁，依勒内继承他担任里昂教会的主教。我们得知，他在年轻时曾是波利卡普的听众。\par

在他所著\WorkTitle{反异端}一书的第三卷中，他收录了一份罗马主教的名录，一直带到厄莱乌泰鲁斯（我们现正考虑他的时代），他在其任职期间撰写了这部著作。他写道：\par

% structure: p0104 source=h2 target=subsection reason=relative-heading-rank; base=h2

\subsection{第六章　罗马主教名录。}

有福的宗徒建立了教会并使其稳固后，将主教的职务托付给了理诺。保禄在致弟茂德的书信中提到了这位理诺。\BibleRef{弟后 4:21}\par

阿纳克莱图斯继承了他；在阿纳克莱图斯之后，从宗徒算起第三位，克肋孟接受了主教职。他曾亲眼见过并亲耳听过有福的宗徒，他们的宣讲仍在他耳中回响，他们的传统仍在他眼前。他并非独自如此，因为许多受过宗徒教导的人仍活着。\par

在克肋孟的时代，格林多的弟兄们之间发生了一场严重的纷争，罗马教会致信格林多人，一封非常合适的信函，使他们和平复和，更新他们的信德，并宣布那新近从宗徒们接受的道理。\par

稍后他又说：\par

厄瓦雷图斯继承克肋孟；亚历山大继承厄瓦雷图斯。然后西斯笃，从宗徒算起第六位，被任命。在他之后，泰勒斯福尔，他荣耀地殉道；然后希吉努斯；然后庇护；在他之后，安尼塞图斯；索特尔继承安尼塞图斯；而现在，从宗徒算起第十二位，厄莱乌泰鲁斯持有主教职务。\par

按照同一秩序与继承，教会中的传统与真理的宣讲已从宗徒传至我们。\par

% structure: p0111 source=h2 target=subsection reason=relative-heading-rank; base=h2

\subsection{第七章　直至那些时代，奇迹仍由信徒施行。}

依勒内同意我们已有的叙述，记录了这些事，在那部由五卷组成的著作中，他给它取名\WorkTitle{驳斥并推翻所谓伪知识}。在同一著作的第二卷中，他指出神圣与奇迹能力的显现一直延续到他的时代，在一些教会中。\par

他说：\par

\Quote{但他们远未能复活死人，如同主所复活的，以及宗徒藉祈祷所行的。并在弟兄中，常常，当因某种需要，我们整个教会以禁食和许多恳求祈求时，死者的灵魂返回，那人藉着圣人的祈祷得到复原。}\par

接着，在其他评论之后，他又说：\par

\Quote{如果他们要说得，连主也只是表面上行了这些事，我们将指引他们去看先知著作，并从其中向他指明，一切事都已这样预先论及他，并且严格地应验了；而且他唯独是天主子。因此，他的真门徒，从他那里领受恩宠，为其他人的益处，奉他的名行这些事，各人按从他那里所领受的恩赐。}\par

4. 他们中有些人确实而真正地驱逐魔鬼，以致那些从恶灵中被洁净的人常常相信并加入教会。另有人预知未来之事，并有神视和预言性的启示。还有人藉着覆手治愈病人，使他们恢复健康。并且，正如我们所说，甚至死人也曾复活，并与我们在一起多年。\par

5. 但我们何必多说？不可能尽数教会普世从天主在耶稣基督——祂在本丢·比拉多治下被钉十字架——的名下所领受的恩赐，她每日为外邦人的益处施行这些恩赐，从不欺骗任何人，也不为钱财而行。因为她从天主白白领受，也白白施予。\par

6. 同一位作者又在另一处写道：\par

\Quote{正如我们也听说教会中的许多弟兄拥有先知之恩，并藉着圣神说各种方言，揭露人心隐秘之事以助其益处，并宣告天主的奥秘。}\par

关于各种神恩仍存于相称者中间直到那时的事实，就谈这些。\par

% structure: p0122 source=h2 target=subsection reason=relative-heading-rank; base=h2

\subsection{第八章 \Person{依勒内}论圣经之言}

1. 既然在本著作的开端，我们应许在必要时引用古代教会长老和作家的言语，他们曾流传下关于正典书卷的传统，而\Person{依勒内}是其中之一，我们现在将引用他的言语，首先是他关于神圣福音所言：\par

2. \Person{玛窦}以希伯来语向希伯来人发表了他的福音，那时\Person{伯多禄}和\Person{保禄}正在\Place{罗马}宣讲并建立教会。\par

3. 他们离去后，\Person{玛尔谷}，\Person{伯多禄}的门徒和翻译，也将\Person{伯多禄}所宣讲的书面传给了我们；而\Person{路加}，\Person{保禄}的随从，将\Person{保禄}所宣告的福音记录成书。\par

4. 随后，\Person{若望}，主的门徒，他曾靠在主的胸膛上，在\Place{亚细亚}的\Place{厄弗所}居住时发表了福音。\par

5. 这些事是他上面提到的那部著作的第三卷中所述。在第五卷中，他谈及\BibleBook{若望}{默示录}和假基督名字的数目如下：\par

\Quote{既然这些事如此，这数目在所有公认的古抄本中均可找到，并且那些亲眼见过\Person{若望}的人也证实此事，而理性也教导我们，那兽的名字的数目，按希腊人的计算方式，在其字母中出现……}\par

6. 随后他论及同一事说：\par

\Quote{我们不敢对假基督的名字妄加断言。因为若有必要在现今清楚宣告他的名字，早该由那见过启示的人宣告了。因为这启示是在不久前，几乎在我们这一代，\Person{多米提安}统治末期被看见的。}\par

7. 这些是他关于\BibleBook{若望}{默示录}在上述著作中所言。他也提及\BibleBook{若望}{一书}，从中引用许多证据，同样也提及\BibleBook{伯多禄}{前书}。他不仅知道，也接纳\WorkTitle{牧者}，写道如下：\par

\Quote{圣经说得真好：\InnerQuote{首先要相信天主是惟一的，祂创造并完成了万有}，}\par

8. 他几乎逐字引用\BibleBook{}{智慧篇}的话，说：\Quote{看见天主产生不死，而不死使我们亲近天主。}他也提到某位宗徒长老的记述，未曾提其名，并给出他对圣经的解释。\par

9. 他又提及殉道者\Person{犹斯定}和\Person{依纳爵}，也引用他们的著作作证。此外，他应许要在专论中从\Person{马西昂}自己的著作驳斥他。\par

10. 关于七十贤士翻译默感圣经一事，请听他写下的言语：\par

\Quote{天主确实成了人，主亲自拯救了我们，赐予童贞女的标记；但不像某些人所说，他们竟敢将圣经译为\InnerQuote{看，一位年轻女子将怀孕生子}，如同\Place{厄弗所}的\Person{特奥多提翁}和\Place{本都}的\Person{阿基拉}——两人都是犹太皈依者——所译；厄彼翁派随从他们，说祂是由\Person{若瑟}所生。}\par

11. 稍后他补充说：\par

因为在罗马人建立帝国之前，马其顿人仍占据亚细亚时，拉古斯之子\Person{托勒密}渴望用所有杰出人物的著作装饰他在\Place{亚历山大里亚}建立的图书馆，便请求\Place{耶路撒冷}人民将他们的圣经译成希腊文。\par

12. 但由于他们当时臣服于马其顿人，他们派了七十位精通圣经和两种语言的长老到\Person{托勒密}那里。天主就这样成就了祂的旨意。\par

13. 但为了逐一试验他们，恐怕他们通过商议而隐瞒圣经的真理，便把他们彼此隔离，命令他们都写出相同的译文。他对所有书卷都这样做了。\par

14. 但当他们聚集在\Person{托勒密}面前，比较各自的译文时，天主受到光荣，圣经被承认为真正神圣的。因为他们所译的一切从头到尾都是相同的话、相同的词、相同的名称，以致外邦人看出圣经是由天主的默感而译成的。\par

15. 这对天主而言并非奇事。天主曾在\Person{拿步高}掳掠百姓、圣经被毁、犹太人七十年后返乡之后，又于波斯王\Person{阿塔薛西斯}时代，默感肋未支派的司祭\Person{厄斯德拉}，使他讲述所有先前先知的话，并将\Person{梅瑟}的律法恢复给百姓。\par

以上是\Person{依勒内}的话。\par

% structure: p0144 source=h2 target=subsection reason=relative-heading-rank; base=h2

\subsection{第九章 \Person{科莫多斯}治下的主教}

1. \Person{安敦宁}作皇帝十九年后，\Person{科莫多斯}接掌政权。在他登基的第一年，\Person{犹利安}成为\Place{亚历山大里亚}教会的主教，此前\Person{阿格里平}已任职十二年。\par

% structure: p0146 source=h2 target=subsection reason=relative-heading-rank; base=h2

\subsection{第十章 哲学家\Person{庞代努斯}}

1. 大约在那时，一位学识极为卓越的人——\Person{Pantænus}——在\Place{Alexandria}负责信友的学校。那里自古就建立了一所神圣学问的学校，延续至今；据我们所知，该校曾由一些极具才能且热忱于神圣事物的人管理。据说，在这些人中，\Person{Pantænus}当时尤为突出，因为他曾受教于那些被称为\OriginalTerm{斯多葛派}{Stoics}的哲学体系。\par

2. 据说，他对\OriginalTerm{圣言}{Divine Word}表现出如此的热忱，以致被任命为基督福音的宣讲者，前往东方各民族，甚至被派遣至\Place{India}。因为当时仍有许多圣言的福音传播者，他们效法宗徒的榜样，热忱地致力于增进和建立\OriginalTerm{圣言}{Divine Word}。\par

3. \Person{Pantænus}便是其中之一，据说他曾前往\Place{India}。据报道，他在那里遇到一些认识基督的人，并发现了\BibleBook{玛窦}{福音}，这部福音书在他到达之前就已存在。因为宗徒\Person{Bartholomew}曾向他们宣讲过，并将\Person{Matthew}用\OriginalTerm{希伯来文}{Hebrew}写成的著作留给了他们，他们一直保存到那时。\par

4. \Person{Pantænus}在完成了许多善行之后，最终成为\Place{Alexandria}学校的负责人，并以口传和著作阐释神圣教义的宝藏。\par

% structure: p0151 source=h2 target=subsection reason=relative-heading-rank; base=h2

\subsection{第十一章　\Person{Clement}（\Place{Alexandria}）}

1. 此时，\Person{Clement}在\Place{Alexandria}与他一同受训于圣经，因而声名远扬。他与古代那位罗马教会负责人——一位宗徒的门徒——同名。\par

2. 在\WorkTitle{Hypotyposes}中，他明确提到\Person{Pantænus}是他的老师。在我看来，他在\WorkTitle{Stromata}的第一卷中也暗示了同一个人，当时他提到他所遇见的一些较杰出的宗徒继承者时，说道：\par

3. \Quote{这部著作并非为炫耀而精心构筑的作品；而是我为老年积存的笔记，作为对抗遗忘的良药；这是一幅未经雕琢的图像，一个粗浅的草图，记录了我有幸听到的那些充满力量与活力的话语，以及那些有福且真正卓越的人物。}\par

4. \Quote{其中一位——爱奥尼亚人——在\Place{Greece}，另一位在\Place{Magna Græcia}；一位来自\Place{Cœle-Syria}，另一位来自\Place{Egypt}。还有人在东方，一位是\OriginalTerm{亚述人}{Assyrian}，另一位是\Place{Palestine}的\OriginalTerm{希伯来人}{Hebrew}。但当我遇见最后一位——论能力他实为第一——我在\Place{Egypt}将他从隐居处搜寻出来，便得到了安息。}\par

5. \Quote{这些人保存了真福教义的真传，直接来自圣宗徒\Person{Peter}、\Person{James}、\Person{John}和\Person{Paul}，子从父领受（但像父亲那样的人很少），他们因天主的旨意，将这些祖传和宗徒的种子传递到我们这里。}\par

% structure: p0157 source=h2 target=subsection reason=relative-heading-rank; base=h2

\subsection{第十二章　\Place{Jerusalem}的主教们}

1. 当时，\Person{Narcissus}是\Place{Jerusalem}教会的主教，他至今仍为许多人称颂。自\Person{Adrian}统治下犹太人战败以来，他是第十五位继任者。我们曾指出，自那时起，\Place{Jerusalem}教会首先由外邦人组成，继受割礼者之后，而\Person{Marcus}是第一位领导他们的外邦人主教。\par

2. 此后，主教职位的继承顺序为：首先是\Person{Cassianus}；其后是\Person{Publius}；然后是\Person{Maximus}；接着是\Person{Julian}；之后是\Person{Gaius}；其后是\Person{Symmachus}和另一位\Person{Gaius}，以及另一位\Person{Julian}；再后是\Person{Capito}、\Person{Valens}和\Person{Dolichianus}；最后是\Person{Narcissus}，他是自宗徒以来第三十位正常继任者。\par

% structure: p0160 source=h2 target=subsection reason=relative-heading-rank; base=h2

\subsection{第十三章　\Person{Rhodo}及其关于\Person{Marcion}分裂的记述}

1. 此时，\Person{Rhodo}——一个\Place{Asia}本地人，如他所述，他曾受教于我们已经认识的\Person{Tatian}——写了几部著作，其中一部是针对\Person{Marcion}异端的。他说，这个异端在他那个时代已分裂为各种意见；在描述那些导致分裂的人时，他准确地驳斥了每个人所捏造的谬误。\par

2. 但请听他所写的：\par

\Quote{因此，他们也彼此意见不合，坚持不一致的观点。因为\Person{Apelles}——那一群中的一员——以他的生活方式和年岁自夸，承认一个本原，但说预言来自一个敌对的神灵，他是因一个名叫\Person{Philumene}的少女（被鬼附身）的答案而被引导至此观点的。}\par

\Quote{但其他人，其中有\Person{Potitus}和\Person{Basilicus}，持有两个本原，正如那位航海者\Person{Marcion}本人一样。}\par

\Quote{这些人跟随那\Place{Pontus}的豺狼，像他一样不能分辨事物，变得鲁莽，未经任何证明就断言两个本原。另有一些人，则陷入更严重的错误，认为不仅有两个，而是有三个本性。其中\Person{Syneros}是领袖和首领，那些维护他教导的人如此说。}\par

5. 同一作者写道，他曾与\Person{Apelles}交谈。他如此说：\par

\Quote{因为老人\Person{Apelles}与我们交谈时，在许多他虚妄的事上被驳倒；因此他也说，根本不需要审查自己的教义，而每个人都应继续持守他所信的。因为他声称，那些信靠\OriginalTerm{被钉十字架者}{the Crucified}的人将要得救，只要他们被发现行善。但正如我们先前所说，他关于天主的观点是所有人中最模糊的。因为他讲一个本原，正如我们的教义一样。}\par

6. 然后，在充分陈述他自己的意见之后，他补充道：\par

\Quote{当我对他说：\InnerQuote{告诉我，你怎么知道这个，或者你怎么能断言只有一个本原？}他回答说，预言自相矛盾，因为它们没有说任何真实的事；因为它们不一致、虚假且自相矛盾。至于如何只有一个本原，他说他不知道，但他如此确信。}\par

7. \Quote{后来当我恳求他说实话时，他发誓他确实那样做，当他说他不知道如何有一位无始的天主时，但他相信。于是我笑了，责备他，因为他自称教师，却不知道如何确证他所教导的。}\par

8. 在同作品中，这位作者致\Person{卡利斯提奥}时，承认自己曾于罗马受教于\Person{塔提安}。他说\Person{塔提安}编写了一部\WorkTitle{问题集}，在其中应许要解释圣经中晦涩隐秘的部分。\Person{洛多}本人也承诺在自己的著作中提供\Person{塔提安}问题的解答。此外，还有他对\WorkTitle{六日创造}的注释存世。\par

9. 但这\Person{阿佩莱斯}以不敬的方式撰写了关于梅瑟法律的许多内容，在其诸多作品中亵渎天主的话，似乎极其热衷于反驳和推翻它们。\par

论及这些，已足够。\par

% structure: p0174 source=h2 target=subsection reason=relative-heading-rank; base=h2

\subsection{第十四章 弗里吉亚人的假先知}

天主教会之敌人，即那极端憎恶善、喜爱恶、对人无所不用其极的，再次活动，使离奇异端在教会中萌生。因为有些人，如同毒蛇，爬遍\Place{亚细亚}和\Place{弗里吉亚}，夸耀说\Person{蒙塔努斯}是\OriginalTerm{护慰者}{Paraclete}，且跟随他的妇女\Person{普黎史拉}和\Person{玛克西米拉}是蒙塔努斯的女先知。\par

% structure: p0176 source=h2 target=subsection reason=relative-heading-rank; base=h2

\subsection{第十五章 罗马的布拉斯都斯裂教}

其他以\Person{弗洛里努斯}为首的人，在罗马活跃。他脱离了教会的司铎职，\Person{布拉斯都斯}也陷入类似的背离。他们还引诱教会中许多人归向他们的见解，各自力图在真理方面引入自己的新说。\par

% structure: p0178 source=h2 target=subsection reason=relative-heading-rank; base=h2

\subsection{第十六章 有关蒙塔努斯及其假先知的情况}

1. 为对抗所谓弗里吉亚异端，那常为真理争战的大能兴起了一位强大无敌的武器，即我们之前提过的\Place{希拉波利}的\Person{阿波利纳里}，与他一起的还有许多有才干的人，他们为我们留下了丰富的历史材料。\par

2. 其中一位，在其反对他们的著作开头，首先指出他曾与他们进行过口头辩论。\par

3. 他以这样的方式开始其著作：\par

\Quote{亲爱的\Person{阿维基乌斯·玛尔凯禄}，你长久且充分地催促我写一篇论著，反对那些以\Person{米提亚德}为名的异端，我一直犹豫至今，并非因无力驳斥虚谎或为真理作证，而是担心和顾虑，恐怕有人会以为我在增补新约福音的教义或规诫，而凡选择按照福音生活的人，既不能增加也不能减少这些。}\par

\Quote{4. 但近日在\Place{迦拉达}的\Place{安基拉}，我发现那里的教会因这个新奇之事大受搅动，这不是他们所谓的预言，而是假预言，如下将证明。因此，我们尽己所能，借主的助佑，在教会中多日辩驳他们分别提出的这些和其他事项，以致教会喜乐并在真理中得到坚固，对方一时懊悔，仇敌悲伤。}\par

\Quote{5. 该地的长老们，以及我们同是长老的\Place{奥特鲁斯}的\Person{佐提库斯}也在场，请求我们留下驳斥真理反对者的话语记录。我们没有这样做，但许诺一旦主允许，就尽快写出寄给他们。}\par

6. 他在著作开头说完这些以及其他内容后，继续陈述上述异端的起因如下：\par

\Quote{他们的敌对以及那使他们与教会分离的新近异端，源于如下原因。}\par

\Quote{7. 据说在\Place{密细亚}与\Place{弗里吉亚}交界处有一个村子叫\Place{阿达堡}。他们称，最初当\Person{格拉图斯}任\Place{亚细亚}总督时，一个名叫\Person{蒙塔努斯}的新皈依者，因对领导的不可遏止的渴望，给了仇敌攻击他的机会。他变得癫狂，突然陷入一种疯狂和出神状态，胡言乱语，开始口吐奇怪的话，以违反教会从起初因传统而传承的恒定习惯的方式说预言。}\par

\Quote{8. 当时听到他虚假言论的一些人感到愤慨，斥责他是附魔者、被魔鬼控制、被欺骗的灵引导、涣散群众，并禁止他讲话，因为他们记得主的区别和警告，要谨慎提防假先知的来到（\BibleRef{玛 7:15}）。但另一些人自以为拥有圣神和先知之恩，沾沾自喜，骄傲不已，忘记了主的区别，竟向那疯狂、阴险、诱惑的灵发出挑战，结果被它欺骗和迷惑。从此，他再也无法被抑制而保持沉默。}\par

\Quote{9. 这样，魔鬼藉着诡计，更确切地说，藉着这种邪恶的巧计，为悖逆者策划毁灭，却反被他们不配地尊崇，它暗中激动并点燃了他们已背离真信仰的理解力。它还另外煽动两个女人，用假神充满她们，使她们像先前提到的那样，狂乱、无理、怪异地说起话来。那灵见它们因它而欢喜夸耀，便称它们有福，用其大应许抬举它们。但有时它也以明智和可信的方式公开斥责它们，使自己看似一个责备者。但弗里吉亚人中被欺骗的只是少数。}\par

\Quote{那骄傲的灵教导他们辱骂整个天上的普世教会，因为假预言的神既未从教会得到尊荣，也未获准进入。}\par

\Quote{10. 因为亚细亚的忠信者多次在亚细亚各地聚集商议此事，审视那些新奇言论，判定其为亵渎，并拒绝这异端，于是这些人被逐出教会，被剥夺共融。}\par

11. 在开头叙述了这些事，并在全书中继续反驳他们的谬误后，他在第二卷书中如下论及他们的结局：\par

12. 既然他们称我们为杀害先知的人，因为我们不接受他们的多言先知——他们说这些先知就是主应许要派遣到人民那里去的——\BibleRef{玛 23:34}，就让他们在天主面前回答：朋友们，在这些开始说话的人中，从\Person{Montanus}和那些女人以下，有谁受过犹太人的迫害，或被不法之徒杀害？没有。有谁曾为那名而被拘捕并钉在十字架上？确实没有。这些女人中有谁曾在犹太人的会堂里受过鞭打，或被石头砸死？没有；从未在任何地方有过。\par

13. 然而，据说\Person{Montanus}和\Person{Maximilla}是死于另一种死亡。因为传言说，他们被疯狂之灵所激动，两人都上吊自尽；并非同时，而是在谣言所传的各自死亡之时。他们就这样死了，像叛徒\Person{Judas}一样结束了自己的生命。\par

14. 同样，正如一般传言所说，那个非凡的人物，可以说是他们所谓预言的首位管家，一个名叫\Person{Theodotus}的人——他好像曾被提到天上并被接纳，陷入恍惚状态，将自己交托给那欺骗之灵——结果被像铁饼一样抛掷出去，悲惨地死去。\par

15. 他们说这些事就是以这种方式发生的。但我们并未亲眼所见，朋友，我们不敢自称知道。或许是这样，或许不是，\Person{Montanus}、\Person{Theodotus}和上述那位女人就是这样死的。\par

16. 他再次在同一本书中说，当时的主教们试图驳斥\Person{Maximilla}里的那灵，但被其他人阻止，这些人显然是与那灵合作的。\par

17. 他写道：\par

\Quote{并且，让那灵在\Person{Asterius Urbanus}的同一部作品中，不要通过\Person{Maximilla}说：\InnerQuote{我像狼一样被从羊群中赶走。我不是狼。我是言语、灵和能力。}而是让他清楚地显示并证明那灵中的能力。并且通过那灵，让他迫使当时在场、为证明和辩驳那多言之灵的人承认他，——那些杰出的人物和主教，来自\Place{Comana}村的\Person{Zoticus}，和来自\Place{Apamea}的\Person{Julian}，\Person{Themiso}的追随者封住了他们的口，拒绝允许那虚假和迷惑之灵被他们驳斥。}\par

18. 在同一部作品中，在说了其他驳斥\Person{Maximilla}虚假预言的话之后，他指明了他写这些记述的时间，并提及她的预言，其中她预言了战争和无政府状态。他以下列方式谴责了它们的虚假：\par

19. \Quote{这一点不是清楚表明是假的吗？因为那女人去世至今已超过十三年，世界上既没有局部战争也没有全面战争；反而，因天主的仁慈，甚至对基督徒而言也有持续的和平。}这些内容摘自第二卷。\par

20. 我还要添加第三卷中的简短摘录，其中他这样反对他们关于许多人殉道的夸口：\par

因此，当他们无言以对，在所说的一切中被驳倒时，他们试图以他们的殉道者为庇护，声称他们有许多殉道者，并且这是与他们同在的所谓预言之灵的能力的确凿证据。但这看来完全是谬误的。\par

21. 因为有些异端有许多殉道者；但我们当然不会因此同意他们，或承认他们拥有真理。首先，那些被称为\OriginalTerm{马西翁派}{Marcionites}的，来自\Person{Marcion}的异端，声称他们有为基督而殉道的大量殉道者；然而他们并不真实地承认基督自己。\par

稍后他继续说道：\par

22. \Quote{当教会中为信德的真理而被召殉道的人，遇到弗里吉亚异端的所谓殉道者时，就与他们分离，并在没有与他们有任何共融的情况下死去，因为他们不愿同意\Person{Montanus}和那些女人的灵。而这是真实的，并且发生在我们这个时代，在迈安德河畔的\Place{Apamea}，与\Person{Gaius}和\Person{Alexander}（来自\Place{Eumenia}）一同殉道的人中，这是众所周知的。}\par

% structure: p0207 source=h2 target=subsection reason=relative-heading-rank; base=h2

\subsection{第十七章 \Person{Miltiades}及其著作。}

1. 在这部作品中，他提到一位作家\Person{Miltiades}，说他曾写了一部书来反对上述异端。在引用他们的一些话之后，他补充道：\par

\Quote{我在他们的一部作品中发现了这些内容，这部作品反对兄弟\Person{Alcibiades}的著作，其中他表明先知不应在神魂超拔中说话，我于是做了一个摘要。}\par

2. 在同一部作品中稍后，他列出了在新盟约下预言的人，其中他列举了一位\Person{Ammia}和\Person{Quadratus}，说道：\par

但假先知陷入神魂超拔状态，在此状态中他不知羞耻或恐惧。从有意的无知开始，如所述，他转至不自主的灵魂疯狂。\par

3. 他们不能指出旧约或新约先知中有谁曾如此被灵带走。他们也不能以\Person{Agabus}、\Person{Judas}、\Person{Silas}、\Person{Philip}的女儿们、在\Place{Philadelphia}的\Person{Ammia}、\Person{Quadratus}，或任何不属于他们的人夸口。\par

4. 稍后他又说道：\Quote{因为如果在\Place{Philadelphia}的\Person{Quadratus}和\Person{Ammia}之后，如他们所声称的，那些与\Person{Montanus}在一起的女人领受了预言之恩，就让他们指出他们中谁是从\Person{Montanus}和那些女人领受的。因为宗徒认为预言之恩应当在全体教会中持续直到最终的来临。但他们不能指出，尽管这是自\Person{Maximilla}去世以来的第十四年。}\par

5. 他如此写道。但他所指的\Person{Miltiades}留下了自己对神圣圣经的热忱的其他纪念，在他所撰写的反对希腊人和反对犹太人的讲论中，分别用两卷书来回答他们。此外，他还向现世的统治者发表了一篇护教书，为他所信奉的哲学辩护。\par

% structure: p0215 source=h2 target=subsection reason=relative-heading-rank; base=h2

\subsection{第十八章 \Person{Apollonius}驳斥弗里吉亚人的方式，以及他所提及的人物。}

1. 由于所谓的弗里吉亚异端在他那个时代仍在弗里吉亚盛行，\Person{Apollonius}，一位教会作家，也着手驳斥它，并写了一部专门的著作反对它，详细纠正了他们中流行的虚假预言，并谴责了异端创始人的生活。但请听他自己关于\Person{Montanus}的话：\par

2. \Quote{他的行动和教训显明了这位新教师是谁。就是他教导婚姻的解除；制定了关于禁食的法规；将夫黎基雅的两个小镇\Place{佩普札}和\Place{提米翁}命名为\Place{耶路撒冷}，想要从各处聚集人到那里；设立了收取钱财的人；以奉献为名安排接受礼物；为传布他教义的人提供薪水，好使那教导藉着贪饕得以盛行。}\par

3. \Quote{我们表明这些首批女先知自己，一旦被圣神充满，就遗弃了她们的丈夫。所以那些称\Person{普黎斯加}为童贞女的人，说的是何等虚假的话。}\par

4. \Quote{难道全部圣经不似乎禁止先知接受礼物和金钱吗？因此，当我看见那女先知接受金、银和贵重衣物时，我怎能不谴责她呢？}\par

5. 稍后他又论及他们的一位宣信者，这样说：\par

\Quote{同样，那个\Person{特米索}，披着貌似有理的贪财，不能忍受宣信的标记，却为大量的财产抛开了锁链。然而，他本应因此谦卑，却胆敢以殉道者自夸，并效法宗徒，写了一封所谓公函，去教训那些信德比他更好的人，为虚浮空谈争辩，亵渎主、宗徒和圣教会。}\par

6. 又论及他们中间被尊为殉道者的其他人，他这样写道：\par

不必多说许多人，让那女先知亲自告诉我们关于\Person{亚历山大}，就是那自称为殉道者的人，她常与他一同宴乐，且被许多人敬拜。我们无需提及他的抢劫和其他胆大妄为的行径——他因此受过罚，但档案里都有记载。\par

7. 这两人中谁赦免另一个的罪？是先知赦免殉道者的抢劫，还是殉道者赦免先知的贪财？因为主曾说：\BibleQuote{不要带金、银，也不要带两件内衣。}\BibleRef{玛 10:9} 但这些人完全相反，在拥有禁止之物的方面犯了过。因为我们将要显明，他们称为先知和殉道者的人，不仅从富人，也从穷人、孤儿和寡妇那里聚敛他们的收益。\par

8. 但他们若自信，就让他们起来讨论这些事，好叫他们若被定罪，以后可以不再犯过。因为先知的果子必须被试验；\BibleQuote{因为树是由它的果子认出来的。}\BibleRef{玛 12:33}\par

9. 但为了使那些想知道关于\Person{亚历山大}的人了解，他曾被厄弗所的总督\Person{艾米略·弗隆提诺}审判；不是因这名，而是为他所犯的抢劫罪——他那时已经是个背教者。后来，他冒称为主的名宣告，被释放了，欺骗了那里忠信的人。他自己的教区，就是他来的地方，没有接纳他，因为他是个强盗。想要知道关于他的人，有亚细亚的公共档案可查。然而，与他共处多年的那位先知却对此一无所知！\par

10. 揭露他的同时，我们也藉着他揭露那先知的伪诈。我们也能就许多其他人证明同样的事。但他们若自信，就让他们经受考验。\par

11. 在他的作品的另一部分，他又论及他们所夸耀的先知如下：\par

\Quote{如果他们否认他们的先知接受过礼物，就让他们承认这点：如果他们被证实接受了礼物，他们就不是先知。我们将为此带来许多证据。但必须审查先知的一切果子。告诉我，先知会染发吗？先知会染眼睑吗？先知会喜爱装饰吗？先知会玩牌和骰子吗？先知会放贷取利吗？让他们承认这些事情是否合法；但我要显明这些事他们已经做了。}\par

12. 同一位\Person{阿颇罗尼}在同一作品中陈述，在他写作之时，自\Person{蒙丹}开始他自称的预言以来，已是第四十年。\par

13. 他还提及，当\Person{玛克西弥拉}在\Place{佩普札}假装说预言时，前面那作者提到过的\Person{佐提古}曾抵挡她，并试图驳斥在她内运行的那神，但被那些赞同她的人阻止了。他还提及那时代的一位殉道者\Person{特拉塞亚}。\par

他还提到一个传统，说救主命令他的宗徒十二年内不要离开\Place{耶路撒冷}。他也引用了\WorkTitle{\BibleBook{若望}{默示录}}的见证，并叙述了一个死人曾在天主大能下在\Place{厄弗所}被\Person{若望}本人复活。他还补充了其他事情，藉此充分而丰富地揭露了我们所谈论的异端之错误。这些都是\Person{阿颇罗尼}所记录的事。\par

% structure: p0233 source=h2 target=subsection reason=relative-heading-rank; base=h2

\subsection{第十九章 \Person{色拉彼雍}论夫黎基雅人的异端。}

1. \Person{色拉彼雍}——据报告，他那时继\Person{玛克西米诺}之后任\Place{安提约基雅}教会的主教——提及\Person{阿波利纳里}反对上述异端的著作。他在写给\Person{卡黎古}和\Person{彭提}的私人信函中也提及此人，在信中他本人揭露了同一异端，并补充了以下话语：\par

2. \Quote{为使你们看见这所谓新预言的撒谎团伙的所作所为，对所有普世弟兄团体是多么可憎，我已给你们送来了有福的\Person{克劳狄·阿波利纳里}的著作，他是\Place{亚细亚}\Place{希耶拉波利}的主教。}\par

3. 在\Person{色拉彼雍}同一封信中，有几位主教的签名，其中一位这样署名：\par

\Quote{\Person{奥勒里·基勒尼}，见证人，祝你健康。}\par

另一人这样署名：\par

\Quote{\Person{艾略·普布里·尤利}，色雷斯殖民地\Place{德贝尔图姆}的主教。正如天主在天上活着，有福的\Person{索塔斯}在\Place{安基雅罗}想要从\Person{普黎斯加}身上驱赶那恶魔，但那些伪善者未允许他。}\par

4. 同一封信中还载有许多其他赞同他们的主教的亲笔签名。\par

关于这些人，就说到这里。\par

% structure: p0242 source=h2 target=subsection reason=relative-heading-rank; base=h2

\subsection{第二十章 \Person{依肋内}反对\Place{罗马}分裂者的著作。}

1. \Person{依勒内}写了几封信，反对那些扰乱罗马教会健全规矩的人。其中一封致\Person{布拉斯都}的《\WorkTitle{论分裂}》；另一封致\Person{弗洛里努}的《\WorkTitle{论君主政体}》，或《\WorkTitle{天主不是灾祸的制造者}》。因为\Person{弗洛里努}似乎拥护此观点。又因他被\Person{瓦伦提努}的错谬引入歧途，\Person{依勒内}写了《\WorkTitle{论八元质}》之作，其中显示他本人曾认识宗徒的首批继承人。\par

2. 在该论的结尾处，我们找到了一段极美的注释，我们不得不将其插入本著。其文如下：\par

\Quote{我恳求凡誊写此书者，以我们的主\Person{耶稣基督}，并以他光荣的来临审判生者死者的名，请将你所写的内容与这份手稿仔细校对并更正，且将此誓词誊写并放置于副本中。}\par

3. 这些事情在他的著作中值得一读，我们在此提及，为使我们能将那些古圣先贤当作细心谨慎的最佳典范。\par

4. 在致\Person{弗洛里努}的书信中（我们已提及），\Person{依勒内}再次提到他与\Person{波利卡普}的亲密关系，说道：\par

\Quote{这些道理，哦，\Person{弗洛里努}（说得温和些），并非健全的判断。这些道理与教会不合，并将接受它们的人推入极大的不敬中。这些道理，甚至连教会外的异端者都不敢宣扬。这些道理，我们之前的司铎们（他们是宗徒的同伴）并没有传授给你。}\par

5. 因为当我还是孩童时，我在\Place{下亚细亚}见过你与\Person{波利卡普}在一起，在皇宫中光彩照人，并试图博得他的赞许。\par

6. 我记得那时的事情比近年的事情更清楚。因为童年所学，随着心智成长，与之结合；因此我能描述那有福的\Person{波利卡普}讲道时所坐的地点，他的进出，他的生活方式，他的外貌，他对众人的演讲，以及他讲述与\Person{若望}及其他见过\Person{主}的人交往的情况。当他记起他们的话，以及他从他们那里听到的关于\Person{主}、关于\Person{主}的奇迹和教导，他从\BibleQuote{生命的圣言}的目击者那里领受后，\newline{}\BibleRef{若一 1:1}\Person{波利卡普}讲论一切与圣经相符。\par

7. 这些事因天主的仁慈告诉我，我留心听记，不是记在纸上，而是记在心里。并常常因天主的恩宠，忠实地回忆它们。我能在天主面前作证，如果那位有福的宗徒的司铎听见了任何这类事情，他必会喊叫，掩耳，并如常喊说：\Quote{善哉天主，你竟容忍我到何时，使我忍受这些事？}他必会逃离那地方，无论他坐或站，听见了这样的话。\par

8. 这可以从他写给邻邦教会以坚固他们，或写给某弟兄劝诫和鼓励的信中清楚显示。至此为\Person{依勒内}之言。\par

% structure: p0253 source=h2 target=subsection reason=relative-heading-rank; base=h2

\subsection{第二十一章 \Person{阿颇罗尼乌斯}如何在\Place{罗马}殉道。}

1. 大约在同一时期，\Person{科摩多斯}在位时，我们的处境更为有利，借着天主的恩宠，普世各处的教会享受和平，救恩之言引导着各种族各民中的每个灵魂虔敬崇拜宇宙的天主。以至于在\Place{罗马}，许多因财富和家族显赫的人，带着全家和亲属归向他们的救恩。\par

2. 但那嫉善的魔鬼，本性恶毒，不能容忍此事，便再准备争斗，设计许多计谋攻击我们。他鼓动\Place{罗马}城的一位仆人（此人很适合这样的目的）来控告\Person{阿颇罗尼乌斯}，他是\Place{罗马}城人，在信友中以学识和哲学闻名。\par

3. 但这可怜的人不合时宜地提出了指控，因为根据敕令，此类告密者不得存活。法官\Person{佩雷尼乌斯}立即判决打断他的双腿。\par

4. 但那深受天主喜爱的殉道者，经法官恳切请求他向元老院陈词，在众人面前为他所见证的信德作了雄辩的辩护。仿佛由元老院下令，他被斩首处死；一条古老的法律要求，凡被带到审判席前且拒绝撤回信仰者，不得释放。凡愿知晓他在法官前的辩论、他对\Person{佩雷尼乌斯}提问的回答，以及他在元老院前的全部辩护者，可从我们所收集的古代殉道记录中找到。\par

% structure: p0258 source=h2 target=subsection reason=relative-heading-rank; base=h2

\subsection{第二十二章 此时闻名的诸位主教。}

在\Person{科摩多斯}在位第十年，\Person{维克托}接续\Person{厄琉特鲁斯}，后者担任主教十三年。同年，\Person{尤利安}完成十年任期后，\Person{德默特琉斯}在\Place{亚历山大}接受了各堂区的管理权。此时，上述的\Person{塞拉皮雍}，即自宗徒以来第八位，仍以\Place{安提约基亚}教会主教闻名。\Person{德奥斐禄}在\Place{巴勒斯坦}的\Place{凯撒勒雅}主持教务；而先前我们提过的\Person{纳尔基索斯}仍负责\Place{耶路撒冷}教会。\Person{巴基卢斯}同时在\Place{希腊}的\Place{科林多}担任主教，\Place{厄弗所}的\Person{波利克拉特斯}也是如此。此外，可能还有许多其他显要人物。但我们只列出了那些信仰纯正已通过文字流传给我们的人。\par

% structure: p0260 source=h2 target=subsection reason=relative-heading-rank; base=h2

\subsection{第二十三章 当时关于逾越节的争议。}

1. 当时发生了一个不小的问题。因为全亚洲的堂区，按更古老的传统，认为应遵守月之第十四天，即犹太人奉命献羔羊的那一天，作为救主逾越节的庆日。因此，不论那天是星期几，都必须在那一日结束斋戒。但世界其他地方的教会却不在此时结束斋戒，因为他们所遵守的习俗，源自宗徒传统，并一直延续至今，即除救主复活日之外，不在其他日子终止斋戒。\par

2. 为此，召开了主教会议和集会，全体一致通过相互通信，制定了教会法令，即主的复活奥迹只应在主日庆祝，并且我们只应在这一天结束逾越节斋戒。当时在巴勒斯坦集会的人，由\Place{凯撒勒雅}的主教\Person{德奥斐罗}和\Place{耶路撒冷}的主教\Person{纳齐苏斯}主持，他们保存了一份文件。此外，还有在\Place{罗马}集会讨论同一问题的人所保存的另一份文件，以\Person{维克多}主教的名义；还有在\Place{本都}的主教，以最年长的\Person{帕尔马斯}为首；以及\Place{高卢}各堂区，由\Person{依勒内}担任主教；还有\Place{奥斯若恩}及其城市的堂区；以及\Place{格林多}教会的主教\Person{巴希鲁斯}的私人信件，和许多其他主教，他们都发表了相同的意见和判断，投了相同的票。\par

3. 上述内容就是他们一致的决定。\par

% structure: p0264 source=h2 target=subsection reason=relative-heading-rank; base=h2

\subsection{第24章 亚洲的分歧}

1. 但亚洲的主教们，以\Person{波利克拉特斯}为首，决定遵守传给他们古老习俗。他本人写信给\Person{维克多}和\Place{罗马}教会，用以下的话陈述了他所承袭的传统：\par

2. 我们遵守准确的日期，不增不减。因为在亚洲，也有伟大的光体安息了，他们将在主来临的日子复活，那时他以荣耀从天降临，并寻找所有的圣人。其中有\Person{斐理伯}，十二宗徒之一，他在\Place{耶辣颇里}安息；还有他的两位年老的童贞女儿，以及另一位女儿，她生活在圣神内，如今安息于\Place{厄弗所}；此外，还有\Person{若望}，他既是见证人又是导师，曾靠在主的胸膛，并且作为司铎，佩戴着司祭牌。\par

3. 他在\Place{厄弗所}安息。\par

4. 还有在\Place{斯米纳}的\Person{波利卡普}，他是主教兼殉道者；以及从\Place{欧墨尼亚}来的\Person{特辣塞亚斯}，也是主教兼殉道者，他在\Place{斯米纳}安息。\par

5. 我又何必提及在\Place{劳狄刻雅}安息的主教兼殉道者\Person{撒加里}，或\Person{真福帕彼里}，或终生生活在圣神内的太监\Person{梅利托}，他躺在\Place{撒尔德}，等待天上主教的职位，当他从死者中复活时？\par

6. 所有这些人都按照福音遵守逾越节第十四天，毫不偏离，而是遵循信仰的规则。我，\Person{波利克拉特斯}，你们中最小的，也遵循我亲属的传统，其中几位我曾密切追随。因为我的亲属中有七位是主教，我是第八位。我的亲属总是遵守人民除去酵母的那一天。\par

7. 因此，弟兄们，我在主内活了六十五年，曾与普世的弟兄们会面，并通读过全部圣经，并不被威吓的话所惊恐。因为比我更伟大的人说过：\Quote{我们应当服从天主，而不是人。}\BibleRef{宗 5:29}\par

8. 然后他写信提到所有与他同在并赞成他的主教们。他的话如下：\par

\Quote{我可以提及那些在场的主教，我是按你们的要求召请了他们；如果我写下他们的名字，那将是一个庞大的数目。他们看到我的卑微，就同意了这封信，知道我没有白白长出白发，而是始终以\Person{主耶稣}来治理我的生活。}\par

9. 于是，主持\Place{罗马}教会的\Person{维克多}立即试图将全\Place{亚洲}的堂区及其同意他们的教会，当作异端者，从共同合一中切除；他写信宣布那里的所有弟兄完全被逐出教会。\par

10. 但这并不使所有主教都满意。他们恳求他考虑和平的事，以及邻里的合一与爱德。他们有一些话存留下来，严厉地责备\Person{维克多}。\par

11. 其中就有\Person{依勒内}，他以他所主持的\Place{高卢}弟兄们的名义写信，坚持认为主的复活奥迹只应在主日庆祝。他恰当地告诫\Person{维克多}，不应将遵守古老习俗传统的整个天主教会都切除。在说了许多其他话之后，他接着写道：\par

12. 因为争议不仅关乎日期，也关乎禁食的方式本身。有些人认为应禁食一天，有些人两天，还有些人更多；此外，有些人将一天算作四十小时的昼夜。\par

13. 这种遵守方式的多样性并非始于我们时代，而是早在我们的祖先时代。很可能他们没有坚持严格精确，从而根据他们自己的简朴和特有方式，为后代形成了一种习俗。然而，所有这些人都照样生活在和平中，我们之间也彼此和平相处；禁食上的分歧反而证实了信仰上的一致。\par

14. 他接着叙述了以下内容，我理应插入此处：\par

在这些长老中，有\Person{索特尔}之前主持你们如今所统治的教会的那些人。我指的是\Person{阿尼塞图斯}、\Person{庇护}、\Person{希吉诺}、\Person{德莱斯福禄}和\Person{西斯笃}。他们自己既不遵守，也不允许以后的人遵守。然而，尽管他们不遵守，却照样与从遵守此习俗的堂区来到他们那里的人和平相处；虽然这种遵守对不遵守的人来说是更反对的。\par

15. 但从未有人因这种形式而被驱逐；相反，你们之前不遵守此习俗的长老，曾将圣体圣事送给其他遵守此习俗的堂区的人。\par

16. 当真福波利卡普在阿尼塞特任教宗期间来到罗马，他们在某些其他事上稍有分歧，便立即彼此和好，不愿为这事争吵。因为阿尼塞特无法说服波利卡普不遵守他一直遵守的、来自主耶稣基督的门徒若望以及其他与他交往的宗徒的传统；波利卡普也无法说服阿尼塞特遵守这传统，因他说应遵循他前任诸位长老的惯例。\par

17. 但事情虽如此，他们仍彼此共融，阿尼塞特许诺由波利卡普在教会中主持圣体圣事，显然作为尊重的表示。他们彼此和平分手，无论是遵守者还是不遵守者，都保持了整个教会的和平。\par

18. 这样，依勒内——他确实名副其实——在这事上成了和平使者，以这种方式为教会的和平劝诫和协商。他就这个有争议的问题，不仅与维克托，也与教会其他多数领袖通信。\par

% structure: p0285 source=h2 target=subsection reason=relative-heading-rank; base=h2

\subsection{第二十五章  众人如何就逾越节达成一致}

1. 我们最近提及的巴勒斯坦那些人，即纳尔奇苏斯、特奥斐卢斯，以及与他们一同的提洛教会的主教卡西乌斯、托勒迈依斯教会的克拉鲁斯，以及与他们集会的人，详细陈述了关于逾越节从宗徒代代相传的传统，在书信结尾加上这些话：\par

2. \Quote{务请将我们的信函副本送至每一教会，以免给那些轻易欺骗自己灵魂的人提供机会。我们向你们表明，在亚历山大里亚，他们也与我们同一天守节。因为书信从我们这里送到他们那里，从他们那里送到我们这里，以致我们在同一方式、同一时间守这圣日。}\par

% structure: p0288 source=h2 target=subsection reason=relative-heading-rank; base=h2

\subsection{第二十六章  依勒内流传至今的优雅著作}

除了我们提及的依勒内的著作和书信外，还存有一卷他写的\WorkTitle{论知识}，反驳希腊人，言简意赅，极为有力；还有一卷他献给弟兄马尔奇安的\WorkTitle{宗徒宣讲之论证}；以及一卷包含各种\WorkTitle{论说}，其中他提及\WorkTitle{希伯来人书}和所谓的\WorkTitle{撒罗满的智慧}，并从中引文。这些就是我们所知的依勒内的著作。\par

科莫杜斯在位十三年后结束统治，他死后不到六个月，塞维鲁成为皇帝，其间佩蒂纳克斯曾短暂统治。\par

% structure: p0291 source=h2 target=subsection reason=relative-heading-rank; base=h2

\subsection{第二十七章  当时其他杰出人士的著作}

当时古代教会人士信仰热忱的众多记录至今仍被许多人保存。其中我们特别注意到赫拉克利特论宗徒的著作，以及马克西默斯关于异端者中广为讨论的问题——罪恶的起源和物质的受造——的著作。还有坎迪杜斯论\WorkTitle{六日创世}的著作，以及阿皮翁论同一主题的著作；同样，塞克斯图斯论复活的著作，以及阿拉比亚努斯的另一论著，还有许多其他人的著作。对这些著作，由于我们缺乏资料，无法在我们的作品中说明他们生活的年代或提供任何历史记载。还有许多其他人的著作流传下来，我们无法列出其名字，这些人是正统的教会人士，正如他们对圣经的诠释所表明的，但我们不知其名，因为他们的著作中未提及名字。\par

% structure: p0293 source=h2 target=subsection reason=relative-heading-rank; base=h2

\subsection{第二十八章  首先倡导阿尔特蒙异端的人；他们的生活方式，以及他们如何胆敢篡改圣经}

1. 在这些作者之一的一篇辛劳之作中，反对阿尔特蒙的异端——这异端在我们的时代由萨莫萨塔的保禄试图复兴——有一段记述适合我们正在考察的历史。\par

2. 因为他批评上述异端是晚近的革新，该异端教导救主仅是凡人，因为他们企图将其夸大如古老。他在作品中列举了许多论据驳斥他们亵渎的谎言，随后加上这些话：\par

3. 因为他们说，所有早期的教师和宗徒都接受并教导他们现在所宣称的，福音的真理一直保存到维克托的时代——他是自\Person{伯多禄}以来的第十三位罗马主教——但从他的继任者则斐利诺开始，真理就被败坏了。\par

4. 他们所说的话也许看似合理，但首先，圣经就与他们相悖。还有比维克托时代更早的某些弟兄的著作，他们为真理写过反对异教徒和当时存在的异端的文章。我是指犹斯定、米尔提亚德、塔提安、克莱孟和许多其他人，在他们的所有著作中，基督被称为天主。\par

5. 因为谁不知道依勒内、默利托以及其他人的著作教导基督既是天主又是人呢？有多少由虔诚弟兄们从初期写下的圣咏和赞美诗，歌颂基督——天主的圣言，称他为天主。\par

6. 那么，既然教会所持的见解已宣讲了这么多年，怎能如他们所宣称的那样，宣讲被延迟到维克托的时代呢？他们明知维克托将鞋匠特奥多图斯——这个否认天主的背道运动的领袖和父亲，也是第一个宣称基督仅是凡人的人——逐出共融，怎能不羞于如此诽谤维克托呢？如果维克托同意他们的意见，像他们的诽谤所声称的那样，他怎么会驱逐这异端的发明者特奥多图斯呢？\par

7. 以上是关于维克托的叙述。他的主教任期持续了十年，则斐利诺在塞维鲁统治的第九年左右被任命为他的继任者。上述那本书的作者关于这异端的创始人，叙述了另一起发生在则斐利诺时代的事件，用了这些话：\par

8. 我要提醒许多弟兄一件发生在我们时代的事，这事若发生在索多玛，我想或许会对他们构成警告。有一个名叫纳塔留的宣信者，不是很久以前，而是在我们自己的时代。\par

9. 此人一度被\Person{阿斯克勒庇奥多图斯}和另一位\Person{戴奥杜笃}（钱商）所欺骗。这两人都是鞋匠\Person{戴奥杜笃}的门徒，正如我所说，后者因这种观念（或更确切地说，这种愚昧）而被当时的主教\Person{维克多}首次逐出教会。\par

10. \Person{纳塔留斯}被他们说服，允许自己被选为这异端的主教，按月领取他们支付的一百五十迪纳里的薪俸。\par

11. 当他如此与他们联合后，主多次通过异象警告他。因为慈悲的天主和我们的主\Person{耶稣基督}不愿祂自己苦难的见证人被逐出教会而丧亡。\par

12. 但他并未重视这些异象，因他被他们中的首位地位和那种毁灭许多人的可耻贪欲所缠绕，便在整夜被圣天使鞭打，受到严厉惩罚。清晨起来，他披上麻衣，撒上灰烬，极为匆忙地流着泪俯伏在主教\Person{则斐琳}面前，不仅匍匐在圣职人员的脚下，也匍匐在平信徒的脚下；他用自己的眼泪感动了仁慈的\Person{基督}的慈悲教会。虽然他极力恳求，并展示了所受鞭打的伤痕，却仍几乎不被重新接纳进入共融。\par

13. 我们将从同一作者那里再添加一些关于他们的其他摘录，内容如下：\par

\Quote{他们轻率而无畏地对待圣经。他们废弃了古代信仰的准则，并不认识\Person{基督}。他们不努力学习圣经所宣告的，却辛苦钻研任何可用来支持他们不虔敬的三段论形式。若有人向他们提出一段圣经经文，他们便看能否从中得出一个结合或分离的三段论。}\par

\Quote{他们既属于地，也谈论地，且不认识那从上面来的那一位，便离弃天主的圣书，投身于几何学。他们中有些人辛苦地丈量\Person{欧几里得}；\Person{亚里士多德}和\Person{泰奥弗拉斯托斯}备受赞赏；\Person{盖伦}甚至可能被某些人崇拜。}\par

\Quote{但那些利用不信者的技艺来支持他们异端观点，并用无神者的诡计玷污圣经纯朴信仰的人，已远离信仰，这又何须多言？因此，他们胆敢擅自修改圣经，声称是自己校正了它们。}\par

\Quote{在这事上我并非说谎，谁若愿意便可知道。因为若有人收集他们各自的抄本，彼此比较，便会发现它们差异甚大。}\par

\Quote{例如，\Person{阿斯克勒庇亚德斯}的抄本与\Person{戴奥杜笃}的抄本不一致；其中许多抄本可以获取，因为他们的门徒殷勤地抄写了各人所称的校正（实为败坏）。再者，\Person{赫尔莫菲鲁斯}的抄本与这些也不一致，\Person{阿波罗尼德斯}的抄本自身也不统一。因为你可以将他们早先准备的抄本与后来败坏的抄本相比较，便会发现它们大不相同。}\par

\Quote{至于这罪行何等胆大，他们自己恐怕并非不知。因为他们要么不相信圣经是由圣神所启示的，从而成为不信者；要么自以为比圣神更聪明，若如此，他们岂非着魔之人？因为他们无法否认这罪行，因抄本出自他们亲手。他们既未从导师那里接受这样的圣经，也拿不出任何从中转抄的抄本。}\par

\Quote{但他们中有些人觉得不值得败坏圣经，而是干脆否认法律和先知，藉此以恩宠为借口，通过他们无法无天和不虔敬的教导，沉沦于毁灭的至极深渊。}\par

关于这些事，到此为止。\par

\SourceRecord{Translated by Arthur Cushman McGiffert.；Nicene and Post-Nicene Fathers, Second Series；Vol. 1.；Edited by Philip Schaff and Henry Wace.；Buffalo, NY: Christian Literature Publishing Co.,；1890.；<http://www.newadvent.org/fathers/250105.htm>.}

% structure: p0001 source=h1 target=section reason=page-title

\section{《\WorkTitle{教会史}（第六卷）》}

% structure: p0002 source=h2 target=subsection reason=relative-heading-rank; base=h2

\subsection{第一章 塞维鲁教难}

当塞维鲁开始迫害教会时，各地宗教的斗士都作出了光荣的见证。在亚历山大城尤其如此，这座城就像一座最显赫的竞技场，来自埃及和整个提巴依德的、按功绩选出的天主的斗士被带到这里，因着他们在众多酷刑和各种死亡方式下所表现出的极大忍耐，从天主那里赢得了冠冕。其中有一位名叫雷奧尼德（Leonides），他被称作奥力振（Origen）的父亲，在他儿子尚年幼时被斩首。关于此子对天主圣言有何等非凡的偏爱——这源于他父亲的教导——简要叙述一下并无不妥，因为他的名望已被多人极隆重地传扬。\par

% structure: p0004 source=h2 target=subsection reason=relative-heading-rank; base=h2

\subsection{第二章 奥力振自幼所受的教育}

1. 若要描述此人求学时期的生活，可说的很多，但仅此题目就需要一部单独的著作。然而，眼下我们删繁就简，尽可能简短地陈述关于他的一些事实，这些事实是从某些书信和至今仍在世、认识他的人的口述中收集的。\par

2. 他们关于奥力振的记述，在我看来，即便从他蹒跚学步时说起，也是值得一提的。\par

那是塞维鲁执政的第十年，莱图斯（Lætus）担任亚历山大城和埃及其余地区的总督，德梅特里乌斯（Demetrius）刚刚接替尤利安（Julian）担任当地教区的主教。\par

3. 既然教难的火焰已被大大点燃，众多人已获得殉道的冠冕，那么对殉道的渴望也抓住了奥力振的灵魂——尽管他还是个孩子——他趋近危险，满怀热忱地冲上前去，投身于战斗。\par

4. 确实，他的生命几乎就要终结，若非神圣的天主眷顾为着许多人的益处，通过他母亲的行动阻止了他的渴望。\par

5. 起初，她哀求他，恳求他怜悯她作为母亲对他的感情；但当她发现，他在得知父亲已被逮捕监禁后反而更加坚决，完全被对殉道的热忱冲昏头脑时，她便将他所有的衣服都藏了起来，从而迫使他留在家中。\par

6. 然而，既然他别无他事可做，而他超越年龄的热忱又不容他安静，他便给父亲寄去一封劝勉信，论及殉道，信中鼓励他说：\Quote{当心不要因我们的缘故而改变心意。}这可以说是奥力振幼年智慧和他对虔敬真诚之爱的第一个证据。\par

7. 因为即使在那时，他已经在信仰的言辞中积累了不小的资源，自幼受过《圣经》的训练。他并非漫不经心地研读它们，因为他的父亲除了给予他通常的博雅教育外，还使这些成为头等大事。\par

8. 首先，在引导他进入希腊科学之前，父亲先让他接受神圣研究的训练，要求他每天学习并背诵。\par

9. 这孩子并不觉得这是负担，反而对这些研究充满热忱和勤奋。他不满足于学习《圣经》话语中简单浅显的内容，而是寻求更深的东西，甚至在那个年龄就忙于更深入的思辨。以致他用种种追问来困扰他的父亲，想要知道受天主启示的《圣经》的真正含义。\par

10. 他的父亲表面上当面责备他，告诉他不要探寻超出年龄或明显含义之外的东西。但私下里却非常欢喜，感谢一切美善的施与者——天主，因祂认为他配作这样一个孩子的父亲。\par

11. 据说他常常站在熟睡的孩子身旁，掀开他的胸膛，仿佛天主圣神居住在其中，并恭敬地亲吻它；他因自己有这样一个优秀的后裔而自认有福。关于年幼的奥力振，人们还传述了这些以及类似的事情。\par

12. 但他的父亲在殉道中结束了生命时，他留下母亲和六个弟弟，他本人还不满十七岁。\par

13. 他父亲的财产被没收归皇室国库，他和家人缺乏生活必需品。但他被证明配得天主的看顾。他得到一位非常富有、在生活方式和其他方面都很杰出的妇人的接待和安顿。这位妇人当时厚待一位著名的异端者，他当时正在亚历山大城；不过这异端者生在安提约基雅。他被她视同义子，倍受恩待。\par

14. 尽管奥力振被迫与他交往，但从那时起他却对自己的信仰正统性给出了强有力的证据。因为由于保禄（Paul）——此人是那异端者的名字——在辩论中表现出的明显技巧，一大群人来到他那里，不仅有异端者，也有我们教内的人。但奥力振从未被诱导与他一起祈祷；因为尽管是个孩子，他仍持守教会的规则，并如他在某处所说，憎恨异端的教导。他父亲已教过他希腊科学，父亲去世后，他更加刻苦专攻文学研究，从而在文献学方面获得了相当准备，并能在父亲去世后不久，通过致力于该学科，挣得一份足以满足他那个年龄需求的丰厚报酬。\par

% structure: p0020 source=h2 target=subsection reason=relative-heading-rank; base=h2

\subsection{第三章 虽极年轻，仍殷勤教授基督的圣言}

1. 然而，当他在学校讲学时——据他本人告诉我们——亚历山大城没有人在信仰上施教，因为所有人都因教难的威胁而被驱散了，于是一些教外人到他那里来听天主的话。\par

2. 他说，其中第一位是普路塔克（Plutarch），他善度此生后，获得了神圣殉道者的荣耀。第二位是赫拉克拉斯（Heraclas），普路塔克的兄弟；他在与前者一起提供了大量哲学和苦修生活的证据后，被认定配得继承德梅特里乌斯担任亚历山大城主教之职。\par

3. 他十八岁时接管了要理学校。他此时也很突出，在\Person{阿奎拉}任\Place{亚历山大}总督的迫害期间，由于他对所有圣殉道者（无论相识与否）所表现出的仁慈和善意，他的名字在信仰领袖中变得著名。\par

4. 因为他不仅在他们被囚期间，直到他们最后被判刑时都与他们同在，而且当圣殉道者被带去处死时，他非常勇敢，与他们一同进入危险。因此，他勇敢地行动，并以极大的勇气亲吻殉道者向他们致意，常常周围的异教群众被激怒，几乎要冲向他。\par

5. 但由于天主的助佑之手，他奇妙地完全逃脱了。这同样的神圣和天上的能力，一次又一次地，数不胜数，因他对基督的言语所怀的巨大热忱和勇敢，在如此危险中保护了他。异教徒对他的敌意如此之大，因为有许多人通过他接受了神圣信仰的教导，以至于他们派兵包围了他居住的房屋。\par

6. 就这样，迫害日益加剧地焚烧着他，以至于整个城市都无法容纳他；但他从一个房屋搬到另一个房屋，因众多追随他神圣教导的人而到处被驱赶。因为他的生活也根据真正哲学实践展现出正确和令人敬佩的品行。\par

7. 因为据说他的生活方式如同他的教导，他的教导也如同他的生活。因此，借着与他同在的天主的能力，他激发了许多人拥有同样的热忱。\par

8. 但当他看到更多的人前来向他求教，而要理学校已由当时主持教会的\Person{德默特琉}单独托付给他时，他认为语法科学的教学与神圣主题的训练不相符，便立即放弃了他的语法学校，认为它无益且阻碍神圣学问。\par

9. 然后，出于适当的考虑，为了不依赖他人的帮助，他卖掉了自己拥有的所有有价值的古代文学书籍，满足于每天从买主那里得到四个奥波。他这样哲学地生活了许多年，摒弃了所有青年欲望的刺激。整整一天，他忍受着不小的锻炼；在夜晚的大部分时间里，他致力于研究圣经。他以最哲学的生活方式尽可能地克制自己；有时通过禁食的锻炼，有时通过有限时间的睡眠。出于热忱，他从不睡在床上，而是睡在地上。\par

10. 最重要的是，他认为应该遵守救主在福音中的话，其中他劝诫不要有两件外衣，不要穿鞋，也不要为未来忧虑。\par

11. 他怀着超越年龄的热忱，忍受寒冷和赤身；并且，他走向极端的贫穷，使周围的人都大为惊讶。事实上，他使许多朋友伤心，因为他们想与他分享自己的财产，但他们看到他为了教导神圣事物而忍受的辛劳。\par

12. 但他并没有放松他的坚忍。据说他有许多年走路从不穿鞋，并且许多年戒酒，以及所有超出必要食物之外的东西；以至于他面临身体崩溃和损坏体质的危险。\par

13. 通过向看到他的人展示如此哲学生活的证据，他激发了许多学生拥有同样的热忱；因此，甚至异教徒中杰出的人物和追求学问与哲学的人也被引导到他的教导之下。他们中有些人从他那里在灵魂深处接受了神圣圣言的信德，在当时的迫害中变得突出；有些人被逮捕并殉道。\par

% structure: p0034 source=h2 target=subsection reason=relative-heading-rank; base=h2

\subsection{第四章 成为殉道者的奥利振学生}

1. 这些人中第一个是\Person{普鲁塔克}，就是刚才提到的。当他被带去处死时，我们所谈论的这个人在他生命尽头与他在一起，几乎被他的同胞杀死，好像他是他死亡的原因。但天主的眷顾这次也保全了他。\par

2. 在\Person{普鲁塔克}之后，奥利振学生中第二位殉道者是\Person{塞雷努斯}，他通过火证明了他所接受的信德。\par

3. 来自同一学校的第三位殉道者是\Person{赫拉克利德}，之后第四位是\Person{赫罗}。前者尚是望教者，后者则刚受洗。两人都被斩首。之后，来自同一学校的第五位是另一位\Person{塞雷努斯}，他被视为虔诚的斗士，据报道，在长时间忍受酷刑后被斩首。在妇女中，\Person{赫莱丝}在尚是望教者时死去，由火接受了洗礼，正如奥利振本人在某处所说的。\par

% structure: p0038 source=h2 target=subsection reason=relative-heading-rank; base=h2

\subsection{第五章 波塔米厄娜}

1. \Person{巴西利德}可算作这些人中的第七位。他带领著名的\Person{波塔米厄娜}去殉道，她因为了保持贞潔和童贞而忍受的许多事情，至今在本地人中仍很出名。因为她心智和身体之美都达到了完美。为基督的信德受了许多苦之后，最后经过可怕且难以言说的酷刑，她与她的母亲\Person{玛尔塞拉}被火处死。\par

2. 据说，名叫\Person{阿奎拉}的法官在她全身施加了严厉的酷刑后，最后威胁要将她交给角斗士进行身体凌辱。她稍加考虑，被问及她的决定时，她作出了一个被认为不敬的回答。\par

3. 于是她立即被宣判，\Person{巴西利德}，一位军队官员，带她去处死。但当人们试图用辱骂的话语骚扰和侮辱她时，他驱赶了侮辱她的人，对她表现出极大的怜悯和仁慈。觉察到这个人对她的同情，她鼓励他鼓起勇气，因为她会在离世后为主为他祈求，他很快就会因他对她所行的仁慈而得到赏报。\par

4. 说了这话，她英勇地忍受了酷刑，从脚底到头顶，身体各处被慢慢浇上滚烫的沥青。这就是这位著名贞女所经历的斗争。\par

5. 此后不久，巴西利德（\OriginalTerm{巴西利德}{Basilides}）因某缘故被同伴要求起誓，他声称自己根本不能起誓，因为他是基督徒，并公开承认了这一点。起初他们以为他在开玩笑，但他坚持如此宣称，便被带到法官面前，并在法官面前承认了自己的信念，于是被投入监狱。但主内的弟兄们到他那里，询问这突然而显著的决心之缘由，据传他说：波塔米埃纳（\OriginalTerm{波塔米埃纳}{Potamiæna}）在她殉道后的三天内，每夜站在他旁边，将一顶冠冕放在他头上，并且说她已为他祈求了主，且已得偿所愿，不久她就会带他同去。\par

6. 于是弟兄们给他盖上了主的印；第二天，他为主作了荣耀的见证后，被斩首。据说在那段时期，亚历山大里亚有许多人迅速接受了基督的圣言。\par

7. 因为波塔米埃纳（\OriginalTerm{波塔米埃纳}{Potamiæna}）在梦中向他们显现并劝勉他们。关于此事，这样说就足够了。\par

% structure: p0046 source=h2 target=subsection reason=relative-heading-rank; base=h2

\subsection{第六章 亚历山大里亚的克莱孟（\OriginalTerm{克莱孟}{Clement}）。}

克莱孟（\OriginalTerm{克莱孟}{Clement}）继承潘塔诺（\OriginalTerm{潘塔诺}{Pantænus}）之后，当时负责亚历山大里亚的教理讲授，以致奥力振（\OriginalTerm{奥力振}{Origen}）尚在童年时便是他的学生之一。在克莱孟所著的《杂记》（\WorkTitle{\OriginalTerm{杂记}{Stromata}}）第一卷中，他列出了一份年表，将事件下推至科莫德（\OriginalTerm{科莫德}{Commodus}）去世之时。因此很明显，该书是在塞维鲁（\OriginalTerm{塞维鲁}{Severus}）统治期间写成的，我们现在所记述的正是他统治的时代。\par

% structure: p0048 source=h2 target=subsection reason=relative-heading-rank; base=h2

\subsection{第七章 著作家犹达斯（\OriginalTerm{犹达斯}{Judas}）。}

此时，另一位著作家犹达斯（\OriginalTerm{犹达斯}{Judas}）论及达尼尔（\OriginalTerm{达尼尔}{Daniel}）中的七十个星期，将年代下推至塞维鲁（\OriginalTerm{塞维鲁}{Severus}）统治的第十年。他认为那时广为谈论的假基督（\OriginalTerm{假基督}{Antichrist}）的来临已经临近。当时对我们子民的迫害所引起的骚动，极大地扰乱了众人的心。\par

% structure: p0050 source=h2 target=subsection reason=relative-heading-rank; base=h2

\subsection{第八章 奥力振（\OriginalTerm{奥力振}{Origen}）的勇敢行为。}

1. 此时，当奥力振（\OriginalTerm{奥力振}{Origen}）在亚历山大里亚进行教理讲授时，他做了一件事，这件事显示了一种不成熟和幼稚的心智，但同时也给出了信德和节制的最高证明。因为他按过分字面和极端的意义理解了\Quote{有些阉人，是为了天国而自阉的}\BibleRef{玛 19:12}这话。同时为了成就救主的话，也为了除去不信者的一切绊脚石——因为他虽年轻，却与男女一同研读神圣之事——他用行动实现了救主的话。\par

2. 他以为此事不会被许多相识的人知道。但他虽想隐瞒，却不能将如此行动保密。\par

3. 当管理该教区的德默特琉（\OriginalTerm{德默特琉}{Demetrius}）终于得知此事时，他极为赞叹这行动的勇敢性质，又看出他的热忱和他信德的真诚，便立即劝勉他勇敢，并更加鼓励他继续他的教理讲授工作。\par

4. 他当时就是这样。但不久之后，见他兴盛起来，在众人中变得伟大而卓越，同一位德默特琉（\OriginalTerm{德默特琉}{Demetrius}）为人性的软弱所胜，便向全世界的诸位主教写信，将他的行为说成是极其愚蠢的。但当时在巴勒斯坦（\OriginalTerm{巴勒斯坦}{Palestine}）各位主教中尤为著名和卓越的凯撒勒雅（\OriginalTerm{凯撒勒雅}{Caesarea}）和耶路撒冷（\OriginalTerm{耶路撒冷}{Jerusalem}）的主教们，认为奥力振（\OriginalTerm{奥力振}{Origen}）极配得荣耀，便祝圣他为司铎。\par

5. 于是他的名声大增，他的名望传遍各处，他在德行和智慧上获得了不小的声誉。但德默特琉（\OriginalTerm{德默特琉}{Demetrius}）除了他年少时的这一行为之外，别无他话可说，便苦毒地控告他，甚至敢于将那些擢升他至司铎职的人一同牵入这些控告之中。\par

6. 这些事发生在稍后。但在当时，奥力振（\OriginalTerm{奥力振}{Origen}）在亚历山大里亚（\OriginalTerm{亚历山大里亚}{Alexandria}）无畏地继续向所有来见他的人日夜讲授神圣之事，将他全部的空闲时间不间断地用于神圣研究和他的学生们。\par

7. 塞维鲁（\OriginalTerm{塞维鲁}{Severus}）执政十八年后，由他的儿子安托宁（\OriginalTerm{安托宁}{Antoninus}）继位。在那次迫害中勇敢忍受、并蒙天主眷顾（\OriginalTerm{天主眷顾}{Providence}）通过宣认的斗争得以保全的人中，有亚历山大（\OriginalTerm{亚历山大}{Alexander}），我们已说过他是耶路撒冷（\OriginalTerm{耶路撒冷}{Jerusalem}）教会的主教。由于他在宣认基督方面的卓越，他被认为配得那主教职位，尽管他的前任纳齐苏斯（\OriginalTerm{纳齐苏斯}{Narcissus}）仍在世。\par

% structure: p0058 source=h2 target=subsection reason=relative-heading-rank; base=h2

\subsection{第九章 纳齐苏斯（\OriginalTerm{纳齐苏斯}{Narcissus}）的奇迹。}

1. 该教区的居民根据继任他的弟兄们的传统，提到纳齐苏斯（\OriginalTerm{纳齐苏斯}{Narcissus}）的许多其他奇迹；其中他们叙述了他所行的以下奇事。\par

2. 他们说，有一次在伟大的逾越节守夜礼中，当执事们整夜守望时，油用尽了。于是全体会众大为沮丧，纳齐苏斯（\OriginalTerm{纳齐苏斯}{Narcissus}）吩咐管理灯烛的人取水来给他。\par

3. 他们立即照办，他就在水上祈祷，并怀着对主的坚定信德，命令他们将水倒入灯中。他们照办后，出于完全意想不到的奇妙而神圣的能力，水的性质竟变成了油。其中一小部分被那里的许多弟兄保存至今，作为这奇事的纪念。\par

4. 他们还讲述了许多关于此人生活值得注意的其他事迹，其中如下：某些卑鄙之人不能忍受他生活的刚强和坚定，又因意识到自己许多恶行而惧怕惩罚，便设计抢先诽谤他，散布了一个可怕的谣言攻击他。\par

5. 为了说服听闻的人，他们用起誓来证实他们的控告：一个招致自己被火烧死，一个招致自己的身体被恶疾侵蚀，第三个招致双眼失明。他们虽如此起誓，却不能动摇信徒们的心，因为纳齐苏斯（\OriginalTerm{纳齐苏斯}{Narcissus}）的节制和道德生活是众所皆知的。\par

6. 但他无论如何不能忍受这些人的恶行；由于他已长期度哲学的生活，他便避开整个教会团体，隐居在荒野和隐秘之处，并在那里住了多年。\par

7. 但那审判的大眼并未对此无动于衷，反而很快俯视这些不虔敬的人，并使那他们用以束缚自己的咒诅降在他们身上。第一个人的住所因一点火星落入而夜间完全烧毁，他与其全家一同灭亡。第二人迅速被他自己咒诅的疾病覆盖，从脚底到头顶。\par

8. 但第三人察觉了其他两人所遭遇的事，惧怕万有的天主那不可避免的审判，便公开承认了他们一同策划的事。他因极大的哀恸而悔改，并持续哭泣，以致双目尽毁。这些人因他们的虚伪而受的惩罚就是如此。\par

% structure: p0067 source=h2 target=subsection reason=relative-heading-rank; base=h2

\subsection{第十章　耶路撒冷的主教们。}

\Person{纳齐苏斯}离开后，无人知道他在何处，邻近教会的主持者认为最好另祝圣一位主教。他名叫\Person{迪乌斯}。他主持了很短时间，\Person{格尔马尼奥}继任。他之后是\Person{戈尔迪乌斯}，在\Person{戈尔迪乌斯}任内，\Person{纳齐苏斯}再次出现，如同从死里复活。弟兄们立刻恳求他接受主教职位，因为众人因他的退隐和哲学，尤其因天主为他报复的惩罚，更加钦佩他。\par

% structure: p0069 source=h2 target=subsection reason=relative-heading-rank; base=h2

\subsection{第十一章　\Person{亚历山大}。}

1. 但由于\Person{纳齐苏斯}因年事已高，无法再履行其公务，天主的眷顾藉着夜间异象给他的启示，召唤上述的\Person{亚历山大}与他一同任职，\Person{亚历山大}当时是另一个教区的主教。\par

2. 于是，他仿佛受天主指引，从最初担任主教的\Place{卡帕多细亚}地方出发，前往\Place{耶路撒冷}，这是出于一个誓愿，并为了解该地的情况。他们非常热诚地接待他，并因他们夜间所见另一个异象，不肯让他回去；那异象向最热心的他们发出了最清晰的信息，因为它告诉他们，如果他們出到门外，就会接收到天主为他们预定的主教。他们这样做了，并在邻近各教会主教的一致同意下，强留他住下。\par

3. \Person{亚历山大}本人在他写给\Place{安提诺依}人的私人信件中（这些信件仍在我们中间保存着），提及\Person{纳齐苏斯}与他共同的主教职务，在信尾这样写道：\par

\Quote{\Person{纳齐苏斯}问候你们，他是在我之前主持这里主教职务的人，现与我一同祈祷，一百一十六岁；他像我一样劝勉你们要同心合意。}\par

这些事就这样发生了。但\Person{塞拉比翁}逝世后，\Person{阿司克莱彼阿得斯}继任\Place{安提约基雅}教会的主教，他本人曾在迫害期间于明认信仰者中备受尊敬。\Person{亚历山大}在写给\Place{安提约基雅}教会的信中提到了他的任命，这样写道：\par

\Quote{\Person{亚历山大}，耶稣基督的仆人和囚犯，致书\Place{安提约基雅}蒙福的教会，在主内问安。主使我在被囚期间轻松愉快，因我得知，藉着天主的眷顾，\Person{阿司克莱彼阿得斯}在真信德方面极有资格，已在\Place{安提约基雅}你的圣教会担任主教之职。}\par

6. 他表明他是通过\Person{克莱孟}寄出这封信的，在信末这样写道：\par

\Quote{我可敬的弟兄们，我通过蒙福的司铎\Person{克莱孟}，一位有德且被认可的人，将这封信寄给你们；你们自己也认识他，也必认出他来。他在这里，藉着主人的眷顾和照管，坚固并建立了主的教会。}\par

% structure: p0078 source=h2 target=subsection reason=relative-heading-rank; base=h2

\subsection{第十二章　\Person{塞拉比翁}及其现存著作。}

1. 其他人很可能保存了\Person{塞拉比翁}文学生涯的其他纪念作品，但传到我们手中的，只有那些写给某个\Person{多米尼诺斯}的信件——他在迫害期间从对基督的信德堕落至犹太人私自的崇拜；以及写给\Person{庞提乌斯}和\Person{卡里库斯}这些教会人士的信件，还有其他给不同人的信；另有一部他关于所谓\WorkTitle{伯多禄福音}的作品。\par

2. 他写这最后一部作品，为驳斥该福音所含的虚谎，因为\Place{罗苏斯}堂区有些人被它引入异端的概念。也许值得从他的作品中摘录一些简短的片段，以显示他对该书的意见。他这样写道：\par

\Quote{3. 因为我们，弟兄们，接纳\Person{伯多禄}和其他宗徒如同基督；但我们明智地拒绝那些假托他们名写的著作，因为知道这些并未传给我们。}\par

\Quote{4. 当我访问你们时，我以为你们都持守真信德，并且因为我未曾读过他们以\Person{伯多禄}之名提出的那部福音，我说：如果这是唯一引起你们争议的事，就让人读它吧。但现在从所告知我的情况得知，他们的心思陷入某种异端，我将快快再回到你们那里。因此，弟兄们，请期待我很快到来。}\par

\Quote{5. 但是，弟兄们，你们将从写给你们的信中得知，我们察觉到\Person{马西安努斯}异端的本质，并且由于他不明白自己所说的话，他自相矛盾。}\par

\Quote{6. 因为我们从其他勤于研究这部福音的人那里获得了它，即从最初使用它的人的后继者——我们称他们为\OriginalTerm{幻象论者}{Docetæ}（因为他们大多数意见与那一学派的教导有关）——我们得以通读全书，并发现许多符合救主真道的内容，但有些内容加添了那道，我们已在后面为你们指出过。}关于\Person{塞拉比翁}就说到这里。\par

% structure: p0085 source=h2 target=subsection reason=relative-heading-rank; base=h2

\subsection{第十三章　\Person{克莱孟}的著作。}

1. \Person{克莱孟}的八卷\WorkTitle{杂记}都保存在我们中间，他给予它们以下标题：\Quote{\WorkTitle{\Person{提图斯·弗拉维乌斯·克莱孟}的关于真哲学的诺斯替注释杂记}}。\par

2. 题为\WorkTitle{雏形}的书卷数目相同。在其中他提及\Person{庞泰努斯}的名字作为他的老师，并陈述了他的意见和传承。\par

3. 此外，尚有他的\WorkTitle{致希腊人的劝勉言论}，三卷\WorkTitle{教师}一书，另一部题为\WorkTitle{富者如何得救？}，论逾越节的作品，论斋戒和论毁谤的讨论，\WorkTitle{论忍耐的劝勉言论，或致新受洗者}，以及题为\WorkTitle{教会法规，或反犹化主义者}的作品——他将后者献给上述的主教亚历山大。\par

4. 在\WorkTitle{杂论}中，他不仅广泛论述了圣经，而且凡希腊作家所说对他似乎有益之处，他也加以引用。\par

5. 他阐明了众多希腊人和外邦人的意见。他也驳斥了异端首领的虚假教义，除此以外，还审查了很大一部分历史，给我们展示了多种学问的范例；他将哲学家的观点与所有其他内容相融合。很可能正因如此，他才给自己的作品取了\WorkTitle{杂论}这一恰当的标题。\par

6. 他在这些作品中也使用了来自有争议的经典的书证，即所谓的\BibleBook{撒罗满}{智慧书}、\BibleBook{}{耶稣·息辣之子}、\BibleBook{}{致希伯来人书}，以及巴拿巴、格肋孟和犹达的书信。\par

7. 他也提到了塔提安\WorkTitle{致希腊人的言论}，并称卡西安诺斯为一部年代学著作的作者。他引用了犹太作家斐罗、阿里斯托布鲁斯、若瑟夫、德默特琉和欧波勒摩斯，指出他们全部在其著作中表明，梅瑟与犹太民族在希腊人最早起源之前就已存在。\par

8. 这些书还包含许多其他学识。在其中的第一卷，作者自称仅次于宗徒的继承者。\par

9. 在这些著作中，他还应许要撰写一部\WorkTitle{论《创世纪》的注释}。在他的论逾越节的作品里，他承认自己曾受朋友们催促，将自古代长老那里听到的传统写下来传给后世；在同一作品中，他提到了默里托和依肋内以及某些其他人，并摘录了他们的著作。\par

% structure: p0095 source=h2 target=subsection reason=relative-heading-rank; base=h2

\subsection{第十四章　他所提及的圣经}

1. 简言之，他在\WorkTitle{纲要}中给出了所有正典圣经的摘要叙述，并未忽略有争议的书卷——我指的是\BibleBook{犹达}{书}和其他公函，以及\BibleBook{巴拿巴}{书}和所谓的\BibleBook{伯多禄}{默示录}。\par

2. 他说\BibleBook{}{致希伯来人书}是保禄的作品，并且是用希伯来语写给希伯来人的；但路加仔细地翻译并发表给希腊人，因此这封书信与\BibleBook{宗徒}{大事录}具有相同的表达风格。\par

3. 但他又说，\Quote{保禄宗徒}几个字可能没有放在开头，因为写给对保禄存有偏见和怀疑的希伯来人时，他明智地不愿一开始就用自己的名字令他们反感。\par

4. 再往下他说道：\Quote{但如今，正如那位有福的长老所说，既然主是全能者的宗徒，被派遣到希伯来人那里，保禄作为被派遣到外邦人那里的人，由于谦虚，并没有自称为希伯来人的宗徒，出于对主的尊敬，并且因为他作为外邦人的宣讲者和宗徒，出于丰富而写信给希伯来人。}\par

5. 此外，在同一著作中，格肋孟按以下方式给出了最早长老关于福音书次序的传统：\par

6. 他说，包含族谱的福音书是首先写成的。按马尔谷所记载的福音有这样一段缘由：当伯多禄在罗马公开宣讲圣言，并藉圣神宣报福音时，许多在场的人请求长期跟随他、记住他言论的马尔谷将这些记录下来。马尔谷写完福音后，就把它交给了那些请求的人。\par

7. 伯多禄得知此事，既没有直接禁止，也没有鼓励。但最后，若望，因他看到外在的事实已在福音中显明，受朋友们催促并被圣神感动，写了一部属灵的福音。这就是格肋孟的叙述。\par

8. 此外，前面提到的亚历山大在给奥力振的一封信中提到了格肋孟，同时也提到了庞代诺，视之为他的亲密朋友。他写道：\par

\Quote{因为，如你所知，这是天主的旨意，要我们之间祖传的友谊坚定不移，甚至更加温暖和坚固。}\par

\Quote{因为我们深知那些在我们前面行路、我们不久将与之同在的有福父亲：庞代诺——真正有福的人和老师，以及圣格肋孟——我的老师和恩人，还有其他像他们一样的人，通过他们我认识了你——万事中的最好者，我的老师和兄弟。}\par

10. 关于这些事，就谈到这里。但阿当曼提——这也是奥力振的一个名字——在则斐琳担任罗马主教时访问了罗马，\Quote{渴望，}如他自己在某处所说，\Quote{见到那最古老的罗马教会。}\par

11. 在那里短暂停留后，他返回了亚历山大里亚。在那里，他以极大的热忱履行教理讲授的职责；当时那里的主教德默特琉催促并甚至恳请他勤勉工作，以裨益弟兄们。\par

% structure: p0108 source=h2 target=subsection reason=relative-heading-rank; base=h2

\subsection{第十五章　赫辣克拉斯}

1. 但当他看到自己没有时间更深入研究神圣事物、查考并诠释圣经，同时也指导前来见他的人——因为他们一个接一个地从早到晚来向他学习，几乎不给他喘息的机会——于是他分开了群众。他从他所熟知的人中挑选了赫辣克拉斯——此人热切钻研神圣事物，并且在其他方面学识渊博，精通哲学——让他成为自己在教导工作上的同伴。他将初学者的基础训练托付给他，而为自己保留了更进深者的教导。\par

% structure: p0110 source=h2 target=subsection reason=relative-heading-rank; base=h2

\subsection{第十六章　奥力振对圣经的勤勉研习}

一、奥力振研究天主圣言如此热切和勤奋，以至他学习了希伯来语，并将犹太人手中的原始希伯来圣经据为己有。他也查考了除七十贤士译本之外的其他圣经译者们的著作。除了众所周知的阿奎拉、息玛谷和德敖多西翁的译本之外，他还发现了某些其他译本——这些译本自远古以来一直被隐藏——我不知道在哪些偏僻的角落里——通过他的搜寻，他将它们公诸于世。\par

二、由于他不知道作者，他仅仅指出，这一个是在亚克兴附近的尼科波利发现的，那一个是在别处发现的。\par

三、在圣咏集的\OriginalTerm{六文本合参}{Hexapla}中，除了四个著名的译本外，他还加入了第五、第六乃至第七种译文。他提到其中有一种是在塞维鲁之子安敦尼时代，在耶里哥的一个坛子里找到的。\par

四、他将所有这些译本收集起来，分成段落，彼此相对排列，并附上希伯来文本本身。这样，他就为我们留下了所谓的\OriginalTerm{六文本合参}{Hexapla}的抄本。他还另外编排了阿奎拉、息玛谷和德敖多西翁与\OriginalTerm{七十贤士译本}{Septuagint}合在一起的版本，即\OriginalTerm{四文本合参}{Tetrapla}。\par

% structure: p0115 source=h2 target=subsection reason=relative-heading-rank; base=h2

\subsection{第十七章　译者息玛谷。}

关于这些译者，应当说明，息玛谷是一个\OriginalTerm{厄彼翁派}{Ebionite}。但是，所谓的厄彼翁派异端主张基督是若瑟和玛利亚的儿子，视他为纯粹的人，并强烈坚持以犹太方式遵守法律，正如我们在本历史中已经看到的那样。息玛谷的注释至今尚存，在其中他似乎通过攻击玛窦福音来支持这一异端。奥力振说，他是从一位名叫犹利安纳的妇女那里得到这些注释以及息玛谷对圣经的其他注释的，他说，她是从息玛谷本人那里继承这些书的。\par

% structure: p0117 source=h2 target=subsection reason=relative-heading-rank; base=h2

\subsection{第十八章　盎博罗削。}

一、大约此时，信奉华伦提奴异端的盎博罗削被奥力振所展示的真理说服，心灵仿佛被光照亮，接受了教会的正统教义。\par

二、还有许多其他人，被奥力振到处传扬的学识声望所吸引，来到他那里，考验他在圣经文学方面的技巧。许多异端者以及不少最杰出的哲学家勤奋地跟随他学习，不仅接受他关于神圣事物的教导，也接受世俗哲学的教导。\par

三、因为当他发现某些人具有出众的才智时，他还教导他们哲学分支——几何学、算术以及其他预备学科——然后进展到哲学家的体系并解释他们的著作。他对每一种都加以观察和评注，以至于即使在希腊人中间，他也成为一位伟大的哲学家而闻名。\par

四、他教导许多学问较少的人学习普通学科，并说这对于他们学习和理解\OriginalTerm{圣经}{Divine Scriptures}会有不小的帮助。因此他认为自己特别有必要精通世俗和哲学学识。\par

% structure: p0122 source=h2 target=subsection reason=relative-heading-rank; base=h2

\subsection{第十九章　与奥力振有关的情节。}

一、他那个时代的希腊哲学家可以证明他在这些学科上的造诣。我们在他们的著作中频繁地看到他。有时他们将作品献给他；有时又将自己的劳动成果作为导师呈交给他评判。\par

二、我们何必说这些呢？就连与我们同时代、居住在西西里并著书攻击我们的波斐利，也试图借此诋毁\OriginalTerm{圣经}{Divine Scriptures}，他提到了那些解释圣经的人；由于他无法从教义本身找到任何可谴责的借口，便转向辱骂和诽谤这些解释者，尤其试图中伤奥力振，他说自己年轻时认识他。\par

三、但实际上，他无意中赞扬了这个人：在有些他不得不说真话的地方说了真话；但在认为不会被察觉的地方说了谎话。有时他指责他是基督徒，有时又描述他在哲学学识上的精湛。但请听他自己的话：\par

四、\Quote{有些人渴望为犹太人的圣经的粗鄙找到一种解决办法，而不是放弃它们，便求助于与所写的文字不一致且不协调的解释，这些解释非但不能为那些外国人提供辩护，反而包含了对他们自己的认可和赞美。因为他们吹嘘梅瑟的简单话语是谜语，并将它们视为充满隐藏奥秘的神谕；他们用愚昧迷惑了心灵的判断，并做出自己的解释。}接着他又说：\par

五、这种荒谬的一个例子就是我年轻时遇到的那个人，他当时因他所留下的著作而大受赞扬，至今仍如此。我指的是奥力振，他深受这些教义的教师们尊敬。\par

六、因为这个人曾是安莫尼的学生，安莫尼在当时任何人中哲学造诣最深，奥力振从这位老师那里获益良多，获得了科学知识；但在正确的人生选择上，他却走上了与他相反的道路。\par

七、因为安莫尼本是基督徒，由基督徒父母抚养长大，当他投身研究和哲学时，立即遵循法律所要求的生活。但奥力振，作为在希腊文学中受教育的希腊人，却转向了野蛮人的狂妄。他将所学得的学问带去贩卖，在生活中行为如同基督徒，违背法律；但在对物质事物和神的看法上，却如同希腊人，将希腊的教导与外邦寓言混合在一起。\par

八、因为他不断地研读柏拉图，并埋头于努默纽、克洛纽、阿颇罗芳、隆基诺、莫德拉图和尼各马可的著作，以及那些在毕达哥拉斯派中著名的人物。他还使用斯多葛派凯勒孟和科尔努图的书籍。通过他们，他熟悉了希腊奥秘的比喻解释，并将其应用于犹太人的圣经。\par

9. 这些话是波斐利在其著作\WorkTitle{反对基督徒}的第三卷中说的。他如实地说出了此人的勤勉和学识，但当他声称此人背离了希腊人，并且阿摩尼乌斯从虔诚生活堕入异教习俗时，他显然说了假话（因为反对基督徒的人还有什么事做不出来呢？）。\par

10. 因为基督的教义是由他的父母传授给奥力振的，正如我们上面所示。而阿摩尼乌斯一生保持神圣的哲学坚定不移、纯正无杂。他流传至今的作品证明了这一点，因为他留下的著作使他在许多人中备受赞誉。例如，题为\WorkTitle{摩西与耶稣的和谐}的作品，以及其他一些为学者所拥有的著作。\par

11. 这些事足以证明假告发者的诽谤，也足以证明奥力振在希腊学问上的造诣。他在一封书信中为自己在这方面勤奋学习辩护，反驳了一些指责他的人，信中写道：\par

\Quote{12. 当我投身于圣言时，我造诣的名声传开了，当异端者和精通希腊学问、特别是哲学的人来到我这里时，似乎有必要检验异端的教义，以及哲学家关于真理的说法。}\par

\Quote{13. 在这方面，我们效法了庞泰努斯，他因在诸如此类的事上充分的准备而在我们之前造福了许多人；也效法了赫拉克拉斯，他如今是亚历山大城司铎团的一员。我发现他与哲学老师在一起，在我开始听那些科目的讲座之前，他已经跟从这位老师五年了。}\par

\Quote{14. 尽管他先前穿着普通服装，但他脱下了它，穿上并至今仍穿着哲学家的外衣；他继续热切地研究希腊著作。}\par

他说这些话是为自己学习希腊文学辩护。\par

15. 大约此时，奥力振仍在亚历山大城，一名士兵前来，将一封阿拉伯总督的信交给堂区主教狄米特里和当时的埃及总督，请求他们尽快派奥力振去与他会面。二人派遣他前往阿拉伯。他在短时间内完成了访问的目的，便回到了亚历山大城。\par

16. 但不久以后，城中爆发了一场大战，他便离开了亚历山大城。他以为留在埃及不安全，于是去了巴勒斯坦，住在凯撒勒雅。在那里，该国教会的众主教请求他公开宣讲并阐释圣经，尽管他尚未被祝圣为司铎。\par

17. 这一点从耶路撒冷主教亚历山大和凯撒勒雅的德奥克提斯图斯就此事写给狄米特里信中的自我辩护可以明显看出，他们写道：\par

\Quote{他在信中声称，这样的事从未听说过，也从未发生过，即平信徒在主教面前讲道。我不明白他如何说出如此明显虚假的话。}\par

\Quote{18. 因为每当发现有能力教导弟兄的人时，圣洁的主教们就会劝勉他们向民众讲道。例如在拉兰达，内昂劝勉欧厄尔皮斯；在依科尼雍，策尔苏斯劝勉保利努斯；在叙纳达，亚弟古斯劝勉德奥多鲁斯，这些是我们的有福弟兄。可能在其他我们所不知道的地方也是如此。}\par

他以这种方式在年轻时便受到尊敬，不仅来自同国人，也来自外国主教。\par

19. 但狄米特里写信派人去请他，并通过教会成员和执事敦促他返回亚历山大城。于是他返回并恢复了惯常的职务。\par

% structure: p0145 source=h2 target=subsection reason=relative-heading-rank; base=h2

\subsection{第二十章 当时作者现存的著作}

1. 当时教会中有许多博学之士，他们彼此之间的书信得以保存且易于获取。这些书信一直保存在埃利亚的图书馆中，该图书馆由当时管理该教会的亚历山大建立。我们能够从该图书馆收集到当前工作的材料。\par

2. 其中，贝里鲁斯除了书信和论文外，还为我们留下了各种优美的著作。他是阿拉伯博斯特拉的主教。同样，管理另一教会的希波吕托斯也留下了一些著作。\par

3. 我们还有一位博学之士卡尤斯的对话录，这是在罗马在则斐林时期与主张弗里吉亚异端的普罗克路斯进行的。在其中，他遏制了对手在提出新圣经时的鲁莽和大胆。他只提到圣宗徒的十三封书信，没有将希伯来人书计算在内。直到今日，罗马人中还有人不认为这是宗徒的著作。\par

% structure: p0149 source=h2 target=subsection reason=relative-heading-rank; base=h2

\subsection{第二十一章 当时知名的众主教}

1. 安东尼努斯在位七年六个月后，马克里努斯继位。他执政仅一年，便由另一位安东尼努斯继位。在他（指后一位安东尼努斯）的第一年，罗马主教则斐林在位十八年后去世，加里斯都接受了主教职。\par

2. 他（指加里斯都）继续任职五年，由乌尔巴诺继任。此后，安东尼努斯在位仅四年，亚历山大成为罗马皇帝。此时，斐肋托也在安提约基雅教会继阿斯克肋皮亚德斯之后任职。\par

3. 皇帝的母亲名叫玛梅亚，是一位极其虔诚、度宗教生活的妇人，若曾有这样的人的话。当奥力振的名声传遍各处，甚至传到她耳中时，她非常想见这个人，尤其是想验证他对神圣事物的著名理解。\par

4. 她在安提约基雅逗留了一段时间，便派一支军事护卫去请他。他与她在一起待了一段时间，向她展示了许多为了主的荣耀和神圣教导卓越性的事，然后急忙回到他惯常的工作。\par

% structure: p0154 source=h2 target=subsection reason=relative-heading-rank; base=h2

\subsection{第二十二章 希波吕托斯流传至今的著作}

1. 当时，希波吕托斯除了许多其他论文外，还写了一部关于逾越节的著作。他在其中给出了一份年表，并提出了一个十六年的逾越节规范，将时间降至亚历山大皇帝的第一年。\par

2. 他的其他著作中，以下作品流传到我们手中：\WorkTitle{六日工程}、\WorkTitle{六日工程后作品}、\WorkTitle{驳马尔基翁}、\WorkTitle{雅歌释义}、\WorkTitle{厄则克耳书部分释义}、\WorkTitle{论逾越节}、\WorkTitle{驳一切异端}；你还能找到许多其他为多人保存的著作。\par

% structure: p0157 source=h2 target=subsection reason=relative-heading-rank; base=h2

\subsection{第23章 \Person{奥力振}的热忱及其晋升司铎}

1. 那时，\Person{奥力振}开始注释圣经，受到\Person{盎博罗削}的催促，后者以无数激励敦促他，不仅以言语劝勉，还提供充足的方法。\par

2. 因为他对七个以上的笔录员口授，他们按指定时间轮流替换。他还雇用了同样多的抄写员，此外还有擅长优美书写的女青年。所有这些，\Person{盎博罗削}都提供了充足的费用，表现出难以言表的勤奋和热忱，为神圣的圣言，以此特别敦促他准备注释。\par

3. 在这些事进行期间，\Person{乌尔巴诺}在担任罗马教会主教八年后，由\Person{彭提安}继任；\Person{则彼尼}继\Person{斐理托}在\Place{安提约基雅}任职。\par

4. 此时，\Person{奥力振}因教会事务的紧迫需要被派往\Place{希腊}，途经\Place{巴勒斯坦}，并在\Place{凯撒勒雅}由该地区的主教祝圣为司铎。为此事对他引起的争议，以及教会首领们对此事的裁决，连同他在盛年时出版的关于圣言的其他著作，都需要单独论述。我们在为他所写的\WorkTitle{辩护书}第二卷中已有所叙述。\par

% structure: p0162 source=h2 target=subsection reason=relative-heading-rank; base=h2

\subsection{第24章 他在\Place{亚历山大里亚}准备的注释}

1. 不妨补充一点，他在解释\WorkTitle{\BibleBook{若望}{福音}}的第六卷中说，前五卷是在\Place{亚历山大里亚}准备的。关于整部福音的著作，流传至今的只有二十二卷。\par

2. 在\WorkTitle{\BibleBook{创}{世纪}}注释第九卷中（全书共十二卷），他说不仅前八卷是在\Place{亚历山大里亚}完成的，而且关于前二十五篇圣咏和\WorkTitle{\BibleBook{哀}{歌}}的注释也是如此。后面五卷流传至今。\par

3. 在这些书中，他还提及他的\WorkTitle{论复活}两卷。在离开\Place{亚历山大里亚}之前，他还写了\OriginalTerm{论首要原理}{De Principiis}；以及题为\OriginalTerm{杂集}{Stromata}的十篇论文，是在\Person{亚历山大}在位期间在同一座城市写成的，正如各卷前的亲笔批注所示。\par

% structure: p0166 source=h2 target=subsection reason=relative-heading-rank; base=h2

\subsection{第25章 他对正典圣经的评述}

1. 在解释第一篇圣咏时，他列出旧约正典圣经的目录如下：\par

\Quote{应当指出，正典书卷，按希伯来人所传，共有二十二卷，与他们的字母数目相对应。} 往下他又说：\par

\Quote{2. 希伯来人的二十二卷书如下：我们称之为\OriginalTerm{创世纪}{Bresith}，意为\InnerQuote{起初}；\OriginalTerm{出埃及记}{Welesmoth}，即\InnerQuote{这些是名字}；\OriginalTerm{肋未记}{Wikra}，\InnerQuote{他呼唤}；\OriginalTerm{户籍记}{Ammesphekodeim}；\OriginalTerm{申命记}{Eleaddebareim}，\InnerQuote{这些话是}；\Person{若苏厄}\OriginalTerm{纳维之子}{Josoue ben Noun}；\OriginalTerm{民长纪与卢德传}{Saphateim}（二者合成一卷）；\OriginalTerm{撒慕尔纪上下}{Samouel}，即\InnerQuote{天主所召}；\OriginalTerm{列王纪上下}{Wammelch David}，即\InnerQuote{达味的王国}；\OriginalTerm{编年纪上下}{Dabreïamein}，即\InnerQuote{日录}；\OriginalTerm{厄斯德拉上下}{Ezra}，即\InnerQuote{助手}；\OriginalTerm{圣咏集}{Spharthelleim}；\OriginalTerm{箴言（撒罗满的）}{Meloth}；\OriginalTerm{训道篇}{Koelth}；\OriginalTerm{雅歌}{Sir Hassirim}（不像有些人以为的\OriginalTerm{歌中之歌}{Song of Songs}）；\OriginalTerm{依撒意亚}{Jessia}；\OriginalTerm{耶肋米亚}{Jeremia}（连同哀歌和书信合成一卷）；\OriginalTerm{达尼尔}{Daniel}；\OriginalTerm{厄则克耳}{Jezekiel}；\OriginalTerm{约伯传}{Job}；\OriginalTerm{艾斯德尔传}{Esther}。除此之外还有\OriginalTerm{玛加伯}{Maccabees}，题为\OriginalTerm{?}{Sarbeth Sabanaiel}。}他在上述著作中列出这些。\par

3. 在他的\WorkTitle{\BibleBook{玛窦}{福音}注释}第一卷中，他坚持教会的正典，证明他只知晓四部福音，写道如下：\par

4. 在天下天主的教会中唯一无可争议的四部福音中，我由传统得知，第一部是由\Person{玛窦}所写，他曾是税吏，后来成为耶稣基督的宗徒，是为犹太归化者准备的，并用希伯来语出版。\par

5. 第二部是由\Person{马尔谷}，他依照\Person{伯多禄}的指示写成，\Person{伯多禄}在其公函中承认他为儿子，说：\BibleQuote{在巴比伦与你们一同被选的教会问候你们，我的儿子\Person{马尔谷}也问候你们。}\BibleRef{伯前 5:13}\par

6. 第三部是由\Person{路加}，这部福音受到\Person{保禄}的称赞，是为外邦归化者撰写的。最后一部是由\Person{若望}。\par

7. 在他的\WorkTitle{\BibleBook{若望}{福音}注释}第五卷中，他论及宗徒书信如此说：但那\BibleQuote{曾使我们适于作新约的仆役，不是凭文字，而是凭圣神}\BibleRef{格后 3:6}，即\Person{保禄}，他\BibleQuote{从耶路撒冷直到依里黎雍，到处传了基督的福音}\BibleRef{罗 15:19}，并没有给他所教导的所有教会写信，而给他写信的那些教会，他也只写了寥寥数行。\par

8. 至于\Person{伯多禄}，基督的教会建立在他身上，\BibleQuote{地狱的权势不能胜过她}\BibleRef{玛 16:18}，他留下了一封公认的书信；或许还有第二封，但这一点存疑。\par

9. 何必再提那曾靠在耶稣怀中的\Person{若望}呢？他留给我们一部福音，尽管他承认他可能写许多，以致世界容不下。他还写了\WorkTitle{默示录}，但被命令保持沉默，不写出七声雷响的话语。\par

10. 他还留下一篇极短的书信；或许还有第二和第三篇，但并非所有人都认为它们是真作，而且它们总共不到一百行。\par

11. 此外，他在关于\WorkTitle{致希伯来人书}的讲道中作了如下陈述：那封题为\WorkTitle{致希伯来人书}的书信，其言辞风格并不像那位自称\BibleQuote{拙于言词}\BibleRef{格后 11:6}的宗徒那样粗鲁，而是希腊文更为纯净，任何能辨别措辞差异的人都会承认。\par

12. 再者，这封书信的思想令人钦佩，并不逊于公认的宗徒著作，任何仔细审视宗徒文本的人都会承认。\par

13. 他在后面又说：\InnerQuote{如果我说出我的看法，我该说，思想属于宗徒，但措辞与语句属于某个记住宗徒教导的人，他闲暇时记下了他老师所说的话。因此，若有教会认为这封书信出自保禄，就让它为此受到称赞。因为古人并非没有理由就把它当作保禄的流传下来。}\par

14. 但究竟是谁写了这封书信，实际上，只有天主知道。据我们之前一些人的说法，罗马主教克莱孟写了这封书信，另一些人所说是路加写的，他是\BibleBook{}{福音}和\BibleBook{宗徒}{大事录}的作者。关于这些事，就到此为止。\par

% structure: p0182 source=h2 target=subsection reason=relative-heading-rank; base=h2

\subsection{第二十六章  海辣克拉斯成为亚历山大里亚主教}

在所述统治的第十年，奥力振从亚历山大里亚迁往凯撒勒雅，将当地的要理学校交给海辣克拉斯。不久之后，亚历山大里亚教会的主教德默特琉去世，他担任此职务整整四十三年，海辣克拉斯接替了他。此时，卡帕多细亚的凯撒勒雅主教斐米利亚诺声名显赫。\par

% structure: p0184 source=h2 target=subsection reason=relative-heading-rank; base=h2

\subsection{第二十七章  主教们如何看待奥力振}

他对奥力振如此热心，以致力劝他来那地区以有益于教会，而且他还去犹太拜访奥力振，与他同住一段时间，以增进对神圣事物的认识。耶路撒冷主教亚历山大和凯撒勒雅主教德奥克提斯督常跟随他，视他为自己唯一的老师，并允许他讲解圣经，以及履行其他属于教会讲论的职务。\par

% structure: p0186 source=h2 target=subsection reason=relative-heading-rank; base=h2

\subsection{第二十八章  马克西米诺统治下的迫害}

罗马皇帝亚历山大在位十三年后结束统治，由马克西米诺凯撒继位。由于他对亚历山大家族（其中包含许多信徒）的仇恨，他开始了一场迫害，命令只处死教会的统治者，认为他们应对福音教导负责。于是，奥力振撰写了\WorkTitle{论殉道}，并将其献给盎博罗削和凯撒勒雅堂区的司铎普洛托克特督斯，因为在迫害中，他们两人都遭遇了异常的艰难，据说他们在马克西米诺（在位仅三年）统治期间以宣认著称。奥力振在他的\WorkTitle{若望福音注释}第二十二卷以及几封书信中指出了这是迫害的时间。\par

% structure: p0188 source=h2 target=subsection reason=relative-heading-rank; base=h2

\subsection{第二十九章  法比亚诺，被天主奇妙地指定为罗马主教}

1. 戈尔狄亚诺继马克西米诺之后成为罗马皇帝；担任罗马教会主教六年的彭提亚诺由安泰洛接替。安泰洛任职一个月后，法比亚诺继任。\par

2. 他们说，在安泰洛去世后，法比亚诺和其他人从乡下来到罗马，并留在那里。在那里，他通过一种极为奇妙的、来自天主和天上恩宠的显现而被选为这一职务。\par

3. 当所有弟兄聚集投票选举谁将接替教会主教职务时，许多人心中想到了几位著名而可敬的人，但法比亚诺虽然在场，却不在任何人心中。然而，他们叙述说，突然一只鸽子飞下，落在他的头上，类似于圣神以鸽子的形状降临在救主身上。\par

4. 于是，全体民众仿佛受同一圣神感动，以极大的热忱和一致高呼他配得，就毫不犹豫地把他带走，安置在主教的座位上。\par

5. 大约在那时，安提约基雅主教则比诺去世，巴比拉继任。在亚历山大里亚，海辣克拉斯在德默特琉之后接受了主教职务，由狄约尼削接替他要理学校的职责，狄约尼削也是奥力振的学生之一。\par

% structure: p0194 source=h2 target=subsection reason=relative-heading-rank; base=h2

\subsection{第三十章  奥力振的学生}

当奥力振在凯撒勒雅继续他惯常的职责时，许多学生不仅从附近，也从其他地方来到他那里。其中，德奥多禄（他在我们时代的主教中以额我略之名著称）和他的兄弟雅典诺多禄，我们知道他们特别著名。发现他们深感兴趣于希腊和罗马学问后，他向他们灌输对哲学的热爱，并引导他们把从前的热忱转向研究神学。他们与他同住了五年，在神圣事物上取得了如此进步，以至于尽管他们还很年轻，两人都被授予了本都各教会的主教职位。\par

% structure: p0196 source=h2 target=subsection reason=relative-heading-rank; base=h2

\subsection{第三十一章  阿非利加诺}

1. 这时，名为\WorkTitle{百夜战}的著作的作者阿非利加诺也很有名。现存有他写给奥力振的一封书信，对\BibleBook{达尼尔}{}中苏撒纳的故事表示怀疑，认为是伪造的和虚构的。奥力振非常详尽地回答了这封信。\par

2. 同一阿非利加诺的其他著作流传到我们手中的是五卷\WorkTitle{年代记}，这是一部准确而费力的著作。他在其中说，他去了亚历山大里亚，是因为海辣克拉斯享有巨大声誉——海辣克拉斯在哲学研究和其他希腊学问方面尤其卓越，我们已提到他被任命为那里教会的主教。\par

3. 现存还有同一阿非利加诺写给雅利斯提德的另一封书信，关于\BibleBook{玛窦}{福音}和\BibleBook{路加}{福音}中基督族谱的所谓矛盾。在这封信中，他根据一个流传到他手中的记载，清楚地表明了两位福音作者的一致，我们已在本著作第一卷的适当位置给出了这个记载。\par

% structure: p0200 source=h2 target=subsection reason=relative-heading-rank; base=h2

\subsection{第三十二章  奥力振在巴勒斯坦凯撒勒雅编著的注释}

1. 大约在这时，奥力振编写了关于\BibleBook{依撒意亚}{先知书}和关于\BibleBook{厄则克耳}{先知书}的注释。关于前者，流传到我们手中的有三十卷，直到\BibleBook{依撒意亚}{先知书}的第三部分，即到旷野中野兽的异象；关于\BibleBook{厄则克耳}{先知书}有二十五卷，他关于这位完整先知所写的全部就是这些。\par

2. 他当时在雅典，完成了关于\BibleBook{厄则克耳}{先知书}的著作，并开始编写关于\BibleBook{}{雅歌}的注释，一直进行到第五卷。他返回凯撒勒雅后，也完成了这些注释，共十卷。\par

3. 但我们为何要在这部史书中提供此人作品的精确目录，这原本需要一部单独的论著？我们在本时代的圣殉道者庞菲禄的生平记述中已提供了这一目录。在展示了庞菲禄在神圣事物上的勤奋之巨后，我们在该书中也给出了他所收集的奥力振及其他教会作家的著作目录。凡愿了解者，可轻易从中得知奥力振的哪些作品已传至我们。但我们现在必须继续我们的史书。\par

% structure: p0204 source=h2 target=subsection reason=relative-heading-rank; base=h2

\subsection{第三十三章 贝里禄斯的错误}

1. 我们最近提到的阿拉伯博斯拉主教\Person{贝里禄斯}，背离了教会准则，试图引入与信仰相异的思想。他竟敢断言，我们的救主和主在居于人间之前，并不以他自己独特的存在形式先存，并且他并不拥有自己的神性，唯有父住在他内。\par

2. 许多主教就此问题与他进行查考和讨论；奥力振也应邀与其他主教一同前往，起初他下去与他会谈，以查明他的真实意见。但当奥力振了解了他的观点，并发现它们是错误的之后，便通过论证说服他，并以证明使他信服，将他带回真道，恢复了他从前正统的意见。\par

3. 至今仍存有贝里禄斯的著作以及为他的缘故召开的主教会议的记录，其中包含奥力振向他提出的问题、在他教区进行的讨论，以及当时所做的一切事情。\par

4. 我们中的年长弟兄传下了许多关于奥力振的其他事迹，我认为应略去不提，因与本书无关。但凡是看来有必要记载的关于他的事，均可在我们与圣殉道者庞菲禄为他所写的\WorkTitle{护教篇}中找到。我们精心准备了这篇护教文，并因着挑剔者的缘故共同完成了此项工作。\par

% structure: p0209 source=h2 target=subsection reason=relative-heading-rank; base=h2

\subsection{第三十四章 斐理伯凯撒}

\Person{戈尔迪安}作罗马皇帝六年之后，斐理伯与其子斐理伯继位。据说，他是一位基督徒，曾在最后一次逾越节守夜礼那日，渴望与众人一起参与教会的祈祷，但未被当时的主教准许进入，直到他告解并算作那些被视为犯罪者、处于补赎地位的人中。因为若他不这样做，因他所犯的许多罪行，他将永远不会被主教接纳。据说他欣然服从，在其行为中表现出对天主的真诚而虔诚的敬畏。\par

% structure: p0211 source=h2 target=subsection reason=relative-heading-rank; base=h2

\subsection{第三十五章 狄奥尼修斯继赫拉克拉斯任主教}

在该帝在位的第三年，赫拉克拉斯去世，他任职十六年；狄奥尼修斯接受了亚历山大里亚教会的主教职。\par

% structure: p0213 source=h2 target=subsection reason=relative-heading-rank; base=h2

\subsection{第三十六章 奥力振的其他著作}

1. 此时，随着信仰的扩展，我们的道理在众人面前被大胆宣讲；奥力振，据说已年逾六旬，且因长期实践而获得了极大的便利，很恰当地允许他的公开讲道由速记员记录下来，这是他此前从未允许过的事。\par

2. 他此时也撰写了一部八卷的著作，回应伊壁鸠鲁派凯尔苏斯为反对我们所写的题为\WorkTitle{真言}的著作，以及二十五卷关于\BibleBook{玛窦}{福音}的注释，此外还有关于\BibleBook{}{十二先知书}的注释，我们只找到了二十五卷。\par

3. 尚存有他写给斐理伯皇帝的一封书信，以及另一封致其妻塞维拉的信，还有致不同人士的几封信。我们将所能收集到的、分散在不同人手中的许多书信按不同类别整理成卷，共一百卷，以使它们不再散失。\par

4. 他还写信给罗马主教法比亚诺，以及许多其他教会领袖，论及他的正统性。这些例证见于我们为他所写的\WorkTitle{护教篇}第八卷中。\par

% structure: p0218 source=h2 target=subsection reason=relative-heading-rank; base=h2

\subsection{第三十七章 阿拉伯人的分歧}

大约同时，在阿拉伯有其他人兴起，倡导一种与真理相异的道理。他们说，在现今时期，人的灵魂与身体一同死亡并消灭，但在复活时，它们将一同更新。当时也召开了一次相当规模的主教会议；奥力振再次应召前往，公开就这一问题发言，其效果如此之大，以致先前陷于错误之人的意见得以改变。\par

% structure: p0220 source=h2 target=subsection reason=relative-heading-rank; base=h2

\subsection{第三十八章 厄尔克赛派的异端}

另一谬误也在此时兴起，称作厄尔克赛派的异端，但它从一开始就被扑灭了。奥力振在关于第八十二篇圣咏的公开讲道中如此论及它：\par

\Quote{刚才有一个人前来，极其自负，宣扬那最近在教会中出现的、被称为\InnerQuote{厄尔克赛派}的无神且不敬的意见。我要向你们显示那意见教导何等邪恶之事，使你们不致被它迷惑。它拒绝每部圣经的某些部分。它也使用旧约和福音的部分，却完全拒绝宗徒。它说否认基督是一件无关紧要的事，并且领悟的人必在必要时用口否认，但心里却不。它们拿出某一本书，说是从天上掉下来的。它们认为，凡听见并相信这书的人，必得赦罪，即另一种罪赦，不同于耶稣基督所赐予的。}这就是关于这些人的记述。\par

% structure: p0223 source=h2 target=subsection reason=relative-heading-rank; base=h2

\subsection{第三十九章 德西乌斯时期的迫害及奥力振的苦难}

1. 斐理伯统治七年之后，德西乌斯继位。由于他对斐理伯的仇恨，他开始了对教会的迫害；在此期间，法比亚诺在罗马殉道，科尔乃略继任主教。\par

2. 在巴勒斯坦，耶路撒冷教会主教亚历山大因基督的缘故，再次被带到凯撒勒雅总督的审判台前；他在第二次宣认中表现高尚，被投入监狱，戴上了可敬年岁的白发冠冕。\par

3. 他在总督法庭前作了光荣而杰出的宣认之后，在狱中安眠，\Person{玛匝巴内斯}接替他成为耶路撒冷的主教。\par

4. 安提约基雅的\Person{巴比拉}，也像\Person{亚历山大}一样在宣认之后死于狱中，\Person{法比乌斯}接替他成为该教会的主教。\par

5. 至于\Person{奥力振}在迫害中遭遇了多少、多么重大的事，以及最终结果如何——那邪恶的魔鬼调集所有兵力，以最大的诡计和能力攻击这个人，超过当时他与之争斗的所有人——他为基督的圣言承受了哪些、多少苦难，捆绑、身体折磨、铁枷和地牢中的苦刑；以及他如何许多天被足枷拉开四孔，忍受敌人的火刑威胁和一切加害；他的苦难如何终结，因审判官竭尽全力不愿结束他的生命；以及他在这些事后留下了哪些充满安慰的话，给需要援助的人——他的许多书信都真实而准确地记述了这些。\par

% structure: p0229 source=h2 target=subsection reason=relative-heading-rank; base=h2

\subsection{第四十章 发生在\Person{狄约尼削}身上的事}

1. 我从\Person{狄约尼削}致\Person{革尔玛诺}的书信中，引用一段关于前者的遭遇。他谈及自己时写道：我在天主面前说话，他知道我没有说谎。我并非出于自己的冲动或没有神性指引而逃跑。\par

2. 但在此之前，就在\Person{德西乌斯}迫害令颁布之时，\Person{撒彼诺}派了一名军需官来搜寻我，我在家里等了四天，等待他的到来。\par

3. 但他到处搜查——道路、河流、田野——凡是他以为我可能藏匿或正在途中的地方。但他被盲目所击打，没有找到房子，因为他没料到，在被追捕的情况下，我会待在家里。第四天之后，天主命令我离开，并以奇妙的方式为我开辟了道路；我、我的随从和许多弟兄一同离开。天主眷顾此事，后来发生的事证实了这一点，在这些事中也许我们对某些人有益。\par

4. 接着他叙述了逃跑后发生的事：\par

\Quote{大约日落时分，我和随行的人一起被士兵抓住，带到\Place{塔波西里斯}，但天主的眷顾，\Person{弟茂德}不在场，没有被捕。但他后来到来，发现房子空无一人，有士兵把守，而我们沦为奴隶。}\par

5. 稍后他说：\par

他那奇妙的安排是怎样的呢？因为要讲真话。有一个乡下人遇见\Person{弟茂德}正在逃跑，惊惶失措，问他为何匆忙。\Person{弟茂德}告诉了他实情。\par

6. 那人听了之后（他正要去参加一场婚宴，因为当时习惯在这样的聚会中通宵达旦），他走进去告诉了宴席上的人。他们好像预先约定好似的，同时起身，迅速冲了出来，呐喊着闯入我们所在之处。看守我们的士兵立刻逃跑，他们来到我们躺着的光秃秃的长榻前。\par

7. 但我，天主知道，起初以为他们是来掠夺的强盗。所以我躺在床上一动不动，只穿一件麻布衣，把旁边剩下的衣服给了他们。但他们叫我起来赶快离开。\par

8. 于是我明白了他们的来意，便呼喊起来，恳求他们离开，让我们独自留下。我请求他们，若想给我任何好处，就赶在那些带走我的人之前，亲手砍下我的头。我这样呼喊时，我的同伴和一切参与者都知道，他们用力将我扶起。但我仰面倒在地上；他们抓住我的手和脚，把我拖走。\par

9. 这些事件的目击者跟在后面：\Person{加约}、\Person{法乌斯托}、\Person{伯多禄}和\Person{保禄}。但抓住我的人急忙把我带出村庄，把我放在一头没有鞍的驴上，载我离去。\par

\Person{狄约尼削}叙述了关于他自己这些事。\par

% structure: p0242 source=h2 target=subsection reason=relative-heading-rank; base=h2

\subsection{第四十一章 亚历山大里亚的殉道者}

1. 同一位作者在致安提约基雅主教\Person{法比乌斯}的书信中，叙述了\Person{德西乌斯}时期亚历山大里亚殉道者的遭遇，如下：\par

我们中间的迫害并非始于皇室法令，而是提前了整整一年。对这城市作恶的先知和始作俑者，不论他是谁，事先煽动并激起了异教群众反对我们，重新点燃了他们本国的迷信。\par

2. 他们被他激起，并找到一切作恶的机会，认为杀我们就是对他们魔鬼唯一的敬虔事奉。\par

3. 他们首先抓住一位名叫\Person{默特拉斯}的老人，命令他说亵渎的话。他不肯服从，他们就用棍棒打他，用尖利的棍子撕裂他的脸和眼睛，把他拖出城外，用石头砸死。\par

4. 然后他们把一位名叫\Person{昆塔}的忠心妇女带到他们的偶像庙里，逼她敬拜。她厌恶地转过身去，他们便捆住她的双脚，拖着她在全城的石板路上走，把她撞在磨石上，同时又鞭打她；然后把她带到同一地点，用石头砸死。\par

5. 于是众人同时冲向虔诚者的住所，凡有人知道是邻居的，就拖出来，抢夺他们的财物。他们把值钱的东西据为己有；把较便宜的木制物品散落在街上焚烧，以致城市好像被敌人占领了一样。\par

6. 但弟兄们退避走开，\Quote{欢喜地接受了被抢掠的财物}，就像\Person{保禄}所见证的那些人一样。据我所知，直到此时，除非有人落入他们手中，否则没有人否认主。\par

7. 然后他们也抓住了那位最可敬的童贞女\Person{阿波罗尼亚}，一位老妇人，打她的颚骨，打掉了她所有的牙齿。他们在城外生火，威胁要烧死她，如果她不加入他们亵渎的呼喊。她稍作恳求，得到释放，随即欣然跳入火中被烧死。\par

8. 随后他们在\Person{塞拉彼翁}自己家中抓住了他，用残忍的酷刑折磨他，打断了他所有的肢体，将他从楼上头朝下扔了下去。没有一条街道、一条公共道路、一条小巷对我们开放，无论昼夜；因为随时随地，他们所有人都呼喊说，若有人不重复他们不敬的话语，就立刻被拖走烧死。\par

9. 事情就这样持续了相当一段时间。但一场暴动和内乱降到了这些悲惨的人民身上，使他们对我们的暴行转而彼此相向。因此，当他们停止对我们的狂怒时，我们得以稍事喘息。但不久，从那次较温和的统治转变的消息传到了我们耳中，对所威胁之事的巨大恐惧攫住了我们。\par

10. 因为谕令到达了，几乎像我们的主所预言的那个最可怕的时刻，若可能，连被选者也要受迷惑（\BibleRef{玛 24:24}）。\par

11. 所有人确实都惊慌失措。许多较显赫的人因恐惧立刻上前；其他在公务部门的人被他们的职责所牵引；另一些人被他们的熟人所催促。当他们的名字被呼喊时，他们走近了不洁和不敬的祭献。有些人脸色苍白、颤抖，仿佛他们不是要去献祭，而是自己成为偶像的牺牲和供物；以致周围的人群嘲笑他们，因为每个人都很清楚，他们既害怕死亡，也害怕献祭。\par

12. 但有些人更加轻易地走向祭台，大胆宣称自己从未做过基督徒。我们的主的预言对他们最真实，即他们\Quote{\InnerQuote{难以}}得救（\BibleRef{玛 19:23-24}）。其余人有的模仿这一类，有的模仿另一类；有的逃跑，有的被抓。\par

13. 其中有些人直至被捆绑和监禁仍保持忠信；有些虽已被囚多日，却在受审前背弃了信仰。另有些人一时忍受了巨大的折磨，最终却退缩了。\par

14. 但主坚定而蒙福的柱石，因他而坚强，并获得了适于他们所有坚强信德的力量和勇力，成了他王国的可敬见证人。\par

15. 其中第一位是\Person{尤利安}，一个患严重痛风以致无法站立或行走的人。他们把他和两个抬他的人一起带上来。其中一人立刻背弃了信仰。但另一人，名叫\Person{克罗尼翁}，别号\Person{欧努斯}，以及老\Person{尤利安}本人，两人都承认了主，被放在骆驼上游遍全城（你们知道这是一座很大的城），在高处被鞭打，最后在熊熊大火中被烧死，周围站满了民众。\par

16. 但一个名叫\Person{贝萨斯}的士兵，在他们被带走时站在旁边，斥责那些侮辱他们的人。民众对他大喊，于是这位最勇敢的天主战士被审讯，在敬虔的伟大争战中英勇作战，被斩首。\par

17. 另有一人，利比亚人，名字和真福上真是\Person{玛卡尔}，法官极力劝他反悔，但他不肯屈服，被活活烧死。在他们之后，\Person{埃庇玛库斯}和\Person{亚历山大}，被囚已久，忍受了无数的刮刑和鞭刑，也被猛火焚烧。\par

18. 与他们一起还有四位妇女。\Person{阿摩纳里乌姆}，一位圣洁的贞女，法官对她无情地过度施刑，因为她从一开始就宣告她不会说出任何他所命令的话；她真诚地持守了诺言，被拖走了。其他的是\Person{梅库里亚}，一位非常杰出的老妇人，以及\Person{狄奥尼西亚}，许多孩子的母亲，她爱主胜过爱自己的儿女（\BibleRef{路 14:26}）。由于总督因施刑无效而感到羞愧，且总是被妇女打败，她们未经刑讯就被刀剑处死。因为勇士\Person{阿摩纳里乌姆}替大家承受了这些。\par

19. 埃及人\Person{赫伦}、\Person{阿特尔}和\Person{伊西多鲁斯}，以及与他们一起的\Person{狄奥斯库鲁斯}，一个约十五岁的男孩，被交了出来。起初法官试图用甜言蜜语欺骗这个少年，以为他能轻易被说服，然后施以酷刑，以为他会轻易屈服。但\Person{狄奥斯库鲁斯}既未被说服，也未被强迫。\par

20. 由于其他人坚定不移，他残酷地鞭打他们，然后将他们投入火中。但钦佩\Person{狄奥斯库鲁斯}公开表现出的杰出态度以及他对规劝的智慧回答，他释放了他，说因他年幼，给他时间悔改。这位最敬虔的\Person{狄奥斯库鲁斯}如今在我们中间，等待更长的争战和更严峻的竞赛。\par

21. 但一个名叫\Person{奈梅西翁}的人，也是埃及人，被指控为强盗同伙；但当他在百夫长面前澄清了这个与真理毫不相干的指控后，他被举报为基督徒，被捆绑带到总督面前。这位最不义的法官对他施加的酷刑和鞭打是加给强盗的两倍，然后将他夹在两个强盗之间烧死，以此藉基督的相似性尊荣这位蒙福者。\par

22. 一队士兵，\Person{阿蒙}、\Person{泽农}、\Person{托勒密}和\Person{英根尼斯}，以及一位老人\Person{特奥斐卢斯}，紧挨着站在法庭前。当一个正被作为基督徒受审的人似乎想要否认时，他们站在一旁咬牙切齿，用面容示意，伸出手来，用身体做手势。当所有人的注意力都转向他们时，在别人抓住他们之前，他们冲上法庭，说他们是基督徒，以致总督及其议会都惊慌失措。那些受审的人在面对苦难时显得极其勇敢，而审判者却战栗。他们欢欣地从法庭走出，因他们的见证而喜乐；是天主亲自使他们荣耀地得胜。\par

% structure: p0266 source=h2 target=subsection reason=relative-heading-rank; base=h2

\subsection{第四十二章 \Person{狄奥尼修斯}所记述的其他人}

1. 许多其他人在城市和乡村中被异教徒撕裂，我将举一个例子来说明。\Person{伊斯基里昂}被一位长官雇佣为管家。他的主人命令他祭献，他拒绝后，主人侮辱他，他坚持不动摇，便虐待他。当他仍然坚持时，主人拿起一根长棍，刺穿了他的内脏，杀死了他。\par

2. 我何必提及那许多在沙漠和山间流浪，因饥饿、干渴、寒冷、疾病、强盗和野兽而丧生的人呢？他们中幸存的人是他们蒙拣选和得胜的见证。\par

3. 但我要将一件事作为例子叙述。\Person{凯勒蒙}年纪很大，是\Place{尼罗}城的主教。他带着妻子逃到\Place{阿拉伯山}，再也没有回来。虽然弟兄们仔细搜寻，却找不到他们或他们的尸体。\par

4. 许多逃到同一\Place{阿拉伯山}的人被野蛮的\Person{撒拉逊人}掳为奴隶。其中一些人费了很大周折、花了大价钱才被赎回；另一些人直到如今仍未赎回。我述说这些事，我的兄弟，并非无的放矢，而是为叫你知道我们遭遇了多少、多大的苦难。那些亲身经历最多的人，最了解这些苦难。\par

5. 稍后他补充说：\Quote{我们中间这些神圣的殉道者，现今与基督同坐，共享祂的王国，有分于祂的审判，并与祂一同审判。他们接纳了一些曾跌倒、背负祭献之罪的弟兄。当他们看到这些人的归化和悔改，足以被那绝不渴望罪人死亡、而是渴望其悔改的天主所接纳时，便考验了他们，然后重新接纳他们，召集他们，与他们相聚，在祈祷和筵席中与他们共融。}\par

6. 那么，弟兄们，关于这些人你们给我们什么建议？我们该做什么？我们是否应有与他们相同的判断和规矩，遵守他们的决定和爱德，向他们所怜悯的人显示慈悲？或者，我们应宣称他们的决定是不义的，将自己设为对他们意见的审判者，使慈悲忧伤，颠覆秩序？这些是\Person{迪奥尼修斯}在提及那些在迫害时期软弱的人时非常恰当地补充的话。\par

% structure: p0273 source=h2 target=subsection reason=relative-heading-rank; base=h2

\subsection{第四十三章 \Person{诺瓦图斯}，他的生活方式与他的异端}

1. 此后，罗马教会的一位司铎\Person{诺瓦图斯}，因傲慢而自高，反对这些人，似乎他们再无救恩的希望，即使他们行一切与真诚纯洁的归化有关的事，他成了那些因想象而骄傲、自称为\OriginalTerm{洁净派}{Cathari}之异端的领袖。\par

2. 于是，在罗马召集了一个非常大的主教会议，有六十位主教，以及更多的司铎和执事；其余各省的牧者则在他们当地私下商议应对之策。所有主教一致确认一项法令：\Person{诺瓦图斯}和那些与他联合的人，以及那些采纳他憎恨弟兄、不人道意见的人，应被教会视为外人；但他们应医治那些遭难跌倒的弟兄，并以悔改的良药服侍他们。\par

3. 我们有\Person{哥尔乃略}（罗马主教）写给\Place{安提约基雅}教会\Person{法比乌斯}的书信，显示了罗马主教会议的决定，以及在意大利、非洲和邻近地区的所有人都认为最好的做法。还有其他书信，用拉丁文写成，是\Person{西彼连}和他在非洲的同工所写，表明他们一致同意必须援助那些曾受诱惑的人，并将异端的领袖和所有与他联合的人从\OriginalTerm{天主教会}{Catholic Church}中剪除。\par

4. \Person{哥尔乃略}另一封关于主教会议决议的书信附于这些之后；还有其他关于\Person{诺瓦图斯}行为的书信，我们从中选取合适的内容，以便任何看到这部作品的人可以了解他。\par

5. \Person{哥尔乃略}以下列话语告知\Person{法比乌斯}\Person{诺瓦图斯}是怎样的人：\par

\Quote{但为叫你知道，很久以前这位杰出的人物就渴望主教职，但他将这野心隐藏起来——以那些从一开始就跟随他的宣信者为掩护，掩盖他的反叛——我愿意述说。}\par

\Quote{6. \Person{玛克西穆}，我们的一位司铎，和\Person{乌尔巴诺}，曾两次因宣认获得最高荣誉，与\Person{西多尼乌斯}，以及\Person{塞莱里努斯}，一个藉着天主的恩宠最英勇地忍受了各种酷刑，并因信德的力量克服了肉体的软弱，大力战胜了仇敌的人——这些人发现了他，察觉了他的诡诈和虚伪，他的伪证和谎言，他的孤僻和残忍的友谊。他们回到了圣教会，在众多主教、司铎和平信徒面前，揭露了他长期隐藏的一切诡诈和邪恶。他们这样做时带着哀叹和悔改，因为他们曾被那诡诈邪恶的野兽的说服而暂时离开了教会。稍后他说：}\par

\Quote{7. 亲爱的弟兄，多么显著的变化和转变，我们在短时间内在他身上看到。因这位最杰出的人，曾以可怕的誓言约束自己绝不寻求主教职，却突然出现成为主教，仿佛被某种机器抛到我们中间。}\par

\Quote{8. 因这位教条主义者，这位教会教义的捍卫者，试图攫取并非从上赐予的主教职，选择了两位放弃自己救恩的同伙。他将他们派到意大利一个微不足道的小角落，在那里用某种虚假的论证欺骗三位粗鲁且非常单纯的主教。他们肯定且强烈地断言，必须尽快来到罗马，以便藉着他们与其他主教共同调解，平息那里发生的一切纷争。}\par

9. 当他们来到时，正如我们所说，他们在邪恶的诡计和伎俩上非常天真，他们与一些像他自己一样被选定的人一起被关起来。到了第十时辰，当他们喝醉并生病时，他强迫他们通过一种虚假而又徒劳的覆手礼，将主教职赋予他。因为主教职没有临到他，他便用诡计和背叛报复。\par

10. 这些主教中的一位不久后回到教会，哀伤并承认他的过犯。我们与他共融，如同与一位平信徒一样，在场的全体民众都为他求情。我们祝圣了其他主教的继承人，并派他们到他们所在的地方。\par

11. 这位福音的报复者那时不知道在一个公教会中应有一位主教；然而他并非不知道（因他如何能不知道呢？）在教会中有四十六位司铎、七位执事、七位副执事、四十二位辅祭、五十二位驱魔员、读经员和守门员，以及超过一千五百位寡妇和困苦的人，所有这些都是由主的恩宠和慈爱所滋养的。\par

12. 但即便是这如此在教会中不可或缺的庞大群众，以及那些因天主的眷顾而富足满足的人，连同许许多多甚至数不胜数的人，都不能使他从这般的绝望和狂妄中回转，召唤他回到教会。\par

13. 再次，他在后面又加上了这些话：请允许我们进一步说：他究竟凭着什么作为或品行，竟有把握争夺主教职？难道是因为他从起初就在教会中长大，为她忍受了许多冲突，并为宗教经历了许多巨大的危险？事实并非如此。\par

14. 但是撒殚，进入并在他内住了很久，成了他相信的缘由。他被驱魔员释放后，患了重病；当他似乎快要死时，在他躺卧的床上接受了注水圣洗；如果我们真能说这样一个接受了洗礼的话。\par

15. 当他疾病痊愈后，他没有接受按照教会法典所必须有的其他事物，即由主教封印。既然他没有接受这个，他怎能接受圣神呢？\par

16. 不久后他又说道：\par

\Quote{在迫害时期，由于怯懦和贪生，他否认自己是司铎。因为当执事们请求并恳求他走出他自我囚禁的房间，并按司铎合法和可能的方式援助那些身处危险并需帮助的弟兄时，他如此不尊重执事们的恳求，竟生气地走开并离开了。因为他说他不再愿意做司铎，因为他仰慕另一种哲学。}\par

17. 略过几件事后，他又加上了以下的话：\par

\Quote{因为这位显赫的人物离弃了天主的教会，在他相信时，他曾因祝圣他担任司铎职务的主教的恩宠，被判断为堪当司铎职。这件事曾遭到全体圣职人员和许多平信徒的反对；因为像他那样因疾病在床上被注水圣洗的人，进入任何圣职都是不合法的；但那位主教请求，只允许他祝圣这个人。}\par

18. 在这些之外，他又加上了这人的另一件，也是最恶劣的过犯，如下：\par

因为当他献祭后，把一部分分给每个人，在递送时，他强迫这个可怜的人以发誓代替祝福。他双手握住对方的手，直到对方以这种方式发誓（因为我要引述他自己的话）才放手：\par

\Quote{你向我指着我主耶稣基督的体血发誓，你永远不会离弃我而转向科尔乃略。}\par

19. 这个不幸的人直到对自己发出诅咒后才品尝；当他拿起饼时，他不说\Quote{亚孟}，却说\Quote{我绝不会回到科尔乃略那里}。再往后他又说：\par

20. \Quote{但你们要知道，他现在已被剥去一切，孤寂无援；因为弟兄们每天都离开他，回到教会。还有真福殉道者梅瑟，他最近在我们中间承受了光荣而可钦佩的殉道，当他还在世时，看到他的胆大和愚妄，拒绝与他以及与他一同脱离教会的五位司铎共融。}\par

21. 在他书信的结尾，他列出了来到罗马并谴责诺瓦督愚行的主教名单，包括他们的名字以及各自所主持的堂区。\par

22. 他也提到了那些没有来到罗马，但通过书信表达同意这些主教投票的人，列出了他们的名字以及他们各自寄信的城市。科尔乃略将这些事写给了安提约基雅主教法比乌。\par

% structure: p0301 source=h2 target=subsection reason=relative-heading-rank; base=h2

\subsection{第四十四章 狄约尼削关于色拉平的记述。}

1. 亚历山大的狄约尼削也曾写了一封书信给这位似乎有些偏向此分裂的法比乌。他在信中写了许多其他关于悔改的事，并叙述了最近在亚历山大殉道者的争战。在另一段叙述之后，他提到了一件奇妙的实事，值得在本著作中占一席之地。内容如下：\par

2. 我要给你一个发生在我们中间的例子。我们中间有一位名叫色拉平的年长信徒，他长期以来生活无可指摘，却在试炼中跌倒了。他多次恳求，但没有人理会他，因为他曾祭献过。但他生病了，连续三天不能说话，没有知觉。\par

3. 第四天稍有好转后，他叫来女儿的儿子，说：\Quote{我的孩子，你要拘留我多久？我恳求你，赶快，快快赦免我。给我叫一位司铎来。}说了这话后，他又不能说话了。孩子跑去找司铎。但那时是夜里，司铎病了，所以不能来。\par

4. 但是，因为我曾下令，临终的人如果请求，尤其若是先前已曾要求，应得赦免，使他们能怀着好希望离去，于是他给了这孩子一小块圣体，吩咐他浸湿后让滴液滴入老人的口中。\par

5. 孩童带着它回来，当他走近、尚未进门时，\Person{塞拉皮雍}再次苏醒，说：\Quote{你来了，我的孩子，司铎不能来；但快照他吩咐的做，好让我离去。}于是孩童将圣体浸湿，放入他口中。他咽下少许，便立刻断了气。\par

6. 这岂不明显表明，他得以保全、生命延续，直到他被赦免，罪过被涂抹，并因他所行的许多善事而蒙认可？\par

\Person{狄约尼削}记述了这些事情。\par

% structure: p0309 source=h2 target=subsection reason=relative-heading-rank; base=h2

\subsection{第四十五章 狄约尼削致诺瓦督斯书函}

1. 但我们来看同一位是如何致函正在扰乱罗马弟兄团体的\Person{诺瓦督斯}。由于他声称某些弟兄是他背教和分裂的起因，仿佛他是在被迫之下才如此行事，请看他写信给他时所用的措辞：\par

2. \Quote{\Person{狄约尼削}致其弟兄\Person{诺瓦督斯}，问候。如果你如所说是被勉强引导的，你若自愿退去，便能证明。因为忍受一切比分裂天主的教会更好。即使为阻止分裂而殉道，其光荣也不亚于拒绝敬拜偶像。不，在我看来更大。因为前者是为自己的灵魂殉道，后者是为整个教会。现在，如果你能说服或促使弟兄们达致同心，你的义德将大于你的过犯；这过犯将不被记念，而那义德将受赞扬。但若你无法说服那些不顺从的人，至少救你自己的灵魂。我祈求你顺利，在主内持守和平。}这是他写给\Person{诺瓦督斯}的。\par

% structure: p0312 source=h2 target=subsection reason=relative-heading-rank; base=h2

\subsection{第四十六章 狄约尼削的其他书函}

1. 他也写了一封关于补赎的书函给在埃及的弟兄们。在其中，他阐述了他认为对于那些失足者所应行的事，并描述了罪的类别。\par

2. 另有一封关于补赎的私人书信，是他写给\Place{赫尔莫波利斯}教区的主教\Person{科农}的，还有一封劝诫性的信，是写给他\Place{亚历山大里亚}的羊群的。其中有写给\Person{奥力振}关于殉道的，以及写给\Place{劳狄刻雅}的弟兄们，\Person{太利米德勒斯}是当地的主教。他也寄了一封关于补赎的信给\Place{亚美尼亚}的弟兄们，\Person{梅洛赞尼斯}是当地的主教。\par

3. 除所有这些之外，当他从\Person{科内略}（罗马的）收到一封反对\Person{诺瓦督斯}的书函时，他也写信给\Person{科内略}。他在信中说明，他曾受\Place{基里基雅}的\Place{塔尔索}主教\Person{赫勒努斯}以及与他同行的各位，即\Place{卡帕多细雅}的主教\Person{非尔弥利亚诺斯}和\Place{巴勒斯坦}的\Person{特奥克提斯托斯}的邀请，前往\Place{安提约基雅}参加教务会议，当时有些人在那里企图建立\Person{诺瓦督斯}的分裂。\par

4. 除此之外，他写道，他得知\Person{法比乌斯}已安眠，\Person{德默特利亚努斯}被任命为\Place{安提约基雅}主教的继任者。他也以这些话论及\Place{耶路撒冷}的主教：\Quote{因为真福\Person{亚历山大}被囚在狱中，安详离世。}\par

5. 除此之外，另有\Person{狄约尼削}的一封执事性的书函存留，是通过\Person{依玻里托斯}寄给罗马的。他还写了另一封关于和平及同样关于补赎的信给他们；又写了一封给当地仍坚持\Person{诺瓦督斯}意见的宣信者。在这些人回归教会之后，他又寄了两封信给同一批人。他还通过书信与许多人交流，这些书信作为多方面的益物留给了如今勤学他著作的人。\par

\SourceRecord{Translated by Arthur Cushman McGiffert.；Nicene and Post-Nicene Fathers, Second Series；Vol. 1.；Edited by Philip Schaff and Henry Wace.；Buffalo, NY: Christian Literature Publishing Co.,；1890.；<http://www.newadvent.org/fathers/250106.htm>.}

% structure: p0001 source=h1 target=section reason=page-title

\section{\WorkTitle{教会史（第七卷）}}

% structure: p0002 source=h2 target=subsection reason=relative-heading-rank; base=h2

\subsection{引言}

在这部\WorkTitle{教会史}的第七卷中，亚历山大的伟大主教\Person{狄约尼削}将再次以他自己的话协助我们，叙述他那个时代的若干事件，这些叙述保存在他留下的书信中。我将从这些开始。\par

% structure: p0004 source=h2 target=subsection reason=relative-heading-rank; base=h2

\subsection{第一章 \Person{德西乌斯}与\Person{加卢斯}的邪恶}

当\Person{德西乌斯}统治不到两年，便连同他的子女被杀，\Person{加卢斯}继位。此时\Person{奥力振}去世，享年六十九岁。\Person{狄约尼削}写信给\Person{赫尔马蒙}，论及\Person{加卢斯}如下：\par

\Quote{加卢斯既不认识德西乌斯的邪恶，也不考虑是什么毁灭了他，反而撞在同一块石头上，尽管那石头就在他眼前。因为当他的统治顺利、事情按他的心意进行时，他攻击那些为他的和平与福祉向天主代祷的圣人们。因此，他与他们一起迫害了他们为他所作的祈祷。} 关于他，就说这些。\par

% structure: p0007 source=h2 target=subsection reason=relative-heading-rank; base=h2

\subsection{第二章 当时的罗马主教}

\Person{科尔乃略}在罗马城担任主教约三年，\Person{路求}继位。他在不到八个月后去世，将职务传给\Person{斯德望}。\Person{狄约尼削}就圣洗问题给他写了第一封信，因为当时出现了一场不小的争论：是否应当以圣洗来洁净那些从任何异端归正的人。因为关于这类人，古代的惯例是只给他们行覆手礼并祈祷。\par

% structure: p0009 source=h2 target=subsection reason=relative-heading-rank; base=h2

\subsection{第三章 \Person{西彼廉}与他同道的诸位主教首先主张，必须藉圣洗洁净那些从异端归正的人。}

首先是迦太基教区的牧者\Person{西彼廉}主张，除非他们藉圣洗从错误中得以洁净，否则不应接纳他们。但\Person{斯德望}认为，没有必要违背从起初就保持的传统而增加任何新事物，对此他非常愤慨。\par

% structure: p0011 source=h2 target=subsection reason=relative-heading-rank; base=h2

\subsection{第四章 \Person{狄约尼削}就此事所写的书信}

因此，\Person{狄约尼削}通过书信就此问题与他进行了广泛的交流，最终向他表明，既然迫害已平息，各地的教会都已拒绝了\Person{诺瓦图斯}的革新，彼此之间和平相处。他写道：\par

% structure: p0013 source=h2 target=subsection reason=relative-heading-rank; base=h2

\subsection{第五章 迫害之后的和平}

1. 但是如今，我的弟兄们，要知道整个东方以及更远地方的所有教会，从前是分裂的，如今已联合为一。各处的主教都同心合意，为这出乎意料的来临的和平而大大欢喜。例如\Place{安提约基雅}的\Person{德默特利亚努斯}，\Place{凯撒勒雅}的\Person{德奥克提斯图斯}，\Place{埃里亚}的\Person{玛匝巴内斯}，\Place{提洛}的\Person{玛里努斯}（\Person{亚历山大}已经安息），\Place{劳狄刻雅}的\Person{赫略多鲁斯}（\Person{特利米德勒斯}已死），\Place{塔尔索}的\Person{海肋努斯}，以及整个\Place{基里基雅}的教会，\Person{菲尔米利亚努斯}，以及整个\Place{卡帕多细雅}。我只列举了更杰出的主教，以免使我的书信冗长，言语累赘。\par

2. 还有整个\Place{叙利亚}，以及你们在需要时予以援助、你们刚刚写信去的\Place{阿拉伯}，\Place{美索不达米亚}，\Place{本都}，\Place{比提尼亚}，简言之，各处都在为这同心合意与兄弟之爱而欢欣，并光荣天主。\Person{狄约尼削}的话到此为止。\par

3. 但\Person{斯德望}任职两年后，\Person{息斯笃}继位。\Person{狄约尼削}又给他写了一封关于圣洗的第二封信，在信中他同时向他表明\Person{斯德望}与其他主教的意见和判断，并论及\Person{斯德望}如此说：\par

4. 因此他先前曾写信论及\Person{海肋努斯}和\Person{菲尔米利亚努斯}，以及\Place{基里基雅}、\Place{卡帕多细雅}、\Place{迦拉达}和邻近地区的所有人，说他不会因这同一原因与他们共融；即他们为异端者重新施洗。但请思量此事的重要性。\par

5. 因为我得知，确实在最大的主教会议上，已就此事通过了法令：从异端归正的人应先接受教导，然后从旧而不洁的酵母的污秽中被洗涤和洁净。而我写信恳求他有关这一切。接着他说：\par

6. \Quote{我也曾写信给我们亲爱的同司铎\Person{狄约尼削}和\Person{斐里蒙}，起先寥寥数语，最近则洋洋洒洒，他们先前曾与\Person{斯德望}抱同样见解，并就同样的事写信给我。} 以上是关于上述争论的话。\par

% structure: p0020 source=h2 target=subsection reason=relative-heading-rank; base=h2

\subsection{第六章 \Person{撒伯里乌}的异端}

在同一封信中，他也提及\Person{撒伯里乌}的异端教导，这些教导在他那个时代日益显著，并说道：\par

\Quote{因为关于目前在\Place{五城区}的\Place{托勒迈斯}所鼓噪的教义——这教义是不虔诚的，带有对全能天主圣父、我们的主耶稣基督的大亵渎，并且包含许多关于他的独生子、万物首生者、成为人的圣言的不信，以及对圣神缺乏认识——由于我收到双方及讨论此事的弟兄们的来信，我便写了若干信函，尽天主助我之能，以教导的方式处理这主题。现将这些信函副本寄给你们。}\par

% structure: p0023 source=h2 target=subsection reason=relative-heading-rank; base=h2

\subsection{第七章 异端者的可憎谬误；\Person{狄约尼削}的神视；以及他所接受的教会规条}

1. 在同一个\Person{狄约尼削}写给罗马司铎\Person{斐里蒙}的第三封论圣洗的书信中，他叙述如下：我审查了异端者的著作和传统，使我的心灵暂时被他们可憎的意见所玷污，但从中得到了这样的益处：我亲自驳斥了他们，并且更加憎恶他们。\par

2. 当司铎中的一位弟兄阻止我，担心我会被他们邪恶的污秽所吞噬（因为这将会玷污我的灵魂）——在我看来，他说的也是实话——这时，一个来自天主的异象临到我，并坚固了我。\par

3. 那临到我的话明确地命令我说：\Quote{尽你所能拿到的，阅读一切，因为你能纠正并验证一切；这从起初就是你信德的缘由。} 我接受这异象，视之为与宗徒的话相符，那话对较强的人说：\Quote{要作精明的兑换银钱者。}\par

4. 然后，在论及一切异端之后，他补充道：\Quote{我从我们可敬的教父赫拉克拉斯领受了这条规则和条例。因为那些从异端归正的人，尽管他们曾背弃教会——或者更确切地说，并未背弃，却似乎与异端交往，并被指控投靠某个错误的教师——当他将他们逐出教会后，虽经他们恳求，也不重新接纳，直到他们公开报告了所有从敌对者那里听到的事；但那时他接纳他们，却不要求另一次圣洗。因为他们先前已从他那里领受了圣神。}\par

5. 再次，在深入讨论这个问题之后，他补充道：\Quote{我也得知这并非仅在非洲新近实行的做法，而是甚至在很久以前，在我们之前的诸位主教时代，这一意见已在最兴旺的教会中、在伊科尼翁和辛纳达的弟兄们的会议中，以及许多其他会议中被采纳。我无法容忍推翻他们的决议，使他们陷入纷争和争论。因为经上说：\InnerQuote{你不可挪移你邻舍的地界，那是你的祖先所设立的。}}\BibleRef{申 19:14}\par

6. 他的第四封论圣洗的书信是写给罗马的狄约尼削的，当时他还是一位长老，但不久后便接受了该教会的主教职。从亚历山大里亚的狄约尼削对他的叙述可以清楚看出，他也是一位博学且可敬的人。他在信中关于诺瓦特写道如下：\par

% structure: p0030 source=h2 target=subsection reason=relative-heading-rank; base=h2

\subsection{第八章 诺瓦特的异端}

\Quote{因为我们完全有理由憎恨诺瓦西安，他分裂了教会，使一些弟兄陷入不虔敬和亵渎，并引入关于天主的邪恶教导，诽谤我们最慈悲的主耶稣基督为无情。除此之外，他还拒绝神圣的圣洗，推翻圣洗前的信仰和宣认，并彻底从他们中驱逐圣神，如果真有任何希望祂会留在或回到他们中的话。}\par

% structure: p0032 source=h2 target=subsection reason=relative-heading-rank; base=h2

\subsection{第九章 异端者的不虔敬圣洗}

1. 他的第五封书信是写给罗马主教西斯都的。在信中，他针对异端者说了很多话之后，讲述了他那个时代发生的一件事，如下：因为，弟兄，我确实需要指导，我请你判断我所遇到的一件事，恐怕我会犯错。\par

2. 因为聚会的一位弟兄，长久以来被认为是信徒，并且在我被祝圣之前——我想也在可敬的赫拉克拉斯被任命之前——他已是会众的一员，当时他与那些刚受洗的人在一起。当他听到问与答时，他来到我面前，哭泣哀叹，俯伏在我脚前，承认并郑重声称他在异端者中所受的圣洗并非如此，也毫不类似，因为它充满了不虔敬和亵渎。\par

3. 他说他的灵魂如今被痛苦刺透，他不敢举目仰望天主，因为他始于那些不虔敬的言语和行为。为此，他恳求能领受这最完美的净化、接纳和恩宠。\par

4. 但我未敢这样做；并说他长久的共融便足以做到这一点。因为我不敢重新开始一位曾听过感恩诵、并一同应答\Quote{阿们}的人；他曾站在祭台旁，伸双手领受祝福的食物；他曾领受，并长久分享我们主耶稣基督的体和血。但我劝慰他鼓起勇气，以坚定的信德和良好的望德去领受圣体。\par

5. 但他不停地哀叹，颤抖着不敢靠近祭台，即使被恳求，也几乎不敢出席祈祷。\par

6. 除此之外，还有同一人物另一封关于圣洗的书信留存，由他和他的堂区写给西斯都和罗马教会。在这封信中，他以详尽的论证思考了当时所争论的问题。在此之后，还有一封写给罗马的狄约尼削，关于路济安。关于这些事，就说到这里。\par

% structure: p0039 source=h2 target=subsection reason=relative-heading-rank; base=h2

\subsection{第十章 瓦莱里安及其治下的迫害}

1. 加卢斯与其他统治者执政不足两年便被推翻，瓦莱里安与其子加里恩努斯接掌帝国。狄约尼削在致赫玛孟的书信中叙述了关于他的情形，我们可从中了解如下：\par

2. 同样，这向若望启示了；因为他说：\Quote{又赐给它说夸大和亵渎话的口；又赐给它权柄，可以任意而行四十二个月。}\BibleRef{默 13:5}\par

3. 奇妙的是，这两件事都发生在瓦莱里安统治时期；而考虑他先前的行为，此事尤为显著，因为他曾对天主的子民温和友善，在他之前的皇帝中没有一人像他那样仁慈友好地对待他们；甚至那些公开自称基督徒的人，在他统治之初也没有像他那样以明显的盛情和友善接待他们。因为他的全家充满了虔诚的人，且是天主的一个教会。\par

4. 但来自埃及的术士会堂的导师和统治者说服他改变了做法，催促他杀戮和迫害纯洁圣洁的人，因为他们反对并阻碍那些邪恶可憎的咒语。因为如今和过去都有一些人在场，被看见，仅凭呼吸和言语，就能驱散邪恶魔鬼的谋划。他还引诱他实行入教仪式和可憎的巫术，献上不被悦纳的祭品；杀害无数儿童，牺牲不幸父亲的后代；剖开新生婴儿的内脏，凌辱切割天主的受造物，仿佛通过这些行为他们可以获得幸福。\par

5. 他接着补充道：马克里安为他们带来的帝国感恩祭品确实辉煌，这是他希望的目标。据说他从前是皇帝的财务总管；但他没有做出任何值得称赞或对大众有益的事，反而应验了先知的这句话：\par

6. \BibleQuote{那些凭自己的心说预言，却不考虑公众福祉的人，有祸了。} \BibleRef{则 13:3} 因为他没有察觉普遍的天主眷顾，也没有寻求那在万有之前、贯穿万有、统御万有者的审判。因此他成了他的天主教会（\OriginalTerm{天主教会}{Catholic Church}）的仇敌，疏远并隔绝了天主的慈悲，尽己所能地逃离自己的救恩。由此，他显出了自己名字的真实含义。\par

7. 再往下他又说：瓦莱里安受此人煽动而作出这些行为，便被交付于侮辱和责骂，正如依撒意亚所说的：\BibleQuote{他们拣选了自己的道路，他们的灵魂喜悦可憎之事；我也要拣选他们的迷惑，将他们的罪报应给他们。} \BibleRef{依 66:3}\par

8. 但此人虽不配，却疯狂渴求王国，因无法将王袍穿在自己残疾的身体上，便推出两个儿子，让他们承担父亲的罪。关于他们，天主所说的宣告是明确的：\BibleQuote{恨我的，我必追讨他的罪，自父及子，直到三四代。} \BibleRef{出 20:5}\par

9. 因他将自己已得逞的邪恶欲望堆积在儿子们头上，便将自己对天主的邪恶和仇恨涂抹在他们身上。\par

狄约尼削就瓦莱里安述说了这些事。\par

% structure: p0050 source=h2 target=subsection reason=relative-heading-rank; base=h2

\subsection{第十一章  此时发生在狄约尼削及埃及众人身上的事件。}

1. 至于在他统治期间如此猛烈肆虐的迫害，以及狄约尼削与他人因对宇宙天主的虔诚而忍受的苦难，他写给当时试图诽谤他的主教盖尔马努的回信中的原话将予以表明。他的陈述如下：\par

2. 我确实有陷入极大愚昧和愚蠢的危险，因被迫叙述天主对我们奇妙的眷顾。但既然经上说：\BibleQuote{隐藏君王的秘密是好事，但揭示天主的作为是光荣的。} \BibleRef{多 12:7} 我将与盖尔马努的暴力相争。\par

3. 我不是独自去见埃米利亚努斯；而是我的同僚司铎马克西穆斯，以及执事福斯图斯、欧瑟比和凯雷蒙，还有一位从罗马来的弟兄与我同去。\par

4. 但埃米利亚努斯起初并没有对我说：\Quote{不得举行集会；} 因为这对他是多余的，对于寻求完成首要之事的人来说，这是最后的事。他关心的不是我们集会，而是我们本人不应是基督徒。他命令我放弃这个；以为如果我背离了，其他人也会跟随我。\par

5. 我既未不恰当地也未用许多话回答他：\BibleQuote{我们应当服从天主，而不是人。} \BibleRef{宗 5:29} 我公开作证，我只崇拜独一的天主，别无他神；我不会背离这一点，也绝不会停止做基督徒。于是他命令我们前往沙漠附近一个叫克弗罗的村庄。\par

6. 但请听双方所说的话，正如所记载的：狄约尼削、福斯图斯、马克西穆斯、玛尔策禄和凯雷蒙被提审时，总督埃米利亚努斯说：\par

7. \Quote{我曾口头与你们论及我们的统治者对你们显示的宽仁；因为他们给了你们自救的机会，只要你们转向合乎本性的事物，崇拜那些保护他们帝国的神祇，并忘记那些违反本性的事物。对此你们有何话说？因为我不认为你们会对他们的仁慈忘恩负义，因为他们是要引导你们走向更好的道路。}\par

8. 狄约尼削回答说：\Quote{并非所有人都崇拜所有的神祇；而是各人崇拜自己所认可的神。因此，我们敬畏并崇拜独一的天主，万有的创造者；祂将帝国赐予蒙受神恩、尊贵的瓦莱里安和加里恩努斯；我们不断为他们的帝国向祂祈祷，愿它屹立不摇。}\par

9. 总督埃米利亚努斯对他们说：\Quote{但如果他是神，谁禁止你们与那些本性是神的神祇一同崇拜他呢？因为你们奉命要敬畏神祇，就是所有人都认识的那些神。} 狄约尼削回答说：\par

10. \Quote{我们不崇拜别的神。} 总督埃米利亚努斯对他们说：\Quote{我看你们既忘恩负义，又对君主的恩惠无动于衷。因此你们不得留在这座城。你们将被遣送到利比亚地区一个叫克弗罗的地方。因为我在君主的命令下选择了此地，绝不允许你们或任何其他人举行集会，或进入所谓的墓地。}\par

11. \Quote{但如果有人在我命令的地方之外被看见，或在任何集会中被发现，他将自招危险。因为适当的惩罚不会落空。所以，去吧，到你们被命令去的地方。}\par

\Quote{他催促我离开，尽管我生病，连一天的喘息都不给。那我有什么机会举行集会，或不举行呢？}\par

12. 再往下他说：但靠着主的帮助，我们没有放弃公开的集会。我更努力地召集城中的那些人，如同我与他们同在；可谓 \BibleQuote{身体不在，精神却在。} \BibleRef{格前 5:3} 但在克弗罗，一支由从城中跟随我们的弟兄和从埃及加入我们的人组成的大教会与我们聚集；在那里 \BibleQuote{天主为我们敞开了圣言的门。} \BibleRef{哥 4:3}\par

13. 起初我们受到迫害和石击；但此后不少外邦人离弃偶像，归向天主。因为直到那时，他们尚未听到圣言，因那是我们首次播种。\par

14. 仿佛天主为此目的将我们带到他们那里，当我们完成了这项职务，祂就将我们转移到另一地方。因为埃米利亚努斯似乎想要将我们运到更荒凉、更似利比亚的地方；于是他命令他们从四面八方聚集到玛勒奥提斯，并给他们分配了该国不同的村庄。但他命令将我们安置在更靠近大路的地方，以便我们首先被逮捕。因为他显然安排并准备好了，以便一旦他想逮捕我们，就能毫不费力地将我们全部抓获。\par

15. 当我最初奉命前往\Place{克弗罗}时，我不知道那地方在哪里，几乎从未听说过这个名字；但我仍欣然前往。然而，当我被告知要迁往\Place{科卢提翁}地区时，在场的人都知道我当时是如何感动的。\par

16. 因为在此我要责备自己。起初我悲伤并极为不安；虽然这些地方更熟悉、更靠近我们，但据说该地缺少弟兄和有品性的人，并且容易受到旅行者的骚扰和强盗的侵袭。\par

17. 但弟兄们提醒我，那里更靠近城市，这安慰了我；而\Place{克弗罗}虽然使我们能与来自埃及的弟兄多有交往，从而能更广泛地扩展教会，但这个地方更靠近城市，我们将更频繁地见到那些真正被爱、关系最密切、最亲爱的人。因为他们会来并留下，并且可以像在更偏远的郊区那样举行特别聚会。结果也确实如此。在其他事情之后，他接着写下了关于他所遭遇之事的以下内容：\par

\Quote{18. 日耳曼努斯确实夸耀许多宣认。他当然可以说他自己忍受了许多逆境。但他能列举出和我们一样多的宣判、没收、剥夺、财产掠夺、尊严丧失、对世俗荣耀的轻视、对总督和议员阿谀的藐视，以及对反对者威胁、呼喊、危险、迫害的忍耐忍受，以及流浪、困苦和各种患难，就像我在德西乌斯和萨比努斯手下所遭遇的，以及现在在埃米利安努斯手下仍然持续的那样吗？但日耳曼努斯在哪里见过？关于他又有何记述？}\par

\Quote{19. 但我还是远离这极大的愚昧，我因日耳曼努斯而陷入了其中。出于同样的原因，我也停止向知道这一切的弟兄们讲述所发生的一切。}\par

20. 同一作者在致多米提乌斯和狄第穆斯的书信中也提到了迫害的一些细节，如下：\Quote{我们的人很多，你们不认识，列出他们的名字是多余的；但要明白，男女老少、少女和主妇、军人与平民，各种种族和年龄的人，有些通过鞭打和火刑，有些通过刀剑，在争斗中得胜并接受了他们的冠冕。}\par

21. 但在某些人的情况下，很长时间不足以上他们看来蒙主悦纳；确实，在我看来也是如此，尚不足够的时间还没有过去。因此他按他所知道合适的时候保留了我，说：\BibleQuote{在悦纳的时候我应允了你，在救恩的日子我帮助了你。}\BibleRef{依 49:8}\par

22. 既然你们询问我们的事务，并希望我们告知你们我们的处境，你们已完全听到，当我们——即我、加尤斯、福斯托斯、伯多禄和保禄——被一名百夫长和官员们，连同他们的士兵和仆人作为囚犯带走时，来自\Place{马雷奥提斯}的某些人强迫我们跟从他们，尽管我们不愿跟随。\par

23. 但现在只有我、加尤斯和伯多禄，被剥夺了其他弟兄，被囚禁在\Place{利比亚}一处荒凉干燥的地方，离\Place{帕雷托尼翁}有三天的路程。\par

24. 他接着说道：\Quote{司铎马克西穆斯、狄奥斯库鲁斯、德默特琉和路济安藏在城里，秘密探望弟兄们；因为福斯蒂诺斯和阿奎拉，他们在世上更显赫，正在\Place{埃及}流浪。但执事福斯托斯、欧瑟伯和凯勒蒙，在瘟疫中幸存下来。欧瑟伯是那位起初就被天主加强和赋予能力，以充沛精力完成对监禁中宣认者的服务，并承担为已完成的蒙福殉道者埋葬遗体这一危险任务的人。}\par

25. \Quote{正如我先前所说，总督至今仍以残忍的方式处死那些受审的人。他用酷刑消灭一些人，用监禁和锁链折磨另一些人；他不允许任何人接近他们，并调查是否有人这样做。尽管如此，天主还是通过弟兄们的热心和坚持，赐予受苦者安慰。}\par

26. 以上是狄约尼削的记述。但应当知道，他所称为执事的欧瑟伯，不久后成为\Place{叙利亚}\Place{劳狄刻雅}教会的主教；而他提及当时为司铎的马克西穆斯，后来继狄约尼削本人成为\Place{亚历山大里亚}的主教。但与他同在的福斯托斯，当时因宣认而闻名，一直存留到我们这时代的迫害中，那时他年事已高，寿数满足，被斩首而殉道结束生命。这就是当时发生在狄约尼削身上的事情。\par

% structure: p0078 source=h2 target=subsection reason=relative-heading-rank; base=h2

\subsection{第12章　\Place{凯撒勒雅}（\Place{巴勒斯坦}）的殉道者}

在瓦勒里安治下的上述迫害期间，\Place{巴勒斯坦}的\Place{凯撒勒雅}有三人，因宣认基督而卓著，被野兽吞噬，获得了神圣的殉道。其中一人名叫普里斯库斯，另一人叫玛尔古斯，第三人的名字是亚历山大。据说这些居住在农村的人，起初表现得怯懦，仿佛粗心大意、不加思考。因为当那些以属天渴望寻求奖赏的人获得机会时，他们却轻看它，恐怕过早夺取殉道的冠冕。但经过商议后，他们急忙前往\Place{凯撒勒雅}，来到审判官面前，并得到了我们提到的那结局。据传，此外，在同一迫害中和同一城市里，一位妇女也经历了类似的争斗。但据说她属于马尔基翁的派别。\par

% structure: p0080 source=h2 target=subsection reason=relative-heading-rank; base=h2

\subsection{第13章　加里恩努斯下的和平}

1. 此后不久，瓦勒里安被蛮族所掳，沦为奴隶；他的儿子成为唯一统治者，更谨慎地处理政务。他立即通过公开谕令停止了对我们的迫害，并以如下诏书指示主教们自由履行他们惯常的职务：\par

\Quote{皇帝凯撒普布里乌斯·李锡尼·加里恩努斯，虔诚的、幸福的、奥古斯都，致狄约尼削、平纳斯、德默特琉及其他主教。我已命令向全世界宣布我恩赐的慷慨，使他们得以离开宗教崇拜的场所。为此目的，你们可使用我这道诏书的副本，使无人骚扰你们。而你们现在能合法做的事，早已由我允准。因此，最高行政官奥勒留·基勒尼乌斯将遵守我所颁布的这条命令。}\par

我以拉丁文译出此文，以便更易理解。另有他的一道诏书存世，致其他主教，允许他们重新占领所谓的墓地。\par

% structure: p0084 source=h2 target=subsection reason=relative-heading-rank; base=h2

\subsection{第十四章 当时兴盛的主教们}

那时，西斯笃仍主持罗马教会；德默特里亚诺继法比之后主持安提约基雅教会；费尔米利亚诺主持卡帕多细雅的凯撒勒雅教会；此外，奥利金的友人额我略及其弟雅特诺多若主持本都的教会；巴勒斯坦凯撒勒雅的德奥克提斯托去世后，多木诺在那里接受主教职，但任期甚短，我们的同代人德奥特克诺继任。他也是奥利金学派的一员。但在耶路撒冷，玛匝巴内斯去世后，多年以来在我们中备受赞誉的希默乃乌斯接替了他的席位。\par

% structure: p0086 source=h2 target=subsection reason=relative-heading-rank; base=h2

\subsection{第十五章 马里诺在凯撒勒雅殉道}

1. 此时，各地教会恢复和平，巴勒斯坦凯撒勒雅的马里诺，因其军功受到尊敬，又因家世和财富显赫，为基督作了见证而被斩首，缘由如下。\par

2. 葡萄枝是罗马人中一种荣誉标志，据说得到它的人就成为百夫长。一个职位空缺，依顺序轮到马里诺。但他正要接受荣誉时，另一个人到法庭声称，根据古代法律，马里诺作为基督徒，不向皇帝献祭，不应接受罗马尊位；该职位应属于他。\par

3. 于是法官（名叫阿凯乌斯）不安，先问马里诺的信仰。见他始终承认自己是基督徒，就给他三小时考虑。\par

4. 他走出法庭后，当地主教德奥特克诺带他到一旁与他谈话，拉着他的手领他进入教堂。与他站在教堂内的圣所，德奥特克诺稍微掀起他的外衣，指着他腰间挂的剑；同时在他面前摆上天主福音的经卷，告诉他选择两者中哪一个。马里诺毫不犹豫伸出右手，取了神圣的经卷。德奥特克诺对他说：\Quote{那么，抓紧，抓紧天主，靠祂加强，愿你得到你所选择的，平安去吧。}\par

5. 他立即返回时，传令官喊他上法庭，因为规定的时间已到。他站在法庭前，显示对信仰更大的热忱，随即被带走，以死亡完成他的路程。\par

% structure: p0092 source=h2 target=subsection reason=relative-heading-rank; base=h2

\subsection{第十六章 关于阿斯提里的事迹}

\Emph{阿斯提里}也因与此事相关的虔诚勇敢而受人纪念。他是罗马元老院成员，得皇帝宠信，因出身高贵和财富而为众人所知。他当时在场目睹殉道者去世，将尸体扛在肩上，用华丽昂贵的衣服为他装扮，以隆重方式准备坟墓，给予合适的安葬。他的朋友中留存至今的，还讲述许多关于他的其他事实。\par

% structure: p0094 source=h2 target=subsection reason=relative-heading-rank; base=h2

\subsection{第十七章 在帕尼亚斯显示我们救主大能的记号}

其中包括以下奇事。在斐尼基人称为帕尼亚斯的凯撒勒雅斐理伯，帕尼乌斯山脚下有泉源，约但河由此流出。据说在某节日，他们将祭牲投入泉中，因魔鬼的力，祭牲奇妙地消失，此事成为在场者的著名奇观。阿斯提里有一次在场目睹此事，见众人惊异，便怜悯他们的迷惑；他举目望天，透过基督恳求统御万有的天主斥责迷惑民众的魔鬼，终结众人的迷惑。据说他如此祈祷后，祭牲立即漂浮在泉面。于是奇迹消失，此后该地再未显任何奇迹。\par

% structure: p0096 source=h2 target=subsection reason=relative-heading-rank; base=h2

\subsection{第十八章 血漏妇人建立的雕像}

1. 既然我提到此城，我认为不应省略一件值得为后代记录的事。据说患血漏的妇人，如我们从神圣福音所知的，从我们的救主那里获得治愈，她来自此地，她的房子在该城仍可看到，救主对她的仁慈有显著的纪念物留存那里。\par

2. 因为在她家门口的一块高石上，立着一尊铜像，是一个跪着的妇人，双手伸出，如在祈祷。对面是另一尊直立的男子像，用同样材料制成，穿着得体的双斗篷，向妇人伸手。在他的脚边，靠近雕像本身，有一种奇特的植物，攀爬到铜斗篷的边缘，能够治疗各种疾病。\par

3. 他们说这尊雕像是耶稣的像。它一直保存到我们时代，我们自己在城里时也曾见过。\par

4. 外邦人中那些古时蒙我救主施恩的人做出这等事，并不奇怪，因为我们也得知祂的宗徒\Person{保罗}和\Person{彼得}，以及基督本人的肖像被保存在绘画中，古人惯于，诚如所料，依照外邦人的习惯，不加区分地将这种敬意给予那些被他们视为救星的人。\par

% structure: p0101 source=h2 target=subsection reason=relative-heading-rank; base=h2

\subsection{第十九章 雅各伯的主教座}

雅各伯的主教座，他由救主本人及宗徒们首先接受了耶路撒冷教会的主教职，并如神圣记载所显示，被称为基督的兄弟，一直保存至今，那些在此相继跟随他的弟兄们，清楚地向所有人展示了古时及现今的人们因虔诚而对圣人所持有并仍持有的敬意。就此而言，已经足够。\par

% structure: p0103 source=h2 target=subsection reason=relative-heading-rank; base=h2

\subsection{第二十章 狄约尼修的节期书信，其中他也颁定了逾越节法则}

\Person{狄约尼修}，除已提及的书信外，当时还写了流传至今的节期书信，其中他使用颂扬之词谈论逾越节庆典。他将其中一封致\Person{弗拉维}，另一封致\Person{多弥提}与\Person{狄第慕}，在其中他提出了一个八年周期，主张不应在春分之前庆祝逾越节。此外，在迫害仍盛行之时，他还致信亚历山大里亚的司铎同道，以及其他不同人士。\par

% structure: p0105 source=h2 target=subsection reason=relative-heading-rank; base=h2

\subsection{第二十一章 亚历山大里亚发生的事}

1. 和平刚恢复，他便回到\Place{亚历山大里亚}；但因叛乱和战争再次爆发，使他无法照管因骚乱而分散在各处的所有弟兄，在逾越节庆典时，仿佛他仍被流放在亚历山大里亚之外，他再次致信他们。\par

2. 在后来写给埃及主教\Person{希辣克斯}的另一封节期书信中，他提到了当时亚历山大里亚盛行的叛乱，如下：\par

\Quote{有什么奇怪，我很难与远处的人通信，因为我甚至无法与自己论理，或为自己的性命筹谋？}\par

\Quote{3. 诚然，我需要寄信给那些如同我自身的一部分、同住一室、同心同德的弟兄，以及同一教会的公民；但我不知如何寄出。因为一个人不仅走出行省的边界，甚至从东方到西方，也比从亚历山大里亚到亚历山大里亚本身更容易。}\par

\Quote{4. 因为城市的中心比以色列两代人所穿越的那广大无路的沙漠更加曲折难行。我们平静无波的海港已变得像被分开并筑墙的海，以色列曾从中驱驰，埃及人即在其大道上被淹没。因为常因那里发生的屠杀，它们看起来像红海。}\par

\Quote{5. 流经城市的河流有时似乎比无水的沙漠更干，比以色列人穿越时因口渴而向\Person{梅瑟}呼求、水从峻峭的磐石流出（藉着那位独行奇事者）时更焦干。}\par

\Quote{6. 它又曾暴涨，淹没四周的乡野、道路和田园；威胁要重现\Person{诺厄}时代发生的大洪水。它流淌着，总是被血、杀戮和溺毙污染，如同\Person{梅瑟}将水变为血时\Person{法郎}所遭遇的那样，发出恶臭。}\par

\Quote{7. 还有什么水能净化那净化一切的水？那浩瀚的、人不可跨越的海洋，若倾入其中，怎能洁净这苦海？或者那流出伊甸的大河，若将其分为四支的一支倾入基红，怎能洗去这污秽？}\par

\Quote{8. 或者，被这些有害瘴气污染的空气何时能变纯净？因为从地上升起这样的蒸汽，从海上吹来这样的风，从河中飘来这样的微风，从港口腾起这样的雾气，使露珠如同从我们周围所有元素中腐烂的尸体中排泄出来的。}\par

\Quote{9. 然而人们惊奇，不能理解这些不断的瘟疫从何而来；这些严重的疾病从何而来；这些各种致命的疾患从何而来；这各种各样大范围的人类毁灭从何而来；为何这座大城不再包含像从前那样多的居民，从幼婴到年寿极高者。但从前四十到七十岁的人如此众多，以致现在即使将十四到八十岁的人登记名册以领取公共粮食，也无法填补其数目。}\par

\Quote{10. 而外表最年轻的人，实际上变得与从前最年老的人同龄。但尽管他们看见人类种族如此持续减少和耗损，尽管他们完全的毁灭在增加和逼近，他们并不战栗。}\par

% structure: p0117 source=h2 target=subsection reason=relative-heading-rank; base=h2

\subsection{第二十二章 降于他们的瘟疫}

1. 这些事件之后，一场瘟疫疾病接着战争而来；在节期临近时，他又写信给弟兄们，描述这场灾难所带来的苦难。\par

\Quote{2. 在其他人看来，目前可能不是举行节期的适当时候。确实，这对他们而言也不是适当时候；无论是悲伤的时候，还是可能被认为特别快乐的时候。现在，确实，一切都是眼泪，人人都在哀悼，哀号每日响彻城中，因为死者与垂死者的众多。}\par

\Quote{3. 因为正如关于埃及长子的记载，现在\BibleQuote{有极大的哀号，因为没有一家没有一具死尸}（\BibleRef{出 12:30}）。}\par

\Quote{4. 因为已经发生了许多可怕的事。首先，他们驱逐我们；当我们孤立无援，被所有人迫害和处死时，我们仍举行了节期。每一个受苦之地对我们都是节庆之所：田野、旷野、船只、客栈、监狱；但完美的殉道者举行了最喜乐的节期，在天上享受盛宴。}\par

5. 在此之后，战争与饥荒接踵而至，我们与外邦人一同承受。然而，那些他们加诸于我们的苦难，却唯有我们独自承担；同时，我们也经历了他们彼此加害与受苦的后果；另一方面，我们也在基督所赐予我们的平安中喜乐，这平安他只赐给了我们。\par

6. \Quote{但在我与他们双方都享受了极为短暂的休息之后，这场瘟疫袭击了我们；对外邦人而言，它比任何恐惧更可怕，比任何灾祸更难忍受；正如他们的一位作家所言，这是唯一超越一切希望的灾祸。然而对我们来说，并非如此，它与其他事情一样，是一种锻炼和考验。因为它并未远离我们，但对外邦人的打击更为严重。}\par

7. 他又接着写道：\par

我们大多数弟兄在超乎寻常的爱德和兄弟情谊中毫不吝惜。他们彼此紧握，无畏地探望病人，不断服侍他们，在基督内为他们服务。他们极其喜乐地与病人一同死去，承担他人的痛苦，将邻人的疾病引到自己身上，甘心接受他们的痛苦。许多照顾病人、给予他人力量的人，自己却死去，将死亡转嫁到自己身上。那句常被视为礼貌客套的流行语，他们当时确实以实际行动实现了，作为他人的\Quote{\OriginalTerm{废物}{offscouring}}而离去。\par

8. 确实，我们中最优秀的弟兄就这样离世，其中包括一些司铎、执事以及民众中最有声望的人；因此，这种死亡形式，因其所彰显的极大虔诚和坚定信德，似乎丝毫不逊于殉道。\par

9. 他们用敞开的双手和怀抱接过圣人们的遗体，合上他们的眼睛和嘴巴；他们用肩膀抬起他们，铺放好；他们紧抱他们，拥抱他们；他们用洗濯和衣裳妥善地预备他们。不久之后，他们自己也受到同样的待遇，因为幸存者不断追随先他们而去的人。\par

10. \Quote{但对外邦人而言，一切完全不同。他们遗弃那些开始生病的人，逃离最亲密的朋友。他们将半死的人扔到街上，将死者像垃圾一样丢弃，无人埋葬。他们躲避任何与死亡有份或共融的事；然而，尽管他们处处小心，仍难以逃脱死亡。}\par

11. 在这封信之后，当城市恢复和平后，他又写了一封节期信函给埃及的弟兄们，此外又写了几封。此外还存有一封论安息日的信，以及另一封论操练的信。\par

12. 此外，他又写信给赫尔马蒙和埃及的弟兄们，详细描述了德西乌斯及其继任者的邪恶，并提及加里恩努斯治下的和平。\par

% structure: p0131 source=h2 target=subsection reason=relative-heading-rank; base=h2

\subsection{第二十三章 加里恩努斯的统治}

1. 但没有什么比亲耳聆听他自己的话更好了，其言如下：\par

然后，他背叛了在他之前的一位皇帝，并与另一位作战，迅速且彻底地与其全家一同灭亡。但加里恩努斯被宣告并被普遍承认为既是旧皇帝又是新皇帝，他在他们之前，并继续存在于他们之后。\par

2. 因为正如依撒意亚先知所言：\InnerQuote{看，起始的事已成就，新事将要发生。}\BibleRef{依 42:9} 因为如同乌云遮蔽太阳的光芒片刻，隐藏了它并取而代之；但当乌云过去或消散，先前升起的太阳再次显现；同样，马克里亚努斯自荐并接近加里恩努斯现有的帝国，他不存在，因为他从未存在。但另一位正如他原有。\par

3. 他的王国，仿佛抛弃了衰老，并从先前的邪恶中得到净化，如今更加生机勃勃地绽放，被更远地看见和听见，并向各方扩展。\par

4. 然后他用以下话语指明他写此信的时期：\par

\Quote{我再次想起回顾那些帝国年岁的日子。因为我发现那些最不虔敬的人，虽然曾闻名一时，却在短时间内变得无名。但更神圣、更虔诚的君主，已度过第七年，如今正完成第九年，我们将在其中庆祝这节日。}\par

% structure: p0138 source=h2 target=subsection reason=relative-heading-rank; base=h2

\subsection{第二十四章 奈颇斯及其分裂}

1. 除这些之外，他还准备了两卷论应许的书。其起因是埃及的一位主教奈颇斯，他教导说，圣经中对圣人们的应许应按照更犹太的方式理解，并且在这地上将有一个身体享乐的千年王国。\par

2. 由于他认为自己可以通过若望默示录确立其私人意见，他写了一本关于此主题的书，题为《反寓意论者》。\par

3. 狄约尼削在他的《论应许》书中对此提出反对。在第一卷中，他陈述了自己关于此教义的意见；在第二卷中，他论述若望默示录，并在开头提及奈颇斯，如此写道：\par

4. 但既然他们提出奈颇斯的某部著作，并自信地依赖它，仿佛它无可争议地证明了基督将在地上为王，我承认，在许多其他方面，我赞同并喜爱奈颇斯，因他的信德、勤勉以及钻研圣经的热诚，也因他广泛的圣咏创作，至今仍有许多弟兄从中得到喜乐；我更加敬重他，因为他已先我们而安息。但真理应被最爱并最受尊崇。虽然我们应毫不吝啬地赞美和赞同那些说得正确的话，但我们也应审视和纠正那些似乎写得不够纯正的。\par

5. 倘若他亲临现场陈述意见，仅凭口头的讨论，通过问答来说服并调和反对者，便已足够。但有些人认为他的著作甚为可信，某些教师视法律和先知为无物，不遵从福音，轻看宗徒书信，却对这部著作的教导许下诺言，仿佛它是某种伟大的隐秘奥秘；他们不许我们较为单纯的弟兄对主荣耀而真正神圣的显现、我们从死者中复活、我们被聚集到他身边并肖似他怀有任何崇高伟大的思想，反而引导他们在天主之国中寄望于微小而必朽的事物，以及如今存在的事物——既然如此，我们必须如同亲临一般，与我们的弟兄内波斯进行辩论。接着他说：\par

6. 当我在\Place{阿尔西诺埃}地区时——如你所知，这一学说已在那里盛行多时，以致导致整个教会分裂和背弃——我召集了村庄中的长老和教师弟兄们，并邀请愿意参加的弟兄一同出席，劝勉他们公开审查这一问题。\par

7. 因此，当他们把这部书当作攻不可破的武器和堡垒带到我面前时，我连续三天从早到晚与他们坐在一起，努力纠正其中所写的内容。\par

8. 我因弟兄们的坚定、真诚、受教和明智而欣喜；我们有序而适度地思考了问题、困难和一致之处。我们并不坚决或争辩地维护先前所持的意见，除非它们显得正确。我们也不回避反对意见，而是尽力把握并确认摆在我们面前的事物；如果给出的道理令我们满意，我们并不羞于改变观点而同意他人；相反，我们诚心诚意、心怀坦荡地在天主面前敞开，接受凡由圣经的证明和教导所确立的一切。\par

9. 最后，这一学说的创始者和推动者——人称\Person{科拉西翁}的——在所有在场的弟兄面前承认并向我们作证，他不再坚持这一观点，也不再讨论、提及或教导它，因为他已完全被反驳此说的论证所说服。另一些弟兄则对会议以及全体所表现的和解与和谐精神表示满意。\par

% structure: p0148 source=h2 target=subsection reason=relative-heading-rank; base=h2

\subsection{第二十五章 若望的默示录}

1. \Emph{随后}，他以这种方式谈论\WorkTitle{若望的默示录}。\par

在我们以前，有些人完全弃绝并拒绝此书，逐章批评，称其缺乏意义和逻辑，并断言其标题是伪造的。\par

2. 因为他们说，这不是若望的作品，也不是启示，因为它被一层浓厚的晦涩之纱厚厚覆盖。他们断言，没有一位宗徒、一位圣人，或教会中的任何人是其作者，而是塞林托——他建立了以其命名的塞林托派——为了给其虚构之作寻求可信的权威，便冠以若望之名。\par

3. 因为他所教导的教义是：基督的国度将是地上的国度。由于他本人沉溺于肉体的享乐，本性完全属感官，他便梦想那国度在于他所贪求的事物，即口腹之欲和性欲；也就是说，吃喝嫁娶，以及节庆、祭献和宰杀牺牲——他以为在这些幌子下可以更好地满足自己的欲望。\par

4. 但我不能冒然拒绝此书，因为许多弟兄对它怀有崇高的敬意。然而我猜想它超出我的理解，其中每一部分都有某种隐藏而更加奇妙的含义。因为即使我不理解，我也怀疑在字面之下蕴含着更深的意义。\par

5. 我不以自己的理性来衡量和判断它们，而是将更多的事归于信德，视它们过于高深，非我所能把握。我不拒绝我所不能理解的，反而因不理解而惊叹。\par

6. 此后，他审阅了整部\WorkTitle{默示录}，并证明不可能按其字面意义来理解，接着写道：\par

可以说，先知在结束全部预言后，宣告那些遵守它的人有福，也包括他自己。因为他\BibleQuote{说：\Quote{那遵守这书卷预言的话的，是有福的。我若望，就是听见并看见这些事的。}}\par

7. 因此，我不否认他被称为若望，也不否认此书是一位若望的作品。我也同意它是出自一位圣洁而受默感的人之手。但我不能轻易承认他就是那位宗徒，载伯德的儿子，雅各伯的兄弟，即写了若望福音和公函的那一位。\par

8. 因为我从二者的风格、表达方式以及整部书的构成来判断，它不是他的作品。因为福音作者无论在福音或公函中都未提及自己的名字，也未宣告自己。\par

9. 接着他又补充道：\par

但若望从未以提及自己或提及他人的方式说话。而默示录的作者在开头就介绍自己：\BibleQuote{耶稣基督的启示，就是天主赐给他，叫他把必要快发生的事指示给他的众仆人；他打发天使，藉着他的使者，指示给他的仆人若望；若望便为天主的话并为耶稣基督的见证，凡他所看见的，都作了见证。}\BibleRef{默 1:1}\par

10. 接着他还写了一封信：\BibleQuote{若望致书给亚细亚的七个教会：愿恩宠与平安归于你们。}\BibleRef{默 1:4}但福音作者即使在其公函中也未冠以己名；而是不加序言，直接从天主启示奥秘本身开始：\BibleQuote{论到那从起初就有的，我们听见、亲眼看见过的。}\BibleRef{若一 1:1}因为对于这样的启示，主也曾祝福伯多禄说：\BibleQuote{西满巴尔约纳，你是有福的，因为不是肉和血启示了你，而是我在天之父。}\BibleRef{玛 16:17}\par

11. 然而，在那些被称为《若望》第二或第三封的书信中，虽然它们很短，却没有出现若望的名字；而是写着匿名的短语\Quote{长老}。但这位作者并不满足于只提一次自己的名字就继续写作；他再次提及：\BibleQuote{我，若望，也就是你们的弟兄，在耶稣基督的患难、国度与忍耐里一同有分的，曾为天主的圣言和耶稣的见证，在那名叫帕特摩的海岛上。}（\BibleRef{默 1:9}）并且接近结尾时他这样说：\Quote{那遵守这书上预言的话语的人是有福的，而我，若望，是看见并听见这些事的。}\par

12. 然而，我们必须相信他所说的，即写这些事的人名叫若望；但他究竟是谁，并不清楚。因为他并不像在福音中常常那样说，他是主所爱的门徒，或是靠在他胸膛的那一位，或是雅各伯的兄弟，或是主的目击者和亲耳听闻者。\par

13. 因为他若想清楚地表明自己，就会提及这些事。但他丝毫没有提及；而是称自己为我们的弟兄和同伴，是耶稣的见证人，并因目睹并听见那些启示而蒙福。\par

14. 但我认为，有许多人与宗徒若望同名，他们因爱慕若望，并因仰慕和效法他，且渴望像他一样蒙主宠爱，便也取了同样的姓氏，正如许多信友的子女被称为保禄或伯多禄一样。\par

15. 例如，在\WorkTitle{宗徒大事录}中也提到另一个若望，别号马尔谷，巴尔纳伯和保禄曾带他同行；关于他，也记着说：\BibleQuote{他们也有若望作助手。}（\BibleRef{宗 13:5}）但我不会说就是他写了此书。因为经上并没有记载他随他们去了亚细亚，而是说：\BibleQuote{保禄和他的同伴从帕佛开船，来到旁非里雅的培尔革，若望却离开他们，回了耶路撒冷。}（\BibleRef{宗 13:13}）\par

16. 但我认为，他是亚细亚的另外一个人；因为有人说在厄弗所有两座墓碑，每座都刻有若望的名字。\par

17. 而且，从思想、用词及其组织来看，可以合理地推断此人与那人不同。\par

18. 因为《福音》和《书信》彼此相符，以同样的方式开始。前者说：\BibleQuote{在起初已有圣言}（\BibleRef{若 1:1}）；后者说：\BibleQuote{那从起初就有的}（\BibleRef{若一 1:1}）。前者说：\BibleQuote{圣言成了血肉，住在我们中间，我们看见了他的光荣，正如父独生者的光荣}（\BibleRef{若 1:14}）；后者以稍加改动的方式说同样的事：\BibleQuote{我们所听见、所看见的、我们亲眼见过、我们亲手摸过的生命之圣言——那生命已显示出来了}（\BibleRef{若一 1:1}）。\par

19. 因为他一开始就引入这些事，坚持它们，正如从下文可见，是为了反驳那些说主没有降生成人的人。因此他也谨慎地补充说：\BibleQuote{我们看见了，并为此作证，也将那与父同在、且显示给我们的永生，传报给你们。我们将所见所闻的也传报给你们。}（\BibleRef{若一 1:2}）\par

20. 他坚持这一点，并不偏离主题，而是在相同的标题和名称下讨论一切，其中一些我们将简要提及。\par

21. 任何人仔细查考，都会发现这些短语：\Quote{生命}、\Quote{光明}、\Quote{脱离黑暗}常常出现在两者中；还有连续出现的\Quote{真理}、\Quote{恩宠}、\Quote{喜乐}、\Quote{主的血肉}、\Quote{审判}、\Quote{罪过的赦免}、\Quote{天主对我们的爱}、\Quote{彼此相爱的命令}、我们应当\Quote{遵守一切诫命}、\Quote{对世界、魔鬼、反基督的谴责}、\Quote{圣神的恩许}、\Quote{天主的义子之恩}、\Quote{不断要求于我们的信德}、\Quote{父与子}，到处出现。事实上，显而易见的是，整部《福音》和《书信》都带有同一特征。\par

22. 但《默示录》与这些著作不同，与它们格格不入；毫不触及，也不丝毫沾边；几乎可以说，连一个共同的音节都没有。\par

23. 不仅如此，那封《书信》——因为我不提《福音》——没有提及，也没有暗含《默示录》的任何内容；《默示录》也没有提及《书信》。但保禄在他的书信中，曾透露一些他的启示，尽管他没有单独写出来。\par

24. 此外，还可以表明，《福音》和《书信》的措辞与《默示录》不同。\par

25. 因为前者不仅在希腊语上没有错误，而且在表达、推理和整体结构上都很优雅。它们远没有暴露出任何\OriginalTerm{生硬措辞}{barbarism}或\OriginalTerm{语法错误}{solecism}，或任何粗俗之处。因为写作者似乎拥有言谈的两项要素，即知识的恩赐和表达的恩赐，正如主将二者都赐予了他。\par

26. 我不否认另一位作者看见了启示，并领受了知识和预言。然而，我察觉他的方言和语言并非准确的希腊语，而是使用\OriginalTerm{生硬习语}{barbarous idioms}，并且在某些地方还有\OriginalTerm{语法错误}{solecisms}。\par

27. 这里不必指出这些，因为我不愿让任何人以为我说这些话是出于嘲弄的意思；我之所以这样说，只是为了清楚地显示这些著作之间的不同。\par

% structure: p0179 source=h2 target=subsection reason=relative-heading-rank; base=h2

\subsection{第二十六章 狄约尼削的书信}

1. 除此以外，狄约尼削还有许多其他书信存世，例如那些驳斥撒伯里乌的书信，致贝尔尼切教会的主教阿孟的信，致特莱斯福禄的信，致尤弗拉诺尔的信，以及另一封致阿孟和尤波若的信。他还写了四部关于同一主题的其他著作，并致送给了他在罗马的同名者狄约尼削。\par

2. 除此以外，我们还有他许多书信，以及用书信体写成的宏篇巨著，例如\WorkTitle{论自然}一书，致青年弟茂德；以及\WorkTitle{论诱惑}一书，他也奉献给了尤弗拉诺尔。\par

3. 此外，在一封致\Place{彭塔波利}各堂区主教\Person{巴西利德}的信中，他说自己曾写过一篇关于\BibleBook{}{训道篇}开头的注释。他还留下了写给同一人的多封书信。关于\Person{狄约尼削}，就说到这里。\par

但我们既已完成了对这些事的记述，请允许我们向后世展示我们这时代的面貌。\par

% structure: p0184 source=h2 target=subsection reason=relative-heading-rank; base=h2

\subsection{第二十七章 \Person{撒摩沙塔的保禄}及其在\Place{安提约基雅}所倡导的异端}

1. \Person{西斯笃}治理罗马教会十一年后，与他同名、亚历山大城的\Person{狄约尼削}接任。大约同时，\Person{德默特里亚努斯}在\Place{安提约基雅}去世，\Person{撒摩沙塔的保禄}接受了该主教职位。\par

2. 由于他持有与教会教导相悖的、对基督低下卑贱的看法，即认为基督的本性只是一个普通人，于是\Person{亚历山大的狄约尼削}受邀前来参加教务会议。但因年迈体衰未能成行，他通过书信就讨论的问题发表了意见。然而，各教会所有其他牧者从四面八方急速聚集到\Place{安提约基雅}，如同一起来对付基督羊群的掠夺者。\par

% structure: p0187 source=h2 target=subsection reason=relative-heading-rank; base=h2

\subsection{第二十八章 当时的卓越主教}

1. 其中最卓越的是：\Person{斐尔米连}，\Place{卡帕多细亚}省\Place{凯撒勒雅}的主教；\Person{额我略}与\Person{阿特诺多鲁斯}兄弟，\Place{本都}地区教会的牧者；\Place{塔尔索}教区的\Person{赫肋努斯}和\Place{依科尼雍}的\Person{尼各马}；此外，\Place{耶路撒冷}教会的\Person{希默纳尤斯}及邻近\Place{凯撒勒雅}教会的\Person{德奥特克诺}；还有\Person{玛息莫}，他卓越地治理着\Place{波斯特拉}的弟兄们。若有人计算，他必能注意到许多其他人，此外还有当时为同一事由聚集在上述城市的司铎和执事。但以上是最卓越的。\par

2. 当这些人多次在不同时间聚会讨论这些事时，每次会议都在辩论和提问；撒摩沙塔派的追随者竭力掩盖和隐瞒他的异说，而其他人则热切地揭露和彰显他的异端以及对基督的亵渎。\par

3. 与此同时，\Person{狄约尼削}在\Person{加里恩努斯}在位第十二年去世，担任亚历山大主教十七年，\Person{玛息莫}接替他。\par

4. \Person{加里恩努斯}统治十五年后，\Person{克劳狄}继位，他在两年后将政权交给了\Person{奥勒良}。\par

% structure: p0192 source=h2 target=subsection reason=relative-heading-rank; base=h2

\subsection{第二十九章 \Person{保禄}被一位来自诡辩学派的司铎\Person{马尔基翁}驳斥后，被绝罚}

1. 在他统治期间，召集了一次由众多主教组成的最后教务会议，\Place{安提约基雅}异端的首领被揭露，他的虚假教义在众人面前显明，他被绝罚，逐出普世的天主教会。\par

2. \Person{马尔基翁}尤其将他从隐藏之处揪出来，驳斥了他。这是一个博学之人，在\Place{安提约基雅}担任希腊学问的诡辩学派校长；但因他信仰基督的卓越高贵，被任命为该教区的司铎。此人曾与他进行辩论，并由速记员记录，我们知道这份记录至今尚存。只有他识破了那个伪装和欺骗他人的人。\par

% structure: p0195 source=h2 target=subsection reason=relative-heading-rank; base=h2

\subsection{第三十章 主教们反对\Person{保禄}的书信}

1. 为此事聚集的牧者一致同意，起草了一封致罗马主教\Person{狄约尼削}和亚历山大的\Person{玛息莫}的书信，并发送到各省。信中向所有人表明了他们自己的热忱和\Person{保禄}的错谬，以及他们与他进行的论辩和讨论，并显示了这个人的全部生活和行为。此时不妨将他们的书信摘录如下：\par

2. \Quote{致\Person{狄约尼削}与\Person{玛息莫}，以及我们在世界各地的所有同工——主教、司铎、执事，并致普世的天主教会：\Person{赫肋努斯}、\Person{希默纳尤斯}、\Person{德奥斐洛}、\Person{德奥特克诺}、\Person{玛息莫}、\Person{普洛克鲁斯}、\Person{尼各马}、\Person{埃利安努斯}、\Person{保禄}、\Person{波拉努斯}、\Person{普罗托格内斯}、\Person{希拉克斯}、\Person{欧提基乌}、\Person{德奥多鲁}、\Person{马尔基翁}、\Person{路基乌}，以及所有与我们同住在邻近城邦和民族的其他主教、司铎、执事，以及各天主的教会，在主内问候亲爱的弟兄们。}\par

3. 稍后他们继续写道：\Quote{我们派人并邀请了远处许多主教来解救我们脱离这致命的教义，例如\Person{亚历山大的狄约尼削}和\Place{卡帕多细亚}的\Person{斐尔米连}，这些可敬的人。前者认为这谬误的首领不配得到致信，便将一封信寄往\Place{安提约基雅}，不是写给他，而是写给整个教区，其副本附于下方。}\par

4. \Quote{但\Person{斐尔米连}两次前来，谴责了他的新论，正如我们这些在场者所知并作证，许多人也了解。然而因他承诺改变观点，\Person{斐尔米连}相信了他，并希望，在不辱没圣言的情况下，能成就必要的事。于是他拖延了此事，被那甚至否认自己的天主和主、且未持守从前的信德的人所欺骗。}\par

5. \Quote{当时\Person{斐尔米连}正再次前往\Place{安提约基雅}，已到达\Place{塔尔索}，因为他已亲身经验到那否认天主的恶行。但正当我们聚集一起召唤他、等待他到来时，他却去世了。}\par

6. 在其他内容之后，他们如下描述他的生活方式：\par

7. \Quote{如今他既已偏离信德准则，转而追求卑劣虚假的教导，既然他不在我们当中，我们便不必论断他的行为：例如，他原本穷困贫乏，既未从父辈继承财富，也未通过交易或经商获利，如今却因不法与渎神行为，以及从弟兄们那里强索的东西而拥有大量财富——他剥夺受屈者的权利，应许以酬报相助，却又欺骗他们，掠夺那些在困境中为求与压迫者和解而乐意付出的人，\BibleQuote{以为得利是敬虔}；见：\BibleRef{弟前 6:5}——}\par

8. 或者因为他骄傲自大，得意洋洋，贪图世俗的尊荣，宁愿被称为\OriginalTerm{杜凯纳里乌斯}{ducenarius}而不愿被称为主教；他在市场上昂首阔步，边走边朗读并背诵书信，前呼后拥，随从护驾，以致因他心中的骄傲自大而使人嫉妒并憎恶信仰——\par

9. 或者因为在教会会议上玩弄诡计，设法自我荣耀，用外表欺骗人，震惊单纯之人的心，为自己预备一个审判台和高高的宝座——不像基督的门徒——并且拥有一个\OriginalTerm{密室}{secretum}——如同世界的统治者——这样称呼它，用手拍大腿，用脚跺审判台——或者因为他斥责和侮辱那些不鼓掌、不像在剧场里那样摇手绢、呼喊跳跃、如同围绕他的男男女女那样以不体面的方式听他讲道，而是像在天主的殿里一样虔诚有序地聆听的人——或者因为他公开粗暴地攻击已经去世的圣言阐释者，抬高自己，不像一个主教，而像一个诡辩家和变戏法的人，\par

10. 并且阻止向吾主\Person{耶稣基督}咏唱圣咏，声称它们是现代人的现代作品，还训练妇女在教会中，在逾越节的大日，向他本人咏唱圣咏，任何人听了都会战栗，并说服邻近地区和城市中阿谀奉承他的主教和长老，在向百姓讲道时宣扬同样的观点。\par

11. 为了预告我们即将记述的一些事，他不肯承认圣子从天降下。这并非只是断言，而是从我们已寄给你们的记录中得到充分证明；尤其在他声称\Quote{\Person{耶稣基督}是从下面来的}这句话中可见一斑。但那些向他咏唱并赞扬他的人，在百姓中间说他们那邪恶的师傅是一个从天降下的天使。他也不禁止这样的事；这傲慢的人甚至在有人这样说时也在场。\par

12. 还有那些被\Place{安提约基雅}人称为\OriginalTerm{同居女}{subintroductæ}的妇女，属于他和与他在一起的长老和执事。虽然他知晓并已定这些人的罪，却对这些人和他们其他不可救药的罪愆视而不见，好使他们与自己捆绑在一起，并因畏惧自身而不敢控告他的恶言恶行。他还使他们致富；因此，那些贪图此类事物的人爱戴并钦佩他。\par

13. 亲爱的诸位，我们知道，主教和全体圣职人员应当成为百姓行一切善工的榜样。我们也并非不知有多少人因他们所引进的这些妇女而跌倒或招致嫌疑。因此，即使我们允许他没有犯任何罪行，他也应当避免由此产生的嫌疑，免得使某人跌倒，或引诱他人效仿他。\par

14. 因为，如果他自己已经打发走了一个，如今身边有两个年轻貌美的女子，无论去哪里都带着她们，同时过着奢侈放纵的生活，他又怎能责备或劝戒别人不要与妇女过于亲近，以免如经上所记载的那样跌倒呢？\par

15. 因这些事，众人都私下哀痛叹息；但他们如此惧怕他的暴虐和权势，竟不敢控告他。\par

16. 然而，正如我们所说，如果此人持守天主教会教理，并与我们共融，人们本可因这种行为追究他的责任；但因为他藐视奥迹，在可憎的\Person{阿尔特玛斯}异端中（因为我们为何不该提及其父呢？）趾高气扬，我们认为无须要求他就这些事作出解释。\par

17. 随后，在书信的结尾处，他们加上了以下这些话：\par

\Quote{因此，我们不得不将他绝罚，因为他敌对\Person{天主}，且不肯服从；并另选一位主教代替他，为\Place{天主教会}服务。我们相信，是遵照\Person{天主}的指引，我们选定了\Person{多姆努斯}，他具备主教应有的一切美德，是已故\Person{德默特里亚努斯}之子，后者曾卓越地管理过同一教区。我们已将此事通知你们，好使你们写信给他，并从他那领受共融信函。至于此人，让他写信给\Person{阿尔特玛斯}罢；凡与\Person{阿尔特玛斯}想法相同的人，愿他们与他共融。}\par

18. 既然\Person{保禄}已从主教职和正统信仰上堕落了，\Person{多姆努斯}如前所述，便成了\Place{安提约基雅}教会的主教。\par

19. 但当\Person{保禄}拒绝交出教堂建筑时，人们便向皇帝\Person{奥勒良}呈请；皇帝极其公正地裁决了此事，命令将教堂交给\Place{意大利}和\Place{罗马城}的主教们所判定的人。于是，此人以极大的耻辱被世俗权力逐出了教会。\par

20. 那时，\Person{奥勒良}就是这样对待我们的；但在他统治期间，他改变了对我们的看法，受某些谋士的鼓动，开始对我们发动迫害。对此，各处都有许多议论。\par

21. 但当他正要付诸实施，可以说是正在签署迫害我们的法令时，天主的审判临到他，在他着手行事的那一刻制止了他，以一种人人都能清楚看见的方式表明：除非那护卫的手，为了管教和纠正，在祂认为合适的时候，按神圣的天上审判允许，否则这世界的统治者永远不能找到机会攻击\Person{基督}的教会。\par

22. 统治六年之后，\Person{奥勒良}为\Person{普罗布斯}所继任。他统治了相同的年数，随后是\Person{卡鲁斯}及其子\Person{卡里努斯}和\Person{努梅里亚努斯}。他们统治不足三年后，政权归于\Person{戴克里先}及其同僚。在他们统治期间，发生了我们时代的迫害，以及随之而来的教堂被毁。\par

23. 在此之前不久，\Place{罗马}的主教\Person{狄奥尼修斯}在任职九年之后去世，\Person{斐理克斯}继任。\par

% structure: p0220 source=h2 target=subsection reason=relative-heading-rank; base=h2

\subsection{第三十一章 此时开始的摩尼教谬误的异端。}

1. 此时，那因其魔鬼般的异端而得名之狂人，已以其理智之扭曲武装自己，如同魔鬼撒殚——那亲自与天主对抗者——推动他走向许多人的毁灭。他一生如同蛮夷，言语和行事皆然；其本性则如魔鬼且疯狂。因此他试图以基督自居，并被自己的疯狂所膨胀，宣称自己就是护慰者和圣神本身；之后，如同基督一样，他选了十二个门徒作为其新学说的同伙。\par

2. 他汇集了许多早已绝迹的邪恶中假神无神的道理，如同致命的毒药，从波斯传播到我们这一方世界。由此，\OriginalTerm{摩尼教人}{Manicheans}这一不敬的名号至今仍在许多人中盛行。这就是那\Quote{假冒为善的知识}的基础，它在那个时代兴起。\par

% structure: p0223 source=h2 target=subsection reason=relative-heading-rank; base=h2

\subsection{第三十二章 当代杰出的教会人物，以及其中哪些人活到教会遭受毁灭之时}

1. 此时，斐理克斯治理罗马教会五年之后，由欧提齐亚努斯接任；但他在不到十个月后，便将席位让给加约，加约活在我们的时代。他担任此职约十五年，随后由玛策利诺斯接替，玛策利诺斯在迫害中被抓获。\par

2. 大约同时，蒂迈欧在多米诺之后接受了安提约基雅的主教职务，接替他的是济利禄，济利禄活在我们的时代。在他的时代，我们认识了多罗塞奥，他当时是博学之人，在安提约基雅被授予司铎之职。他热爱神圣事物中的美，并致力于希伯来语言，能流利阅读希伯来文圣经。\par

3. 他属于那些特别博学之人，并且对希腊预备学科并不陌生。此外，他生来就是阉人。因此，仿佛奇迹一般，皇帝将他纳入家族，并委托他管理\Place{提洛}的紫色染料工坊，以示尊荣。我们曾听他在教会中智慧地讲解圣经。\par

4. 济利禄之后，提拉诺斯接受了安提约基雅教区的主教职务。在他的时代，教堂遭到了毁灭。\par

5. 欧瑟伯，来自亚历山大里亚，在苏格拉底之后治理了劳狄刻雅教区。他迁往那里的原因与保禄的事件有关。他为此前往叙利亚，并因当地热心神圣事物的人而无法返回家乡。在我们同时代的人中，他是宗教的美好榜样，这从我们所引用狄约尼削的话中容易看出。\par

6. 接替他的是阿纳托利乌斯，正如人们所说，一位好人接续另一位好人。他也是亚历山大里亚人。在学识和希腊哲学技能方面，如算术、几何、天文学和一般辩证法，以及在物理学理论方面，他在我们时代最有能力的人中首屈一指，并且在修辞学方面也名列前茅。据说，为此亚历山大里亚市民请求他在那里建立一所亚里士多德哲学学校。\par

7. 他们讲述他在亚历山大里亚的皮鲁奇乌姆围城期间许多其他杰出的行为，为此他特别受到所有高级官员的尊敬；但我只举以下例子。\par

8. 据说，面包已经耗尽，被围困者难以抵挡饥荒更甚于外敌；但他到场，以这种方式为他们提供食物。由于城市的另一部分与罗马军队结盟，因此并未被围，阿纳托利乌斯派人去找欧瑟伯——那时他还在去叙利亚之前，仍留在未被围困者中，并且享有盛名，其名声甚至传到了罗马将军那里——向他告知了围城中因饥荒而濒死的人。\par

9. 当欧瑟伯得知此事，他请求罗马统帅赐予最大的恩惠，即准许从敌人那边投奔过来的人安全。获得许可后，他告知了阿纳托利乌斯。阿纳托利乌斯一接到消息，便召集了亚历山大里亚的元老院，起初提议所有人都应与罗马人和解。但当看到他因这个建议而恼怒时，他说：\Quote{但我认为你们不会反对我，如果我建议你们将多余的人、对我们毫无用处的人，如老人、妇女和孩童，送出城门，任凭他们去往所愿之处。因为何必徒然保留这些无论如何都快死的人呢？何必用饥饿摧毁那些残废和肢体不全的人呢？我们只应供养男人和青年，将必要的面包分给那些需要守城的人。}\par

10. 他用这样的论据说服了集会，并率先起身投票，主张所有对军队无用的人，无论男女，都应离开城市，因为如果留下并毫无必要地继续留在城中，他们就没有得救的希望，而会死于饥饿。\par

11. 由于元老院所有其他人都同意，他几乎拯救了所有被围困的人。他首先安排属于教会的人，然后是城中其他人，无论年龄大小，都得以逃脱，不仅包括法令所列的类别，还包括许多其他人，偷偷穿上妇女的衣服；通过他的安排，他们在夜间出了城门，逃到罗马军营。在那里，欧瑟伯如同父亲和医生，接待了所有因长期围困而虚弱的人，并以各种谨慎和关怀使他们恢复。\par

12. 劳狄刻雅教会接连享有如此两位牧者，他们在天主的眷顾下，在上述战争之后从亚历山大里亚来到该城。\par

13. 阿纳托利乌斯著作不多；但在流传于我们手中的那些著作中，我们可以看出他的口才和学识。在这些著作中，他特别陈述了他对逾越节的意见。以下摘录似乎重要，特此列出。\par

14. \Emph{从\WorkTitle{阿纳托利的逾越节典制}}有一年第一个月的月朔，即每一个十九年周期的开始，是在埃及月份\OriginalTerm{法墨诺特}{Phamenoth}的第二十六日；但按马其顿人的月份，则是\OriginalTerm{狄斯特罗斯}{Dystrus}的第二十二日，或如罗马人所说，四月\OriginalTerm{朔日}{Kalends}前第十一日。\par

15. 在上述\OriginalTerm{法墨诺特}{Phamenoth}第二十六日，太阳不仅被发现已进入第一个宫段，而且已在该宫段中过了第四天。他们惯常称此宫段为第一个\OriginalTerm{十二宫区}{dodecatomorion}、\OriginalTerm{春分}{equinox}、月份的起始、周期的首端，以及行星运行的起点。但他们称在此之前的那个为最后之月、第十二宫段、最后一个\OriginalTerm{十二宫区}{dodecatomorion}，以及行星运行的终点。因此我们坚持认为，那些将第一个月置于其中，并以此确定\OriginalTerm{逾越节}{passover}第十四日的人，所犯的不是轻微或寻常的错误。\par

16. 这并非我们自己的见解；而是古时的犹太人，甚至在基督之前，就已知晓，并为他们所谨慎遵守。这可以从\Person{斐罗}、\Person{约瑟夫}和\Person{穆赛乌斯}所说的话中得知；不仅从他们，也从更早的两位被称为\Quote{大师}的\Person{阿伽托布利}，以及著名的\Person{阿里斯托布鲁}那里得知，他被\Person{斐拉德尔斐斯的托勒密}和他的父亲选为神圣\OriginalTerm{希伯来圣经}{Hebrew Scriptures}的\OriginalTerm{七十子}{seventy interpreters}之一，并把他关于梅瑟律法的\WorkTitle{释义}书籍献给同一些国王。\par

17. 这些作者在解释关于\OriginalTerm{出谷记}{Exodus}的问题时说，所有人都应当在春分之后，第一个月的中间，献上\OriginalTerm{逾越节祭献}{passover offerings}。而这事发生在太阳经过太阳的（或如他们中某些人所称的\OriginalTerm{黄道圈}{zodiacal circle}）第一个宫段之时。\Person{阿里斯托布鲁}补充说，为了逾越节，不仅太阳必须经过\OriginalTerm{春分段}{equinoctial segment}，月亮也必须如此。\par

18. 因为有两个\OriginalTerm{昼夜平分线段}{equinoctial segments}，春分和秋分，彼此直接相对，而逾越节的日子被定为该月的第十四日，从傍晚开始，月亮将处于与太阳直接相对的位置，正如在\OriginalTerm{满月}{full moons}中所见；太阳将在\OriginalTerm{春分}{vernal equinox}的宫段，而月亮则必然在\OriginalTerm{秋分}{autumnal equinox}的宫段。\par

19. 我知道他们说了许多其他事情，有些似乎有道理，有些近乎绝对的证明，他们借此试图证明，必须在春分之后守\OriginalTerm{逾越节}{passover}和\OriginalTerm{无酵节}{feast of unleavened bread}。但我避免对此类事情要求这种证明，因为\OriginalTerm{梅瑟法律}{Mosaic law}的面纱已经除去，以至于现在我们终于得以揭开脸面，如同在镜中不断瞻仰\Person{基督}以及基督的教导和苦难。\BibleRef{格后 3:18}但希伯来人的第一个月临近春分，\WorkTitle{以诺书}的教导也显示出来。\par

20. 同一位作者还留下了十卷的\WorkTitle{算术原理}，以及其他证明他在神圣事物上的经验和造诣的著作。\par

21. \Person{忒奥提克诺斯}，\Place{巴勒斯坦}地\Place{凯撒勒雅}的主教，首先祝圣他为主教，意图使他死后在自己的教区接替他。有一小段时间，他们两人共同管理同一教会。但为审理\Person{保禄}案件而召开的会议将他召到\Place{安提约基雅}，当他经过\Place{劳狄刻雅}城时，\Person{欧瑟比}已去世，他被那里的弟兄们留住。\par

22. 在\Person{阿纳托利}去世之后，迫害之前该教区的最后一位主教是\Person{斯德望}，他以其哲学和其他希腊学识的知识而受到许多人的钦佩。但他对神圣信德并不同样热衷，正如迫害的进程所表明的；因为它显示他是一个怯懦而不勇敢的伪善者，而不是一个真正的哲学家。\par

23. 但这并未严重损害教会，因为\Person{忒奥多图斯}立即整顿了他们的局面，由天主、万物的\OriginalTerm{救主}{Saviour}亲自立为该教区的主教。他以行为证明了他那与主相关的名字和主教职分。因为他在身体的医术和灵魂的医治艺术上都很卓越。在仁爱、真诚、同情和帮助需要他援助之人的热忱方面，没有其他人能与他相比。他也非常致力于神圣的学问。他就是这样的人。\par

24. 在\Place{巴勒斯坦}的\Place{凯撒勒雅}，\Person{阿伽皮乌}接替了\Person{忒奥提克诺斯}，后者极其热忱地履行了其主教职务。我们也知道阿伽皮乌勤奋工作，在他对人民的监督中表现出最真挚的天主眷顾，特别以慷慨的手照顾所有穷人。\par

25. 在他的时代，我们认识了\Person{庞费禄}，这位最雄辩的人，真正哲学性的生活，被认为配得该教区\OriginalTerm{司铎}{presbyter}的职分。要展示他是怎样的人以及他从何而来，并非小事。但我们在关于他的专题著作中描述了他生活的一切细节，他所建立的学校，他在\OriginalTerm{迫害}{persecution}中多次\OriginalTerm{宣认}{confessions}所受的试炼，以及他最终获得荣誉的\OriginalTerm{殉道}{martyrdom}冠冕。但在那里所有的人中，他的确是最可钦佩的。\par

26. 在我们最近的时代的人中，我们认识了\Place{亚历山大}的司铎中的\Person{皮里乌斯}，以及\Place{本都}各教会的主教\Person{梅利提乌}——稀有之人。\par

27. 前者以其极端贫困的生活和哲学学识而著称，在沉思和解释神圣事物以及教会公开讲道方面极为勤奋。\Person{梅利提乌}，学者们称他为\Quote{阿提卡的蜜糖}，是每一个人都会描述为在所有学问上最有成就的人；要足够钦佩他的修辞技巧是不可能的。可以说他在这方面是天赋的；但在其他方面，谁又能超越他伟大经验与博学的卓越呢？\par

28. 因为在所有知识分支中，你若曾试他一试，就会说他是最有技巧和最有学问的。此外，他生活的德行也同样引人注目。我们在\OriginalTerm{迫害}{persecution}时期仔细观察过他，当时他整整七年逃避迫害的狂怒，逃亡在\Place{巴勒斯坦}地区。\par

29. 继我们刚才提到的主教希默纳乌斯之后，\Person{Zambdas}接受了耶路撒冷教会的主教职。他不久即去世，\Person{Hermon}——这时代迫害之前的最后一位主教——继承了那一直保存至今的宗座。\par

30. 在亚历山大，继狄约尼削去世后担任了十八年主教的玛克西穆，由\Person{Theonas}继任。在他任职期间，与彼埃里乌同时被任命为亚历山大司铎的\Person{Achillas}声名显赫。他被委以神圣信仰的学府，展现出极为罕见、无与伦比的哲学果实，以及真正符合福音的行为。\par

31. 特奥纳斯任职十九年后，\Person{Peter}接受了亚历山大主教职，在整十二年间甚为卓越。其中，他在迫害之前治理教会不到三年，此后余生则恪守更严格的纪律，并不隐晦地关怀教会公共福祉。因此，他在迫害第九年被斩首，荣膺殉道之冠。\par

32. 我在这些书中，从救主的诞生直到敬拜场所被毁——三百零五年的时期内——写下了传承的记述，现在请允许我转向我们这时代为信仰英勇作战者的斗争，并将这些冲突的规模与重大性笔之于书，留待后人知晓。\par

\SourceRecord{Translated by Arthur Cushman McGiffert.；Nicene and Post-Nicene Fathers, Second Series；Vol. 1.；Edited by Philip Schaff and Henry Wace.；Buffalo, NY: Christian Literature Publishing Co.,；1890.；<http://www.newadvent.org/fathers/250107.htm>.}

% structure: p0001 source=h1 target=section reason=page-title

\section{\WorkTitle{教会史（卷八）}}

% structure: p0002 source=h2 target=subsection reason=relative-heading-rank; base=h2

\subsection{引言。}

我们已在七卷书中描述了自宗徒时代以来的事件，故在此第八卷中，我们认为适当记录我们时代的一些最重要事件，以供后代知晓，这些事件值得永久记载。我们的叙述将由此开始。\par

% structure: p0004 source=h2 target=subsection reason=relative-heading-rank; base=h2

\subsection{第一章 在我们这个时代迫害之前的事件。}

1. 在我们时代的迫害之前，对宇宙天主的虔敬之道——通过基督向世界宣扬——在所有民族中，无论是希腊人还是蛮族，所受的荣耀与自由的广度与性质，非我们所能恰当地描述。\par

2. 统治者对我们子民的恩宠可作为证据：因为他们将省份的治理托付给他们，并因对其教义怀有极大友谊，免去了他们对献祭的忧虑。\par

3. 何需提及那些在皇宫中的人，以及所有统治者，他们允许自己的家眷、妻子、儿女和仆人在他们面前公开谈论神圣的圣言与生命，甚至容忍他们几乎夸耀其信仰的自由？\par

4. 事实上，他们非常看重他们，视他们优于同僚。例如那位\Person{多罗特乌斯}，对他们最为忠诚和忠实，因此在他们当中，尤其被那些拥有最尊贵职位和治理权的人所尊敬。与他一起的还有著名的\Person{哥尔哥尼乌斯}，以及所有因天主的圣言而被视为配得上同样殊荣的人。\par

5. 人可以看到，每座教会的统治者都受到所有官员和总督的极大青睐。但谁能描述那些巨大的集会，以及每个城市中聚集的群众，还有祈祷之所中著名的聚会？他们不满足于古老的建筑，因此从根基起在所有城市中建造了大教堂。\par

6. 没有嫉妒阻碍这些事务的进展，它们逐渐前进，日益增长和扩大。也没有任何恶鬼能通过人类计谋诽谤或阻碍它们，只要神圣而天上的手看顾并护卫自己的子民，视他们为堪当的。\par

7. 然而，由于过度的自由，我们陷入了松懈和怠惰，互相嫉妒和辱骂，几乎如同拿起武器彼此攻击：统治者如投枪般的言语相互攻击，民众结党反对民众，骇人的虚伪和掩饰升至邪恶的顶峰。神圣的审判带着宽容，照其所悦，在群众仍然聚集之时，温和而适度地困扰了主教职。\par

8. 这场迫害始于军队中的弟兄。但我们如同没有感觉一般，并不急于使神祇悦纳和平息；有些人如同无神论者，以为我们的事务无人关注、无人治理；于是我们罪上加罪。而那些被视为我们牧者的人，抛弃了虔敬的束缚，激动地彼此冲突，除了堆积争吵、威胁、嫉妒、敌意和仇恨外，无所作为，如同暴君急切地想维护自己的权力。那时，真如耶肋米亚的话：\Quote{主在愤怒中使熙雍女子昏暗，将以色列的光荣从天上抛到地上，在发怒之日不记念他的脚凳。主也倾覆了以色列一切华美之物，推倒了他的一切堡垒。} \BibleRef{哀 2:1}\par

9. 又正如圣咏中所预言的：\Quote{他废弃了他仆人的盟约，将他的圣地亵渎至地——在教堂的毁灭中——推倒了他的一切堡垒，使他的要塞成为懦弱。凡路过的人都掠夺了这百姓的众多财物；他又成为邻人的羞辱。因为他高举了他敌人的右手，使他刀剑的援助退缩，在战争中未支持他。但他剥夺了他的净化，将他的宝座抛到地上。他缩短了他日子的时间，此外，还将羞辱倾注在他身上。}\par

% structure: p0014 source=h2 target=subsection reason=relative-heading-rank; base=h2

\subsection{第二章 教堂的毁灭。}

1. 这一切都在我们身上应验了，因为我们亲眼目睹祈祷之所被彻底拆毁，神圣的圣经在市场中付诸烈焰，教会的牧者各处卑怯地躲藏，其中一些被羞辱地捕获，受敌人嘲弄。那时，又正如另一个先知的话：\Quote{藐视倾倒在统治者身上，他使他们走在无人踏足、无路可走的荒野。}\par

2. 但描述最终降临于他们的悲惨不幸并非我们的职责，而且我们认为记录他们在迫害之前的彼此分裂和反自然行为也不合适。因此，我们决定除了那些可以辩护神圣审判的事情外，不叙述任何关于他们的事。\par

3. 因此，我们不提及那些被迫害动摇的人，也不提那些在一切关乎救恩的事上遭遇船难、因自己的意志沉入深渊的人。我们只在此历史中一般地介绍那些首先对我们自己、其次对后代有益的事件。让我们接着简要描述神圣圣言的见证人的神圣争战。\par

4. 在\Person{戴克里先}统治的第十九年，Dystrus月——罗马人称为三月——救主受难节将近之时，皇家敕令各处颁布，命令将教堂夷为平地，圣经被火销毁，并命令那些拥有荣誉职位的人被降级，家仆若坚持宣认基督信仰，则剥夺自由。\par

5. 这是针对我们的第一道敕令。但不久之后，又颁布了其他法令，命令先将各地教会的统治者投入监狱，然后用尽一切伎俩强迫他们献祭。\par

% structure: p0020 source=h2 target=subsection reason=relative-heading-rank; base=h2

\subsection{第三章 在迫害中所忍受的冲突性质}

1. 那时，确实有许多教会的领袖热切地忍受了可怕的苦难，并提供了崇高斗争的榜样。但另有一大群人，精神因恐惧而麻木，在最初的冲击下轻易地软弱了。其余的人各自承受了不同形式的折磨。一人的身体被棍棒鞭打。另一人受到难以忍受的折磨和刮削，其中有些人悲惨地死去。\par

2. 另一些人经历了不同的斗争。这样，有一个人，当周围人用强力逼迫他、拖他去行那可憎不洁的祭献时，他却像已献了祭一样被释放，尽管他并没有献祭。另一个人，虽然根本未曾靠近，也没有触碰任何污秽之物，但当别人说他已献祭时，他默然承受了控告，离开了。\par

3. 另一个人被半死不活地抓起，像已死的人一样被抛弃；又有一个躺在地上，被拖着双脚走了很长一段路，被算在已献祭的人中。一个人大声喊叫，以响亮的声音声明他拒绝献祭；另一个人高喊他是基督徒，因着对那救恩之名的宣认而光彩照人。另一个人声明他从未献祭，也永远不会献祭。\par

4. 但他们被一大队为此而列阵的士兵击打嘴巴，并被禁声；他们被打脸和面颊，被强行赶走；不虔敬的仇敌认为，无论如何，使他们的目的看似达成，是何等重要。但这些事对圣殉道者毫无用处；对于他们，我们有什么言语足以精确描述呢？\par

% structure: p0025 source=h2 target=subsection reason=relative-heading-rank; base=h2

\subsection{第四章 天主的著名殉道者，他们以记忆充满各处，并为宗教赢得各种冠冕}

1. 因为我们可以讲述许多人，他们不仅从大迫害开始之时，而且远在那之前，当和平尚存之际，就为宇宙天主的宗教显示了令人钦佩的热忱。\par

2. 因为那获得权力者，虽然看似从沉睡中醒来，但从\Person{德西乌斯}和\Person{瓦勒里安}以后，他一直暗中且不被察觉地密谋反对教会。他并非立刻与我们全体作战，而是首先只对军队中的人进行试探。因为他以为，如果他先攻击并制服这些人，其他人就容易被拿下。于是，许多士兵极其乐意地选择了退隐生活，以免否认他们对宇宙创造者的虔诚。\par

3. 因为当那首领，无论他是谁，开始迫害士兵时，将士兵分成部落，清洗那些在军中登记的人，给他们选择：要么服从以得到属于他们的荣誉，要么因不服从命令而被剥夺荣誉。基督王国的大量士兵毫不迟疑，立刻选择了对祂的宣认，胜过他们正享有的表面荣耀和繁荣。\par

4. 他们中偶尔有人，因其虔诚的坚忍，不仅失去了职位，而且失去了生命。但至今，这个阴谋的煽动者还以温和的方式行事，仅在某些情况下才敢流到血；因为信徒的众多，可能使他害怕，并阻止他立刻向所有人开战。\par

5. 但当他更大胆地发动攻击时，在所有城市和乡村的居民中，可以看到多少、何种的天主殉道者，那就无法述说了。\par

% structure: p0031 source=h2 target=subsection reason=relative-heading-rank; base=h2

\subsection{第五章 在\Place{尼科美底亚}的人}

1. 在\Place{尼科美底亚}，针对教会的敕令刚一公布，有一个人，并非无名之辈，而是因着显赫的世俗尊位而备受尊敬，出于对天主的热情，被炽热的信德所激励，趁敕令公开张贴之际，把它撕成碎片，视之为亵渎不敬之物；此事发生时，当时两位君主正同在城中——最年长的那位，以及在他之后掌第四位权柄的那位。\par

2. 但这个人，在那一地首先如此显扬自己之后，承受了那种勇敢所必然带来的事，并保持心境愉快、镇定，直至死亡。\par

% structure: p0034 source=h2 target=subsection reason=relative-heading-rank; base=h2

\subsection{第六章 在宫廷中的人}

1. 这个时期产生了神圣而光荣的殉道者，尤其是那些历来被颂扬、被赞为勇敢的人，无论在希腊人或野蛮人中，以\Person{多罗特乌斯}和在宫廷中与他在一起的仆人们为代表。虽然他们从主人那里得到最高的荣誉，并被主人视为自己的儿女，但他们将为了宗教而受的辱骂、考验，以及被发明出来的各种死亡形式，看作比今生的荣耀和奢华更大的财富。\par

2. 我们将描述其中一人结束生命的方式，让读者从他的情况推断其他人的苦难。有一个人，在上述城市中，被带到我们所说的统治者面前。他奉命献祭，但他拒绝了，于是被命令剥去衣服，高高吊起，用棍棒鞭打全身，直到他被征服，即使违背自己的意愿，也去做所命令的事。\par

3. 但他不为这些痛苦所动，他的骨头已经露出，他们将醋和盐混合，浇在他被撕裂的身体部位上。当他蔑视这些剧痛时，一个铁架和火被拿来。他身上剩余的肉，像准备食用的肉一样，被放在火上，不是一下子，以免他立刻断气，而是一点一点地。那些把他放在火刑堆上的人，直到他在这样的痛苦后同意所命令的事，才被允许停止。\par

4. 但他坚守他的决心，在酷刑仍在进行时，胜利地交出了自己的生命。宫廷中一位仆人的殉道就是这样，他确实无愧于他的名字，因为他被称为\Person{伯多禄}。\par

5. 其余者的殉道，虽不亚于他，但为简略起见，我们略过不予记述，只记\Person{多洛泰乌斯}和\Person{戈尔贡尼乌斯}，以及许多其他皇室人员，在经历各种苦刑后，被绞杀而终，并赢得了天主所赐的胜利之荣冠。\par

6. 此时，\Person{安提默}——当时主持\Place{尼苛默底亚}教会——因给基督作证而被斩首。一大群殉道者加入了他，因为在那些日子里，\Place{尼苛默底亚}的皇宫不知何故发生火灾，这被错误地怀疑到我们的人身上。该地所有虔诚的家族都奉皇命集体处死，有的死于刀下，有的死于火中。据传，男女老少怀着一种神圣而难以形容的热忱冲入火中。刽子手又将大批其他人捆绑起来，放到船上，投入深海。\par

7. 而那些曾被尊为他们的主人的人，认为必须将已按适当葬礼埋葬于地的皇室仆人的尸体挖出，投入海中，以免有人——如他们所想——视他们为神明，敬拜躺在坟墓中的他们。\par

8. 这些事发生在迫害初期于\Place{尼苛默底亚}。但不久之后，由于在称为\Place{默利泰内}的乡间以及\Place{叙利亚}全境有人企图篡夺政权，一道皇室敕令发出，命各地教会首领被投入监狱并受捆绑。\par

9. 此后所见之事，笔墨难以尽述。各处有大批人被囚禁；各地监狱——早先为杀人犯和盗墓者预备——如今都充满了主教、司铎、执事、读经员和驱魔员，以致再无空间容纳那些因罪行被定罪的人。\par

10. 当第一道敕令之后又跟随后续法令，命令狱中之人若肯献祭便可获释，而拒绝者则受许多酷刑折磨时，又有谁能再次数清各省殉道者的众多数目？尤其是在\Place{阿非利加}、\Place{毛里塔尼亚}、\Place{忒拜斯}和\Place{埃及}的殉道者？从这最后一块地方，许多人前往其他城市和省份，因殉道而成名。\par

% structure: p0045 source=h2 target=subsection reason=relative-heading-rank; base=h2

\subsection{七、在\Place{腓尼基}的埃及人}

1. 我们知晓其中在\Place{巴勒斯坦}显赫者，以及在\Place{提洛}者。看见他们的人，谁不惊讶于无数的鞭打，以及这些真正奇妙的信仰斗士在鞭打下的坚韧？又惊讶于他们在鞭打后立即与嗜血的野兽搏斗——被丢到豹子、各种熊、野猪和受火与炽铁刺激的公牛面前？又惊讶于这些高贵的人在面对各种野兽时所表现的惊人耐力？\par

2. 我们本人亲临这些事件发生之时，并记录了我们的殉道救主耶稣基督的神圣权能，这权能临在于殉道者身上并有力地彰显出来。很长一段时间，那些食人的野兽不敢触碰或靠近这些天主所爱之人的身体，却冲向那些从外面刺激、驱赶它们的人。它们丝毫不敢碰这些圣洁的运动员，尽管他们独自站着，赤身露体，向野兽挥手以吸引它们——因为他们奉命这样做。但每当野兽冲向他们时，它们就像被某种神圣力量抑制而退却。\par

3. 这种情况持续了很久，令旁观者惊奇不已。当第一只野兽无所作为时，第二只和第三只又被放出来对付同一个殉道者。\par

4. 人们不能不惊叹这些圣者不可战胜的坚韧，以及那些身体年轻者持久不变的恒心。你能看见一个不到二十岁的青年，未被捆绑，双手伸出呈十字形，怀着无畏不颤的心智，热切地向天主祈祷，丝毫没有从他站立之处后退或退缩，而熊和豹子呼吸着愤怒与死亡，几乎触到他的皮肉。然而，它们的口被某种神圣而不可理解的力量抑制——不知何故——又跑回原处。此人正是如此。\par

5. 你又能看见其他的人——因为他们共五人——被丢到一头野公牛面前。这公牛用角将那些从外面靠近的人抛向空中，将他们撕裂，任其半死地被抬走；但当它愤怒而威胁地冲向独自站着的圣洁殉道者时，却不能靠近他们；尽管它跺脚，到处用角顶撞，因烧红的铁器刺激而呼吸愤怒与威胁，却仍被神圣的上智所阻止。因它丝毫未能伤害他们，人们便放出其他野兽对付他们。\par

6. 最后，在经历这些可怕而多样的攻击之后，他们全被刀剑所杀；而他们未被葬于土中，而是被抛入海浪之中。\par

% structure: p0052 source=h2 target=subsection reason=relative-heading-rank; base=h2

\subsection{八、在\Place{埃及}的人}

1. 那些在\Place{提洛}为信仰而英勇奋战的埃及人所经历的争战便是如此。但我们也必须钦佩那些在本土殉道的人；在那里，成千上万的男女和儿童，为了救主的教导而轻视今生，忍受了各种死亡。\par

2. 他们中有些人，在经受刮削、拉肢、最严厉的鞭打以及无数其他酷刑——甚至听来都可怕——之后，被投入火焰；有些人溺死于海中；有些人勇敢地将头伸给斩首者；有些人在酷刑中死去；另一些人则饿死。还有一些人被钉在十字架上；有些按通常对待恶犯的方式；另一些则更残酷地被倒钉在十字架上，头朝下，并一直被留活直到在十字架上饿死。\par

% structure: p0055 source=h2 target=subsection reason=relative-heading-rank; base=h2

\subsection{九、在\Place{忒拜斯}的人}

1. 在底比斯，殉道者所受的凌辱和酷刑难以尽述。他们被用贝壳刮遍全身，直至死亡，这酷刑可代替铁钩。妇女们被绑住一只脚，用机械吊在空中，全身赤裸无蔽，呈现在众人眼前，成了最可耻、残酷、无人性的景象。\par

2. 另一些人被绑在树枝和树干上而死去。因为人们用机械把最粗壮的树枝拉拢，将殉道者的四肢绑在上面，然后让树枝恢复原位，顷刻间就把他们的肢体撕裂了。\par

3. 这一切事并非只进行了数日或短短一段时间，而是持续许多年。有时一天处死十余人，有时二十余人；又有一次不少于三十人，约六十人，甚至一百名男子连同幼童和妇女，被处以各种不同的酷刑。\par

4. 我们本人也身临其境，曾目睹一天之内大批人群被处决：有的被斩首，有的受火刑；以致杀人的刀都钝了，因乏力而折断，连刽子手也疲惫不堪，彼此替换。\par

5. 我们看见了信从天主基督者所表现的最令人惊叹的热忱和真正神圣的精力与热诚。因为第一个人被判刑之后，接二连三有人冲向审判台，承认自己是基督徒。他们对可怕的事和种种酷刑视若无睹，勇敢无畏地宣认了宇宙天主的信仰。他们喜乐、欢笑、愉快地接受了最终的死刑判决，歌唱着、向上主歌颂赞美，直至最后一息。\par

6. 这些固然令人惊叹，但更令人惊叹的，是那些因财富、高贵出身、荣誉、学识和哲学而卓著的人，将这一切置于次要地位，而把真宗教和对我们救主耶稣基督的信仰放在首位。\par

7. 这样的人中有\Person{斐洛罗穆斯}，他在\Place{亚历山大}的帝国政府中担任高职，每日执行审判，并有与他身份和罗马尊严相称的卫兵随行。同样还有\Person{斐肋亚斯}，他是\Place{特穆伊斯}教会的\OriginalTerm{主教}{Episcopus}，因其爱国心和对祖国的服务，以及其哲学学识而著称。\par

8. 这些人虽然有许多亲戚和其他朋友恳求他们，甚至有身居高位的人和审判官本人恳求他们怜悯自己、顾惜妻儿，但他们丝毫不被这些事所动，以致选择贪生怕死、藐视我们救主关于宣认与否认的规诫。相反，他们以刚毅和哲人的心智，更确切地说，以虔诚和爱天主的心灵，坚持面对审判官的一切威胁和辱骂；于是二人都被斩首。\par

% structure: p0064 source=h2 target=subsection reason=relative-heading-rank; base=h2

\subsection{第10章　殉道者斐肋亚斯记述亚历山大事件的著作}

1. 既然我们提到\Person{斐肋亚斯}在世俗学问上享有盛誉，那就让他在以下摘录中作自己的见证吧。在这段文字中，他向我们表明了他是什么人，同时比我们更准确地描述了当时发生在\Place{亚历山大}的殉道事迹：\par

2. 摆在我们面前的，是神圣圣经中赐予我们的这一切榜样、模式和崇高标记，与我们一起的真福殉道者毫不迟疑，以真诚的灵魂之眼注视那统御万有的天主，心志已定，为宗教而死，坚守了自己的召叫。因为他们知道，我们的主耶稣基督为我们而成了人，为要除尽一切罪恶，并为我们提供进入永生的途径。因为祂\BibleQuote{不以自己与天主同等为应当把持不舍的，反而空虚自己，取了奴仆的形体，与人相似，形状也一见如人；祂贬抑自己，听命至死，且死在十字架上。}\BibleRef{斐 2:6}\par

3. 因此，这些怀有基督的殉道者切慕更大的恩赐，忍受了各种考验和种种折磨人的酷刑；不但一次，甚至有人二次受刑。尽管守卫们以各种方式，不仅用言语，而且用行动竞相恐吓他们，但他们并未放弃决心，因为\BibleQuote{完美的爱德驱除恐惧}\BibleRef{若一 4:18}。\par

4. 谁能以言语描述他们面对每种酷刑时的勇气和刚毅呢？因为凡愿意虐待他们的，都得到了许可：有人用棍棒，有人用木杖，有人用鞭子，还有人用皮带，有人用绳索。\par

5. 凌辱的景象千奇百怪，显示出极大的恶毒。有些人双手反绑，被吊在木架上，每个肢体都被某些机械拉长。然后行刑者奉命用刑具撕裂他们的全身；不仅像对待杀人犯那样撕裂肋部，还撕裂腹部、膝盖和面颊。另一些人被吊在门廊下，只用一只手悬挂，忍受着最剧烈的痛苦，因为关节和四肢被拉伸。又有人面对面被绑在柱子上，脚不着地，身体的重量全压在被捆绑的地方，将其拉得紧紧的。\par

6. 他们忍受这些，并非仅限总督与他们谈话或闲暇之时，而是几乎一整天。因为当他转向别人时，他留下自己属下的官员监视先前的那些人，观察是否有人因酷刑而屈服。他又下令毫不留情地将他们投入锁链，随后当他们奄奄一息时，把他们扔倒在地拖走。\par

7. 因为他说，他们不应对我们有丝毫怜悯，而应认为我们已不复存在，这样——我们的仇敌除了鞭打之外，又设计了这第二种折磨方式。\par

8. 在此等暴行之后，有些人被置于足枷上，双脚被拉到四个孔洞中，以致他们被迫仰卧在足枷上，因全身布满鞭打新伤而无法支撑自己。另一些人被抛在地上，因累积的酷刑折磨而躺在那里，向旁观者展示出更为可怕的严厉景象，因为他们身上带着各种不同刑罚的印记，这些刑罚是被人发明的。\par

9. 如此继续，有些人在酷刑中死去，以他们的坚毅羞辱了敌手。另一些人半死不活地被关进监狱，在痛苦中挣扎，数日后死去；其余的人则因所受照料而康复，随着时间的推移和在监狱中的长期拘禁而重获信心。\par

10. 因此，当他们奉命选择：要么触摸那污秽的祭品，从而免于骚扰，从他们那里获得可诅咒的自由；要么拒绝献祭而被判处死刑；他们毫不迟疑，欣然赴死。因为他们知道圣经早已宣告的话。经上记载：\BibleQuote{凡向其他神明献祭的，必被彻底消灭。}\BibleRef{出 22:20} 以及 \BibleQuote{你不可有别的神在我面前。}\BibleRef{出 20:3}\par

11. 这就是那位真正热爱哲学且热爱天主的殉道者的言辞，在最终判决之前，他尚在狱中时，以此致信他教区的弟兄们，向他们说明自己的处境，同时劝勉他们在他即将去世之后，仍要持守基督的宗教。\par

12. 但我们何必细述这些事，不断增添神圣殉道者在普世战斗的新事例呢？尤其是他们已不再按普通法律受审，而是如战争中的敌人般被攻击。\par

% structure: p0077 source=h2 target=subsection reason=relative-heading-rank; base=h2

\subsection{第十一章 在弗里吉亚的人}

1. 弗里吉亚有一个完全由基督徒居住的小镇，当镇上的男子都在镇内时，士兵们将其完全包围。他们放火焚烧，将男子、妇女和儿童一起烧死，而此时他们正呼求基督。他们这样做，是因为该城的所有居民，连同市长、总督以及所有官员和全体民众，都承认自己是基督徒，丝毫不服从那些命令他们崇拜偶像的人。\par

2. 另有一位名为\Person{阿道克图斯}的罗马显贵，出身于意大利贵族家庭，在皇帝治下历经所有官职，以至他清白地担任过所谓的执政官及财政大臣等高级职务。除此之外，他在虔敬的行为和对天主基督的宣认上出类拔萃，并以殉道者的冠冕为装饰。他仍在担任财政大臣期间，为宗教而忍受了战斗。\par

% structure: p0080 source=h2 target=subsection reason=relative-heading-rank; base=h2

\subsection{第十二章 其他许多人，无论男女，以各种方式受苦}

1. 何必再逐一提名，或列举众多男子的数量，或描绘基督可敬的殉道者所受的各种苦难呢？有些人在阿拉伯被斧头砍死。有些人在卡帕多西亚四肢被打断。有些人在美索不达米亚，脚朝上被高高吊起，头朝下，身下燃着微火，被燃烧木柴升起的烟雾窒息而死。另一些人在亚历山大港被割掉鼻子、耳朵和手，其他肢体和部位也被切断而遭肢解。\par

2. 何必再提及安提约基雅的那些人，他们被放在烤架上烤，不是要杀死他们，而是使他们遭受持久的刑罚？或者另一些人，宁愿将右手伸入火中，也不愿触摸那邪恶的祭品？有些人畏惧考验，宁可从高处跳下，也不愿被抓住落入敌人之手，认为死亡胜过不敬神之人的残酷。\par

3. 有一位圣洁的妇人——心灵因德性而卓越，身体上是一位女子——她在安提约基雅因财富、家族和名声而超群出众，以宗教原则教养了她的两个女儿，她们正值青春年华。由于她们引起极大嫉妒，人们用尽一切手段搜寻她们藏身之处；当得知她们不在时，便用诡计将她们召唤到安提约基雅。就这样，她们落入了士兵的网罗。当这妇人看到自己和女儿们如此无助，并知道人们将对她们做出令人难以启齿的可怕之事——其中最难忍受的，是她们贞洁受侵犯的威胁——她便劝勉自己和两位少女，她们甚至不应听从这种话。她说，将自己的灵魂交给魔鬼奴役，比任何死亡和毁灭都更糟糕；她给她们指出了从这一切中唯一得救的途径——逃到基督那里。\par

4. 她们听从了她的劝告。于是她们整理好衣服，从路中央走开，向卫兵请求一点时间退避，然后投身于旁边流淌的河中。\par

5. 她们就这样自尽了。但在同一安提约基雅城，另有两位少女，她们在一切事上侍奉天主，是真正的姐妹，家世显赫，生活出众，年轻而朝气，心智严肃，举止虔诚，热忱令人钦佩。仿佛大地不能承受如此卓越，崇拜魔鬼的人命令将她们投入海里。她们就这样被处死了。\par

6. 在本都，另一些人遭受了骇人听闻的苦难。他们的手指被尖锐的芦苇刺入指甲下。熔化的铅，因热而沸腾冒泡，被倒在另一些人的背上，他们身体最敏感的部位被灼烤。\par

7. 另一些人的肚腹和私处遭受了可耻、不人道且难以言表的折磨，这些是高贵而守法的审判官为了显示他们的严厉，作为更可敬的智慧表现而发明的。新的酷刑不断被创造出来，仿佛他们试图在竞赛中相互超越以获得奖品。\par

8. 但在这些灾难的末尾，当他们最终无法想出更残酷的手段，且厌倦了杀害，被流血所充满和饱足时，他们转而采取他们认为仁慈和人道的对待，以至于似乎不再对我们策划可怕的事。\par

9. 因为他们说，城市不应被自己百姓的血所污染，统治者的政府——它对所有人都是仁慈而温和的——也不应因过度的残忍而受诽谤；相反，仁慈而王权的恩惠应普及众人，我们不应再被处死。因为施加这种惩罚于我们，应因统治者的仁爱而停止。\par

10. 因此，命令被下达：要挖出我们的眼睛，并使我们的一条肢体残缺。因为在其眼中，这些事是仁爱的，是对我们最轻的刑罚。因此，由于这不虔敬者所给予我们的这种仁慈对待，已无法数算那些右眼先被剑割出、随后被火灼烧，或左足因烧灼关节而致残，之后被判决到省属铜矿——与其说是服役，不如说是受苦和艰辛——的人，其数目难以估量。除此之外，还有其他人遭遇了其他无法详述的考验；因为他们英勇的忍耐超越了任何描述。\par

11. 在这些争战中，基督的崇高殉道者们闪耀于整个世界，处处令目睹他们英勇的人惊奇；我们救主那真正神圣和不可言喻的权能的明证，藉着他们彰显出来。要一一提名将是漫长的工作，若非不可能的话。\par

% structure: p0092 source=h2 target=subsection reason=relative-heading-rank; base=h2

\subsection{第十三章 教会的主教们以他们的血证明他们所宣讲的宗教的真实性}

1. 至于在主要城市中为主殉道的教会领袖，我们在虔诚者的纪念中首先要提及的是\Place{尼科美底亚}城的主教\Person{安提摩斯}，他被斩首。\par

2. 在\Place{安提约基雅}的殉道者中，有该教区的司铎\Person{路济安}，他的一生最为卓越。在\Place{尼科美底亚}，他在皇帝面前宣讲基督的天国，首先以口头的辩护，随后也以行为。\par

3. 在\Place{腓尼基}的殉道者中，最卓越是基督属灵羊群的那些献身牧者：\Place{提洛}教区的主教\Person{提兰尼翁}、\Place{漆冬}教会的一位司铎\Person{则诺比乌斯}，以及\Place{厄默撒}附近各教会的主教\Person{西尔瓦努斯}。\par

4. 最后一位与其他人一起在\Place{厄默撒}被抛给野兽，就这样被接入殉道者的行列。另外两位在\Place{安提约基雅}藉着至死的忍耐荣耀了天主的圣言。主教被投入深海。\Person{则诺比乌斯}是一位技艺高超的医生，因施加于他肋旁的严酷折磨而死去。\par

5. 在\Place{巴勒斯坦}的殉道者中，\Place{加沙}附近各教会的主教\Person{西尔瓦努斯}与另外三十九人在\Place{斐诺}铜矿被斩首。在那里，\Place{埃及}主教\Person{佩琉斯}和\Person{尼鲁斯}与其他人也遭受火刑而死。\par

6. 在这些殉道者中，我们必须提及司铎\Person{庞斐禄}，他是\Place{凯撒勒雅}教区的极大光荣，在我们这时代的人中最令人钦佩。\par

7. 他英勇事迹的德行，我们已在适当之处记载。在\Place{亚历山大里亚}和整个\Place{埃及}及\Place{忒拜斯}亡故的光荣者中，应首先提及\Place{亚历山大里亚}主教\Person{伯多禄}——他是基督宗教最卓越的导师之一；以及与他同在的司铎\Person{福斯图斯}、\Person{狄乌斯}和\Person{阿摩尼乌斯}，基督的完美殉道者；还有\Place{埃及}各教会的主教\Person{斐理亚}、\Person{赫西基乌斯}、\Person{帕基米乌斯}和\Person{德奥多鲁斯}，此外还有许多其他杰出人士，被他们本国和本地区的教区所纪念。\par

要描述那些为神圣宗教而受苦之人的争战、并精确叙述各人所遇的事，并非我们之责。这乃是那些事件的目击者所当做的工。我将把自己亲眼所见的事，在另一部作品中为后世叙述。\par

8. 但在本书中，我将在已提供的之外，添加我们的迫害者所颁布的撤销令，以及迫害开始时发生的事件——这些对阅读它们的人将极有益处。\par

9. 何等的言语足以描述在反对我们的战争之前，罗马政府——当时统治者对我们友好而和睦——其繁荣的伟大与丰盈？那时，那些位居最高统治地位、任职十年或二十年的人，在安然的和平中度过时光，有节庆、公共比赛和最欢快的愉悦与快乐。\par

10. 正当他们的权威不断增长、日益增强时，他们突然改变了对我们的和睦态度，开始了无情的战争。但这运动的第二年尚未过去，整个政府便发生了革命，颠覆了一切。\par

11. 因为我们所提到的那位首领患上重病，心智错乱；而那位享有第二等尊荣的人，与他一起退入私生活。他刚这样做，整个帝国就分裂了——这是一件前所未有的事。\par

12. 不久之后，\Person{君士坦提乌斯}皇帝——他一生对臣民最为仁慈和友好，对神圣圣言也最友善——按自然规律结束了生命，留下自己的儿子\Person{君士坦丁}作为皇帝和奥古斯都继位。他是他们中第一个被列入诸神之列，死后获得人们所能给予皇帝的一切尊荣。\par

13. 他是众皇帝中最仁慈、最温和的一位，也是我们时代中唯一一位在执政期间始终以符合其职位的方式行事的皇帝。此外，他对所有人都极其友善和恩惠。他丝毫没有参与对我们的战争，反而保护了他治下的虔诚者不受伤害和凌辱。他没有拆毁教堂建筑，也没有对我们采取任何其他行动。他的结局是荣耀且三倍有福的。他是唯一一位在去世时将帝国幸福而光荣地留给自己儿子作为继承人的皇帝——其子各方面都最为明智和虔诚。\par

14. 其子\Person{君士坦丁}立即登基执政，被士兵们宣告为最高皇帝和奥古斯都，并且早在许久之前就已由万王之王的天主亲自宣告。他证明自己是他父亲对我们教义之虔诚的效仿者。他就是这样的人。\par

但此后，\Person{李锡尼}由统治者们共同投票被宣告为皇帝和奥古斯都。\par

15. 这些事使\Person{马克西米努斯}极为苦恼，因为直到那时，所有人都只称他为凯撒。他因此极其专横，擅自夺取了尊号，自立为奥古斯都。与此同时，我们之前提到的那位在退位后重获尊号的人，因被发觉密谋反对\Person{君士坦丁}的生命，而以最可耻的方式丧命。他是第一个因其邪恶与不敬而使其法令、雕像和公共纪念碑被销毁的人。\par

% structure: p0110 source=h2 target=subsection reason=relative-heading-rank; base=h2

\subsection{第十四章　宗教敌人的品性}

1. 其子\Person{马克森提乌斯}获得了罗马的统治权，起初为了讨好和奉承罗马人民，假装信奉我们的信仰。为此，他命令臣民停止迫害基督徒，假装虔诚，以便显得比他前任更仁慈和温和。\par

2. 但他的行为并未证明他是所期望的那种人，反而陷入了各种邪恶，不回避任何不洁或放荡，犯下奸淫，沉溺于各种腐败。因为他将妻子从合法的配偶身边夺走，凌辱她们，然后极不光彩地将她们送回丈夫身边。他不仅对无名小卒这样做，而且尤其侮辱罗马元老院中最显赫和最杰出的人物。\par

3. 他所有的臣民，无论民众还是统治者，显赫还是卑微，都因沉重的压迫而疲惫不堪。尽管他们保持沉默，忍受着痛苦的奴役，但暴君的凶残暴行却没有任何缓解。有一次，他因一点小事就让卫兵屠杀民众；大批罗马平民在城中被杀，使用的不是斯基泰人和蛮族的矛和武器，而是他们自己同胞的。\par

4. 不可能数清为了财富而被处死的元老院议员的人数；无数人因各种借口被杀。\par

5. 为使他的一切邪恶达到顶点，暴君诉诸巫术。他在占卜时剖开孕妇，又检查新生婴儿的内脏。他屠杀狮子，并施行各种可憎的仪式以召唤恶魔和避免战争。因为他唯一的希望是借此确保胜利。\par

6. 无法述说这个罗马暴君压迫臣民的方式，以至于他们陷入前所未有生活必需品的极度匮乏，据我们同时代的人说，无论在罗马还是其他地方都未曾有过。\par

7. 但东方的暴君\Person{马克西米努斯}秘密地与罗马暴君结成了友好的联盟，如同与邪恶中的兄弟一样，长期试图掩盖此事。但最终被发觉，遭受了应有的惩罚。\par

8. 他在邪恶上与罗马暴君何等相似，甚至远超于他，这真是奇妙。因为术士和魔法师的首领被他授予最高官阶。他变得极其胆怯和迷信，极为看重偶像和魔鬼的错误。确实，可以这样说，没有占卜者和神谕，他连一个指头都不敢动。\par

9. 因此，他比前任更猛烈和不间断地迫害我们。他下令在每个城市建造庙宇，并迅速修复因时间久远而毁坏的神圣树林。他在每一个地方和城市任命偶像祭司；并在每个省份任命一位特别在各类服务中出类拔萃的行政官员作为大祭司，配以士兵和卫队。他给所有术士，好像他们是虔诚和众神所爱的一样，授予行省和最大的特权。\par

10. 从这时起，他通过极端征收金银财物、最痛苦的起诉和各种罚款，折磨和骚扰不只是一个城市或地区，而是他治下的所有省份。他夺走富人从祖先继承的财产，并将巨额财富和大笔金钱赐予他身边的谄媚者。\par

11. 他陷入了极度的愚昧和酗酒，以至于在狂欢中精神错乱和疯狂；他在醉酒时发出的命令，清醒后又后悔。他不允许任何人在放荡和堕落上超过他，反而使自己成为周围人（包括统治者和臣民）在邪恶上的导师。他鼓动军队在各种狂欢和无度中放荡生活，并鼓励总督和将领们以贪婪和贪欲虐待臣民，几乎就像与他共同执政一样。\par

12. 何必叙述此人的淫荡无耻行径，或列举他与多少人通奸？因为他每经过一个城市，必连续不断地奸污妇女，玷污处女。\par

13. 在这件事上，他成功了，唯独基督信徒除外。因为他们藐视死亡，所以毫不介意他的权势。这些人忍受了火、刀剑、十字架、野兽、深海、肢解、焚烧、刺穿、挖眼、全身伤残，此外还有饥饿、矿坑和锁链。他们为宗教而显出一切忍耐，胜于将属于天主的崇敬转移给偶像。\par

14. 妇女们也不逊于男子，为了圣言的教导而英勇奋战，与男子一同承受冲突，并获得了同等的德行之冠。当她们被拖去行淫时，她们宁可将生命交于死亡，也不将身体交于不洁。\par

15. 其中，在亚历山大里亚，有一位极其卓越、显赫的基督信徒妇女，被暴君抓去行淫，却以最英勇的坚定制胜了马克西米努斯那激情且不节制的灵魂。她因财富、家世和学识而受人尊敬，却将这一切视为不如贞洁重要。他多次催促她，但她已准备赴死；他却不能处死她，因为他的欲望强于愤怒。\par

16. 因此，他放逐了她，并夺去了她所有财产。还有许多其他人，甚至无法忍受异教统治者强暴的威胁，忍受了各种酷刑、折磨和致命的惩罚。\par

这些人确实值得钦佩。但最值得钦佩的是罗马的那位妇女，她真正是所有妇女中最尊贵、最贞洁的；暴君马克森提乌斯，在行为上与马克西米努斯如出一辙，试图凌辱她。\par

17. 她得知那些为暴君从事此类事务的人已在她家中（她也是一位基督信徒），并且她的丈夫，虽是罗马的行政长官，竟容许他们将带走她，她请求稍作妆扮，便进入内室，独自一人，以剑自刺。她立刻死去，将自己的尸体留给那些前来接她的人。她以行动，比任何言语都更有力地向现在和以后的所有人证明：基督信徒所拥有的美德，是唯一不可战胜、永不毁灭的财富。\par

18. 统治东西方的两位暴君同时推行了如此邪恶的行径。经过仔细考察，谁会犹豫不决，不把如此祸患归咎于针对我们的迫害呢？尤其当这场极度的混乱直到基督信徒获得自由后才止息。\par

% structure: p0130 source=h2 target=subsection reason=relative-heading-rank; base=h2

\subsection{第15章 发生在异教徒身上的事}

1. 在迫害的整整十年间，他们不断彼此谋害、互相争战。海路无法航行，人从任何港口出航都难免遭受各种暴行：被拖上刑架，肋旁被撕裂，为了通过种种酷刑查证他们是否来自敌方；最终受十字架或火刑。\par

2. 此外，盾牌、胸甲在准备，箭矢、长矛及其他军械也在备办，战船和海军装备到处征集。无人不随时预期会遭敌人攻击。除此之外，还有饥荒和瘟疫临到他们；关于这些，我们将在适当之处叙述必要的内容。\par

% structure: p0133 source=h2 target=subsection reason=relative-heading-rank; base=h2

\subsection{第16章 局势的好转}

1. 整个迫害期间的情况就是这样。但在第十年，借着天主的恩宠，迫害完全止息了；它自第八年后开始减弱。因为当天主和天上的恩宠向我们显示慈祥、眷顾的时候，我们的统治者，甚至那些曾热烈发动对我们战争的人，极其惊人地改变了心意，发布了一道撤销令，并以关于我们的仁慈谕令和法规，扑灭了那已被点燃的迫害大火。\par

2. 但这并非出于人为，也非如有人所说，出于我们统治者的怜悯或仁爱——远非如此，因为从一开始直到那时，他们每日都在筹划更严厉的措施对付我们，并使用更多样的刑具不断发明暴行——而是明明出于天主眷顾的照管：一方面与祂的子民和好，另一方面攻击那煽动这些祸患的人，向他显示愤怒，视他为全部迫害暴行的元凶。\par

3. 虽然这些事必须按照天主的审判发生，然而圣言却说：\BibleQuote{祸哉，那因他而使人跌倒的人！}\BibleRef{玛 18:7}因此，来自天主的惩罚降到他身上，从他的肉体开始，延及他的灵魂。\par

4. 因为他的私处突然生出一个脓肿，从中形成一个深陷的溃烂，无法控制地蔓延到他内部脏腑。由此生出不可胜数的蛆虫，散发出死亡般的恶臭；因为他整个身体的体积，由于贪食，在患病之前已变成一团过度柔软的脂肪，现已腐烂，对走近的人呈现出可怕而难以忍受的景象。\par

5. 有些医生完全无法忍受那极其难闻的气味，被杀死；另一些人，因整个躯体肿胀已无恢复希望，又无法给予任何帮助，也被毫不怜悯地处死。\par

% structure: p0139 source=h2 target=subsection reason=relative-heading-rank; base=h2

\subsection{第17章 统治者的撤销令}

1. \Emph{挣扎}于如此众多的灾祸中，他想起自己曾对虔敬者施加的残暴。于是，他将心思转向自己，首先公开向宇宙的天主认罪，然后召集侍从，命令立即停止对基督信徒的迫害，并应通过法律和皇谕敦促他们修建教堂、举行惯常的敬礼，为皇帝祈祷。话语一出，行动随即跟从。\par

2. 谕旨在各城公布，内含撤销反我等行动的诏令，形式如下：\par

3. \Person{Cæsar Galerius Valerius Maximinus}，\OriginalTerm{不败者}{Invictus}，\OriginalTerm{奥古斯都}{Augustus}，\OriginalTerm{大祭司}{Pontifex Maximus}，日耳曼征服者，埃及征服者，底比斯征服者，五次萨尔马提亚征服者，波斯征服者，两次卡帕多西亚征服者，六次亚美尼亚征服者，米底征服者，阿迪亚贝尼征服者，第二十次人民保民官，第十九次皇帝，第八次执政官，祖国之父，\OriginalTerm{代执政官}{Proconsul}；\par

4. 以及\Person{Cæsar Flavius Valerius Constantinus}，\OriginalTerm{虔诚的}{Pius}，\OriginalTerm{幸福的}{Felix}，\OriginalTerm{不败者}{Invictus}，\OriginalTerm{奥古斯都}{Augustus}，\OriginalTerm{大祭司}{Pontifex Maximus}，人民保民官，第五次皇帝，执政官，祖国之父，\OriginalTerm{代执政官}{Proconsul}；\par

5. 以及\Person{Cæsar Valerius Licinius}，\OriginalTerm{虔诚的}{Pius}，\OriginalTerm{幸福的}{Felix}，\OriginalTerm{不败者}{Invictus}，\OriginalTerm{奥古斯都}{Augustus}，\OriginalTerm{大祭司}{Pontifex Maximus}，第四次人民保民官，第三次皇帝，执政官，祖国之父，\OriginalTerm{代执政官}{Proconsul}；向各行省人民问安：\par

6. 在我们为公共福利和利益所颁发的其他事物中，我们先前曾希望将一切恢复为符合罗马人的古老法律和公共纪律，并规定那些背弃祖先宗教的基督徒也应恢复良好的心态。\par

7. 因为某种程度上，如此傲慢已占据他们，如此愚昧已笼罩他们，以致他们不遵行其祖先或许先前已确立的古老制度，反而各随己意，按自己的意愿制定法律并遵守，从而分散在各处聚集成独立的会众。\par

8. 当我们发布这道谕令，命他们返回古时所立的制度时，许多人虽面临危险仍顺从了，但许多人遭受折磨，忍受了各种死亡。\par

9. 既然许多人仍执迷于同样的愚妄，我们看出他们既不向天上诸神献上应得的敬拜，也不尊重基督徒的天主，考虑到我们的仁爱和一贯作风——我们惯于向所有人宽恕——我们决定应最乐意地在此事上亦施予恩免：他们可以再次成为基督徒，并重建他们惯常聚会的会堂，但条件是不得做出任何违反纪律的事。我们将在另一封信中指示地方官应遵守的事项。\par

10. 因此，因着我们的这份恩免，他们应祈求他们的天主，保佑我们、人民以及他们自身的安全，使公共福祉在每一处得以保存，使他们能在各自的家中安然生活。\par

11. 此敕令的大意如此，已尽可能从拉丁文译成希腊文。现在该思考这些事件之后发生的事了。\par

% structure: p0151 source=h2 target=subsection reason=relative-heading-rank; base=h2

\subsection{以下内容见于某些抄本第八卷中。}

1. 敕令的作者在此认罪后不久便脱离了痛苦而死。据说他乃是这场迫害苦难的始作俑者，早在其他皇帝行动之前，就已试图使军中的基督徒背弃信德，并首当其冲地针对自己家中之人，有的被褫夺军职，有的受尽凌辱，更有的面临死亡威胁，最终煽动与其共同执政者发动全面的迫害。对此等皇帝的结局，不可默然不提。\par

2. 当时四人掌握最高权力，年长且位尊者在迫害持续不到两年后，如我们已述，便逊位引退，以普通平民身份度过余生。\par

3. 他们的结局如下：第一位在年资和尊荣上最高，因长久而极其痛苦的肉体疾病而死。第二位则自缢身亡，依某种恶魔的预言，这是因他诸多大胆罪行而受的惩罚。\par

4. 在他们之后的最后一位，即我们说过是整个迫害的发起者，遭受了我们所述之事。但在他之前的\Person{君士坦提乌斯}皇帝，最为仁慈和善，其执政期间始终配其职分。而且，他以最和善慷慨之道对待众人。他对抗我等的战事不曾有丝毫参与，且将治下的虔诚信众完好无损地保全，未加凌辱。他既不拆毁教堂建筑，亦不对我们另施诡计。他一生结局幸福而三倍蒙福。他临终时，其帝业幸福荣耀地传给自己之子，这位儿子在各方面至为睿智与虔诚。他旋即继位，被士兵拥戴为最高皇帝和\OriginalTerm{奥古斯都}{Augustus}；\par

5. 并显出其效法父王对我们教义的虔诚。我们所述四人之死便如此，分别在不同时期发生。\par

6. 其中，唯有我们前面略提的那一位，与后来分享统治权的同僚，最终公开向所有人颁布了以上所述认罪书，即他所发布的那道书面敕令。\par

\SourceRecord{Translated by Arthur Cushman McGiffert.；Nicene and Post-Nicene Fathers, Second Series；Vol. 1.；Edited by Philip Schaff and Henry Wace.；Buffalo, NY: Christian Literature Publishing Co.,；1890.；<http://www.newadvent.org/fathers/250108.htm>.}

% structure: p0001 source=h1 target=section reason=page-title

\section{\WorkTitle{教会史（第九卷）}}

% structure: p0002 source=h2 target=subsection reason=relative-heading-rank; base=h2

\subsection{第一章　所谓的缓解}

1. 前述那道帝王的收回成命之谕，在亚细亚全境及邻近各省均有张贴。此事之后，东方的暴君——若有过此人的话，那便是最不虔敬者，对宇宙天主的宗教极端敌视的\Person{马克西米努斯}——对其中内容毫不满意，非但没有将上述诏令送交其下属总督，反而口谕他们放松对我们的战争。\par

2. 因为他无法以其他方式反对上级的决定，便将已颁布的法律秘而不宣，并留意不使其在所辖地区公布，而是以未成文的方式命令其总督们放松对我们的迫害。他们彼此以书面形式传达了此命令。\par

3. 至少，在他们中享有最高官职荣誉的\Person{萨比努斯}，以一封拉丁文书信将皇帝的旨意传达给各省总督，其译文如下：\par

4. 我们最神圣的君主皇帝陛下，以持续而最虔诚的热忱，曾引导所有人的心遵行神圣而正确的生活方式，使那些似乎过与罗马人不同生活的人，也应向不朽众神献上当有的敬拜。但某些人的顽固和最不可战胜的决心，竟至于既不能因命令的正当理由而背离其目的，也不能因即将来临的惩罚而畏惧。\par

5. 既然因此已发生许多这样的人因这种行径而陷于危险，我们最强大的君主皇帝陛下，以其敬神的崇高高贵，认为因这种原因将人置于如此巨大的危险与皇帝陛下的旨意不符，遂命令其忠仆我本人致书您的智慧：若有基督徒被发现参与其本族人的敬拜，您不应骚扰或危害他，也不应认为有必要以此为借口惩罚任何人。因为经过如此长时间的实践证明，他们无论如何也不能被说服放弃如此顽固的行径。\par

6. 因此，您应负责致函各城市的督办、市政官和区监，使他们知道不必再对此事多加关注。\par

7. 于是，各省的统治者，以为所写之事的意图已真正显明给他们，便以书面形式将皇帝的旨意传达给各城市的督办、市政官和行政长官。但他们不仅限于书面，而且更迅速地在行动中也完成了所认为的皇帝意愿。他们将那些因其对天主的宣认而被囚禁的人释放，又将那些已被送往矿场受罚的人中释放；因为他们错误地认为这正是皇帝的真正意愿。\par

8. 这些事既已完成，立刻，如同黑夜中闪耀的光芒，人们可以在每一城市看到聚集的会众、拥挤的集会，以及按照他们惯例举行的聚会。每一个不信的外教人对此都惊讶不已，惊叹如此奇妙的转变，并大声宣告基督徒的天主是伟大而唯一真实的。\par

9. 我们中的一些人，曾忠信而勇敢地承受了迫害的争战，如今又对所有人坦然而勇敢；但那些在信德上患病、被风暴摇撼了灵魂的人，则急切地寻求医治，恳求并哀求强壮者向他们伸出救助之手，并恳求天主怜悯他们。\par

10. 那时，宗教的高贵斗士们，那些从矿场的苦难中被释放出来的人，返回了自己的家园。他们幸福而欢欣地经过每一座城市，充满难以言喻的喜乐和无法用言语表达的勇气。\par

11. 大群的人沿着公路和市场前行，以赞美诗和圣咏赞美天主。你可能会看到，那些不久之前还在最残酷的判决之下被捆绑着赶出故乡的人，如今带着明亮而喜悦的面容回到自己的炉边；以至于那些从前渴求我们血的人，见到这出乎意料的奇迹，也因所发生的事而向我们祝贺。\par

% structure: p0014 source=h2 target=subsection reason=relative-heading-rank; base=h2

\subsection{第二章　随后的逆转}

1. 但是，那位正如我们所说统治东方地区的暴君——他完全憎恶善，是一切有德之人的仇敌——再也无法忍受此事；实际上，他并未允许事情以这种方式继续下去达六个月之久。他图谋一切可能的手段来摧毁和平，首先试图以借口阻止我们在墓地聚会。\par

2. 然后，通过一些恶人的中介，他派遣使团反对我们，鼓动\Place{安提约基雅}的市民向他请求一个极大的恩惠，即绝不允许任何基督徒居住在他们的地区；另一些人也被暗中鼓动做同样的事。在\Place{安提约基雅}，这一切的始作俑者是\Person{德奥特克诺斯}，一个暴烈而邪恶的人，一个骗子，其品格与其名字不符。他似乎是该城市的督办。\par

% structure: p0017 source=h2 target=subsection reason=relative-heading-rank; base=h2

\subsection{第三章　在安提约基雅新立的雕像}

此人对我们发动了各式各样的战争，并使人像搜捕不洁的贼一样，在我们的隐避处殷勤地搜捕我们的人，又对我们发明了各种诽谤和指控，成为无数人死亡的原因，最后他立了一尊\OriginalTerm{友谊的朱庇特}{Jupiter Philius}雕像，配以某些戏法和巫术仪式。他为此发明了不洁的入教仪式和不祥的神秘礼，以及可憎的净化方法，并通过他假装发出的神谕向皇帝展示其戏法；又通过取悦统治者的谄媚，挑动恶魔反对基督徒，宣称神已命令将基督徒作为他的敌人驱逐出城市及邻近地区的边界。\par

% structure: p0019 source=h2 target=subsection reason=relative-heading-rank; base=h2

\subsection{第四章 反对我们的陈情书}

1.\ 此人率先推动此事，并成功达到目的，这一事实激励了同一辖境内所有其他各城的官员也准备类似的陈情书。各省总督察觉到此事合皇帝心意，便暗示其属下也如法炮制。\par

2.\ 暴君以诏书表示对其举措深为满意，迫害便重新燃起，针对我们。于是各城设立了供像的司祭，此外由\Person{马克西米努斯}本人任命了高级司祭。后者选自那些在公共生活中最为显赫、曾在所任官职中享有盛誉的人，而且他们对于自己所崇拜的神祇怀有极大的热忱。\par

3.\ 简而言之，皇帝那异常的迷信导致他所有臣民——无论是统治者还是平民——为了取悦他，而无所不用其极地反对我们，以为通过密谋杀害我们、并向我们表现出任何新的恶意，就能最好地表达对他所赐恩惠的感激。\par

% structure: p0023 source=h2 target=subsection reason=relative-heading-rank; base=h2

\subsection{第五章 伪造的行实}

1.\ 于是，他们伪造了《\WorkTitle{比拉多行实及我们的救主}》，其中充满了对基督的各种亵渎，经皇帝批准后发送给他治下的整个帝国，并附书面命令，要求无论城乡各地，均须公开张贴于众人眼前，而且教师应将其交给学生，取代惯常的课业，供其研读并背诵。\par

2.\ 正在这些事发生之际，另一位罗马人称为\OriginalTerm{公爵}{Dux}的军事指挥官，在腓尼基的\Place{大马士革}市场上抓获了一些恶名昭彰的妇女，以施刑相威胁，迫使她们书面供称自己曾是基督徒，并知晓基督徒的恶行——他们甚至在教堂内行淫乱之事；此外，她们还按他的意愿对我们的宗教作了许多其他诽谤。他将她们的供词记录下来，呈报皇帝，皇帝下令也将这些文件在各处各城公布。\par

% structure: p0026 source=h2 target=subsection reason=relative-heading-rank; base=h2

\subsection{第六章 此时殉道的人}

1.\ 然而不久之后，这位军事指挥官便自戕身亡，为其邪恶付出了代价。但我们不得不再次忍受流放和严酷的迫害，各省总督又一次对我们变得极其凶狠；以致一些在\OriginalTerm{圣言}{Divine Word}上卓著的人也被捕，并被毫不留情地判处死刑。其中三人在腓尼基的\Place{埃梅萨}城，因承认自己是基督徒，而被投入野兽吞食。他们中有一位是主教\Person{西尔瓦努斯}，一位非常年迈的老人，已尽职达四十年之久。\par

2.\ 大约同时，\Person{伯多禄}——他极其杰出地主理\Place{亚历山大}的各堂区，因其卓越的品行和对圣经的研习而成为主教的神圣典范——也无缘无故、完全出乎意料地被捕，仿佛奉了\Person{马克西米努斯}之命，立即被斩首，不予解释。与他一同遇难的，还有\Place{埃及}的许多其他主教。\par

3.\ \Person{路济安}，安提约基雅堂区的一位司铎，各方面都是一位极其卓越的人，生活节制，以对神圣事物的学识著称，他被带到\Place{尼苛美底亚}城——当时皇帝恰好驻跸于此——在向统治者为其所信奉的教义作了辩护之后，被投入监狱并被处死。\par

4.\ 在短时间内，我们遭受了如此多的磨难，来自美德的敌人\Person{马克西米努斯}，以致这场对我们掀起的迫害似乎比先前的更为残酷。\par

% structure: p0031 source=h2 target=subsection reason=relative-heading-rank; base=h2

\subsection{第七章 刻在铜柱上反对我们的诏令}

1.\ 反对我们的陈情书以及所发布的帝国敕令副本，被镌刻并竖立在城中的铜柱上——这在别处从未有过。学校里的孩子们每日口中不离\Person{耶稣}和\Person{比拉多}的名字，以及那些狂妄地伪造的《行实》。\par

2.\ 我认为有必要在此插入\Person{马克西米努斯}这份张贴于铜柱上的文件，以便同时显明这位憎恨天主者的自夸与傲慢，以及紧接其后的、神圣的不容忍邪恶者对不敬者的惩罚——在惩罚的压力下，他不久之后便改变了对我们的做法，并以成文法确认之。\par

诏书内容如下：\par

\Emph{取自\Place{提洛}铜柱的\Person{马克西米努斯}针对反对我们的陈情书所作答复的译文副本。}\par

3.\ 如今，人类心智的软弱力量终于能够摆脱并驱散以往包围着那些比不敬者更为可怜之人的感官、并使其陷入黑暗而毁灭性的无知的一切错误迷雾，并且认识到，这世界是由永生的诸神慈爱的眷顾所治理并建立的。\par

4.\ 你们对敬虔的决心给出了最明确的证明，这令我们感到何等感激、何等愉悦、何等快慰，实难置信；因为在此之前，每个人早已知道你们对永生的诸神怀着何等的尊重与敬畏，所展示的并非空洞言辞的信心，而是持续而卓越的丰功伟绩。\par

5.\ 因此，你们的城邦完全可称为永生诸神的宝座和居所。至少，从众多迹象可以看出，它因天上诸神的临在而繁荣。\par

6.\ 看哪，你们的城邦不顾一切私利，搁置了以往为自身所呈的请愿，当它觉察到那可恶的虚妄的信徒又开始蔓延，并燃起最猛烈的火焰——如同被忽略而熄灭的火葬堆，当柴薪被重新点燃时——便立即投奔我们虔诚的怀抱，如同投奔一切宗教的首都，请求某种补救和援助。\par

7.\ 显然，由于你们的信仰与虔诚，诸神赐予了你们这救赎的心志。\par

因此，那位统治你们荣耀之城、保护你们祖先的神明、你们的妻儿、你们的炉灶与家园免于一切毁灭性灾祸的至高至强的\OriginalTerm{朱庇特}{Jove}，已将这有益的决意注入你们的灵魂；祂显示并证明，以合宜的虔敬遵守对不死神明的崇拜与神圣礼仪是多么卓越、荣耀与有益。\par

8. 因为谁能如此无知或如此缺乏一切理解，以至于不察觉，正是由于神明的慈爱眷顾，大地才不拒绝撒入其中的种子，也不以徒然的期望使农夫的希望落空；不敬神的战争并非不可避免地降临地上，腐败空气影响下衰败的躯体也不被拖向死亡；大海不被狂野的风暴吹得汹涌澎湃；意外的飓风不突然爆发，掀起毁灭性的暴风雨；此外，大地——万物的养育者与母亲——不因可怕的震动从最深的地基摇撼，其上的山脉也不沉入裂开的深渊。无人不知，这一切，乃至比这些更糟的灾祸，从前常常发生。\par

9. 所有这些不幸的发生，都是因为那些不敬神之人的虚妄空幻的毁灭性错误，当它充满他们的灵魂时，几乎可以说，使整个世界蒙受耻辱。\par

10. 在其他话语之后，他接着说：让他们观看那已茁壮成长、在广阔田地上摇曳着穗头的庄稼，以及那因丰沛雨水和恢复温和柔顺的气候而闪闪发光的草地与花朵。\par

11. 最后，让所有人都因最强大可畏的\OriginalTerm{玛尔斯}{Mars}之神威被我们的虔敬、祭献和尊崇所平息而欢喜；并因此享有稳固和平的安宁；让那些已经完全抛弃那盲目错误与迷惑、回归正确健全心志的人更加欢喜，如同从意外的风暴或重病中被解救出来，将在余生中收获欢乐的果实。\par

12. 但如果他们仍然坚持那可憎的虚妄，就让他们，如你们所愿，被远远驱逐出你们的城市与领土；这样，依你们在此事上值得称颂的热忱，你们的城市，摆脱一切污秽与不敬，便能按其本性，以合宜的虔敬从事不死神明的神圣礼仪。\par

13. 但为使你们知道，你们在此事上的请求在我们看来多么可悦，以及我们的心意多么愿意无需纪念与请愿而自动施恩，我们准许你们的虔诚，为报答你们这虔敬的心意，可以请求你们所愿的任何重大恩赐。\par

14. 现在你们可以请求，并要得到它；因为你们必立即获得。这恩赐，赐予你们的城市，将永远为对不死神明的虔诚信奉，以及你们因这一选择而从我们方的仁慈中获得应得的奖赏，提供证据；并且将显明给你们的子孙后代。\par

15. 这事在各省公布出来反对我们，剥夺了我们从人那里得到的一切美好希望；因此，按照那句神圣的话语，\BibleQuote{若可能，连被选者也会绊倒} \BibleRef{玛 24:24}。\par

16. 如今，当我们大多数人的希望几乎熄灭时，突然，那些要在某些地方执行上述诏令的人尚未完成旅程，天主——祂自身教会的保护者——就向我们显现了祂天界的代祷，几乎阻止了暴君对我们的夸口。\par

% structure: p0051 source=h2 target=subsection reason=relative-heading-rank; base=h2

\subsection{第八章 与这些事相关的灾祸：饥荒、瘟疫与战争。}

1. 冬季惯常的雨水和阵雨不再按往常的丰沛降在地上，出现了意想不到的饥荒，此外还有瘟疫，以及另一种严重的疾病，即一种溃疡，因其外表如火而得名\OriginalTerm{痈}{carbuncle}。这病蔓延全身，极大危害患者生命；但由于它主要侵袭眼睛，使无数男女老少失明。\par

2. 此外，暴君被迫与亚美尼亚人交战，他们自古就是罗马人的朋友和盟友。由于他们也是基督徒，热心敬拜神，这位天主的敌人曾试图强迫他们向偶像和魔鬼献祭，从而将朋友变成敌人，将盟友变成仇敌。\par

3. 这一切突然同时发生，驳斥了暴君对神明的空幻夸口。因为他曾夸口，由于他热心偶像并敌视我们，在他时代既无饥荒、瘟疫，也无战争。因此，这些灾祸同时齐至，也成为他自身灭亡的先兆。\par

4. 他本人及其军队在与亚美尼亚人的战事中战败，他治下其他城市的居民也惨遭饥荒和瘟疫的折磨，以致一斗小麦卖到两千五百\OriginalTerm{阿提卡德拉克马}{Attic drachms}。\par

5. 城中死者不计其数，乡间及村庄死者更多。以致从前包含大量农村人口的税册几乎完全注销；几乎所有人都迅速死于饥荒和瘟疫。\par

6. 因此，有些人想用自己最珍贵的东西，向那些供应较好的人换取一点点食物；另一些人则一点点变卖家产，陷入最后的贫困。有些人嚼着干草茬，轻率地吞食有害的草药，从而损坏和毁掉了自己的身体。\par

7. 有些城里的贵族妇女，因贫困所迫，沦落到羞耻的境地，外出到市场乞讨，她们面容谦逊、衣着得体，显露出往日良好的教养。\par

8. 有些人衰敗如鬼魂，奄奄一息，東倒西歪地到處蹣跚，因虛弱無法站立，倒在街中心；全身伸直躺臥時，乞求給一小塊食物，並以最後的氣息呼喊：飢餓！只有力氣發出這最痛苦的呼喊。\par

9. 但其他看似供應較好的人，見到乞討者眾多，大為震驚，在大量施捨之後，終究變得冷酷無情，預料自己不久也將遭受與乞討者相同的災禍。因此，市場和巷道中，赤裸的屍體多日無人埋葬，呈現給觀看者最可悲的景象。\par

10. 有些屍體也成了狗的食物，為此，倖存者開始殺狗，以免牠們發瘋去吞噬人。\par

11. 但更糟的是瘟疫，它吞噬了整個家庭和家族，尤其是那些因糧食充足而未被饑荒消滅的人。因此，富人、統治者、地方長官和眾多官員，彷彿被饑荒刻意留下來給瘟疫一般，迅速而突然地死去。每個地方都充滿哀號；每條巷道、市場和街道，除了眼淚、哀悼的樂器和聲音之外，別無所見所聞。\par

12. 如此，死亡以這兩種武器——瘟疫和饑荒——作戰，在短時間內摧毀了整個家族，使人們可以看到兩三具屍體同時被抬出。\par

13. 這就是\Person{马克西米努斯}的誇耀以及各城市對我們所採取之措施的報應。\par

於是，基督徒普遍的熱忱與虔敬的證據，向所有外邦人彰顯了出來。\par

14. 因為惟有他們在如此災禍中，以行動展現了同情和仁愛。每日不斷有人繼續照顧和埋葬死者，因為有許多人無人照顧；另有人將城內各地因饑荒受苦的人聚集一處，給他們所有人麵包；這事便在所有人中傳開，他們讚美基督徒的天主；並因事實本身而確信，承認惟有他們才是真正虔誠和敬畏天主的。\par

15. 這些事既已如此成就，天主——基督徒偉大而天上的護衛者——在已描述的事件中，向所有人顯示了祂對我們所遭受之大罪的忿怒和義憤之後，為我們恢復了祂眷顧的明亮而恩慈的陽光；以致在至深的黑暗中，一道平安之光從祂那裡最奇妙地照耀著我們，並向所有人顯明，天主自己一直是我們事務的統治者。祂時常管教祂的子民，以祂的懲罰來糾正他們，但當管教足夠之後，祂又向那些仰望祂的人顯示仁慈和恩惠。\par

% structure: p0068 source=h2 target=subsection reason=relative-heading-rank; base=h2

\subsection{第九章 蒙天主所愛的皇帝們的勝利}

1. 因此，\Person{君士坦丁}——我們已提及他為皇帝，生為皇帝之子，一位最虔誠最明智父親的虔敬兒子——和\Person{李錫尼}，僅次於他——兩位蒙天主所愛的皇帝，因其智慧和虔敬同受尊敬——因著天主、萬有的絕對統治者與救主的激勵，起來反對兩個最不虔敬的暴君，以天主為盟友，與他們正式交戰，結果\Person{马克森提乌斯}在\Place{罗马}被\Person{君士坦丁}以非凡的方式擊敗，而東方的暴君也未久存，在\Person{李錫尼}（當時尚未瘋狂）手下遭遇最恥辱的死亡。\par

2. \Person{君士坦丁}，無論在尊嚴和皇帝等級上都更為優越，首先憐憫\Place{罗马}受壓迫的人，並祈禱懇求天上的天主、祂的聖言，以及耶穌基督——萬有的救主——作為援助，然後率領全軍前進，決心恢復羅馬人祖傳的自由。\par

3. 但\Person{马克森提乌斯}信賴巫術，多於臣民的忠誠，不敢出城門一步，反而用龐大的軍隊和無數兵士，將他在\Place{罗马}附近及整個\Place{意大利}轄下的每個地方、區域和城鎮都設防加固。然而皇帝依靠天主的幫助，攻擊了暴君的第一、第二和第三軍，將他們全部擊敗；並且穿過\Place{意大利}大部分地區，已非常接近\Place{罗马}。\par

4. 然後，為了不讓他為了暴君而被迫與羅馬人作戰，天主親自將後者（彷彿用鎖鏈捆綁）拉出城門一段距離，並證實了那些古時記載在圣经中對不虔敬者的威脅——大多數人以為是神話而不信，但信徒卻相信——一句話，藉著行動本身，向所有親眼見到這奇事的人，無論信徒或非信徒，證實了這些威脅。\par

5. 因此，如同在\Person{梅瑟}本人和古時蒙天主所愛的希伯來民族時期一樣，\BibleQuote{他把法郎的戰車和軍兵投進海中，使他的精選戰車長沉入紅海，用洪水淹沒他們}，同樣，\Person{马克森提乌斯}也與他的士兵和護衛一同\BibleQuote{像石頭一樣沉入深淵}\BibleRef{出 15:5}，當他在天主的權能面前逃跑時，那權能與\Person{君士坦丁}同在；他穿過橫在路上的河流，他曾用船隻搭成橋樑，從而為自己的毀滅準備了工具。\par

6. 關於他，可以說：\Quote{他挖了坑，又挖深了，卻掉進自己所挖的坑裡；他的勞苦必歸到他頭上，他的不義必落在他的頂上。}\par

7. 这样，河上的桥被毁，通道下沉，船与船上的人立即消失在深渊中，那个最不虔敬的人自己首先，然后与他一起的持盾者，正如神圣的神谕所预言的，\BibleQuote{像铅沉在大水之中}；\BibleRef{出 15:10} 因此，那些从天主获得胜利的人，虽不是用言语，却用行动，如同天主的伟大仆人梅瑟和与他同在的人，恰当地唱着他们曾对古代不虔敬的暴君所唱的，说：\BibleQuote{我们要歌颂主，因为他大获全胜；将马和骑手投进海中；他成了我的救助者和保护者，为了我的救恩；}和\BibleQuote{主啊，众神之中谁能像你？谁能像你，神圣可畏，荣耀显赫，施行奇事？}\BibleRef{出 15:11}\par

8. 君士坦丁因他实际的行动，向天主，即宇宙的统治者、他胜利的赐予者，咏唱这些和类似的赞歌，当他胜利地进入罗马时。\par

9. 所有元老院成员和其他最显赫的人，连同全体罗马人民，包括儿童和妇女，以闪耀的目光和整个灵魂，以欢呼的喜悦和无限的快乐，接待他作为他们的解救者、救主和恩人。\par

10. 但他，作为天生对天主有虔敬的人，没有因欢呼而得意，也没有因赞美而自满；反而看出他的援助来自天主，立刻命令在他自己的雕像手中安放一个救主受难的标记。\par

11. 当他将它，即带有十字架救恩的记号，放在罗马最公开的地方，在其右手上，他命令用拉丁文刻上以下铭文：\Quote{藉此救恩的记号，即勇气的真实证明，我拯救并解放了你们的城市脱离暴君的桎梏，并且，释放了罗马元老院和人民，我恢复了他们古代的光荣和辉煌。}\par

12. 此后，君士坦丁本人和与他一起的皇帝李锡尼（他尚未陷入后来他堕入的那种疯狂）赞颂天主为一切福乐的赐予者，同心合意地制定了一项完全而最完整的诏令，以利于基督徒，并将天主为他们所做的奇妙之事和战胜暴君的消息连同诏令的副本，寄给仍在统治东方国家并假装对他们友好的马克西米努斯。\par

13. 但他，如同一个暴君，对所闻感到非常痛苦；但既不愿向别人让步，又不敢由于命令者的原因而压制所命令的事，仿佛出于他自己的权柄，他迫不得已地向他的属下发出了这份有利于基督徒的第一个文告，虚构一些从未被他做过的事来反对自己。\par

\Emph{暴君马克西米努斯书信的翻译副本。}\par

14. 约维乌斯·马克西米努斯·奥古斯都致萨比努斯。我相信以下事实对你和所有人都是明显的：我们的先辈戴克里先皇帝和马克西米安努斯皇帝，当他们看到几乎所有人都抛弃了众神的敬拜而依附基督徒的党派时，他们正确地下令，所有放弃对同一不朽众神敬拜的人，应当通过公开的惩罚和刑罚被召回对众神的敬拜。\par

15. 但是，当我第一次在吉祥的征兆下来到东方时，我得知在某些地方，许多能够提供公务服务的人因上述原因被法官放逐，我命令每一位法官，从今以后，他们不应再严厉对待各省居民，而应由奉承和劝诫召回他们对众神的敬拜。\par

16. 然后，当根据我的命令，这些命令被法官遵守时，结果，那些居住在东方地区的人中，没有人被放逐或侮辱，反而由于没有对他们施加严厉手段，他们被带回到对众神的敬拜。\par

17. 但后来，去年我在吉祥的征兆下去到尼科米底亚并暂居那里时，同一城市的市民带着众神的像来见我，恳切地请求，决不允许这样一群人居住在他们国家。\par

18. 但是当我得知有许多同一宗教的人居住在这些地区时，我回答说，我很高兴地感谢他们的请求，但我看出这请求并非由所有人提出，因此，如果有任何人坚持同一迷信，他有权按照自己的意愿行事，即使他希望承认对众神的敬拜。\par

19. 尽管如此，我认为有必要对尼科米底亚居民和其他同样向我提出这一请求的城市给予友好答复，即，他们的城中不应有基督徒居住——因为所有古代皇帝都采取了同样的做法，并且这也为众神所喜悦，一切人和国家政权本身都依靠他们——并批准他们代表他们神祇的崇拜所提出的请求。\par

20. 因此，尽管此前已发送特别信件给你的忠诚，并且同样已发出命令，不应采取严厉措施对待那些希望遵循这种做法的各省居民，而应温和适度地对待他们——然而，为防止他们受到受益人或其他人的侮辱或勒索，我认为应当在这封信中也提醒你的坚定，你应当以奉承和劝诫引导我们的各省居民承认对众神的关心。\par

21. 因此，如果有人出于自己的选择决定接受对众神的敬拜，应当欢迎他；但如果有人希望追随自己的宗教，你就留给他们权力。\par

22. 因此，你的忠诚理应守护所托付于你的事，确保无人有权以侮辱和勒索压迫我们的属民，因为如上所述，用劝诫和谄媚来将我们的属民召回敬拜众神，更为适宜。然而，为使我们的这道命令为所有属民所知，你有责任以你颁发的谕令，宣告所吩咐的事。\par

23. 因为他乃是出于被迫而行此事，并非出于本意下令，故无人视他为真诚或可信，因他先前做出类似让步之后，已显出其反复无常和欺诈的本性。\par

24. 因此，我们中无人敢举行聚会，甚至不敢出现在公共场合，因为他的通告并未涉及这些，只是命令避免加害于我们，并未下令我们可以聚会、建造教堂，或履行我们惯常的礼仪。\par

25. 然而君士坦丁和利西尼乌斯——和平与虔诚的拥护者——曾致信于他要求准许此事，并以谕旨和法令将之赐予其所有臣民。但这位最邪恶之人在这件事上不肯让步，直到被天主审判所迫，才最终被迫违背己意而行。\par

% structure: p0095 source=h2 target=subsection reason=relative-heading-rank; base=h2

\subsection{第十章 暴君的倾覆及其临终之言}

1. 迫使他采取这一行动的原因如下。由于他缺乏审慎和帝王之智，再也无法承担那本不配交托给他的庞大政权，他以卑鄙的方式处理政务，且心中毫无理性地因骄傲自大而高抬自己，甚至对那在出身、教养、教育、品德和才智，尤其是对真天主的节制与虔诚方面皆远超于他的帝国同僚，他也开始胆大妄为，自居首位。\par

2. 他因愚妄而疯狂，撕毁了与利西尼乌斯所立的条约，发动了一场无情的战争。于是，在短时间内，他使万事陷入混乱，搅动了每一座城市，集合了全部兵力，包括不可计数的士兵，前去与他交战，心中得意洋洋，寄望于那被他当作鬼神的魔鬼，以及他士兵的众多。\par

3. 交战时，他被剥夺了天主的眷顾，胜利因着那唯一的天主、万有的主宰，而归于当时统治的利西尼乌斯。\par

4. 首先，他所信赖的军队被摧毁；当所有卫兵都抛弃他，留下他独自一人，逃向胜利者时，他尽快秘密脱下那本不属于他的帝国服饰，以懦弱、卑怯、不男人的方式混入人群，然后逃窜，躲藏在田野和村庄中。然而，尽管他如此谨慎自保，仍几乎未能逃脱敌人之手，以此行为显明神圣的神谕是信实真实的，其中说道：\BibleQuote{君王不能因大军得救，巨人不能因大力得救；马是为获救而虚妄的，他也不能因力量之大而得解脱。} \BibleRef{咏 33:16-17}\par

5. \BibleQuote{看，主的眼目注视那些敬畏祂的人，注视那些期望祂仁慈的人，为救他们的灵魂脱离死亡。} \BibleRef{咏 33:18-19}\par

6. 如此，这位暴君满面羞惭，回到了自己的国度。首先，在狂怒中，他杀死了许多他先前所钦佩的诸神的司祭和先知——那些以其神谕煽动他发动战争的人——将他们作为巫师和骗子，此外还当作他安全的背叛者。然后，他将荣耀归于基督徒的天主，并颁布了一部关于他们自由的最完整全面的法令，随即被一种致命的疾病抓住，未得喘息，便离世了。他所颁布的法令如下：\par

\Emph{暴君为基督徒所颁谕旨之副本，译自罗马语文。}\par

7. 皇帝凯撒卡尤斯·瓦勒留·马克西米努斯，日耳曼征服者、萨尔马提亚征服者、虔诚者、幸运者、无敌者、奥古斯都。我们相信，这是显而易见的，无人不知，凡是回顾过去的人都知道并意识到，我们始终以各种方式关心属民的福祉，并愿向他们提供对所有人特别有益、符合共同利益和好处的事物，以及一切有助于公共福祉且合乎各人意愿的事。\par

8. 因此，在此之前，我们心中清楚，在奉我们的父母、最神圣的戴克里先和马克西米安的命令——即废除基督徒集会的命令——的借口下，官员们实行了许多勒索和掠夺；而且这些邪恶不断增长，损害了我们的属民——我们尤其渴望对他们施以适当关怀——结果他们的财产正在丧失。为此，去年我们向各省总督发送了信件，其中我们谕令：若有人愿意遵循此种习例或遵守此同一宗教，应允许其不受阻碍地追求其目的，无人可阻挠或禁止；并且所有的人都应享有自由，无需任何恐惧或猜疑，去做各人所偏好之事。\par

9. 但即使现在，我们也不禁觉察到，有些法官误解了我们的命令，使我们的百姓有理由怀疑我们法令的含义，并导致他们在履行那些合其心意的宗教礼仪时过于勉强。\par

10. 因此，为使今后一切畏惧疑虑的猜疑得以消除，我们已命令颁布此谕旨，以便使所有人明确：凡愿意皈依此派别和宗教者，均凭此恩准获得许可；并且每个人，照其意愿或喜好，均可按其习惯奉行他所选择遵守的宗教。我们还准许他们建造主的圣所。\par

11. 为使我们的这项恩赐更为宏大，我们以为还应颁布命令：凡是先前合法属于基督徒的房屋和土地，若因我们父母的命令而没入国库或为任何城市所没收——无论它们已被出卖或赠与他人——所有这些都应归还原主，即基督徒，以便在此事上众人也得知我们的虔诚与关怀。\par

12. 以上是那暴君的话，在他刻于石柱的反对基督徒的法令颁布后不到一年，便发布了这些谕令。而此前不久，我们还被视为不敬神的恶徒、无神论者和一切生命的毁灭者，以致我们不被允许在任何城市甚至乡村或荒野居住——如今，他却为基督徒颁布法令和谕令；那些不久前还在这暴君面前被火与剑、野兽和飞鸟所毁灭，受尽各种酷刑和惩罚，并作为无神论者和不敬神的恶徒而遭受最悲惨死亡的人，现在竟被他承认为拥有宗教的人，并被允许建造教堂；那暴君亲自作证并承认他们享有某些权利。\par

13. 他作了这样的承认后，仿佛因之得了些好处，或许所受的痛苦比应受的少；他被天主突如其来的鞭笞击打，在战争第二次战役中丧命。\par

14. 但他的结局不像那些军事首领，他们在战场上为美德和朋友勇敢作战，常光荣赴死；他如同一个不敬神的仇敌，当他的军队还在战场上列阵时，他却留在家里躲藏，遭受了应得的惩罚。因为他全身被天主突然的鞭笞击打，被可怕的疼痛和折磨所苦，扑倒在地，因饥饿而消瘦，他的皮肉被一种无形而天主所遣的火焚烧，以致他整个形貌改变，只剩下一副因岁月久远而干枯的骨架；因此，在场的人只能将他的身体看作他灵魂的坟墓——灵魂埋葬在一具已死且完全消融的身体里。\par

15. 当那热力更猛烈地燃烧他骨髓深处时，他的眼珠迸出，从眼眶脱落，使他失明。那时他仍在呼吸，向主自由承认，并呼喊死亡；最后，他承认因自己对基督的暴行而公正地遭受这些，于是断气身亡。\par

% structure: p0112 source=h2 target=subsection reason=relative-heading-rank; base=h2

\subsection{第十一章　宗教仇敌的最终毁灭}

1. 这样，当\Person{马克西米努斯}——那位曾是宗教仇敌中唯一存留且被视为其中最恶者——被除掉后，教会的重建从根基开始，藉着万有主宰天主的恩宠，基督的圣言照耀归于宇宙之天主的荣耀，获得了比先前更大的自由，而不敬神的宗教仇敌则蒙受极度的耻辱和羞惭。\par

2. 因为\Person{马克西米努斯}本人首先被皇帝们宣布为公敌，并通过公开告示被宣告为最不敬神、可憎恶和憎恨天主的暴君。那些为纪念他或其子女而在各城竖立的肖像，有的被从原处推倒在地，撕成碎片；其余的则被涂上黑漆而毁掉了面容。那些为他立起的雕像也同样被推倒、打碎，暴露在那些有意侮辱和嘲弄它们的人的嘲笑和戏弄之中。\par

3. 此外，所有其他宗教仇敌的尊荣也被剥夺，所有与\Person{马克西米努斯}同流合污的人都被处死，尤其是那些因谄媚而被他授予高官厚禄、并对我们的教理举止傲慢的人。\par

4. 这样的人有\Person{佩乌凯提乌斯}，他是他最亲密的伙伴，受他尊崇和奖赏超过众人，曾担任第二和第三任执政官，并被任命为首相；还有\Person{库尔基阿努斯}，他也同样经历了各级官职的晋升，并因在埃及无数次处决基督徒而闻名；此外还有不少人，正是通过他们，\Person{马克西米努斯}的暴政得以巩固和扩展。\par

5. 还有\Person{塞奥特克努斯}，也被那绝不放过他对基督徒所行之事的正义所追索。因为当他在安提约基雅竖立那雕像时，他似乎处于最幸福的状态，且已被\Person{马克西米努斯}任命为总督。\par

6. 但\Person{李锡尼乌斯}下到\Place{安提约基雅}城，搜查骗子，并拷问那新立雕像的先知和祭司，问他们为何施行欺骗。他们在拷问之下无法再隐瞒，便道出整个骗局是\Person{塞奥特克努斯}的诡计。因此，李锡尼乌斯对他们全体施行公正的判决：首先处死了\Person{塞奥特克努斯}本人，然后以最严厉的酷刑处死了他的同谋者。\par

7. 此外，还有\Person{马克西米努斯}的子女——他已通过将他们的名字刻在匾牌和雕像上，使他们分享皇权。那暴君的亲属，之前曾趾高气扬、傲慢地压迫众人，现在与上述之人同受惩罚，并蒙受极度的羞辱。因为他们未曾受教，也不明白圣经中的劝诫：\par

8. \BibleQuote{你们不要信赖权贵，也不要信赖世人，因为他们毫无救援；他的气一断，就要归于尘土；那一天，他的策划尽归乌有。}\par

不敬神的人就这样被清除后，政府得以稳固且无可争辩地归属于\Person{君士坦丁}和\Person{李锡尼}，他们理应享有此权。他们首先清除了世界对天主的敌视，因念及天主所赐的恩惠，便通过为基督徒颁布的法令，表现出他们对美德与天主的爱，以及对天主的虔敬与感恩。\par

\SourceRecord{Translated by Arthur Cushman McGiffert.；Nicene and Post-Nicene Fathers, Second Series；Vol. 1.；Edited by Philip Schaff and Henry Wace.；Buffalo, NY: Christian Literature Publishing Co.,；1890.；<http://www.newadvent.org/fathers/250109.htm>.}

% structure: p0001 source=h1 target=section reason=page-title

\section{教会史（第十卷）}

% structure: p0002 source=h2 target=subsection reason=relative-heading-rank; base=h2

\subsection{第一章　天主赐予我们的和平}

1. 愿一切感谢归于全能的天主，宇宙的统御者和君王，最大的感谢归于耶稣基督，我们灵魂的救主和救赎者，藉着祂，我们祈求那和平能永远稳固地保全在我们内，不为外在的困扰和内心的困扰所动摇。\par

2. 至圣的\Person{保利诺}，既然依照你们的意愿，我们已在先前各卷之后添加了《教会史》的第十卷，我们要将此卷献给你，宣告你是全书的印证；我们也将以完美的数目，恰当地添加关于教堂重建的完美颂词，顺从那以以下的话劝勉我们的天主圣神：\par

3. \Quote{你们应向主唱新歌，因祂行了奇事。祂的右手和祂的圣臂为祂施行了拯救。主已彰显了祂的救恩，在万民眼前揭示了祂的正义。}\par

4. 依照那命令我们唱新歌的训示，让我们继续表明，在我们所描述的那可怖的阴沉景象之后，我们现在得以看见并庆祝这样的事，即许多真确的义人和天主殉道者在我们之前渴望在地上看见却未曾看见，渴望听见却未曾听见的事。\BibleRef{玛 13:17}\par

5. 但他们匆匆前行，得到了更美好的事物，被带到天堂和神圣喜悦的乐园。我们承认这些事甚至超越了我们的配得，我们对那大恩的赐予者所显示的恩宠感到惊奇，我们理当钦崇祂，以我们灵魂的全部力量朝拜祂，并证实那些记载的言辞的真实性，其中说：\Quote{你们来看主的作为，祂在地上所行的奇事；祂消除战争直到地极，祂折断弓矛，焚烧盾牌。}\par

6. 我们因这些在我们时代明显应验的事而欢喜，让我们继续叙述。\par

7. 天主的全部仇敌种族以上述方式被消灭，如此突然地从人的眼前消失。于是，又一段神圣的言辞应验了：\Quote{我看见那恶人趾高气扬，像黎巴嫩的香柏树一样高举自己；我经过时，看，他已不在了；我寻找他的地方，却找不着。}\par

8. 最后，一个明亮而辉煌的日子，没有乌云遮蔽，以天上光明之光照耀着普世各地的基督教会。甚至那些不与我们共融的人，也未尝被阻止分享同样的福分，或至少受到它们的影响，享受天主赐予我们的恩惠的一部分。\par

% structure: p0011 source=h2 target=subsection reason=relative-heading-rank; base=h2

\subsection{第二章　教堂的重建}

1. 于是，所有的人都摆脱了暴君的压迫，摆脱了先前的苦难，各人以各自的方式承认护佑虔诚者的那一位是惟一真正的天主。我们这些寄望于天主基督的人，更是有着说不出的喜乐，一种灵感洋溢的欢欣为我们所有人绽放，因为我们看见不久前还被暴君的恶行所荒废的每一处地方，仿佛从一场漫长而致命的瘟疫中复苏，圣殿再次从地基上耸立至极高的高度，并获得了远比那些被毁的旧殿更大的辉煌。\par

2. 至高的统治者还通过再三为基督徒颁布的法令，更广泛地将天主的慷慨施恩确认给我们；皇帝的个人信函被送至主教们那里，附有荣誉和金钱的馈赠。将这些从罗马语译成希腊语的文献，恰如其分地插入本书的适当位置，如同刻在神圣的石版上，或许并非不恰当，以便它们作为纪念留存于所有后来者。\par

% structure: p0014 source=h2 target=subsection reason=relative-heading-rank; base=h2

\subsection{第三章　各地的奉献典礼}

1. 此后，我们所有人所渴望和祈求的景象出现了：在诸城中举行了奉献典礼和新建造的祈祷之所的祝圣仪式，主教们聚集，外乡人从远方前来会合，民众与民众之间展现了彼此的爱，基督身体的肢体在完全的和谐中合而为一。\par

2. 于是，那预言性地预示将来之事的先知话得以应验：\Quote{骨与骨相连接，关节与关节相结合，}\BibleRef{则 37:7}以及在那受默感的篇章中以隐晦表达所真实宣告的一切。\par

3. 有一种天主圣神的能力贯穿所有肢体，一个灵魂存在于众人之内，同样的信仰热忱，以及一篇从所有人发出的赞美天主的颂歌。是的，完美的礼仪由主教们主持，神圣的仪式被隆重举行，教会的庄严制度被遵守，此处有圣咏的歌唱和天主所交托给我们的经文的宣读，彼处有神圣而奥妙的礼典的执行；救主受难的奥妙象征得以施行。\par

4. 同时，各年龄的人，无论男女，都以整个心灵的力量，在祈祷和感恩中，以喜乐的心和灵魂，将尊崇归于天主，他们恩惠的赐予者。出席的每一位主教，各尽其能，发表了颂扬的演说，为聚会增添了光彩。\par

% structure: p0019 source=h2 target=subsection reason=relative-heading-rank; base=h2

\subsection{第四章　颂扬事务辉煌的赞词}

1. 有一位才能平庸的人，写成了一篇讲辞，在众多如同教会聚会般聚集的牧者面前站出来；他们安静得体地聆听，他以下面的话向一位在各方面最卓越、受天主爱戴的主教致辞，因他的热忱，在\Place{腓尼基}最为辉煌的\Place{提洛}的圣殿得以建成。\par

\Emph{教堂建造颂词，致提洛主教\Person{保利诺}}\par

2. 披上神圣长袍、佩带天上荣耀之冠的天主的朋友与司铎，以及圣神的灵感傅油与司祭之衣；而祢，啊，天主新圣殿的荣耀，蒙祂赐予年岁的智慧，却展现出青春与蓬勃德行的宝贵工作与事迹，——那位拥抱整个宇宙的天主，亲自授予祢这卓越的尊荣，为基督——祂的独生与首生圣言——以及祂神圣而神圣的净配，建造并更新这地上的殿宇；——\par

3. 有人或可称祢为新贝匝肋耳，神圣会幕的建筑师，或新耶路撒冷更美好的君王撒罗满，抑或新的则鲁巴贝耳，为天主的殿宇增添了远胜从前的荣耀；——\par

4. 而你们，啊，基督神圣羊群的乳儿，善言的居所，智慧的学校，以及庄严虔诚的宗教听众：\par

5. 许久以前，我们便已获准向天主高唱赞美诗歌，那时我们听闻圣经的宣读，得知天主奇妙的记号与主的奇妙作为赐予人的恩惠，受教如此说：\BibleQuote{天主啊！我们亲耳听见了，祖先也对我们讲述了祢在他们昔日、古时所作的事。}\par

6. 但现在，我们不再仅凭传闻或报告感知我们至慈天主与万有之王的高举之臂和天上的右手，而是仿佛亲身目睹古时所记的宣告是信实而真实的，我们获准高唱第二首凯歌，并扬声歌唱说：\BibleQuote{正如我们听见，也亲眼看见；在万军之主的城中，在我们天主的城中。}\par

7. 除了这座新建且由天主建造的城——它是\BibleQuote{永生天主的教会，真理的柱石和基础}，\BibleRef{弟前 3:15}——此外还能是哪座城呢？关于它，另有神圣神谕如此宣告：\BibleQuote{啊，天主的城！人们论及你，说了荣耀的事。}既然至慈的天主藉其独生子的恩宠将我们聚集于此，愿每一位被召的人扬声歌唱说：\Quote{我对那向我说\InnerQuote{我们要进入上主的殿宇}的人，甚觉欢喜}，以及\BibleQuote{上主，我爱慕祢圣殿的美丽，和祢荣耀所居之处。}\par

8. 让我们不仅各人单独，而是一齐同心同德，尊崇祂并高声呼喊说：\BibleQuote{上主伟大，在我们的天主城中，在祂的圣山上应受大赞美。}因为祂确实伟大，祂的殿宇也伟大，崇高宽敞，\BibleQuote{在世人中美丽绝伦}。\BibleQuote{上主伟大，独自施行奇事}；\BibleQuote{祂行伟大而不可测之事，荣耀而奇妙，不可胜数}；\BibleRef{约 9:10}伟大的是祂\BibleQuote{改变时季，废王立王}；\BibleRef{达 2:21}\BibleQuote{祂从尘埃中提拔穷苦人，从粪土中举扬贫困者}。\BibleQuote{祂从高位上推下有权势者，却举扬了卑微的人。祂使饥饿者饱享美物，却打碎了骄傲者的臂膀。}\BibleRef{路 1:52}\par

9. 祂不仅对信徒，也对不信者证实了古老事件的记录；祂是行奇事者，行大事者，万有的主宰，整个世界的创造者，全能、至慈、唯一的天主。让我们向祂唱新歌，在心中补充说：\BibleQuote{唯独行大奇事的天主，因为祂的仁慈永远长存}；\BibleQuote{祂击杀大君王，杀害有名之王，因为祂的仁慈永远长存}；\BibleQuote{上主在我们卑微时记念了我们，从仇敌手中救拔了我们。}\par

10. 让我们永不止息地向万有的圣父这样高声呼求。让我们口中常尊崇那第二个施恩者——神圣知识的导师，真宗教的教师，不敬神者的毁灭者，暴君的诛灭者，生命的革新者，耶稣，我们这些绝望者的救主。\par

11. 因为唯独祂，作为至慈圣父的唯一至慈之子，依照祂父慈爱的旨意，甘愿穿上我们这些卧于腐朽中的本性，如同一位卓越的医师，为拯救病人而察验其苦痛，触摸其恶臭的伤口，并从他人的悲惨中为自己招致痛苦；同样，我们不仅患病、被可怕的溃疡和已坏死的伤口折磨，甚至已卧于死者之中，祂却从死亡的口中将我们拯救归祂自己。因为天上没有其他者能有如此能力，毫不受损地服务于如此多人的救恩。\par

12. 唯独祂触及我们深重的腐朽，唯独祂承担了我们的劳苦，唯独祂承受了因我们不敬神所应受的惩罚，将我们这些不仅半死、且已躺在坟墓与墓穴中、全然污秽可憎者挽回，藉祂的仁慈热忱，出乎任何人甚至我们自己意料地，从古至今拯救我们，并慷慨分施圣父的恩赐——祂是生命与光明的赐予者，我们伟大的医师、君王与主，天主的基督。\par

13. 因为那时，当全人类因有罪恶魔的诡计和憎恨天主的邪灵之权势，被埋葬在幽暗黑夜与深深黑暗之中时，祂仅以显现，便以其光芒之射线，一次而永远地松开了我们罪孽紧绑的绳索，如同蜡被熔化。\par

14. 但是，当恶毒的嫉妒和邪恶的魔鬼对这般的恩宠与仁慈几乎气炸，并转过来向我们施展他所有致死的势力，起初，他像一条疯狗，向投掷它的石头咬牙切齿，将它的狂怒倾泻在投掷它的人身上，对着非生命的投射物发泄；他对准圣所的石头和房屋的无生命材料发泄其凶残的疯狂，摧毁了教堂——至少他以为如此——然后发出可怕的嘶嘶声和蛇一般的声音，时而通过不敬神的暴君的威胁，时而通过亵渎神明的统治者的渎神法令，吐出死亡，并且用他有害且毁灭灵魂的毒药感染被他俘获的灵魂，几乎用他充满死亡的、对死偶像的祭献杀死他们，并驱使每一个具有人形的野兽和每一种野蛮的物种攻击我们——那时，的确，\InnerQuote{伟大议会的天使}，大能的统帅，在他王国的至强士兵们通过忍耐和坚忍在一切事上表现出足够的操练之后，突然再次显现，涂抹并消灭了他的仇敌和对手，以致他们似乎连名字都不曾有过。\par

15. 但他的朋友和亲属，他提升到最高的荣耀，不仅在众人面前，也在天上的掌权者、太阳、月亮、星辰以及整个天地面前，以至于如今，正如前所未有地，至尊的统治者们，意识到从他那里所得的尊荣，唾弃死偶像的面孔，践踏魔鬼的不洁礼仪，嘲笑从他们祖先传下来的古老错觉，只承认一位天主，所有人的共同恩主，也包括他们自己。\par

16. 而且他们承认基督，天主子，万有的普世君王，并在纪念碑上宣称他为救主，以皇帝的诏书永垂不朽，在统治大地的城市中央，记录他的正义作为和对不敬神者的胜利。因此，我们的救主耶稣基督是唯一从永远就被甚至世间最高位者所承认的，不是作为人群中普通的王，而是作为普世天主的真子，并被敬拜为真天主，而且理应如此。\par

17. 因为哪个曾有过的国王达到了这样的德能，以致用他自己的名字充满了地上所有人的耳朵和舌头？哪个国王，在制定了如此虔诚而明智的法律之后，将它们从地极延伸到地极，以致它们被永久地在众人耳中诵读？\par

18. 谁曾用他温和且最爱人法律废除了未开化民族的野蛮和残忍习俗？谁，在全时代中被所有人攻击，却显示了如此超人的德能，以致每日兴盛，并一生保持年轻？\par

19. 谁曾建立一个民族，这民族古时甚至未曾被听闻，但现在并未隐藏在某个角落，而是四处散布在阳光之下？谁曾如此用虔诚的武器坚固他的士兵，以致他们的灵魂比金刚石更坚硬，在与对手的争战中光芒闪耀？\par

20. 哪个王能如此得胜，甚至在死后仍领导他的士兵，为他的敌人设立战利品，用他的王宫——甚至神圣的殿宇及其祝圣的祭品——充满每个地方、每个地区和城市，希腊的和外邦的，就像这座殿宇及其卓越的装饰和还愿供物一样？它们本身如此伟大而威严，值得惊奇和赞叹，是我们救主主权的清晰标志。因为如今，也\InnerQuote{他一发言，万有造成；他一命令，万物生成。}（\BibleRef{咏 33:9}）因为有什么能抗拒普世君王、治理者和天主圣言本人的命令呢？\par

21. 需要专门的论述来准确地审视和解释这一切；也要描述工人们的热情在被称为神圣的那位眼中是何等伟大，他注视着我们所有人所构成的活圣殿，并察看那由活石和移动的石头所建造的房屋，它稳固地建立在宗徒和先知的基础上，而房角石本身是耶稣基督，他不仅被那已不复存在的古老建筑的建造者所弃绝，也被那由众人堆砌的、至今仍存在的结构的建造者——邪恶工程的恶建筑师——所弃绝。但圣父无论在昔日还是今日都赞许了他，并使他成为我们这共同教会的房角石。\par

22. 谁看见这由我们自身形成的、永生天主的活圣殿——我说这最大且真正神圣的圣所，其内殿对群众是看不见的，真是神圣且是至圣所——敢去描述它？谁能甚至窥视圣所之内，除了众人大司祭，唯独他被允许探索所有理性灵魂的奥秘？\par

23. 但或许另一个人，唯一的一个，被允许在这同一事业中居第二位，即这支军队的统帅，第一位的大司祭亲自以第二位尊荣了他，作为你们神性羊群的牧者，他藉着父的分配和判断获得了你们这些子民，如同立他为自己的仆人和解释者，一个新亚郎或默基瑟德，相似于天主子，并按照你们众人同心的祈祷，由他存留并持续保守。\par

24. 因此，唯他被允许——若不是第一，至少是第一和最大司祭之后的第二位——观察和监督你们灵魂的深处；他藉着经验和时间准确地考验了每一个人，并藉着他的热忱和关切，使你们都在虔诚的行为和教义中得到安排，并且他比任何人更能对那藉着天主助佑他自己所完成的事，给予与事实相称的说明。\par

25. 至于我们的首位大司祭，经上说：\Quote{凡看见圣父所做的，圣子也同样做。} \BibleRef{若 5:19} 同样，这位司祭以纯洁的心灵之眼仰望他，视他为首位导师，以他所作的一切为原型，产生出它们的肖像，尽可能按相同的样式制造，丝毫不逊色于那位\OriginalTerm{贝匝肋耳}{Beseleel}，就是天主亲自\Quote{以智慧和聪敏之神充满了他} \BibleRef{出 35:31}，并赋予其他工艺和科学知识，召叫他建造那按天上象征性样式所造的圣殿。\par

26. 同样，这位也按自己灵魂中带有整个基督、圣言、智慧、光明的肖像，为至高天主建造了这座宏伟的圣殿，对应于那更大者，犹如可见者对应于不可见者；其灵魂的伟大、心智的丰富与慷慨，以及你们众人因捐献者的宽宏大量而表现出的竞相仿效，在实现同一目标上毫不逊色，真是难以言表。而这个地方——这值得首先提及——曾被敌人的诡计用各种垃圾覆盖，他没有忽视，也没有屈服于造成那种状况者的邪恶，尽管他本可选择其他地点，因为城中还有许多其他场所，在那里他可以少些劳苦，免于麻烦。\par

27. 但他首先自己投入工作，然后用热忱坚定了全体人民，将他们组成一个巨大的身体，进行了第一场战斗。因为他认为，这座特别被敌人围攻、首先与我们一同并为我们受苦、如同失去儿女的母亲一般的教会，应与我们一起在至仁慈天主特别的恩宠中喜乐。\par

28. 因为当大牧者赶走了野兽、豺狼和一切残忍凶暴的牲畜，并如神圣谕言所说，\Quote{打碎了狮子的颚骨}，他决意在同一地点重新聚集她的子女，并以最正义的方式设立她的羊栈，\Quote{使仇敌和报复者蒙羞}，驳斥天主的敌人不虔敬的狂妄。\par

29. 如今他们已不存在——那些仇恨天主的人——因为他们从未真正存在。在短暂地骚扰他人并自身受扰之后，他们遭受了当得的惩罚，并将自己、朋友和亲属引向完全毁灭，以致自古铭刻在圣书中的宣告已由事实证明为真。在这些宣告中，神圣的话语真实地论及他们，其中这样说：\par

30. \Quote{恶人拔出了剑，他们弯了弓，要射杀心里正直的人；愿他们的剑刺入自己的心，他们的弓被折断。} 又说：\Quote{他们的纪念归于消灭，} 并且 \Quote{你永远涂抹了他们的名号；} 因为当他们也遭难时，\Quote{他们呼求，却无人拯救；呼求上主，他却不听。} 但 \Quote{他们的脚被缠住，他们跌倒，我们却起来，站得正直。} 而那预先宣告的话——\Quote{上主啊，在你的城中，你必藐视他们的像}——已被众人亲眼看见为真。\par

31. 但他们像巨人一样与天主作战，就这样灭亡了。但那被世人遗弃、被世人拒绝的，因她向天主的忍耐，获得了我们所见到的圆满，以致以赛亚的预言论及她说：\par

32. \BibleQuote{干旱的旷野要喜乐，愿旷野欢乐，如百合花开放，沙漠要开花欢乐。} \BibleRef{依 35:1} \Quote{你们要加强软弱的手，稳固颤抖的膝。你们心里懦弱的人，要坚强，不要害怕。看，我们的天主必来报复，他必来拯救我们。} 因为他说：\Quote{在旷野中水已流出，在干渴之地有水池，干地要成为水浇的草地，在干渴之地要有水泉。}\par

33. 这些早已预言的事，记载在圣书中；但如今不只是传闻，而是事实。这旷野，这干地，这寡居被弃的，\Quote{他们用斧头砍伐她的城门，如同在森林中，用斧子和锤子将其毁坏，} 他们又毁坏她的书籍，\Quote{用火焚烧天主的圣所，亵渎你名的居所于地，} \Quote{凡路过的人都摘取她，折断她的篱笆，林中的野猪蹂躏她，田野的野兽吞吃她，} 如今因基督奇妙的权能，当他愿意时，她已变得如同百合花。因为那时她也如慈父般，受他的示意管教；\Quote{因为主所爱的，他必管教，又鞭打凡所收纳的儿子。}\par

34. 然后，在按情节所需受了一定的管教之后，她奉命重新喜乐。她如百合花开放，向众人发出神圣的馨香。因为经上说：\BibleQuote{旷野中水已流出} \BibleRef{依 35:6}，即神圣重生之洗的拯救之泉。如今她，不久前还是旷野，\BibleQuote{已成了水浇的草地，干渴之地有泉涌出。} \BibleRef{依 35:7} 先前\Quote{软弱}的手如今\Quote{确实强壮了}；这些工程是强壮之手的伟大而确凿的证据。先前\Quote{软弱无力}的膝，恢复了惯常的力量，正直地行走在神圣知识的道路上，奔向至慈牧者同类的羊群。\par

35. 若有人因暴君的威胁而心智麻木，拯救的圣言也不视他们为无可救药；他也医治他们，并催促他们接受神圣的安慰，说：\BibleQuote{你们懦弱的人，要受安慰；要加强，不要害怕。} \BibleRef{依 35:4}\par

36. 这位我们的新而卓越的\Person{则鲁巴贝尔}，听到了预先宣告的圣言——那曾因天主而成为荒芜的，在痛苦的掳掠和荒凉的可憎之后，应享有这些事——他没有忽略那死亡的躯体；而是首先以你们的共同共识，藉着祈祷和恳求，平息了圣父的义怒，并呼求那唯一使死者复活的作为他的同盟和同工，在净化并解除其病患之后，扶起了那倒下的。他给她穿上的，不是旧衣裳，而是他从神圣的谕示中再次学到的，那些谕示清楚地宣告：\BibleQuote{这殿后来的光荣，要比先前的大。}\BibleRef{盖 2:9}\par

37. 这样，他圈起更大的空间，用一道环绕整个建筑的围墙巩固了外院，作为整个建筑最安全的堡垒。\par

38. 他竖立并延伸了一个巨大而高耸的门廊，朝向初升的太阳的光芒，使那些站在圣所之外远处的人，能够充分看见里面的事物，几乎使那些信仰的陌生人转向入口，以致无人能不因记起从前的荒凉和现今不可思议的改变而留下印象。他期望这样的人受此感动，可能被吸引，并因这景象而被诱导进入。\par

39. 但当人进入大门后，他不允许他以不洁未洗的脚立即进入圣所；而是在圣殿与外门之间留下尽可能大的空间，用四个横向的回廊环绕并装饰它，构成一个四方的空间，四周立有柱子，并用木制的格子屏风连接，升至合适的高度；并在中间留出一片露天区域，以便看见天空，以及阳光中明亮的自由空气。\par

40. 在这里，他放置了神圣净化的象征，在圣殿对面设立了喷泉，提供充足的水，使进入圣所的人可以洁净自己。这是进入者的第一站；它同时为所有人呈现美丽而灿烂的景象，并为那些仍需要初级教导的人提供一个合适的停留处。\par

41. 但越过这一景象，他制作了通往圣殿的开放入口，在内又有许多其他的门廊，在一侧安装了三扇门，同样面向太阳的光芒。中间的门用铜板装饰，包有铁边，并刻有精美的浮雕；他将其做得比其他的更高更宽，仿佛让它们作为护卫，如同护卫一位王后。\par

42. 同样，他为整个圣殿两侧的走廊安排了许多门廊，并在他之上制作了各种通往建筑的开口，以便引入更多光线，用非常精美的木雕装饰它们。但君王之室，他用了更美丽更辉煌的材料，毫不吝啬地慷慨花费。\par

43. 我觉得在此赘述建筑的長阔、难以形容的辉煌与庄严，以及那耀眼的作品的外观、高耸入云的顶端和其上的昂贵的黎巴嫩香柏，似乎是多余的；这香柏是神圣的谕示未曾忽略提及的，它说：\Quote{上主的树木和祂所栽种的黎巴嫩香柏都要欢喜。}\par

44. 何必现在描述每一部分巧妙的建筑安排和超卓的美丽呢？因为亲眼所见已经使得耳闻的教导成为多余。当圣殿如此完成时，他为那些主礼者提供了高高的宝座，此外在整个建筑中按适当顺序排列了座位，最后在中央放置了至圣所、祭台，并且为使群众不能接近，他用精细雕刻的木格子围起来，呈现给观看者奇妙的景象。\par

45. 甚至连铺地也没有被他忽视；因为他也用各种美丽的大理石装饰了它。然后最后他转向圣殿以外的部分，提供了宽敞的半圆形壁龛和两侧的建筑，这些建筑连接到长方形教堂，并与内部结构的入口相通。这些是由我们最和平的\Person{撒罗满}——天主的圣殿的建造者——所建立的，为那些仍需要藉着水和圣神得到净化和洒扫的人，以便上述的预言不再仅仅是言语，而是事实；因为如今确实应验了\BibleQuote{这殿后来的光荣，比先前的大}。\BibleRef{盖 2:9}\par

46. 因为她是必要且相称的：既然她的牧人和主曾为她尝到死亡，并在他的苦难之后，把她所承受的那卑贱的身体改变成荣耀的身体，使那从朽坏中被救出的肉身进入不朽，那么她也应享受救主的安排。因为她从他那里接受了比这些大得多的应许，她渴望永远不停地与光明天使的歌团分享在重生中更伟大的光荣，\BibleRef{玛 19:28} 在不朽身体的复活中，在超越诸天的天主的宫殿里，与\Person{耶稣基督}自己——公共的恩主和救主——同在。\par

47. 但如今，那从前丧偶和荒芜的，因天主的恩宠被这些花朵所装扮，真正如同百合花，正如预言所说，她接受了婚礼的袍服和美丽的冠冕，并被\Person{依撒意亚}教导跳舞，以敬畏的话语向天主君王献上感恩祭。\par

48. 让我们听她说：\BibleQuote{我的灵魂因上主而欢喜；因为他给我穿上救恩的衣裳，给我披上喜乐的外袍；他像新郎一样给我戴上花冠，又像新娘一样用珠宝装饰我；有如大地生出嫩芽，花园使所种之物发芽；同样，上主天主使正义和赞美在万民面前生发。}\BibleRef{依 61:10}\par

49. 她用这些话欢跃。同样，天上的新郎——圣言耶稣基督也亲自回答她。请听主说：\Quote{不要惧怕，因为你没有蒙羞；不要羞愧，因为你没有受辱；因为你必忘记从前的羞耻，不再记念你寡居的羞辱。}\Quote{主没有召你如同被遗弃、心灰意冷的女子，也没有如同自幼被恨的女子。你的天主说：我离弃你不过片时，却要以大怜悯怜悯你；我曾在怒气中掩面不顾你，却要以永远的慈爱怜恤你。救赎你的主说。}\BibleRef{依 54:6}\par

50. \Quote{醒来，醒来，你喝了上主手中烈怒的杯；你喝了那使人东倒西歪的爵，甚至喝尽了。你从你所生的众子中，没有一个来引导你的；也没有一个搀扶你的手。} \Quote{看哪，我已将那使人东倒西歪的杯，就是我烈怒的爵，从你手中接过来；你必不再喝。我却要将这杯递在苦待你的人手中。}\par

51. \Quote{醒来，醒来，披上你的能力，披上你的荣美。抖下尘土，起来坐下；松开你颈项的锁链。} \Quote{你举目四面观看，他们都聚集来到你这里。我指着我的永生起誓，你必要以他们为妆饰佩戴，以他们为华带束腰，像新妇一样。因为你荒废、败坏、毁坏之处，现在因居民之故，必太狭窄；吞灭你的必远离你。}\par

52. \Quote{因为你所失丧的儿子必在你耳边说：这地方太狭窄，给我地方居住。那时你心里必说：我既丧子无夫，是谁给我生了这些？我独自一人，这些人在哪里呢？}\par

53. 这些是依撒意亚预言的，是古时在圣经中关于我们记载的；我们总有一天要从它们的应验中得知其真实无误。\par

54. 因为当新郎——圣言向他的新娘，就是神圣而圣洁的教会，说出这样的话时，这位伴郎——在她荒凉、如同尸体一般躺卧、在世人眼中毫无希望之际——按照你们众人同心合意的祈祷，理所当然地伸出你们的手，在天主——普世君王的命令下，在耶稣基督大能的彰显中，唤醒她、扶起她；扶起后，又按他从圣经的描述所学到的，将她建立起来。\par

55. 这诚然是一个极伟大的奇事，超越一切赞叹，尤其对于仅注重外表的人而言；但比奇事更奇妙的，是那些原型及其精神模型与神圣典范，我是指我们灵魂中那受感动的、理性的建筑的再现。\par

56. 神圣的圣子亲自按自己的肖像创造了这灵魂，在各方面都赋予她天主的模样：不朽的本性、非物质的、理性的、脱离一切尘世物质，一个具有自身理智的存在；当他一旦将她从无中召唤到有，便使她成为圣洁的新娘，成为他自己及圣父的至圣宫殿。这一点，他也用以下话语清楚宣告并承认：\Quote{我要在他们内居住，要在他们中往来；我要作他们的天主，他们要作我的子民。}\BibleRef{格后 6:16} 这就是那完美而纯洁的灵魂，从起初就这样被造，以承载天上圣言的形象。\par

57. 但因那恶毒的魔鬼的嫉妒和热忱，她出于自己的自由选择，变得感官化、喜爱邪恶，天主就离弃了她；她仿佛失去了保护者，便成了那些长久嫉妒她之人的轻易猎物和唾手可得之物；被那些看不见的敌人和属灵的仇敌用攻城器械和机器攻击，就遭受了可怕的坠落，以至于在她里面没有一块德行的石头留在另一块上，而是完全死在地上，完全失去了她关于天主的自然观念。\par

58. 但这按天主的肖像所造的她如此躺卧时，毁坏她的并非我们所见的林中野猪，而是一种毁灭性的魔鬼和属灵的野兽，它们用情欲如同自己邪恶的火箭欺骗她，烧毁了那真正神圣的天主圣殿，并亵渎了他名号的帐幕直到地上。接着，它们用一堆堆尘土埋葬了那可怜者，毁灭了一切得救的希望。\par

59. 但那光照如神、拯救的圣言，她的保护者，在她为自己的罪受了应得的惩罚之后，又恢复了她，赢得了那极其仁慈的圣父的恩宠。\par

60. 他首先征服了最高统治者的灵魂，然后通过那些最神圣的君王，从整个大地清除了所有不虔敬的毁灭者，以及那些可怕、仇恨天主的暴君。然后，他将那些属于他的朋友——他们早已为生命而奉献给他，但在罪恶的风暴中，仿佛在他的庇护下隐藏着——带到光中，用圣神的大恩赐相称地尊荣他们。并且，又藉着他们，用教义的规劝言语——如同铁锹和锄头——清除了那些不久前还被不虔敬法令的污秽覆盖、被各种质料和垃圾所累的灵魂。\par

61. 当他使你们所有心灵的土地洁净明亮之后，最终将它托付给这位充满智慧、深受天主喜爱的统治者；他既有判断力和谨慎，又具其他恩赐，能准确地考察和辨别托付给他的人的心灵，从第一天起直到今日，从未停止建造。\par

62. 如今祂已在你们众人中提供了光辉的金子，再是精炼无杂的银子，以及珍贵而贵重的宝石，以至于在事实中又为你们应验了神圣而神秘的预言，它说：\InnerQuote{看，我使你的石头成为红玉，你的根基为蓝宝石，你的城垛为碧玉，你的城门为水晶，你的城墙为精选的石头；你的众子都要受上主教导，你的儿女要享完全的平安；你必在正义中被建造。}\par

63. 因此，在正义中建造时，祂按各人的力量分派全体百姓。有些人祂只加固了外围的围墙，以真诚的信德砌成；这就是那大多数不能承受更大建筑的人。另一些人祂允许进入建筑，命令他们站在门口，为进来的人作向导；这些人可以恰当地比作圣殿的门廊。还有一些人祂以第一根柱子支撑，这些柱子安放在四方形大厅外，使他们初习四福音书文字的入门。还有另一些人祂联结在会堂两侧；这些是正在前进和进步的望教者，与忠信者所蒙受的圣事内观相距不远。\par

64. 祂从中取出那些如同金子般被圣洗洗净的纯洁灵魂，然后用比外部柱子更好、由圣经内在而神秘的教导制成的柱子支撑他们，并以窗户光照他们。\par

65. 祂以唯一普世君王和独一天主的荣耀的巨大门廊装饰整个圣殿，并在圣父权威的两侧安置基督和圣神作为第二光芒，祂在整个建筑中充足而荣耀地彰显其余一切细节中真理的清晰与光辉。祂从各处选取活生生的、移动的、预备好的灵魂石头，用它们建造那伟大而王家的房屋，内外都辉煌而充满光明；因为不仅灵魂和理智，连他们的身体也因纯洁和谦逊的盛开装饰而变得荣耀。\par

66. 在这圣殿中也有宝座，以及大量的座位和长凳，即在所有那些灵魂中，圣神的恩赐安坐其上，正如古时神圣的宗徒和与他们一起的人所见的，当时\InnerQuote{有舌头如火焰显现，分开落在他们每人头上}。\BibleRef{宗 2:3}\par

67. 但在众人之首，合理的推测是基督自己以其丰满居住，而在位居第二的那些人中，则按各人所能容纳基督和圣神的能力而定。此外，那些分别托付给每人教导和照顾的灵魂——可能是天使的座位。\par

68. 但那伟大、庄严而独特的祭台，除了一切司铎共有的灵魂的至圣所之外，还能是什么呢？耶稣自己，宇宙的大司祭，天主的独生子，站在这祭台的右边，以明亮的眼睛和伸出的手接受众人献上的馨香，以及他们祈祷中所献的、不流血的、非物质的祭品，并将它们带到天上的父、宇宙的天主面前。祂自己首先朝拜祂，唯独向父献上应得的尊敬，也恳求祂始终对我们众人仁慈和恩惠。\par

69. 这就是那伟大的圣殿，宇宙的伟大造物主——圣言——在整个世界中建造的，使之成为地上那些在天穹之上事物的理智形象，以便在包括地上理性受造物的整个受造界中，祂的父得受尊崇和朝拜。\par

70. 但那在天之上的领域，连同那里地上事物的模型，以及所谓的上天耶路撒冷，天上的熙雍山，永生天主的超世之城，在其中无数天使的歌团和名录于天上的首生者的教会（\BibleRef{希 12:22}），以我们无法言说、不能领悟的赞美诗歌颂扬他们的造物主和宇宙的最高统治者——凡有死的人谁能堪当称颂这事？\InnerQuote{因为眼所未见，耳所未闻，人心所未想到的，就是天主为爱祂的人所预备的。}（\BibleRef{格前 2:9}）\par

71. 既然我们、男人、孩童、妇女，小的和大的，已经部分分享了这些事，让我们全体同心同灵，不停息地宣认和赞美这些巨大恩惠的赋予者，\InnerQuote{祂赦免你的一切罪愆，医治你的一切疾病，救赎你的性命脱离死亡，以慈爱和怜悯为你加冕，以美物满足你的愿望。}\InnerQuote{祂没有按我们的罪恶对待我们，也没有按我们的过犯报复我们；}\InnerQuote{东离西有多远，祂使我们的过犯离我们也有多远。父亲怎样怜恤自己的儿女，上主也怎样怜恤敬畏祂的人。}\par

72. 现在和以后的所有时间，让我们将这些思想重新点燃在记忆中，日夜、每时每刻、每次呼吸，可以说，默观当前节庆以及这光明辉煌之日的创始者和主宰，用灵魂的一切力量爱祂、朝拜祂。现在起来，让我们大声恳求祂保护并保守我们到底在祂的羊栈中，赐予我们祂那永不间断、不可动摇的平安，直到永远，因我们的救主耶稣基督；愿荣耀藉着祂归于祂，直到永远。阿们。\par

% structure: p0093 source=h2 target=subsection reason=relative-heading-rank; base=h2

\subsection{第五章 帝国法律抄本}

1. 最后，让我们附上君士坦丁和李锡尼皇帝敕令的拉丁文译本。\par

% structure: p0095 source=h3 target=subsubsection reason=relative-heading-rank; base=h2

\subsubsection{译自拉丁文的帝国敕令抄本}

2. 很久以前，我们便认识到不应否认宗教自由，而应允许每个人按自己的判断和意愿履行其宗教义务。因此我们曾下令，每个人，包括基督徒和其他人，都应保持其自身派别和宗教的信仰。\par

3. 但既然在那道诏书中，在赐予他们如此自由的同时，似乎明确添加了许多不同的条件，他们当中有些人或许在稍后便退出了这种遵守。\par

\Quote{4. 我，\Person{君士坦丁·奥古斯都}，和我，\Person{李锡尼·奥古斯都}，在吉兆下抵达\Place{米兰}，并全面考虑一切关乎公共福祉和繁荣之事后，我们决定，在其他事项中，或者说首先，要颁布那些似乎对每个人都有益的法令；即，那些应保持对神祇的敬畏和虔诚的法令。我们决定，即，赐予基督徒和所有人选择他们自己宗教的自由，以便无论何种天上的神性都能眷顾我们和所有生活在我们统治下的人。}\par

\Quote{5. 因此，我们以健全正直的意图决定，不得拒绝任何人选择并遵循基督徒宗教仪式的自由，而应给予每个人自由，使其心思专注于他认为适合的宗教，以便神性在所有事物中向我们展示他惯常的关怀和眷顾。}\par

\Quote{6. 我们适宜地写下了这是我们的旨意：完全放弃我们之前关于基督徒写给尔忠诚的信中所包含的那些条件，使一切看来非常严苛且不符合我们温和之政令得以废止，并且现在每个同样渴望遵守基督徒宗教的人都可以不受骚扰地这样做。}\par

\Quote{7. 我们决定将此事最充分地通报尔之忠诚，以便尔知道，我们已经赐予这些基督徒自由和完全的自由去遵守他们自己的宗教。}\par

\Quote{8. 既然我们已自由地赐予他们这一切，尔之忠诚便认识到，自由也已赐予其他可能希望遵循自己宗教仪式的人；这显然符合我们时代的安宁，即每个人应有选择并崇拜他所喜悦的任何神祇的自由。我们这样做，是为了使我们绝不在任何方面歧视任何等级或宗教。}\par

\Quote{9. 关于基督徒，我们还进一步规定：他们从前通常聚集的场所，关于这些场所在之前致尔忠诚的信中已下达不同命令，若发现有人从我们的国库或从任何其他人那里购买了这些场所，应归还给上述基督徒，不得索取金钱或任何其他补偿，不得拖延或犹豫。}\par

\Quote{10. 如果有人碰巧接收了上述场所作为赠礼，应尽快将其归还给这些基督徒；但条件是，如果那些购买或接收这些场所作为赠礼的人要求从我们恩赐中得到任何补偿，他们可以到该地区的法官那里，以便由我们的仁慈为他们提供安排。所有这些事情应由尔之忠诚立即并无任何延迟地授予基督徒团体。}\par

\Quote{11. 既然已知上述基督徒不仅拥有他们通常聚集的场所，还拥有其他场所，这些场所不属于他们中的个人，而属于整个团体，即基督徒团体，尔将命令，根据我们上述律法，所有这些场所应毫不犹豫地归还给这些基督徒；即归还给他们的团体和集会；当然要遵守上述规定，即那些无偿归还者，正如我们之前所说，可望从我们恩赐中获得赔偿。}\par

\Quote{12. 在所有这些事情上，为了上述基督徒团体的利益，尔应竭尽全力，以便我们的命令能迅速执行，并且借此，通过我们的仁慈，也能为共同的和公共的安宁做出安排。}\par

\Quote{13. 因为通过这种方式，正如我们之前所说，我们在许多事情上已经经历的天主眷顾将在整个时间中持续稳固。}\par

\Quote{14. 为了使我们这恩典法令的条款为众人所知，我们希望我们写下的这些内容将由尔在各处公布并告知所有人，以便没有人不知道我们这一恩典法令。}\par

% structure: p0109 source=h3 target=subsubsection reason=relative-heading-rank; base=h2

\subsubsection{另一道皇帝谕令的副本，表明恩赐仅授予天主教会。}

\Quote{15. 致我们最尊敬的\Person{阿努利努斯}，问候。我们仁慈的习惯，最尊敬的\Person{阿努利努斯}，是愿意属于他人的权利不仅不受侵犯，而且应被归还。\Signature{再会，我们最尊敬和最亲爱的阿努利努斯。}}\par

\Quote{16. 因此，我们的旨意是，当你收到此信时，如果任何此类事物属于基督徒的天主教会，在任何城市或其他地方，但现为公民或任何其他人所持有，你应命令将其立即归还给上述教会。因为我们已决定，这些教会先前所拥有的事物应归还给他们。}\par

\Quote{17. 因此，既然你的忠诚见到我们的此命令最为明确，尔应尽快归还一切先前属于上述教会的事物，无论是花园、建筑，还是任何其他物品，以便我们知道你已最仔细地遵守了我们此命令。\Signature{再会，我们最尊敬和最亲爱的阿努利努斯。}}\par

% structure: p0113 source=h3 target=subsubsection reason=relative-heading-rank; base=h2

\subsubsection{一封书信的副本，皇帝在信中命令在\Place{罗马}举行主教会议，以促进教会的合一与和谐。}

\Quote{18. \Person{君士坦丁·奥古斯都}致\Person{弥尔提亚德}，\Place{罗马}主教，以及\Person{马尔库斯}。既然\Person{阿努利努斯}，\Place{非洲}最杰出的总督，已经向我发送了许多这样的信件，其中提到\Place{迦太基}城主教\Person{凯基利安努斯}被他在\Place{非洲}的一些同僚在许多事情上指控；并且在我看来，在那些天主眷顾已自由地交托给我忠诚的省份，且人口众多，民众似乎趋向更恶劣的路线，仿佛分裂为两派，主教们存在分歧，这是一件非常严重的事情——}\par

19. 我认为最好让\Person{采奇利亚努斯}本人，连同十位似乎是控告他的主教，以及另外十位他可能认为对辩护必要的同僚，一起乘船前往\Place{罗马}；在那里，在你们以及你们的同僚\Person{雷泰丘斯}、\Person{马特尔努斯}和\Person{马里努斯}（我已命令他们为此目的急速赶往\Place{罗马}）面前，听取他的陈述，正如你们所能理解的，这是符合至圣法律的。\par

20. 但为使你们能够对这些事情获得最完美的了解，我已将\Person{阿努利努斯}寄给我的文件副本附在我的信函中，并发送给你们上述的同僚。你们的坚定性阅读这些文件后，将考虑如何以最精确的方式调查和公正地裁决上述案件。因为你们的勤勉并不忽视我对合法天主教会怀有如此敬意，以致我不希望你们在任何地方留下分裂或纷争。愿伟大天主的尊威保佑你们，最尊敬的诸位，多年安康。\par

\Emph{一份诏书副本，其中皇帝命令召开另一次主教会议，以消除主教间的所有分歧。}\par

21. \Person{君士坦丁奥古斯都}致\Place{叙拉古}主教\Person{克雷斯特斯}。当有些人开始邪恶乖戾地就神圣敬礼、天上权能和公教教义彼此分歧时，为结束他们之间的争端，我先前曾命令从\Place{高卢}派遣几位主教，并从\Place{阿非利加}传唤那些持续固执争吵的敌对双方；在他们和\Place{罗马}主教在场的情况下，对引起骚乱的事情进行仔细审查和裁决。\par

22. 但由于正如所发生的，有些人既忘记了自己的救恩，也忘记了对于至圣宗教应有的敬意，甚至至今仍不肯终止敌意，不愿服从已经作出的判决，声称那些表达意见和决定的人为数甚少，或者他们在判决时过于仓促和草率，尚未对所有应当准确调查的事情进行审查——因此，那些本应彼此保持兄弟般的和谐关系的人，竟可耻地，更确切地说，可憎地彼此分裂，给那些灵魂远离此至圣宗教的人提供了嘲笑的借口。因此，我认为有必要规定：这种本应在判决已由他们自己自愿协议作出后就应当停止的分歧，如今若有可能，应通过许多人的出席来结束。\par

23. 既然我们已命令许多来自不同地区的主教于八月朔日之前在\Place{阿尔勒}城集合，我们也认为适宜写信给你，让你从\Place{西西里}总督、最杰出的\Person{拉特罗尼亚努斯}那里获得公共交通工具，并带上你自行选择的另外两位次级主教及三名沿途服侍你的仆人，在指定日期之前前往上述地点；以便通过你的坚定性以及在场其他人的明智同心与和谐，经过聆听那些目前彼此不和且我们也同样命令出席之人的陈述后，按照正确的信仰解决这一因某些可耻的争执而可耻地持续至今的争端，并恢复兄弟般的和谐，即使只是逐步实现。\par

24. 愿全能的天主保佑你多年健康。\par

% structure: p0122 source=h2 target=subsection reason=relative-heading-rank; base=h2

\subsection{六、帝国诏书：拨款予各教会}

1. \Emph{\Person{君士坦丁奥古斯都}}致\Place{迦太基}主教\Person{采奇利亚努斯}。既然我们乐意在\Place{阿非利加}、\Place{努米底亚}和\Place{毛里塔尼亚}各省向合法至圣天主教信仰的某些神职人员拨发经费以支付其开支，我已写信给\Place{阿非利加}卓越的财政大臣\Person{乌尔苏斯}，并指示他确保向你——坚定者，支付三千\OriginalTerm{弗利斯}{folles}。\par

2. 因此，当你收到上述款项后，应命令按照\Person{霍西乌斯}送交你的清单，将其分发给上述所有人。\par

3. 但倘若你发现尚有任何不足以实现我这一针对所有人的意图，你可毫不犹豫地向我们的司库\Person{赫拉克利德斯}索取任何你认为必需之物。因为当他出席时我已命令他：若你——坚定者——向他请求任何款项，他应确保立即支付。\par

4. 并且，由于我得知有些心怀不安的人试图以某种可耻的腐化手段使民众偏离至圣的天主教会，你应当知道：我已命令总督\Person{阿努利努斯}以及代理禁军长官\Person{帕特里丘斯}——他们当时在场——不仅要注意其他事务，更尤其要注意此事，且在类似情形发生时不得忽视。因此，你若发现任何人执迷不悟地继续这种疯狂行为，应立即前往上述官员处并向他们报告此事；使他们能够按照我当面对他们的命令加以纠正。愿伟大天主的尊威保佑你多年安康。\par

% structure: p0127 source=h2 target=subsection reason=relative-heading-rank; base=h2

\subsection{七、圣职人员的豁免}

% structure: p0128 source=h3 target=subsubsection reason=relative-heading-rank; base=h2

\subsubsection{帝国诏书：皇帝命令教会首领豁免一切政治义务}

1. 致问候予你，我们最敬重的\Person{阿努利努斯}。从诸多情况看来，当那宗教（其中保存着对至圣天上威权的最崇敬）被轻视时，公共事务便面临重大危险；但当它被合法采纳并遵守时，则藉着神圣的恩惠，为罗马的声名带来最显著的繁荣，为人类一切事务带来非凡的幸福——因此我以为，最敬重的\Person{阿努利努斯}，那些以应有的圣洁和对此律法的恒常遵守，献身于神圣宗教敬拜的人，应当为其辛劳获得酬报。\par

2. 因此，我意即：在你所辖的省区内，在\Person{塞西利亚努斯}所主持的公教会内，那些献身于这神圣宗教、通常被称为圣职人员的人，应完全豁免一切公共义务，以免他们因任何谬误或亵渎的疏忽而偏离对神祇应尽的侍奉，而能无碍地专务自己的律法。因为看来，当他们向神祇表达最大的崇敬时，国家便获得最大的利益。再会，我们最敬重且亲爱的\Person{阿努利努斯}。\par

% structure: p0131 source=h2 target=subsection reason=relative-heading-rank; base=h2

\subsection{第八章 \Person{李锡尼乌斯}后来的邪恶及其死亡}

1. 藉着我们救主的显现，神圣而天上的恩宠赐予我们如此福分，因我们所享有的和平，丰盛的恩惠遍及众人。于是，我们的事务在欢喜和庆节中达到顶峰。\par

2. 但恶毒的嫉妒和那喜爱邪恶的魔鬼，不能忍受看见这些事；而且，我们已提及的那些暴君所遭遇的事件，并不足以使\Person{李锡尼乌斯}归回理智。\par

3. 至于后者，虽然他的统治昌盛，且被尊为仅次于伟大皇帝\Person{君士坦丁}的第二位，并与后者有最紧密的婚姻纽带，却放弃了善行的效法，转而仿效那些他亲眼目睹其结局的不虔暴君的恶行，选择追随他们的原则，而不是与那位比他们更善者保持友好关系。他嫉妒共同的恩主，向他发动了一场不虔且最可怕的战争，不顾自然律、盟约、血亲，也不顾誓约。\par

4. 因为\Person{君士坦丁}，如同一位至为仁慈的皇帝，给予他真实恩宠的凭证，没有拒绝与他结盟，也没有拒绝将他的妹妹嫁给他，反而赐予他荣耀，使他分享先祖的尊贵和古老皇室的血脉，并授予他作为姻亲与共治者参与统治全帝国的权利，将不亚于他自己所拥有的罗马行省部分交给他治理和管辖。\par

5. 然而\Person{李锡尼乌斯}却反其道而行之；每日对上级策划各种阴谋，图谋各种恶事，以便以恶报恩。起初他试图隐藏自己的准备，假装成朋友，屡次施行欺诈和诡计，希望轻易达成目的。\par

6. 但天主是\Person{君士坦丁}的朋友、保护者和守护者，将暗中策划的阴谋暴露于光，挫败了它们。虔诚的伟大盔甲具有如此强大的力量，可以抵御敌人、保守我们的安全。受此护卫，我们最蒙天主眷顾的皇帝逃脱了那可憎之人的无数阴谋。\par

7. 但当\Person{李锡尼乌斯}察觉他的秘密准备丝毫未如他所愿进行——因为天主向那位蒙天主眷顾的皇帝揭露了每一项阴谋和邪恶——他再不能隐藏自己，便发动了公开的战争。\par

8. 并且在他决定与\Person{君士坦丁}开战的同时，也着手与宇宙的天主作战——他知道\Person{君士坦丁}所敬拜的就是这位天主——并开始暂时温和而隐秘地攻击\Person{君士坦丁}的虔敬臣民，这些臣民从未对他的政府造成任何伤害。他这样做是受其内在的邪恶驱迫，这邪恶使他陷入可怕的盲目。\par

9. 因此他没有将那些在他之前迫害基督徒之人的记忆放在眼前，也没有记住他自己曾被指定为那些因所行不虔之人的毁灭者和刽子手。但他偏离了理智，总之是被疯狂所攫，他决心与作为\Person{君士坦丁}盟友的天主自己作战，而不是与受天主助佑的那个人作战。\par

10. 首先，他将每个基督徒逐出他的宫廷，因而使自己——可怜的人——丧失了基督徒为他向天主献上的祈祷，而基督徒们按照他们祖先的教导，习惯为所有人祈祷。接着他下令，城市中的士兵除非选择向魔鬼献祭，否则应被开除并剥夺军衔。然而，这些与随后发生的更大事相比，还是小事。\par

11. 何必详细叙述那憎恨天主者所做的一切，以及这最无法无天的人如何制定不法的法律？他颁布法令：任何人不得向狱中的受苦者施以人道，给予食物；不得向被捆绑而饿死的人显示怜悯；任何人不得以任何方式行善或做好事，即使被大自然本身所感动而同情邻人。这确实是一项公然可耻且最残忍的法律，旨在驱逐一切天然的仁慈。此外，还作为惩罚规定：凡显示怜悯者，应与他们所怜悯之人同受其苦；凡仁慈地照料受苦者的人，应被投入锁链和监牢，与受苦者一同承受惩罚。\Person{李锡尼乌斯}的法令就是如此。\par

12. 何须详述他在婚姻及临终事宜上的革新——竟胆敢废止罗马人久经审慎制定且行之有效的古法，而引入某些野蛮残酷、确属非法且无视法纪的法令？他为了损害其统辖的行省，编造了无数指控及种种勒索金银的手段，施行土地新测度，并向那些已非在世、早已亡故的乡民横征暴敛。\par

13. 又有何必要详述这位人类公敌在以上种种之外，加诸无辜者的流放、他强夺出身高贵、名望卓著之人的年轻妻子，将他们交给其手下某些更为卑劣之徒肆意凌辱，以及众多有夫之妇和童贞女，虽自己年事已高，仍放纵情欲——我说，又有何必要详述这些事，因其最后行径的穷凶极恶，反使先前种种显得微末不足道？\par

14. 终于，他疯狂到极点，竟攻击主教，自以为他们——作为万有天主的仆人——必会敌视其措施。他尚未公开对付他们，因惧怕其上司，却仍如前秘密而诡诈地，利用总督的背信来消灭其中最杰出者。他们被杀害的方式甚是奇特，前所未闻。\par

15. 他在\Place{阿马西亚}及\Place{本都}其他城市行径的残忍超乎任何极端。天主的教会有些再度被夷为平地，有些被关闭，致使惯常前往者无人能进入并尽献于天主的敬拜。\par

16. 因为他邪恶的良心使他以为有人不为他祈祷；但他却深信我们一切所为皆为蒙天主宠爱的皇帝着想，并为他向天主恳求。因此他急匆匆将怒气转向我们。\par

17. 当时那些总督中谄媚他的人，察觉做这些事能取悦不虔暴君，便迫使一些主教承受惯常加于罪犯的刑罚，毫无理由地将无辜者如同杀人犯般拖走并处罚。有些人现在遭受一种新式死亡：身体被刀剑砍成许多碎片，在这野蛮而最可怖的景象之后，被抛入深海以喂鱼。\par

18. 于是天主的敬拜者再度逃亡，田野、荒漠、森林与山岳再次接纳了\Person{基督}的仆人。那不虔暴君在这些措施得逞之后，终于谋划重新掀起对众人的迫害。\par

19. 他本可如愿以偿，工作中无人能阻，若非天主——其子民生命的护卫者——以极速预见将发生之事，在如同黑暗阴沉的深夜中发出大光，为众人兴起一位拯救者，以大能的手臂将其仆人\Person{君士坦丁}领入那些地域。\par

% structure: p0151 source=h2 target=subsection reason=relative-heading-rank; base=h2

\subsection{第9章 \Person{君士坦丁}的胜利，以及他治下罗马帝国臣民所获得的祝福。}

1. 因此，天主从天上赐予他相称虔敬的果实，对不虔者的胜利之荣冠，并将那有罪者连同其所有谋士与朋友伏于\Person{君士坦丁}脚下。\par

2. 因\Person{李锡尼}将狂妄推到极致，这位爱天主之皇帝，认为不能再容忍，基于健全的判断，将公正的坚定原则与仁爱结合，欣然决定前来保护受暴君压迫的人，并着手除去几个毁灭者，以拯救人类的大多数。\par

3. 因为当他从前仅施行仁爱，对那不值得同情的人施以怜悯，却一无所成；\Person{李锡尼}并未放弃其邪恶，反而对治下百姓更加愤怒，受难者无救恩之望，如同被一头凶残野兽所压迫。\par

4. 因此，美善的保护者，混合对恶的憎恨与对善的喜爱，与其子\Person{克里斯普斯}——一位最仁慈的君主——一同出征，向所有垂死之人伸出救援之手。他们父子，正如在万有之君天主的护卫下，以\Person{圣子}——万物救主为统率和盟友，从四面展开军队对抗天主的敌人，赢得轻易的胜利；天主在战斗中事事使他们如愿以偿。\par

5. 于是，骤然间，快过述说，昨日及前日还在吐露死亡与威吓之人不复存在，连他们的名号也无人记念，但他们的碑文和荣誉却遭受应有的羞辱。\Person{李锡尼}曾亲眼目睹先前不虔暴君所遭遇之事，他自身也同样承受，因为他未受教训，也未从邻人的惩罚中学会智慧，反而步上他们所行的不虔之路，公正地被投入同一深渊。如此他倒伏在地。\par

6. 而\Person{君士坦丁}，这位最有力的胜利者，佩戴一切虔敬美德，与其子\Person{克里斯普斯}——一位最蒙天主钟爱、处处肖似其父的君主——收复了属于他们的东方；他们形成一个统一的罗马帝国，如同古时，将整个世界从日出处到对面区域，包括南北，直到日落之极，都置于他们和平的统治之下。\par

7. 因此，一切先前迫害者的恐惧从人心中被除去，他们欢庆盛大而喜庆的日子。万物充满光明，那些先前沮丧的人们面带笑容、目光炯炯地彼此相望。在城中与乡间，他们以舞蹈和讚歌首先光荣天主，万有的君王，因为他们如此受教，然后光荣虔诚的皇帝及其天主所爱的儿女。\par

8. 过去的邪恶被遗忘，每一件不虔敬的行为都被忘却；他们享受当前的恩惠，并期待未来的恩惠。得胜的皇帝在各地颁布充满仁慈的谕令，以及蕴含善意和真虔敬的法律。\par

9. 如此，在一切暴政被清除之后，属于他们的帝国得以稳固且无匹敌地只归君士坦丁及其儿子们。他们消除了前任的不敬神，承认天主所赐的恩惠，以他们在众人眼前所行的作为，显露出他们对德行的爱、对天主的爱，以及他们的虔敬和对天主的感恩。\par

\Emph{在天主助佑下，\Person{欧瑟伯·庞斐利}的\WorkTitle{教会史}第十卷终。}\par

\SourceRecord{Translated by Arthur Cushman McGiffert.；Nicene and Post-Nicene Fathers, Second Series；Vol. 1.；Edited by Philip Schaff and Henry Wace.；Buffalo, NY: Christian Literature Publishing Co.,；1890.；<http://www.newadvent.org/fathers/250110.htm>.}
